\documentclass[11pt,oneside,openany]{book}

% 1151 憲法專題研究（東吳大學法律學系學士班法二Ａ，湯德宗老師，BLW52401）課堂逐字稿。
% 一週一部：Teile/<週>.tex 由 Skripte/gen_tex.py 自 Notizen/<週>/ 各表產生（cue 表、校訂、斷句、說話人、編註）。
% 每個 cue 前有 % cue N 起秒 迄秒，hypertarget 為 w01:cue:N 之形。編註在 Notizen/<週>/編註.tsv。
% 宏包：Translation 倉 LaTeX/course-notes.sty。編譯：./build.sh。

\def\TransLaTeXDir{../../../../1142/Translation/LaTeX}
\input{\TransLaTeXDir/course-notes.sty}
\addbibresource{refs.bib}
\CourseNotesMeta{憲法專題研究（大學部）課堂逐字稿}{王逸帆}
\CourseNotesChinese[部]

% 段首時間碼：排在左邊界外，小字灰色
\newcommand{\tc}[1]{\noindent\llap{\makebox[0pt][r]{\scriptsize\color{gray}#1\hspace{0.8em}}}}
% 說話人標籤（老師以外者，或回到老師時）
\newcommand{\speaker}[1]{{\small\bfseries #1：}}
\setlength{\emergencystretch}{2em}

\newcommand{\documentversion}{v0.1（2026-11-24）}

\begin{document}
\frontmatter
\coursetitle{憲法專題研究}{課堂逐字稿}%
  {東吳大學法律學系學士班　憲法專題研究（權力分立原則）\\115 學年度第 1 學期　湯德宗老師}%
  {\documentversion}
\tableofcontents
\mainmatter
% 自 Notizen/W01_0914/ 各表由 Skripte/gen_tex.py 產生；改字改表後重產，不手改本檔。
\chapter{課程安排與權力分立原則導論（2026-09-14）}

2026 年 9 月 14 日下午第 7、8 節課堂錄音（約 124 分鐘）之逐字稿。左邊界為段首時間碼；未標說話人者為湯德宗老師。


\section{課前事務與教室設備}

\tc{0:00:01}
% cue 1 1.060 45.470
\hypertarget{c01:cue:1}{}第一件事情，幫老師插桌子，這個我給你提包，我給你提包，芝子巾。幫老師來之前的小事情，另一件事情就是，我們等一下課堂上會有一些事情，那它最主要是今天這節課，下課以後，比較辛苦，要幫老師把這個欄上的尊重，尊重，主要是這個，剩下的事情就是聯繫啦，哪一個同學，有什麼初春的，
% cue 2 45.590 63.850
\hypertarget{c01:cue:2}{}這個狀況呢，這個下個禮拜，他比方說，他原來要報告的，他不行啊，什麼狀態啊，他設計要做一些什麼處理啊，會是像這樣的事情，沒有什麼特殊的事情，會想到的是什麼，那你們有清楚嗎？

\tc{0:01:03}
% cue 3 63.850 85.160
\hypertarget{c01:cue:3}{}有，但是沒有，有，好，那你們現在手上，知道上課的做課計畫拿到了嗎？有人，你有看到嗎？你們是不是要放在什麼地方？學校有一個什麼，教室的地方？
% cue 4 85.160 88.350
\hypertarget{c01:cue:4}{}你們是不是等聽不懂，在教，不然你們怎麼連到什麼上課？
% cue 5 88.350 98.490
\hypertarget{c01:cue:5}{}第二，請問，這不是學校規定，老師都要教的嗎？有，亂講嗎？沒有啊？是。
% cue 6 98.490 117.100
\hypertarget{c01:cue:6}{}對啊，那你們要上課，要選課，與否，就是要看這個，不然要看什麼？有人會開這個機器嗎？
% cue 7 117.100 144.990
\hypertarget{c01:cue:7}{}對，謝謝，學校的朋友，不該，不所要出現，哪兩個辦法，是這個事情。
% cue 8 144.990 148.730
\hypertarget{c01:cue:8}{}有一個按鈕，它按下去，它就是對面的男生。
% cue 9 148.730 152.780
\hypertarget{c01:cue:9}{}你現在身邊不太喜歡，找一個男生好不好？
% cue 10 152.780 167.290
\hypertarget{c01:cue:10}{}因為系裡面講一下，說這個機器沒有開，我是要用的。
% cue 11 167.290 169.290
\hypertarget{c01:cue:11}{}你跟誰講，我也不知道，你去跟西班講，我也不知道。
% cue 12 169.290 182.400
\hypertarget{c01:cue:12}{}有，有，有，有，知道我們自己，好，你知道我們人不知道是誰嗎？好。

\tc{0:03:02}
% cue 13 182.400 194.800
\hypertarget{c01:cue:13}{}我有一個問題，我們現在，首先第一電視機是我剛才講的，你們需要開，或者是會進化給他。
% cue 14 194.800 203.630
\hypertarget{c01:cue:14}{}要不然我們就，不知道在講什麼，所以到時開始講比較多，有大堆。
% cue 15 203.630 230.180
\hypertarget{c01:cue:15}{}我現在是沒有，接著，但是我想教室演，是嗎？可以，好，不要講。
% cue 16 230.180 242.030
\hypertarget{c01:cue:16}{}大老師講，還是你們要自己看，不要問，好，比較喜歡，習慣的是聽到老師講嗎？
% cue 17 242.030 253.430
\hypertarget{c01:cue:17}{}但是等一下，有你們講的，想什麼，沒有我們。

\section{指定教科書}

\tc{0:04:13}
% cue 18 253.430 272.110
\hypertarget{c01:cue:18}{}反正學習會有一些書要用，但是這個，蘇來呢，再說，現在還在，老師應該在做最後的修訂。
% cue 19 272.110 278.310
\hypertarget{c01:cue:19}{}大概是 500，現在在我再講的，指定教科書，特別是第一的，卷一\footnote{湯德宗著《憲法結構與動態平衡：權力分立新論 卷一》，新學林，增訂五版，2026 年，課綱所列三部指定教科書之一。課綱第 2、4、8、9、10、12 週的指定閱讀頁碼出自此卷。}的部分。
% cue 20 278.310 285.380
\hypertarget{c01:cue:20}{}卷二\footnote{湯德宗著《違憲審查與動態平衡：權力分立新論 卷二》，新學林，增訂五版，2026 年。課綱第 6 週（10 月 19 日）所列釋字第 3、175、14、319、461、632 號的指定閱讀頁碼出自此卷。}當然也很好，但是卷二對這個學習的上課，不夠遠一點，不要是卷一，卷二比較難。
% cue 21 285.380 297.040
\hypertarget{c01:cue:21}{}那有聽過嗎？OK，那反正蘇來，再說，剩下來，是老師的釋憲學思心路\footnote{湯德宗著，上、中、下冊，新學林，二版，2026 年；初版題為《釋憲學思心路：湯大法官德宗意見書心法尋繹》，時英出版社 2025 年出版。課綱第 13、14、16 週的指定閱讀頁碼出自此書。}，上中下三冊。

\tc{0:04:57}
% cue 22 297.040 301.380
\hypertarget{c01:cue:22}{}請問有看過，有聽過這個書的人舉手，有看過嗎？
% cue 23 301.380 303.640
\hypertarget{c01:cue:23}{}有一點沒有，那是那些聲音嗎？
% cue 24 303.640 306.530
\hypertarget{c01:cue:24}{}你一開始的時候，你會有嘎，聲音，他有一個聲音，不一想。
% cue 25 306.530 309.940
\hypertarget{c01:cue:25}{}有聽過釋憲學思心路，看過學思心路的修成。
% cue 26 309.940 321.840
\hypertarget{c01:cue:26}{}你講，哦，你同一塊，你一生會看過，去年就來參加了。
% cue 27 321.840 325.840
\hypertarget{c01:cue:27}{}去年是五月十五花半年，在學校。
% cue 28 325.840 337.090
\hypertarget{c01:cue:28}{}對，哦，你那個，學做下課，有很多事情。

\tc{0:05:37}
% cue 29 337.090 342.110
\hypertarget{c01:cue:29}{}所以班長，現在是真的要猜想這個朋友園。
% cue 30 342.110 346.030
\hypertarget{c01:cue:30}{}我們要學會自己在開這個東西。
% cue 31 346.030 349.730
\hypertarget{c01:cue:31}{}所以我有點驚訝，因為以前我們每次上課。
% cue 32 349.730 352.330
\hypertarget{c01:cue:32}{}班上總是有個會。
% cue 33 352.330 354.770
\hypertarget{c01:cue:33}{}所以好像很少教導人。
% cue 34 354.770 356.770
\hypertarget{c01:cue:34}{}欸，是不是，那你關上什麼事情。
% cue 35 356.770 359.400
\hypertarget{c01:cue:35}{}怎麼開，要跟同學講一下。
% cue 36 359.400 361.400
\hypertarget{c01:cue:36}{}你現在要好，對，試試看。
% cue 37 361.400 368.440
\hypertarget{c01:cue:37}{}現在不會有聲音，這裡就會有聲音。
% cue 38 368.440 374.180
\hypertarget{c01:cue:38}{}它有一個鈕，我們要按下去，然後有看到一聲的。

\tc{0:06:14}
% cue 39 374.180 376.180
\hypertarget{c01:cue:39}{}那剛才那個同學，老師來請問一下。
% cue 40 376.180 378.240
\hypertarget{c01:cue:40}{}你一次打了五年級嗎。
% cue 41 378.240 383.340
\hypertarget{c01:cue:41}{}三年級，表示你很厲害。
% cue 42 383.340 392.240
\hypertarget{c01:cue:42}{}那你是三級，3CA,3C,3D,3D,3D。
% cue 43 392.240 397.020
\hypertarget{c01:cue:43}{}3D，所以你是朱小語，朱小語。
% cue 44 397.020 400.630
\hypertarget{c01:cue:44}{}朱小語，你可以跟我講一下。
% cue 45 400.630 402.630
\hypertarget{c01:cue:45}{}你看那個書，有什麼感覺。
% cue 46 402.630 406.200
\hypertarget{c01:cue:46}{}好，你經過。

\tc{0:06:46}
% cue 47 406.200 408.160
\hypertarget{c01:cue:47}{}我們現在同學可以做一件事情。
% cue 48 408.160 409.160
\hypertarget{c01:cue:48}{}畫得做他們的。
% cue 49 409.160 412.410
\hypertarget{c01:cue:49}{}我們現在手機打開了，你打到四個名字。
% cue 50 412.410 414.400
\hypertarget{c01:cue:50}{}然後你打四個名字，四個名字。
% cue 51 414.400 415.490
\hypertarget{c01:cue:51}{}看你一定能看到書的樣子。
% cue 52 415.490 423.090
\hypertarget{c01:cue:52}{}書的樣子，能不能再多。
% cue 53 423.090 427.470
\hypertarget{c01:cue:53}{}因為現在書第一版的。
% cue 54 427.470 430.770
\hypertarget{c01:cue:54}{}現在是第二版。
% cue 55 430.770 433.960
\hypertarget{c01:cue:55}{}我猜想第二版的書還沒有出來。
% cue 56 433.960 436.470
\hypertarget{c01:cue:56}{}所以現在可能看不到。
% cue 57 436.470 437.470
\hypertarget{c01:cue:57}{}你最多只能看到一個圖片。
% cue 58 437.470 439.570
\hypertarget{c01:cue:58}{}有沒有看到圖片。
% cue 59 439.570 441.180
\hypertarget{c01:cue:59}{}沒有啊，有。
% cue 60 441.180 442.180
\hypertarget{c01:cue:60}{}你看到圖片，還是有嗎。
% cue 61 442.180 444.270
\hypertarget{c01:cue:61}{}好，那就好。

\tc{0:07:24}
% cue 62 444.270 447.280
\hypertarget{c01:cue:62}{}這什麼回事，我並不懂每天。
% cue 63 447.280 449.530
\hypertarget{c01:cue:63}{}像我有一個人會把插開。
% cue 64 449.530 451.530
\hypertarget{c01:cue:64}{}因為我做同學會用手機的文件。
% cue 65 451.530 455.710
\hypertarget{c01:cue:65}{}所以這個時候是應該等待的。
% cue 66 455.710 457.460
\hypertarget{c01:cue:66}{}對。
% cue 67 457.460 461.140
\hypertarget{c01:cue:67}{}那我現在已經插了。
% cue 68 461.140 465.720
\hypertarget{c01:cue:68}{}那現在這個已經有聲音了嗎。
% cue 69 465.720 469.080
\hypertarget{c01:cue:69}{}喂，喂。
% cue 70 469.080 472.210
\hypertarget{c01:cue:70}{}好，那麻煩你要關小聲一點。
% cue 71 472.210 475.400
\hypertarget{c01:cue:71}{}然後你們去學一下怎麼看。
% cue 72 475.400 479.440
\hypertarget{c01:cue:72}{}太小聲了。
% cue 73 479.440 481.470
\hypertarget{c01:cue:73}{}要比剛才大的聲音好。
% cue 74 481.470 482.470
\hypertarget{c01:cue:74}{}差不多這樣就可以了。
% cue 75 482.470 483.470
\hypertarget{c01:cue:75}{}對吧。
% cue 76 483.470 484.630
\hypertarget{c01:cue:76}{}所以會不會能聽到。
% cue 77 484.630 486.630
\hypertarget{c01:cue:77}{}因為太少了，好多別人。

\section{新書發表會與延伸講次}

\tc{0:08:06}
% cue 78 486.630 488.810
\hypertarget{c01:cue:78}{}好，我問你們一件事啊。
% cue 79 488.810 492.600
\hypertarget{c01:cue:79}{}今年的 11 月 18 日期寫下來。
% cue 80 492.600 495.600
\hypertarget{c01:cue:80}{}星期三下午七八節。
% cue 81 495.600 499.540
\hypertarget{c01:cue:81}{}我希望在這個七樓。
% cue 82 499.540 500.510
\hypertarget{c01:cue:82}{}模擬法庭。
% cue 83 500.510 501.820
\hypertarget{c01:cue:83}{}一年級的小朋友。
% cue 84 501.820 502.820
\hypertarget{c01:cue:84}{}這裡沒有一年級的很好。
% cue 85 502.820 505.160
\hypertarget{c01:cue:85}{}應該沒有人不知道模擬法庭吧。
% cue 86 505.160 507.160
\hypertarget{c01:cue:86}{}這座大樓的樓上。
% cue 87 507.160 509.640
\hypertarget{c01:cue:87}{}那我想要辦一個活動就是。
% cue 88 509.640 510.640
\hypertarget{c01:cue:88}{}新書發表會。
% cue 89 510.640 512.570
\hypertarget{c01:cue:89}{}那我現在名字大體上就這樣子。
% cue 90 512.570 514.800
\hypertarget{c01:cue:90}{}聽一下你們的意見。
% cue 91 514.800 515.800
\hypertarget{c01:cue:91}{}還可以來得及修改。
% cue 92 515.800 520.370
\hypertarget{c01:cue:92}{}總標題大法官談古論今。
% cue 93 520.370 521.790
\hypertarget{c01:cue:93}{}那種範圍。
% cue 94 521.790 522.820
\hypertarget{c01:cue:94}{}談古就是什麼。
% cue 95 522.820 523.810
\hypertarget{c01:cue:95}{}夠夠熟呢。
% cue 96 523.810 526.040
\hypertarget{c01:cue:96}{}論今就是論今朝發生的事情。
% cue 97 526.040 527.350
\hypertarget{c01:cue:97}{}請問你們。
% cue 98 527.350 528.350
\hypertarget{c01:cue:98}{}謝謝。
% cue 99 528.350 530.350
\hypertarget{c01:cue:99}{}假如你們有什麼時事要問的話。
% cue 100 530.350 532.730
\hypertarget{c01:cue:100}{}所以是談古論今。
% cue 101 532.730 533.730
\hypertarget{c01:cue:101}{}那我要講什麼呢。
% cue 102 533.730 534.730
\hypertarget{c01:cue:102}{}我還在考慮中。
% cue 103 534.730 536.520
\hypertarget{c01:cue:103}{}那下面呢。
% cue 104 536.520 537.520
\hypertarget{c01:cue:104}{}副標題。
% cue 105 537.520 540.670
\hypertarget{c01:cue:105}{}釋憲學思心路二版。
% cue 106 540.670 544.120
\hypertarget{c01:cue:106}{}關於新書發表暨簽書會。
% cue 107 544.120 545.680
\hypertarget{c01:cue:107}{}有興趣嗎。
% cue 108 545.680 548.600
\hypertarget{c01:cue:108}{}我們班上可以成為我們演講的。
% cue 109 548.600 549.600
\hypertarget{c01:cue:109}{}基本觀眾嗎。

\tc{0:09:11}
% cue 110 551.750 553.290
\hypertarget{c01:cue:110}{}那現在來講一講。
% cue 111 553.290 554.320
\hypertarget{c01:cue:111}{}有什麼好處。
% cue 112 554.320 557.500
\hypertarget{c01:cue:112}{}如果我們班上有足夠的人買書。
% cue 113 557.500 559.560
\hypertarget{c01:cue:113}{}學校有足夠的人買書。
% cue 114 559.560 560.780
\hypertarget{c01:cue:114}{}我想做一件事。
% cue 115 560.780 564.090
\hypertarget{c01:cue:115}{}但是我要徵求一下你們的意願。
% cue 116 564.090 566.600
\hypertarget{c01:cue:116}{}在臺大這邊我也是這樣做的。
% cue 117 566.600 569.680
\hypertarget{c01:cue:117}{}如果班上有足夠的人買書。
% cue 118 569.680 571.960
\hypertarget{c01:cue:118}{}我就希望今天我們是上到五點半。
% cue 119 571.960 573.960
\hypertarget{c01:cue:119}{}也許我們可以五點二十結束。
% cue 120 573.960 575.960
\hypertarget{c01:cue:120}{}五點半也沒有問題。
% cue 121 575.960 577.370
\hypertarget{c01:cue:121}{}後面也沒有課。
% cue 122 577.370 580.590
\hypertarget{c01:cue:122}{}所以我想加半小時的客服。
% cue 123 580.590 582.000
\hypertarget{c01:cue:123}{}客服吧。
% cue 124 582.000 583.000
\hypertarget{c01:cue:124}{}我們都叫客服吧。
% cue 125 583.000 584.380
\hypertarget{c01:cue:125}{}客服有個標題。
% cue 126 584.380 587.400
\hypertarget{c01:cue:126}{}如何讀懂大法官解釋。
% cue 127 587.400 591.940
\hypertarget{c01:cue:127}{}然後老師就以釋憲學思心路這三冊裡面的案子。
% cue 128 591.940 596.420
\hypertarget{c01:cue:128}{}那有可能會連結到以前的三次的之前。
% cue 129 596.420 597.420
\hypertarget{c01:cue:129}{}我還沒有來的。
% cue 130 597.420 600.420
\hypertarget{c01:cue:130}{}以及三次以後我走掉以後的憲法。
% cue 131 600.420 602.770
\hypertarget{c01:cue:131}{}來的這些案子我把你串起來。
% cue 132 602.770 605.110
\hypertarget{c01:cue:132}{}然後我來跟你講。
% cue 133 605.110 606.780
\hypertarget{c01:cue:133}{}什麼讀懂大法官解釋。
% cue 134 606.780 608.740
\hypertarget{c01:cue:134}{}我會想一個方法。
% cue 135 608.740 610.960
\hypertarget{c01:cue:135}{}這等於算是這集客服延伸。
% cue 136 610.960 612.470
\hypertarget{c01:cue:136}{}所以它叫客服。
% cue 137 612.470 614.600
\hypertarget{c01:cue:137}{}真的是名副其實的客服。
% cue 138 614.600 617.650
\hypertarget{c01:cue:138}{}任務就是幫助你們讀懂大法官解釋。

\section{讀懂大法官解釋的層次}

\tc{0:10:17}
% cue 139 617.650 621.380
\hypertarget{c01:cue:139}{}你們可能會覺得老師你太小看我們呢。
% cue 140 621.380 623.700
\hypertarget{c01:cue:140}{}法官解釋有多少次嗎。
% cue 141 623.700 624.700
\hypertarget{c01:cue:141}{}看不懂嗎。
% cue 142 624.700 627.740
\hypertarget{c01:cue:142}{}那你先不要急著。
% cue 143 627.740 629.710
\hypertarget{c01:cue:143}{}馬上還致富地回答說。
% cue 144 629.710 630.710
\hypertarget{c01:cue:144}{}我硬幹得很。
% cue 145 630.710 632.830
\hypertarget{c01:cue:145}{}等我開始講了。
% cue 146 632.830 633.960
\hypertarget{c01:cue:146}{}你就會知道懂。
% cue 147 633.960 636.500
\hypertarget{c01:cue:147}{}有很多種層次的懂。
% cue 148 636.500 638.450
\hypertarget{c01:cue:148}{}層次。
% cue 149 638.450 642.520
\hypertarget{c01:cue:149}{}認識文字是最起碼的層次。
% cue 150 642.520 643.770
\hypertarget{c01:cue:150}{}他講的那個意思。
% cue 151 643.770 644.770
\hypertarget{c01:cue:151}{}那個意思都看不懂。
% cue 152 644.770 645.770
\hypertarget{c01:cue:152}{}他層次沒有乾淨。
% cue 153 645.770 648.150
\hypertarget{c01:cue:153}{}我再問你第二次。
% cue 154 648.150 649.370
\hypertarget{c01:cue:154}{}他為什麼要這樣講。
% cue 155 649.370 651.910
\hypertarget{c01:cue:155}{}你們回答好了。
% cue 156 651.910 653.930
\hypertarget{c01:cue:156}{}第三次。
% cue 157 653.930 655.220
\hypertarget{c01:cue:157}{}這一個解釋裡面。
% cue 158 655.220 657.410
\hypertarget{c01:cue:158}{}有協同有不同\footnote{指協同意見書與不同意見書。憲法訴訟法第 35 條規定，大法官贊成裁判之主文，而對其理由有補充或不同意見者，得提出協同意見書；對於裁判之主文曾於評議時表示部分或全部不同意見者，得提出部分或全部之不同意見書。\cite{csfa}。}。
% cue 159 657.410 658.810
\hypertarget{c01:cue:159}{}那請問你。
% cue 160 658.810 660.040
\hypertarget{c01:cue:160}{}這個協同跟不同。
% cue 161 660.040 661.040
\hypertarget{c01:cue:161}{}跟這個多數意見。
% cue 162 661.040 663.280
\hypertarget{c01:cue:162}{}差別在哪裡。
% cue 163 663.280 664.930
\hypertarget{c01:cue:163}{}問你人的講。
% cue 164 664.930 666.410
\hypertarget{c01:cue:164}{}那第四。
% cue 165 666.410 667.620
\hypertarget{c01:cue:165}{}那這句比較難回答。
% cue 166 667.620 669.900
\hypertarget{c01:cue:166}{}這下文第四句就更難回答。
% cue 167 669.900 670.900
\hypertarget{c01:cue:167}{}那這兩個。
% cue 168 670.900 671.900
\hypertarget{c01:cue:168}{}不管是協同或不同。
% cue 169 671.900 672.900
\hypertarget{c01:cue:169}{}他們為什麼要協同。
% cue 170 672.900 675.210
\hypertarget{c01:cue:170}{}他們為什麼要在這時候。
% cue 171 675.210 676.210
\hypertarget{c01:cue:171}{}表達協同跟不同。
% cue 172 676.210 677.810
\hypertarget{c01:cue:172}{}他們在爭取。
% cue 173 677.810 679.810
\hypertarget{c01:cue:173}{}或者他們堅持的事實。
% cue 174 679.810 682.540
\hypertarget{c01:cue:174}{}這是表現到這裡。
% cue 175 682.540 684.570
\hypertarget{c01:cue:175}{}如果講到這四個層次。
% cue 176 684.570 685.530
\hypertarget{c01:cue:176}{}你都能懂。
% cue 177 685.530 686.600
\hypertarget{c01:cue:177}{}那大概解釋很真。
% cue 178 686.600 688.540
\hypertarget{c01:cue:178}{}你大概算看得懂。

\tc{0:11:28}
% cue 179 688.540 690.370
\hypertarget{c01:cue:179}{}剩下來的就比較難了。
% cue 180 690.370 693.690
\hypertarget{c01:cue:180}{}那就是老師會在釋憲學思心路裡面。
% cue 181 693.690 694.690
\hypertarget{c01:cue:181}{}講幾則。
% cue 182 694.690 695.850
\hypertarget{c01:cue:182}{}你們知道禪宗。
% cue 183 695.850 697.100
\hypertarget{c01:cue:183}{}公案嗎。
% cue 184 697.100 699.420
\hypertarget{c01:cue:184}{}公案的還是聽懂的。
% cue 185 699.420 699.520
\hypertarget{c01:cue:185}{}據說的。
% cue 186 699.420 700.480
\hypertarget{c01:cue:186}{}有聽懂。
% cue 187 700.480 701.480
\hypertarget{c01:cue:187}{}公案。
% cue 188 701.480 701.580
\hypertarget{c01:cue:188}{}沒有啊。
% cue 189 701.480 703.990
\hypertarget{c01:cue:189}{}應該在五十個人左右。
% cue 190 703.990 704.990
\hypertarget{c01:cue:190}{}有沒有人聽懂公案。
% cue 191 704.990 706.050
\hypertarget{c01:cue:191}{}我們現在班上多少。

\tc{0:11:46}
\speaker{學生}
% cue 192 706.050 707.050
\hypertarget{c01:cue:192}{}四十一個人。

\tc{0:11:47}
\speaker{老師}
% cue 193 707.050 708.050
\hypertarget{c01:cue:193}{}我看。
% cue 194 708.050 709.080
\hypertarget{c01:cue:194}{}我看應該都是四十一個人。
% cue 195 709.080 711.560
\hypertarget{c01:cue:195}{}有沒有人聽不懂公案二字。
% cue 196 711.560 712.560
\hypertarget{c01:cue:196}{}公案聽不懂啊。
% cue 197 712.560 714.190
\hypertarget{c01:cue:197}{}來來來來。
% cue 198 714.190 715.220
\hypertarget{c01:cue:198}{}你們查一下。
% cue 199 715.220 716.210
\hypertarget{c01:cue:199}{}你們手上些太。
% cue 200 716.210 718.280
\hypertarget{c01:cue:200}{}你們不是有不夠嗎。
% cue 201 718.280 719.750
\hypertarget{c01:cue:201}{}查一下查一下。
% cue 202 719.750 733.100
\hypertarget{c01:cue:202}{}中。

\tc{0:12:13}
% cue 203 733.100 734.050
\hypertarget{c01:cue:203}{}公案。
% cue 204 734.050 735.960
\hypertarget{c01:cue:204}{}這個是禪宗。
% cue 205 735.960 736.930
\hypertarget{c01:cue:205}{}在幫助大家。
% cue 206 736.930 740.680
\hypertarget{c01:cue:206}{}理解佛法重要的附近。
% cue 207 740.680 743.720
\hypertarget{c01:cue:207}{}他們對佛法的理解。
% cue 208 743.720 746.160
\hypertarget{c01:cue:208}{}不限於經典。
% cue 209 746.160 747.450
\hypertarget{c01:cue:209}{}比如說啊。
% cue 210 747.450 748.450
\hypertarget{c01:cue:210}{}你誦金剛經\footnote{《金剛般若波羅蜜經》，後秦鳩摩羅什譯，一卷，收於大正藏第 8 冊，經號 235。\url{https://cbdata.dila.edu.tw/stable/works?work=T0235}。}啊。
% cue 211 748.450 749.450
\hypertarget{c01:cue:211}{}你誦什麼普門品\footnote{〈觀世音菩薩普門品〉，鳩摩羅什譯《妙法蓮華經》卷七所收的第二十五品。\url{https://cbdata.dila.edu.tw/stable/works?work=T0262}。}啊。
% cue 212 749.450 750.450
\hypertarget{c01:cue:212}{}你誦什麼。
% cue 213 750.450 750.550
\hypertarget{c01:cue:213}{}普門品啊。
% cue 214 750.450 751.450
\hypertarget{c01:cue:214}{}不是只有這個。
% cue 215 751.450 753.770
\hypertarget{c01:cue:215}{}他們在日常生活當中。
% cue 216 753.770 756.990
\hypertarget{c01:cue:216}{}尤其是有經典的禪師。
% cue 217 756.990 758.370
\hypertarget{c01:cue:217}{}像老師。
% cue 218 758.370 760.460
\hypertarget{c01:cue:218}{}帶離他的弟子部門。
% cue 219 760.460 762.810
\hypertarget{c01:cue:219}{}通常是藉一些事件。
% cue 220 762.810 764.290
\hypertarget{c01:cue:220}{}來敲打學生。
% cue 221 764.290 767.470
\hypertarget{c01:cue:221}{}那就是這就是他們的禪宗。
% cue 222 767.470 768.470
\hypertarget{c01:cue:222}{}禪宗有很多公案哦。
% cue 223 768.470 770.550
\hypertarget{c01:cue:223}{}數量非常多。
% cue 224 770.550 772.160
\hypertarget{c01:cue:224}{}那有點像這個。
% cue 225 772.160 773.410
\hypertarget{c01:cue:225}{}那我會。
% cue 226 773.410 775.380
\hypertarget{c01:cue:226}{}在這個釋憲學思心路裡面。
% cue 227 775.380 776.380
\hypertarget{c01:cue:227}{}找到幾則。
% cue 228 776.380 778.570
\hypertarget{c01:cue:228}{}明顯的公案。
% cue 229 778.570 779.570
\hypertarget{c01:cue:229}{}有些不明顯。
% cue 230 779.570 780.540
\hypertarget{c01:cue:230}{}那就算了。
% cue 231 780.540 781.860
\hypertarget{c01:cue:231}{}明顯的公案。
% cue 232 781.860 783.020
\hypertarget{c01:cue:232}{}那我來問你。

\tc{0:13:03}
% cue 233 783.020 784.050
\hypertarget{c01:cue:233}{}我隨便舉一個例子。
% cue 234 784.050 786.040
\hypertarget{c01:cue:234}{}我說這一段文字。
% cue 235 786.040 786.970
\hypertarget{c01:cue:235}{}十頁左右。
% cue 236 786.970 788.140
\hypertarget{c01:cue:236}{}看完了嗎。
% cue 237 788.140 789.570
\hypertarget{c01:cue:237}{}來。
% cue 238 789.570 791.150
\hypertarget{c01:cue:238}{}現在請回答我。
% cue 239 791.150 792.090
\hypertarget{c01:cue:239}{}這裡面。
% cue 240 792.090 793.370
\hypertarget{c01:cue:240}{}哪一個是白臉。
% cue 241 793.370 794.370
\hypertarget{c01:cue:241}{}哪一個是黑臉。
% cue 242 794.370 799.470
\hypertarget{c01:cue:242}{}那會這樣問你。
% cue 243 799.470 800.470
\hypertarget{c01:cue:243}{}問你這樣問會提出。
% cue 244 800.470 801.470
\hypertarget{c01:cue:244}{}這叫看懂。
% cue 245 801.470 802.470
\hypertarget{c01:cue:245}{}看不懂大。
% cue 246 802.470 806.000
\hypertarget{c01:cue:246}{}那白臉的黑臉有什麼差別。
% cue 247 806.000 807.700
\hypertarget{c01:cue:247}{}那個黑臉。
% cue 248 807.700 808.700
\hypertarget{c01:cue:248}{}為什麼敢跟老師講話。
% cue 249 808.700 809.700
\hypertarget{c01:cue:249}{}講那麼久。
% cue 250 809.700 813.450
\hypertarget{c01:cue:250}{}之類。
% cue 251 813.450 816.150
\hypertarget{c01:cue:251}{}然後為什麼這個解釋。
% cue 252 816.150 820.410
\hypertarget{c01:cue:252}{}跟兩年以後發生的事情。
% cue 253 820.410 822.600
\hypertarget{c01:cue:253}{}要怎麼溝擊在一起。
% cue 254 822.600 824.940
\hypertarget{c01:cue:254}{}這樣子慢慢你就。
% cue 255 824.940 825.940
\hypertarget{c01:cue:255}{}哦。
% cue 256 825.940 827.870
\hypertarget{c01:cue:256}{}原來這套文章的解釋。
% cue 257 827.870 830.560
\hypertarget{c01:cue:257}{}這麼複雜的故事在裡面。
% cue 258 830.560 832.180
\hypertarget{c01:cue:258}{}表面上看不見。
% cue 259 832.180 835.720
\hypertarget{c01:cue:259}{}所以就要花一點時間去做。

\tc{0:13:55}
% cue 260 835.720 835.820
\hypertarget{c01:cue:260}{}好。
% cue 261 835.720 837.920
\hypertarget{c01:cue:261}{}回到剛才講的這一點。
% cue 262 837.920 839.640
\hypertarget{c01:cue:262}{}第一。
% cue 263 839.640 840.640
\hypertarget{c01:cue:263}{}11 月 18 號星期三。
% cue 264 840.640 841.760
\hypertarget{c01:cue:264}{}七八節。
% cue 265 841.760 844.970
\hypertarget{c01:cue:265}{}它保留起來。
% cue 266 844.970 845.970
\hypertarget{c01:cue:266}{}第二。
% cue 267 845.970 846.970
\hypertarget{c01:cue:267}{}老師的這個活動。
% cue 268 846.970 848.350
\hypertarget{c01:cue:268}{}希望你們參加。
% cue 269 848.350 849.350
\hypertarget{c01:cue:269}{}第三。
% cue 270 849.350 850.730
\hypertarget{c01:cue:270}{}你們現在跟我講。
% cue 271 850.730 852.470
\hypertarget{c01:cue:271}{}意向的表達。
% cue 272 852.470 853.470
\hypertarget{c01:cue:272}{}假如這個活動。
% cue 273 853.470 854.400
\hypertarget{c01:cue:273}{}翻得不錯。
% cue 274 854.400 855.530
\hypertarget{c01:cue:274}{}覺得熱戀。
% cue 275 855.530 856.530
\hypertarget{c01:cue:275}{}買書要買的文化。
% cue 276 856.530 858.720
\hypertarget{c01:cue:276}{}那我們就在這個課的後面。
% cue 277 858.720 859.840
\hypertarget{c01:cue:277}{}加上課題。
% cue 278 859.840 861.930
\hypertarget{c01:cue:278}{}那個課題的名稱。
% cue 279 861.930 863.890
\hypertarget{c01:cue:279}{}就是如何讀懂大法官解釋。
% cue 280 863.890 866.110
\hypertarget{c01:cue:280}{}然後老師由細胞的。
% cue 281 866.110 868.620
\hypertarget{c01:cue:281}{}會慢慢的挑案件。
% cue 282 868.620 870.740
\hypertarget{c01:cue:282}{}帶你們進去那個複雜的三冊。
% cue 283 870.740 873.000
\hypertarget{c01:cue:283}{}那三冊都是老師的意見書。
% cue 284 873.000 881.560
\hypertarget{c01:cue:284}{}方言的嘔心瀝血。
% cue 285 881.560 885.340
\hypertarget{c01:cue:285}{}我已經想不出來。
% cue 286 885.340 887.430
\hypertarget{c01:cue:286}{}還有什麼方法可以行動。

\tc{0:14:47}
% cue 287 887.430 891.930
\hypertarget{c01:cue:287}{}但是老師最精華的歲月。
% cue 288 891.930 894.270
\hypertarget{c01:cue:288}{}從 55 歲做到 63 歲。
% cue 289 894.270 896.970
\hypertarget{c01:cue:289}{}人生還有比這個更精華的歲月嗎。
% cue 290 896.970 898.710
\hypertarget{c01:cue:290}{}之前不懂事。
% cue 291 898.710 899.860
\hypertarget{c01:cue:291}{}之後老血。
% cue 292 899.860 901.110
\hypertarget{c01:cue:292}{}再更老一點。
% cue 293 901.110 902.400
\hypertarget{c01:cue:293}{}對整直務人。
% cue 294 902.400 905.380
\hypertarget{c01:cue:294}{}人生歲月很短。
% cue 295 905.380 907.790
\hypertarget{c01:cue:295}{}但我們的歲月跟別人一樣。
% cue 296 907.790 908.790
\hypertarget{c01:cue:296}{}別人快樂的吃香喝辣。

\tc{0:15:08}
% cue 297 908.790 911.940
\hypertarget{c01:cue:297}{}大法官的解釋。
% cue 298 911.940 912.940
\hypertarget{c01:cue:298}{}你不寫意見書。
% cue 299 912.940 915.120
\hypertarget{c01:cue:299}{}會有壓力。
% cue 300 915.120 916.120
\hypertarget{c01:cue:300}{}沒有人逼你寫。
% cue 301 916.120 917.560
\hypertarget{c01:cue:301}{}你要寫意見書。
% cue 302 917.560 919.620
\hypertarget{c01:cue:302}{}那你就給兩倍工作。
% cue 303 919.620 921.450
\hypertarget{c01:cue:303}{}參加大法官會議\footnote{1948 年施行的司法院組織法第 3 條規定，司法院設大法官會議，以大法官十七人組織之，行使解釋憲法並統一解釋法律命令之職權。審理程序先後依 1958 年制定的司法院大法官會議法與 1993 年更名的司法院大法官審理案件法；2019 年 1 月 4 日修正公布並更名的憲法訴訟法自 2022 年 1 月 4 日施行，改由大法官組成憲法法庭審理案件。}。
% cue 304 921.450 923.450
\hypertarget{c01:cue:304}{}本身出一個多數意見。
% cue 305 923.450 925.670
\hypertarget{c01:cue:305}{}這多數意見的過程當中。
% cue 306 925.670 927.670
\hypertarget{c01:cue:306}{}你可能要竭力的反對。
% cue 307 927.670 928.670
\hypertarget{c01:cue:307}{}竭力的去爭取。
% cue 308 928.670 929.670
\hypertarget{c01:cue:308}{}那爭取不到。
% cue 309 929.670 931.670
\hypertarget{c01:cue:309}{}那時候你自己要再來寫意見書。
% cue 310 931.670 932.760
\hypertarget{c01:cue:310}{}意見書你敢寫。
% cue 311 932.760 934.880
\hypertarget{c01:cue:311}{}而且你敢會記成冊。
% cue 312 934.880 936.980
\hypertarget{c01:cue:312}{}那就是可受公評。
% cue 313 936.980 939.390
\hypertarget{c01:cue:313}{}那這是增加非常多的壓力。
% cue 314 939.390 940.720
\hypertarget{c01:cue:314}{}你們會不會有些人。
% cue 315 940.720 941.660
\hypertarget{c01:cue:315}{}你寫完以後。
% cue 316 941.660 945.590
\hypertarget{c01:cue:316}{}你寫完以後。
% cue 317 945.590 946.590
\hypertarget{c01:cue:317}{}就說。
% cue 318 946.590 947.530
\hypertarget{c01:cue:318}{}無聲無事。
% cue 319 947.530 949.590
\hypertarget{c01:cue:319}{}不可遇人也。

\tc{0:15:49}
% cue 320 949.590 950.590
\hypertarget{c01:cue:320}{}這是歐陽修\footnote{歐陽修為北宋文學家，宋仁宗嘉祐六年任參知政事。「平生所為，未嘗有不可對人言者」出自《宋史．司馬光傳》所載司馬光自述：「吾無過人者，但平生所爲，未嘗有不可對人言者耳。」司馬光於宋哲宗時拜尚書左僕射兼門下侍郎。\href{https://zh.wikisource.org/wiki/\%E5\%AE\%8B\%E5\%8F\%B2/\%E5\%8D\%B7336}{https://zh.wikisource.org/wiki/宋史/卷336}。}的靈驗。
% cue 321 950.590 951.650
\hypertarget{c01:cue:321}{}聽得懂我現在。
% cue 322 951.650 952.650
\hypertarget{c01:cue:322}{}聽到他這幾個字嗎。
% cue 323 952.650 956.080
\hypertarget{c01:cue:323}{}這一生。
% cue 324 956.080 958.230
\hypertarget{c01:cue:324}{}沒有一件事情。
% cue 325 958.230 962.430
\hypertarget{c01:cue:325}{}可以對人說。
% cue 326 962.430 967.840
\hypertarget{c01:cue:326}{}就是這個人的無論起。
% cue 327 967.840 969.190
\hypertarget{c01:cue:327}{}歐陽修都有聽過嗎。
% cue 328 969.190 970.170
\hypertarget{c01:cue:328}{}我不知道。
% cue 329 970.170 971.510
\hypertarget{c01:cue:329}{}因為現在。
% cue 330 971.510 973.670
\hypertarget{c01:cue:330}{}108 課綱改立法古文。
% cue 331 973.670 974.860
\hypertarget{c01:cue:331}{}所有的古文都不見了。
% cue 332 974.860 977.010
\hypertarget{c01:cue:332}{}歐陽修寫了很多有名的古文。
% cue 333 977.010 978.390
\hypertarget{c01:cue:333}{}簡單的講。
% cue 334 978.390 979.710
\hypertarget{c01:cue:334}{}他就說什麼。
% cue 335 979.710 980.710
\hypertarget{c01:cue:335}{}一個相當於宰相。
% cue 336 980.710 981.710
\hypertarget{c01:cue:336}{}這樣的職務。
% cue 337 981.710 983.630
\hypertarget{c01:cue:337}{}他也算是。
% cue 338 983.630 984.630
\hypertarget{c01:cue:338}{}人體法俗世的。
% cue 339 984.630 985.790
\hypertarget{c01:cue:339}{}重要人物。
% cue 340 985.790 988.320
\hypertarget{c01:cue:340}{}人家跟他講話。
% cue 341 988.320 989.320
\hypertarget{c01:cue:341}{}他就一句話。
% cue 342 989.320 990.320
\hypertarget{c01:cue:342}{}解決他自己。
% cue 343 990.320 991.600
\hypertarget{c01:cue:343}{}他說我這一生。
% cue 344 991.600 992.600
\hypertarget{c01:cue:344}{}沒有做任何對不起。
% cue 345 992.600 994.080
\hypertarget{c01:cue:345}{}的事情。
% cue 346 994.080 995.080
\hypertarget{c01:cue:346}{}所以你問我。
% cue 347 995.080 996.110
\hypertarget{c01:cue:347}{}我沒有事情。
% cue 348 996.110 999.510
\hypertarget{c01:cue:348}{}改給人說。
% cue 349 999.510 1000.510
\hypertarget{c01:cue:349}{}那這是需要。
% cue 350 1000.510 1001.990
\hypertarget{c01:cue:350}{}有效方向自行。
% cue 351 1001.990 1003.210
\hypertarget{c01:cue:351}{}定解。

\tc{0:16:43}
% cue 352 1003.210 1004.210
\hypertarget{c01:cue:352}{}一個人要始終如一。
% cue 353 1004.210 1005.200
\hypertarget{c01:cue:353}{}是不是很不容易。
% cue 354 1005.200 1007.310
\hypertarget{c01:cue:354}{}我這樣講。
% cue 355 1007.310 1008.250
\hypertarget{c01:cue:355}{}有過分嗎。
% cue 356 1008.250 1009.280
\hypertarget{c01:cue:356}{}你們在座的人。
% cue 357 1009.280 1010.370
\hypertarget{c01:cue:357}{}是不是想說。
% cue 358 1010.370 1011.370
\hypertarget{c01:cue:358}{}很難講。
% cue 359 1011.370 1014.870
\hypertarget{c01:cue:359}{}我這個想法。
% cue 360 1014.870 1015.870
\hypertarget{c01:cue:360}{}應該從 10 歲。
% cue 361 1015.870 1016.870
\hypertarget{c01:cue:361}{}開始卻有了吧。
% cue 362 1016.870 1017.870
\hypertarget{c01:cue:362}{}10 歲也結束不懂事。
% cue 363 1017.870 1018.900
\hypertarget{c01:cue:363}{}10 歲以後說。
% cue 364 1018.900 1019.900
\hypertarget{c01:cue:364}{}我旁邊的人。
% cue 365 1019.900 1020.900
\hypertarget{c01:cue:365}{}每次考試的考度好。
% cue 366 1020.900 1021.820
\hypertarget{c01:cue:366}{}你把你考試的考度好。
% cue 367 1021.820 1023.470
\hypertarget{c01:cue:367}{}你把你考慮。
% cue 368 1023.470 1025.230
\hypertarget{c01:cue:368}{}發奮圖強。
% cue 369 1025.230 1026.170
\hypertarget{c01:cue:369}{}兩天以後。
% cue 370 1026.170 1030.780
\hypertarget{c01:cue:370}{}現在你們現在還有很多夢想。
% cue 371 1030.780 1031.950
\hypertarget{c01:cue:371}{}還在發奮圖強。
% cue 372 1031.950 1033.940
\hypertarget{c01:cue:372}{}不然你不會做到這個班上。
% cue 373 1033.940 1035.000
\hypertarget{c01:cue:373}{}我相信你們聽到很多人。
% cue 374 1035.000 1036.000
\hypertarget{c01:cue:374}{}說。
% cue 375 1036.000 1037.000
\hypertarget{c01:cue:375}{}老師這課很難。
% cue 376 1037.000 1038.000
\hypertarget{c01:cue:376}{}什麼謠言一大堆。
% cue 377 1038.000 1040.000
\hypertarget{c01:cue:377}{}你們還有勇氣都在這裡。
% cue 378 1040.000 1041.000
\hypertarget{c01:cue:378}{}看你是不是有夢想。

\tc{0:17:21}
% cue 379 1041.000 1045.410
\hypertarget{c01:cue:379}{}我做一個老師存在的任務之一。
% cue 380 1045.410 1047.570
\hypertarget{c01:cue:380}{}就是讓你們的夢想成真。
% cue 381 1047.570 1049.500
\hypertarget{c01:cue:381}{}不是給你交流水。
% cue 382 1049.500 1050.760
\hypertarget{c01:cue:382}{}但是夢想要成真。
% cue 383 1050.760 1055.460
\hypertarget{c01:cue:383}{}需要帶你你步一步的走到正確的道路上用力。
% cue 384 1055.460 1056.840
\hypertarget{c01:cue:384}{}用力用的是蠻力。
% cue 385 1056.840 1057.870
\hypertarget{c01:cue:385}{}用錯力。
% cue 386 1057.870 1059.450
\hypertarget{c01:cue:386}{}全面不好。
% cue 387 1059.450 1060.860
\hypertarget{c01:cue:387}{}什麼事情。
% cue 388 1060.860 1062.490
\hypertarget{c01:cue:388}{}那是用錯了力。
% cue 389 1062.490 1063.610
\hypertarget{c01:cue:389}{}用對人力。
% cue 390 1063.610 1065.610
\hypertarget{c01:cue:390}{}而又肯用力。
% cue 391 1065.610 1067.650
\hypertarget{c01:cue:391}{}然後持之以恆的用力。
% cue 392 1067.650 1068.790
\hypertarget{c01:cue:392}{}那你就會成功。
% cue 393 1068.790 1071.110
\hypertarget{c01:cue:393}{}動物的情況差不多就是這樣對不對。
% cue 394 1071.110 1072.370
\hypertarget{c01:cue:394}{}比上不足。
% cue 395 1072.370 1073.370
\hypertarget{c01:cue:395}{}比下有餘。
% cue 396 1073.370 1074.530
\hypertarget{c01:cue:396}{}往上走。
% cue 397 1074.530 1075.530
\hypertarget{c01:cue:397}{}吃力。
% cue 398 1075.530 1076.510
\hypertarget{c01:cue:398}{}但是上得去。
% cue 399 1076.510 1078.840
\hypertarget{c01:cue:399}{}要自由落體。
% cue 400 1078.840 1080.610
\hypertarget{c01:cue:400}{}那也很。
% cue 401 1080.610 1081.600
\hypertarget{c01:cue:401}{}不要兩個月就落得一趟。
% cue 402 1081.600 1082.830
\hypertarget{c01:cue:402}{}就不面部全面。
% cue 403 1082.830 1083.740
\hypertarget{c01:cue:403}{}自己不認識。
% cue 404 1083.740 1086.420
\hypertarget{c01:cue:404}{}我這樣想。
% cue 405 1086.420 1087.990
\hypertarget{c01:cue:405}{}不對嗎。
% cue 406 1087.990 1089.990
\hypertarget{c01:cue:406}{}有沒有有偏僅真實嗎。
% cue 407 1089.990 1091.430
\hypertarget{c01:cue:407}{}沒有啊。
% cue 408 1091.430 1092.750
\hypertarget{c01:cue:408}{}好啦。

\tc{0:18:12}
% cue 409 1092.750 1093.750
\hypertarget{c01:cue:409}{}現在意向調查一下。
% cue 410 1093.750 1095.160
\hypertarget{c01:cue:410}{}你們認為。
% cue 411 1095.160 1096.130
\hypertarget{c01:cue:411}{}如果哪些人夠多。
% cue 412 1096.130 1098.130
\hypertarget{c01:cue:412}{}我們就可以進行這一個群體。
% cue 413 1098.130 1099.220
\hypertarget{c01:cue:413}{}然後老師在這個課之後。
% cue 414 1099.220 1101.180
\hypertarget{c01:cue:414}{}就有一個課譜。
% cue 415 1101.180 1102.340
\hypertarget{c01:cue:415}{}課譜就叫做。
% cue 416 1102.340 1104.370
\hypertarget{c01:cue:416}{}如何不懂大法官解釋。
% cue 417 1104.370 1105.370
\hypertarget{c01:cue:417}{}然後老師去計劃。
% cue 418 1105.370 1106.370
\hypertarget{c01:cue:418}{}又信我的。
% cue 419 1106.370 1107.370
\hypertarget{c01:cue:419}{}南打這三次的裡面。
% cue 420 1107.370 1108.370
\hypertarget{c01:cue:420}{}比較粗暗的案子。
% cue 421 1108.370 1109.280
\hypertarget{c01:cue:421}{}講一講。
% cue 422 1109.280 1111.070
\hypertarget{c01:cue:422}{}對於這種想法。
% cue 423 1111.070 1112.160
\hypertarget{c01:cue:423}{}表示認同。
% cue 424 1112.160 1113.190
\hypertarget{c01:cue:424}{}有興趣。
% cue 425 1113.190 1114.190
\hypertarget{c01:cue:425}{}願意參加的。
% cue 426 1114.190 1116.120
\hypertarget{c01:cue:426}{}舉手我看看。
% cue 427 1116.120 1116.220
\hypertarget{c01:cue:427}{}好 謝謝。
% cue 428 1116.120 1118.120
\hypertarget{c01:cue:428}{}可以到處師父一帶。

\section{簡歷、影片清單與心得分享}

\tc{0:18:38}
% cue 429 1118.120 1120.210
\hypertarget{c01:cue:429}{}好 現在請回到講剛來。
% cue 430 1120.210 1122.330
\hypertarget{c01:cue:430}{}這講到裡面有幾件事情。
% cue 431 1122.330 1123.500
\hypertarget{c01:cue:431}{}第一件事情。
% cue 432 1123.500 1125.460
\hypertarget{c01:cue:432}{}你們找到老師的簡歷\footnote{課綱「參考書資料暨網路資源」列有「湯德宗教授簡歷」的 Google Drive 連結，第 1 週分週進度亦列此項。}了。
% cue 433 1125.460 1127.090
\hypertarget{c01:cue:433}{}好 今天特別請助理。
% cue 434 1127.090 1128.090
\hypertarget{c01:cue:434}{}是要檢查這個連結。
% cue 435 1128.090 1130.090
\hypertarget{c01:cue:435}{}能連結得到嗎。
% cue 436 1130.090 1131.180
\hypertarget{c01:cue:436}{}有人連結到。
% cue 437 1131.180 1132.180
\hypertarget{c01:cue:437}{}舉手我看一下。
% cue 438 1132.180 1133.780
\hypertarget{c01:cue:438}{}好 謝謝。
% cue 439 1133.780 1134.780
\hypertarget{c01:cue:439}{}老師是有的喔。
% cue 440 1134.780 1135.780
\hypertarget{c01:cue:440}{}可以連結的。
% cue 441 1135.780 1136.780
\hypertarget{c01:cue:441}{}等一下請問。
% cue 442 1136.780 1137.780
\hypertarget{c01:cue:442}{}請問。
% cue 443 1137.780 1138.770
\hypertarget{c01:cue:443}{}這是剛老師的簡歷。
% cue 444 1138.770 1139.920
\hypertarget{c01:cue:444}{}吾愛吾師\footnote{語本亞里斯多德《尼各馬科倫理學》第 1 卷第 6 章（1096a），他在批評柏拉圖學派的理型論之前說，友人與真理兩者皆可親，而義務要求以真理為先；H. Rackham 英譯作「Both are dear to us, yet 'tis our duty to prefer the truth.」\url{https://www.perseus.tufts.edu/hopper/text?doc=Perseus:text:1999.01.0054:bekker\%20page=1096a}。}。
% cue 445 1139.920 1141.180
\hypertarget{c01:cue:445}{}更愛真理。
% cue 446 1141.180 1142.510
\hypertarget{c01:cue:446}{}但是你們要先認識。
% cue 447 1142.510 1144.580
\hypertarget{c01:cue:447}{}老師長得什麼樣子。
% cue 448 1144.580 1146.570
\hypertarget{c01:cue:448}{}等一下開放問題給你們提問。

\tc{0:19:06}
% cue 449 1146.570 1147.630
\hypertarget{c01:cue:449}{}第二個。
% cue 450 1147.630 1149.630
\hypertarget{c01:cue:450}{}你們有看過老師的 YouTube\footnote{課綱「參考書資料暨網路資源」列有「YT 頻道影片清單」的 Google Drive 連結。分週表第 1 週列「湯德宗教授 YT 頻道影片清單」，第 5 週（10 月 12 日）與第 15 週（12 月 21 日）的末行標「YT」。} 的舉手。
% cue 451 1149.630 1152.290
\hypertarget{c01:cue:451}{}好 老師的 YouTube 到現在為止。
% cue 452 1152.290 1154.290
\hypertarget{c01:cue:452}{}已經超過三百萬總觀看人數。
% cue 453 1154.290 1155.290
\hypertarget{c01:cue:453}{}三百多萬。
% cue 454 1155.290 1157.830
\hypertarget{c01:cue:454}{}那這在我們國內。
% cue 455 1157.830 1159.990
\hypertarget{c01:cue:455}{}我相信是很難。
% cue 456 1159.990 1161.150
\hypertarget{c01:cue:456}{}有人可以給的。
% cue 457 1161.150 1162.530
\hypertarget{c01:cue:457}{}而且這是不收費的。
% cue 458 1162.530 1164.750
\hypertarget{c01:cue:458}{}那老師每一個學期。
% cue 459 1164.750 1167.780
\hypertarget{c01:cue:459}{}都會給大家一個最新的影片清單。
% cue 460 1167.780 1168.770
\hypertarget{c01:cue:460}{}為什麼。
% cue 461 1168.770 1170.930
\hypertarget{c01:cue:461}{}因為你們上憲法課。
% cue 462 1170.930 1172.150
\hypertarget{c01:cue:462}{}有非常大的可能。
% cue 463 1172.150 1174.440
\hypertarget{c01:cue:463}{}會在那裡面撈到你要看的東西。
% cue 464 1174.440 1176.790
\hypertarget{c01:cue:464}{}不含我們這學期上課。
% cue 465 1176.790 1177.780
\hypertarget{c01:cue:465}{}建議有建議。
% cue 466 1177.780 1179.800
\hypertarget{c01:cue:466}{}所以你可以從檢查來看。
% cue 467 1179.800 1182.340
\hypertarget{c01:cue:467}{}每一次我就跟你講。
% cue 468 1182.340 1186.050
\hypertarget{c01:cue:468}{}你要出行坐。
% cue 469 1186.050 1189.910
\hypertarget{c01:cue:469}{}在睡覺時都要看。
% cue 470 1189.910 1191.060
\hypertarget{c01:cue:470}{}那也太辛苦。
% cue 471 1191.060 1192.060
\hypertarget{c01:cue:471}{}就是沒事的時候就拿來看一看。
% cue 472 1192.060 1193.630
\hypertarget{c01:cue:472}{}這對你們來講。
% cue 473 1193.630 1194.630
\hypertarget{c01:cue:473}{}只有好處沒有壞處。
% cue 474 1194.630 1195.910
\hypertarget{c01:cue:474}{}就是常常聽。
% cue 475 1195.910 1198.070
\hypertarget{c01:cue:475}{}讓你的頭腦可以磨練。
% cue 476 1198.070 1200.360
\hypertarget{c01:cue:476}{}磨練的精準。
% cue 477 1200.360 1201.480
\hypertarget{c01:cue:477}{}磨練的精準。
% cue 478 1201.480 1203.020
\hypertarget{c01:cue:478}{}聽別人講話。
% cue 479 1203.020 1204.980
\hypertarget{c01:cue:479}{}你很快就聽出。
% cue 480 1204.980 1206.980
\hypertarget{c01:cue:480}{}他講話的美感在哪裡。
% cue 481 1206.980 1208.520
\hypertarget{c01:cue:481}{}這比有人聽過。
% cue 482 1208.520 1209.520
\hypertarget{c01:cue:482}{}我要慢聽到他的腳事。
% cue 483 1209.520 1210.810
\hypertarget{c01:cue:483}{}讓你講多了吧。

\tc{0:20:10}
% cue 484 1210.810 1212.970
\hypertarget{c01:cue:484}{}好 另外一個請你們看一下。
% cue 485 1212.970 1215.940
\hypertarget{c01:cue:485}{}上學期心得分享\footnote{課綱「參考書資料暨網路資源」所列「前期上課心得分享」，為 Google Drive 上的檔案。}。
% cue 486 1215.940 1217.940
\hypertarget{c01:cue:486}{}那是你們上學期的同學寫的。
% cue 487 1217.940 1219.190
\hypertarget{c01:cue:487}{}老師把它整理。
% cue 488 1219.190 1220.190
\hypertarget{c01:cue:488}{}沒有全部網易。
% cue 489 1220.190 1222.160
\hypertarget{c01:cue:489}{}你大概挑的是二分之一。
% cue 490 1222.160 1223.160
\hypertarget{c01:cue:490}{}二分之一沒有。
% cue 491 1223.160 1224.160
\hypertarget{c01:cue:491}{}三分之一四分之一大概。
% cue 492 1224.160 1225.540
\hypertarget{c01:cue:492}{}不記得多少。
% cue 493 1225.540 1226.790
\hypertarget{c01:cue:493}{}反正是言之有物的。
% cue 494 1226.790 1228.770
\hypertarget{c01:cue:494}{}我們就把它放一放。
% cue 495 1228.770 1231.070
\hypertarget{c01:cue:495}{}找到了嗎。
% cue 496 1231.070 1233.420
\hypertarget{c01:cue:496}{}有這個藍威嗎。
% cue 497 1233.420 1235.610
\hypertarget{c01:cue:497}{}有。
% cue 498 1235.610 1235.710
\hypertarget{c01:cue:498}{}沒有。
% cue 499 1235.610 1238.230
\hypertarget{c01:cue:499}{}有找不到的人用舉手。
% cue 500 1238.230 1239.230
\hypertarget{c01:cue:500}{}因為我現在在跟你講話。
% cue 501 1239.230 1240.230
\hypertarget{c01:cue:501}{}你如果有困難。
% cue 502 1240.230 1241.290
\hypertarget{c01:cue:502}{}趕快問一下別人。
% cue 503 1241.290 1243.290
\hypertarget{c01:cue:503}{}至少你現在跟得上。
% cue 504 1243.290 1244.700
\hypertarget{c01:cue:504}{}你如果這樣放棄。
% cue 505 1244.700 1245.700
\hypertarget{c01:cue:505}{}那我就不知道該怎麼放。

\tc{0:20:45}
% cue 506 1245.700 1247.820
\hypertarget{c01:cue:506}{}一個老師上課要準備。
% cue 507 1247.820 1248.820
\hypertarget{c01:cue:507}{}這麼多東西給你。
% cue 508 1248.820 1249.820
\hypertarget{c01:cue:508}{}這個老師對你們是和好。
% cue 509 1249.820 1252.170
\hypertarget{c01:cue:509}{}講不好聽。
% cue 510 1252.170 1254.170
\hypertarget{c01:cue:510}{}你這個老師有點像假大修行。
% cue 511 1254.170 1255.610
\hypertarget{c01:cue:511}{}你一般上課就上課。
% cue 512 1255.610 1256.610
\hypertarget{c01:cue:512}{}講那麼多幹什麼。
% cue 513 1256.610 1257.610
\hypertarget{c01:cue:513}{}給這些幹什麼。
% cue 514 1257.610 1258.670
\hypertarget{c01:cue:514}{}給這些都要花力氣的。
% cue 515 1258.670 1260.080
\hypertarget{c01:cue:515}{}給這些都要整理的。
% cue 516 1260.080 1264.460
\hypertarget{c01:cue:516}{}給這些都要花力氣的。
% cue 517 1264.460 1265.460
\hypertarget{c01:cue:517}{}給這些都要整理的。
% cue 518 1265.460 1265.560
\hypertarget{c01:cue:518}{}你看上學期的人。
% cue 519 1265.460 1267.960
\hypertarget{c01:cue:519}{}他們是怎麼樣給你們的經驗。

\tc{0:21:07}
% cue 520 1267.960 1270.120
\hypertarget{c01:cue:520}{}我現在覺得這個事情。
% cue 521 1270.120 1272.340
\hypertarget{c01:cue:521}{}越來越讓人憂慮。
% cue 522 1272.340 1274.650
\hypertarget{c01:cue:522}{}你們好像跟同學之間。
% cue 523 1274.650 1277.750
\hypertarget{c01:cue:523}{}還有跟學長之間很少聯繫。
% cue 524 1277.750 1280.030
\hypertarget{c01:cue:524}{}跟你們的學弟妹也很少聯繫。
% cue 525 1280.030 1282.250
\hypertarget{c01:cue:525}{}為什麼學校越來越冷漠。
% cue 526 1282.250 1284.250
\hypertarget{c01:cue:526}{}越來越個人化。
% cue 527 1284.250 1285.470
\hypertarget{c01:cue:527}{}知道。
% cue 528 1285.470 1286.720
\hypertarget{c01:cue:528}{}沒有辦法回答。
% cue 529 1286.720 1288.200
\hypertarget{c01:cue:529}{}以前有一陣子。
% cue 530 1288.200 1289.200
\hypertarget{c01:cue:530}{}我還擔心說。
% cue 531 1289.200 1291.200
\hypertarget{c01:cue:531}{}學長給了錯誤的訊息。
% cue 532 1291.200 1293.200
\hypertarget{c01:cue:532}{}會影響到你們的學者。
% cue 533 1293.200 1294.450
\hypertarget{c01:cue:533}{}但我們現在發現。
% cue 534 1294.450 1295.450
\hypertarget{c01:cue:534}{}你們根本沒有學長。
% cue 535 1295.450 1296.470
\hypertarget{c01:cue:535}{}連對學弟妹。
% cue 536 1296.470 1297.760
\hypertarget{c01:cue:536}{}你們獨來獨往的一個人。
% cue 537 1297.760 1299.330
\hypertarget{c01:cue:537}{}全世界我最大。
% cue 538 1299.330 1303.030
\hypertarget{c01:cue:538}{}沒有年輕不太接收訊息。
% cue 539 1303.030 1304.540
\hypertarget{c01:cue:539}{}好了。
% cue 540 1304.540 1305.540
\hypertarget{c01:cue:540}{}這是有關。
% cue 541 1305.540 1307.540
\hypertarget{c01:cue:541}{}有關應該要給的東西都給了。

\section{教學目標與憲法必修時數的縮減}

\tc{0:21:47}
% cue 542 1307.540 1309.600
\hypertarget{c01:cue:542}{}現在我們來講一下教學目標。
% cue 543 1309.600 1311.950
\hypertarget{c01:cue:543}{}這個學期的課。
% cue 544 1311.950 1312.950
\hypertarget{c01:cue:544}{}憲法專題。
% cue 545 1312.950 1313.950
\hypertarget{c01:cue:545}{}專題是什麼名字。
% cue 546 1313.950 1314.950
\hypertarget{c01:cue:546}{}權力分立原則\footnote{英文 doctrine of separation of powers。課綱教學目標規定本學期以此原則為核心，先講解其理論與實作（政府體制典範），再檢驗大法官有關「政府體制」（或稱「憲法結構」，constitutional structure）領域的指標解釋。}。
% cue 547 1314.950 1318.120
\hypertarget{c01:cue:547}{}那學校開這個課。
% cue 548 1318.120 1320.230
\hypertarget{c01:cue:548}{}配合學校。
% cue 549 1320.230 1322.190
\hypertarget{c01:cue:549}{}把你們的憲法必修課。
% cue 550 1322.190 1323.190
\hypertarget{c01:cue:550}{}中環部憲法必修課。
% cue 551 1323.190 1326.720
\hypertarget{c01:cue:551}{}從每個學期兩個學分。
% cue 552 1326.720 1328.810
\hypertarget{c01:cue:552}{}一整年四學分。
% cue 553 1328.810 1330.870
\hypertarget{c01:cue:553}{}改成一個學期只有三學分。
% cue 554 1330.870 1335.600
\hypertarget{c01:cue:554}{}我必須要說這個改變不是我樂見的改變。
% cue 555 1335.600 1338.790
\hypertarget{c01:cue:555}{}但是那時候我還沒有到學校來專任。
% cue 556 1338.790 1341.300
\hypertarget{c01:cue:556}{}專任之前老師在東吳。
% cue 557 1341.300 1343.380
\hypertarget{c01:cue:557}{}學校教多少。
% cue 558 1343.380 1344.480
\hypertarget{c01:cue:558}{}二十五。
% cue 559 1344.480 1346.330
\hypertarget{c01:cue:559}{}教我們二屬於憲法。
% cue 560 1346.330 1347.330
\hypertarget{c01:cue:560}{}因為非常多的學長。
% cue 561 1347.330 1349.360
\hypertarget{c01:cue:561}{}教我們爸爸媽媽。
% cue 562 1349.360 1350.430
\hypertarget{c01:cue:562}{}已經經歷了。
% cue 563 1350.430 1354.840
\hypertarget{c01:cue:563}{}教我們。
% cue 564 1354.840 1355.740
\hypertarget{c01:cue:564}{}這個改變是錯誤。

\tc{0:22:35}
% cue 565 1355.740 1357.320
\hypertarget{c01:cue:565}{}錯誤的原因很簡單。
% cue 566 1357.320 1359.510
\hypertarget{c01:cue:566}{}大法官的解釋越來越多。
% cue 567 1359.510 1361.690
\hypertarget{c01:cue:567}{}為什麼我們憲法課的時數越來越少。
% cue 568 1361.690 1362.870
\hypertarget{c01:cue:568}{}請說給我聽聽看。
% cue 569 1362.870 1365.120
\hypertarget{c01:cue:569}{}還是你認為他媽媽都在放屏。
% cue 570 1365.120 1366.440
\hypertarget{c01:cue:570}{}講的話你都沒有聽。
% cue 571 1366.440 1368.660
\hypertarget{c01:cue:571}{}所以他出他的幾十以上的課。
% cue 572 1368.660 1369.660
\hypertarget{c01:cue:572}{}是這樣。
% cue 573 1369.660 1371.130
\hypertarget{c01:cue:573}{}第二。
% cue 574 1371.130 1373.120
\hypertarget{c01:cue:574}{}現在改成一個學期三學分。
% cue 575 1373.120 1375.090
\hypertarget{c01:cue:575}{}普遍的情況就是教不完。
% cue 576 1375.090 1379.250
\hypertarget{c01:cue:576}{}因為兩個學期各兩學分都教不完。
% cue 577 1379.250 1380.880
\hypertarget{c01:cue:577}{}我認為比較合理的是。
% cue 578 1380.880 1382.880
\hypertarget{c01:cue:578}{}兩個學期各三學分。
% cue 579 1382.880 1386.730
\hypertarget{c01:cue:579}{}不然你憲法是學不好。

\tc{0:23:06}
% cue 580 1386.730 1390.200
\hypertarget{c01:cue:580}{}那我聽到之前我不是第一次開這個課。
% cue 581 1390.200 1393.230
\hypertarget{c01:cue:581}{}老師 2019 年從十八年到晚卸了一個。
% cue 582 1393.230 1395.480
\hypertarget{c01:cue:582}{}成功點的學校的好業。
% cue 583 1395.480 1396.480
\hypertarget{c01:cue:583}{}就來這裡專任了。
% cue 584 1396.480 1399.020
\hypertarget{c01:cue:584}{}那我已經開了二年的課。
% cue 585 1399.020 1401.050
\hypertarget{c01:cue:585}{}現在二年都要八年了。
% cue 586 1401.050 1404.370
\hypertarget{c01:cue:586}{}八年了。
% cue 587 1404.370 1405.750
\hypertarget{c01:cue:587}{}所以我很清楚。
% cue 588 1405.750 1407.320
\hypertarget{c01:cue:588}{}每一年我問。
% cue 589 1407.320 1408.220
\hypertarget{c01:cue:589}{}大家都說。
% cue 590 1408.220 1410.350
\hypertarget{c01:cue:590}{}老師我來上你的課。
% cue 591 1410.350 1411.350
\hypertarget{c01:cue:591}{}是因為我讀了憲法學的一趟。
% cue 592 1411.350 1412.350
\hypertarget{c01:cue:592}{}為什麼。
% cue 593 1412.350 1413.990
\hypertarget{c01:cue:593}{}因為我從來講不完。
% cue 594 1413.990 1415.660
\hypertarget{c01:cue:594}{}為什麼講不完。
% cue 595 1415.660 1416.660
\hypertarget{c01:cue:595}{}因為時數太少。
% cue 596 1416.660 1419.060
\hypertarget{c01:cue:596}{}剛剛進來又沒有概念。
% cue 597 1419.060 1420.090
\hypertarget{c01:cue:597}{}好。
% cue 598 1420.090 1421.090
\hypertarget{c01:cue:598}{}這是一個多講的前面意。

\tc{0:23:41}
% cue 599 1421.090 1423.570
\hypertarget{c01:cue:599}{}老師剛開始教書的時候。
% cue 600 1423.570 1426.080
\hypertarget{c01:cue:600}{}憲法二年級三年級。
% cue 601 1426.080 1427.100
\hypertarget{c01:cue:601}{}可能是一年級。
% cue 602 1427.100 1429.450
\hypertarget{c01:cue:602}{}我也不喜歡他放在一年級。
% cue 603 1429.450 1431.550
\hypertarget{c01:cue:603}{}一年級的小朋友。
% cue 604 1431.550 1435.140
\hypertarget{c01:cue:604}{}剛剛大學聯考或者推甄進來。
% cue 605 1435.140 1437.550
\hypertarget{c01:cue:605}{}心裡面只是要求解放。
% cue 606 1437.550 1439.640
\hypertarget{c01:cue:606}{}我要認識男女朋友。
% cue 607 1439.640 1440.640
\hypertarget{c01:cue:607}{}我要出去玩。
% cue 608 1440.640 1442.770
\hypertarget{c01:cue:608}{}我要享受我的大學生活。
% cue 609 1442.770 1445.080
\hypertarget{c01:cue:609}{}你跟我講做國家主人。
% cue 610 1445.080 1447.080
\hypertarget{c01:cue:610}{}教我怎麼做國家主人。
% cue 611 1447.080 1448.210
\hypertarget{c01:cue:611}{}這個好無聊喔。
% cue 612 1448.210 1450.300
\hypertarget{c01:cue:612}{}這等到老點再來消口。
% cue 613 1450.300 1453.450
\hypertarget{c01:cue:613}{}為什麼學生要不用排在一年級。
% cue 614 1453.450 1454.960
\hypertarget{c01:cue:614}{}我就不懂啊。

\tc{0:24:14}
% cue 615 1454.960 1456.470
\hypertarget{c01:cue:615}{}如果是我。
% cue 616 1456.470 1458.690
\hypertarget{c01:cue:616}{}我認為這套課是必修課。
% cue 617 1458.690 1461.610
\hypertarget{c01:cue:617}{}大學四年之內你要修完。
% cue 618 1461.610 1462.530
\hypertarget{c01:cue:618}{}不然不能畢業。
% cue 619 1462.530 1463.840
\hypertarget{c01:cue:619}{}這麼簡單。
% cue 620 1463.840 1465.860
\hypertarget{c01:cue:620}{}那你要哪年修。
% cue 621 1465.860 1466.860
\hypertarget{c01:cue:621}{}按照你心趣。
% cue 622 1466.860 1468.080
\hypertarget{c01:cue:622}{}成熟的年齡。
% cue 623 1468.080 1471.580
\hypertarget{c01:cue:623}{}你是一個找得比較快的人。
% cue 624 1471.580 1473.960
\hypertarget{c01:cue:624}{}就覺得你要當國家主人。
% cue 625 1473.960 1474.960
\hypertarget{c01:cue:625}{}那你就來上課。
% cue 626 1474.960 1475.960
\hypertarget{c01:cue:626}{}我現在要玩。
% cue 627 1475.960 1477.960
\hypertarget{c01:cue:627}{}我不想當國家主人還是孩子。
% cue 628 1477.960 1479.240
\hypertarget{c01:cue:628}{}這一次有爸媽去團黨。
% cue 629 1479.240 1481.240
\hypertarget{c01:cue:629}{}天塌下來有別人的主持。
% cue 630 1481.240 1482.360
\hypertarget{c01:cue:630}{}那你就不要念。
% cue 631 1482.360 1483.580
\hypertarget{c01:cue:631}{}反正大學畢業之前念就好。
% cue 632 1483.580 1485.930
\hypertarget{c01:cue:632}{}但是後來教育部體力。
% cue 633 1485.930 1486.930
\hypertarget{c01:cue:633}{}不一定要大一。
% cue 634 1486.930 1487.930
\hypertarget{c01:cue:634}{}那現在就變成大一的課。
% cue 635 1487.930 1489.890
\hypertarget{c01:cue:635}{}大一現在學校都自己創造說。
% cue 636 1489.890 1490.890
\hypertarget{c01:cue:636}{}我們覺得這個課呢。
% cue 637 1490.890 1491.890
\hypertarget{c01:cue:637}{}太難了。
% cue 638 1491.890 1492.890
\hypertarget{c01:cue:638}{}所以把他砍一半。
% cue 639 1492.890 1494.500
\hypertarget{c01:cue:639}{}讓我們覺得教學。
% cue 640 1494.500 1496.630
\hypertarget{c01:cue:640}{}這是真的造成教學的課。

\tc{0:24:56}
% cue 641 1496.630 1497.620
\hypertarget{c01:cue:641}{}然後說呢。
% cue 642 1497.620 1498.850
\hypertarget{c01:cue:642}{}沒有關係。
% cue 643 1498.850 1499.910
\hypertarget{c01:cue:643}{}我們約一個這樣子的課。
% cue 644 1499.910 1500.910
\hypertarget{c01:cue:644}{}然後呢。
% cue 645 1500.910 1502.910
\hypertarget{c01:cue:645}{}有一個學習憲法案例解析。
% cue 646 1502.910 1504.900
\hypertarget{c01:cue:646}{}有一個學習憲法專題。
% cue 647 1504.900 1510.180
\hypertarget{c01:cue:647}{}讓你們可以自己問一個問題。
% cue 648 1510.180 1512.240
\hypertarget{c01:cue:648}{}如果他的基礎沒有打好。
% cue 649 1512.240 1516.850
\hypertarget{c01:cue:649}{}後面要使用。
% cue 650 1516.850 1518.620
\hypertarget{c01:cue:650}{}讓老師經歷。
% cue 651 1518.620 1519.870
\hypertarget{c01:cue:651}{}那這個學期。
% cue 652 1519.870 1520.930
\hypertarget{c01:cue:652}{}我們就是講。
% cue 653 1520.930 1524.600
\hypertarget{c01:cue:653}{}這個權力分立原則。

\section{憲法第 22 條與隱私權}

\tc{0:25:24}
% cue 654 1524.600 1528.290
\hypertarget{c01:cue:654}{}至於在這個學校裡面。
% cue 655 1528.290 1529.690
\hypertarget{c01:cue:655}{}開過的其他專題。
% cue 656 1529.690 1530.690
\hypertarget{c01:cue:656}{}舉例說。
% cue 657 1530.690 1531.820
\hypertarget{c01:cue:657}{}憲法 22 條\footnote{中華民國憲法第 22 條：「凡人民之其他自由及權利，不妨害社會秩序公共利益者，均受憲法之保障。」司法院釋字第 603 號解釋以本條為隱私權的依據。\url{https://law.moj.gov.tw/LawClass/LawAll.aspx?pcode=A0000001}。}。
% cue 658 1531.820 1532.850
\hypertarget{c01:cue:658}{}現在我講過。
% cue 659 1532.850 1534.850
\hypertarget{c01:cue:659}{}你沒有念下請你翻卡條。
% cue 660 1534.850 1536.140
\hypertarget{c01:cue:660}{}憲法 22 條。
% cue 661 1536.140 1537.140
\hypertarget{c01:cue:661}{}看到這個條問題。
% cue 662 1537.140 1538.140
\hypertarget{c01:cue:662}{}舉手。
% cue 663 1538.140 1539.190
\hypertarget{c01:cue:663}{}請他念一下。
% cue 664 1539.190 1540.550
\hypertarget{c01:cue:664}{}憲法 22 條。
% cue 665 1540.550 1541.550
\hypertarget{c01:cue:665}{}從什麼樣子找到的情緒。
% cue 666 1541.550 1556.440
\hypertarget{c01:cue:666}{}以前你們上課的時候。
% cue 667 1556.440 1557.440
\hypertarget{c01:cue:667}{}抓到這一條。
% cue 668 1557.440 1558.910
\hypertarget{c01:cue:668}{}假如有。
% cue 669 1558.910 1560.910
\hypertarget{c01:cue:669}{}那老師就算是認真了。
% cue 670 1560.910 1562.260
\hypertarget{c01:cue:670}{}那這一條。
% cue 671 1562.260 1563.260
\hypertarget{c01:cue:671}{}我們通常叫做。
% cue 672 1563.260 1564.260
\hypertarget{c01:cue:672}{}未有列舉的基本。
% cue 673 1564.260 1566.600
\hypertarget{c01:cue:673}{}跟前面的。
% cue 674 1566.600 1568.330
\hypertarget{c01:cue:674}{}第七條\footnote{中華民國憲法第 7 條：「中華民國人民，無分男女、宗教、種族、階級、黨派，在法律上一律平等。」第 8 條至第 18 條與第 21 條另列舉人身自由、居住遷徙、言論講學著作出版、秘密通訊、信仰宗教、集會結社、生存權工作權財產權、請願訴願訴訟、參政、應考試服公職及受國民教育等權利。\url{https://law.moj.gov.tw/LawClass/LawAll.aspx?pcode=A0000001}。}。
% cue 675 1568.330 1569.330
\hypertarget{c01:cue:675}{}平等去開始。
% cue 676 1569.330 1570.420
\hypertarget{c01:cue:676}{}每一個權力。
% cue 677 1570.420 1573.190
\hypertarget{c01:cue:677}{}都有一個特定的範圍。
% cue 678 1573.190 1574.590
\hypertarget{c01:cue:678}{}言論自由啦。
% cue 679 1574.590 1575.820
\hypertarget{c01:cue:679}{}什麼什麼。
% cue 680 1575.820 1576.730
\hypertarget{c01:cue:680}{}這條學沒有範圍。
% cue 681 1576.730 1577.790
\hypertarget{c01:cue:681}{}不一樣。
% cue 682 1577.790 1581.320
\hypertarget{c01:cue:682}{}我們的專題裡面。
% cue 683 1581.320 1582.320
\hypertarget{c01:cue:683}{}你們去上。
% cue 684 1582.320 1583.320
\hypertarget{c01:cue:684}{}老師的 YouTube。
% cue 685 1583.320 1584.540
\hypertarget{c01:cue:685}{}我有一個學期。
% cue 686 1584.540 1585.540
\hypertarget{c01:cue:686}{}整個學期。
% cue 687 1585.540 1586.540
\hypertarget{c01:cue:687}{}再講憲法二十二條。
% cue 688 1586.540 1587.540
\hypertarget{c01:cue:688}{}就憲法專題。
% cue 689 1587.540 1588.450
\hypertarget{c01:cue:689}{}我就講。
% cue 690 1588.450 1589.530
\hypertarget{c01:cue:690}{}未有列舉的基本權。
% cue 691 1589.530 1590.460
\hypertarget{c01:cue:691}{}二十二條。
% cue 692 1590.460 1591.750
\hypertarget{c01:cue:692}{}二十二條。
% cue 693 1591.750 1592.750
\hypertarget{c01:cue:693}{}你看是一個條文。
% cue 694 1592.750 1593.750
\hypertarget{c01:cue:694}{}這一個複雜的。
% cue 695 1593.750 1595.510
\hypertarget{c01:cue:695}{}不得了。
% cue 696 1595.510 1596.510
\hypertarget{c01:cue:696}{}非常見證的隱私權\footnote{司法院釋字第 603 號解釋文：「隱私權雖非憲法明文列舉之權利，惟基於人性尊嚴與個人主體性之維護及人格發展之完整，並為保障個人生活私密領域免於他人侵擾及個人資料之自主控制，隱私權乃為不可或缺之基本權利，而受憲法第二十二條所保障（本院釋字第五八五號解釋參照）。」\url{https://cons.judicial.gov.tw/docdata.aspx?fid=100&id=310784}。}。
% cue 697 1596.510 1597.960
\hypertarget{c01:cue:697}{}光光隱私權。
% cue 698 1597.960 1598.960
\hypertarget{c01:cue:698}{}就可以講高級的一半。

\tc{0:26:38}
% cue 699 1598.960 1599.960
\hypertarget{c01:cue:699}{}隱私權有好多種。
% cue 700 1599.960 1601.110
\hypertarget{c01:cue:700}{}舉例說。
% cue 701 1601.110 1603.360
\hypertarget{c01:cue:701}{}我沒有得到你的同意。
% cue 702 1603.360 1604.360
\hypertarget{c01:cue:702}{}可被在你身上。
% cue 703 1604.360 1606.360
\hypertarget{c01:cue:703}{}紮一個真臭無細細的血。
% cue 704 1606.360 1608.800
\hypertarget{c01:cue:704}{}你跟我說。
% cue 705 1608.800 1610.290
\hypertarget{c01:cue:705}{}可以。
% cue 706 1610.290 1611.400
\hypertarget{c01:cue:706}{}有人說不可以。
% cue 707 1611.400 1612.620
\hypertarget{c01:cue:707}{}不可以。
% cue 708 1612.620 1614.920
\hypertarget{c01:cue:708}{}不可以欺貪你手。
% cue 709 1614.920 1615.880
\hypertarget{c01:cue:709}{}這個隱私。
% cue 710 1615.880 1616.910
\hypertarget{c01:cue:710}{}是身體。
% cue 711 1616.910 1618.230
\hypertarget{c01:cue:711}{}物理的隱私。
% cue 712 1618.230 1619.700
\hypertarget{c01:cue:712}{}沒有經過你同意。
% cue 713 1619.700 1620.700
\hypertarget{c01:cue:713}{}可以進你的房間。
% cue 714 1620.700 1621.700
\hypertarget{c01:cue:714}{}可以進你家。
% cue 715 1621.700 1623.490
\hypertarget{c01:cue:715}{}就是這一種概念。

\tc{0:27:03}
% cue 716 1623.490 1624.740
\hypertarget{c01:cue:716}{}第二種隱私。
% cue 717 1624.740 1625.740
\hypertarget{c01:cue:717}{}英文叫做。
% cue 718 1625.740 1626.740
\hypertarget{c01:cue:718}{}decisional privacy\footnote{英文，可譯為自主決定的隱私，指就涉及自身的事項自行決定而不受他人干預的權利。Stanford Encyclopedia of Philosophy 的 Privacy 條目記載，美國法上的隱私權討論以此為重心，所涉為同性婚姻與墮胎；美國聯邦最高法院於 1965 年明確承認一項獨立於資訊隱私與聯邦憲法第四修正案的隱私權。\url{https://plato.stanford.edu/entries/privacy/}。}。
% cue 719 1626.740 1627.760
\hypertarget{c01:cue:719}{}這個對女生。
% cue 720 1627.760 1628.760
\hypertarget{c01:cue:720}{}我去整局。
% cue 721 1628.760 1629.760
\hypertarget{c01:cue:721}{}這個例子特別懂。
% cue 722 1629.760 1632.100
\hypertarget{c01:cue:722}{}你懷孕了。
% cue 723 1632.100 1633.640
\hypertarget{c01:cue:723}{}你要不要繼續懷孕。
% cue 724 1633.640 1634.640
\hypertarget{c01:cue:724}{}還是你要終止懷孕。
% cue 725 1634.640 1635.640
\hypertarget{c01:cue:725}{}就要多灘。
% cue 726 1635.640 1637.400
\hypertarget{c01:cue:726}{}那這個事情。
% cue 727 1637.400 1640.130
\hypertarget{c01:cue:727}{}跟其他人有什麼關係。
% cue 728 1640.130 1641.700
\hypertarget{c01:cue:728}{}路人說。
% cue 729 1641.700 1642.660
\hypertarget{c01:cue:729}{}欸。
% cue 730 1642.660 1644.000
\hypertarget{c01:cue:730}{}我不可以。
% cue 731 1644.000 1645.220
\hypertarget{c01:cue:731}{}終止生命。
% cue 732 1645.220 1646.220
\hypertarget{c01:cue:732}{}這種不道德。
% cue 733 1646.220 1649.110
\hypertarget{c01:cue:733}{}那你是現代女性。
% cue 734 1649.110 1650.450
\hypertarget{c01:cue:734}{}你會跟她講。
% cue 735 1650.450 1651.860
\hypertarget{c01:cue:735}{}孩子懷在你肚子裡。
% cue 736 1651.860 1652.860
\hypertarget{c01:cue:736}{}還是懷在我肚子裡。
% cue 737 1652.860 1654.370
\hypertarget{c01:cue:737}{}懷在我肚子裡。
% cue 738 1654.370 1655.370
\hypertarget{c01:cue:738}{}You shut up。
% cue 739 1655.370 1656.590
\hypertarget{c01:cue:739}{}你懂的屁。
% cue 740 1656.590 1658.940
\hypertarget{c01:cue:740}{}跟你半點關係都沒有。
% cue 741 1658.940 1659.940
\hypertarget{c01:cue:741}{}痛苦是我痛苦。
% cue 742 1659.940 1661.910
\hypertarget{c01:cue:742}{}你講的是我將來的發展。
% cue 743 1661.910 1663.770
\hypertarget{c01:cue:743}{}會影響到的是我的身體。
% cue 744 1663.770 1664.770
\hypertarget{c01:cue:744}{}跟我的人格。
% cue 745 1664.770 1665.830
\hypertarget{c01:cue:745}{}跟你什麼關係。
% cue 746 1665.830 1666.830
\hypertarget{c01:cue:746}{}所以這件事情。
% cue 747 1666.830 1667.830
\hypertarget{c01:cue:747}{}你沒有決定的權。
% cue 748 1667.830 1669.020
\hypertarget{c01:cue:748}{}靠房間站。
% cue 749 1669.020 1670.140
\hypertarget{c01:cue:749}{}這個英文叫做。
% cue 750 1670.140 1671.140
\hypertarget{c01:cue:750}{}decisional privacy。
% cue 751 1671.140 1673.510
\hypertarget{c01:cue:751}{}還有一個。

\tc{0:27:53}
% cue 752 1673.510 1674.510
\hypertarget{c01:cue:752}{}所以我們到。
% cue 753 1674.510 1675.510
\hypertarget{c01:cue:753}{}603 號解釋\footnote{司法院釋字第 603 號解釋（2005 年 9 月 28 日）宣告戶籍法第 8 條第 2 項、第 3 項強制按捺指紋並予錄存，否則不予發給國民身分證之規定，與憲法第 22 條、第 23 條之意旨不符，自解釋公布之日起不再適用。解釋文稱資訊隱私權「乃保障人民決定是否揭露其個人資料、及在何種範圍內、於何時、以何種方式、向何人揭露之決定權，並保障人民對其個人資料之使用有知悉與控制權及資料記載錯誤之更正權」。\url{https://cons.judicial.gov.tw/docdata.aspx?fid=100&id=310784}。}講的。
% cue 754 1675.510 1676.860
\hypertarget{c01:cue:754}{}603 號解釋。
% cue 755 1676.860 1678.860
\hypertarget{c01:cue:755}{}我當你這樣講的時候。
% cue 756 1678.860 1680.080
\hypertarget{c01:cue:756}{}有任何一個人。
% cue 757 1680.080 1681.080
\hypertarget{c01:cue:757}{}不管你幾年級。
% cue 758 1681.080 1682.080
\hypertarget{c01:cue:758}{}你現在可以舉手。
% cue 759 1682.080 1683.050
\hypertarget{c01:cue:759}{}給我。
% cue 760 1683.050 1684.050
\hypertarget{c01:cue:760}{}603 號解釋。
% cue 761 1684.050 1685.230
\hypertarget{c01:cue:761}{}大概在講什麼內容。
% cue 762 1685.230 1687.350
\hypertarget{c01:cue:762}{}我今天跟你講。
% cue 763 1687.350 1689.150
\hypertarget{c01:cue:763}{}跟理事有關。
% cue 764 1689.150 1690.150
\hypertarget{c01:cue:764}{}任何一個人都可以舉手。
% cue 765 1690.150 1691.150
\hypertarget{c01:cue:765}{}講錯人歸過分。
% cue 766 1691.150 1692.210
\hypertarget{c01:cue:766}{}只會加分。
% cue 767 1692.210 1693.590
\hypertarget{c01:cue:767}{}等一下我會講。
% cue 768 1693.590 1694.650
\hypertarget{c01:cue:768}{}這是我們學習。
% cue 769 1694.650 1695.650
\hypertarget{c01:cue:769}{}更好表現的一種。
% cue 770 1695.650 1697.410
\hypertarget{c01:cue:770}{}請說。
% cue 771 1697.410 1698.410
\hypertarget{c01:cue:771}{}舉高一點。
% cue 772 1698.410 1699.410
\hypertarget{c01:cue:772}{}不要害怕。
% cue 773 1699.410 1700.410
\hypertarget{c01:cue:773}{}請講。
% cue 774 1700.410 1701.410
\hypertarget{c01:cue:774}{}603 講什麼。
% cue 775 1701.410 1702.410
\hypertarget{c01:cue:775}{}用大意講。
% cue 776 1702.410 1703.410
\hypertarget{c01:cue:776}{}你現在也不急看。
% cue 777 1703.410 1704.410
\hypertarget{c01:cue:777}{}來不及看。
% cue 778 1704.410 1705.410
\hypertarget{c01:cue:778}{}你不要緊。
% cue 779 1705.410 1706.630
\hypertarget{c01:cue:779}{}你現在有關於什麼。
% cue 780 1706.630 1709.290
\hypertarget{c01:cue:780}{}這件事情。

\tc{0:28:29}
\speaker{學生}
% cue 781 1709.290 1715.630
\hypertarget{c01:cue:781}{}這件事情。

\tc{0:28:35}
\speaker{老師}
% cue 782 1715.630 1716.630
\hypertarget{c01:cue:782}{}牽涉到哪一種隱私呢。
% cue 783 1716.630 1717.630
\hypertarget{c01:cue:783}{}好。

\tc{0:28:37}
% cue 784 1717.630 1720.030
\hypertarget{c01:cue:784}{}那現在回來。
% cue 785 1720.030 1721.030
\hypertarget{c01:cue:785}{}講剛才。
% cue 786 1721.030 1722.030
\hypertarget{c01:cue:786}{}講得對。
% cue 787 1722.030 1723.030
\hypertarget{c01:cue:787}{}你等一下下課。
% cue 788 1723.030 1724.030
\hypertarget{c01:cue:788}{}來找我。
% cue 789 1724.030 1725.510
\hypertarget{c01:cue:789}{}我們要加分。
% cue 790 1725.510 1726.510
\hypertarget{c01:cue:790}{}這個資訊隱私呢。
% cue 791 1726.510 1727.500
\hypertarget{c01:cue:791}{}是另外一種隱私。
% cue 792 1727.500 1728.500
\hypertarget{c01:cue:792}{}剛才我們講的。
% cue 793 1728.500 1729.500
\hypertarget{c01:cue:793}{}第一種是物理上的隱私。
% cue 794 1729.500 1730.500
\hypertarget{c01:cue:794}{}我的房間。
% cue 795 1730.500 1731.440
\hypertarget{c01:cue:795}{}我的家。
% cue 796 1731.440 1733.550
\hypertarget{c01:cue:796}{}第二種是決策隱私。
% cue 797 1733.550 1734.550
\hypertarget{c01:cue:797}{}因為這影響我們。
% cue 798 1734.550 1735.550
\hypertarget{c01:cue:798}{}這個決定影響。
% cue 799 1735.550 1736.520
\hypertarget{c01:cue:799}{}不影響到你們。
% cue 800 1736.520 1737.510
\hypertarget{c01:cue:800}{}不幹。
% cue 801 1737.510 1738.510
\hypertarget{c01:cue:801}{}幹清合適。
% cue 802 1738.510 1739.510
\hypertarget{c01:cue:802}{}幹裹事。
% cue 803 1739.510 1740.800
\hypertarget{c01:cue:803}{}為重大。
% cue 804 1740.800 1742.020
\hypertarget{c01:cue:804}{}這現在講第三種。
% cue 805 1742.020 1744.020
\hypertarget{c01:cue:805}{}有關於個人資料。
% cue 806 1744.020 1745.020
\hypertarget{c01:cue:806}{}每一個人。
% cue 807 1745.020 1746.210
\hypertarget{c01:cue:807}{}每一個主體。
% cue 808 1746.210 1748.260
\hypertarget{c01:cue:808}{}都有權利決定。
% cue 809 1748.260 1750.380
\hypertarget{c01:cue:809}{}有關於你的資料。
% cue 810 1750.380 1751.470
\hypertarget{c01:cue:810}{}要不要。
% cue 811 1751.470 1753.660
\hypertarget{c01:cue:811}{}在什麼範圍內。
% cue 812 1753.660 1755.660
\hypertarget{c01:cue:812}{}在什麼條件下。
% cue 813 1755.660 1756.690
\hypertarget{c01:cue:813}{}要向誰公開。
% cue 814 1756.690 1758.880
\hypertarget{c01:cue:814}{}以及公開以後。
% cue 815 1758.880 1761.140
\hypertarget{c01:cue:815}{}那個資料跑到哪裡去。
% cue 816 1761.140 1764.000
\hypertarget{c01:cue:816}{}然後你要不要收回。
% cue 817 1764.000 1765.000
\hypertarget{c01:cue:817}{}註銷。
% cue 818 1765.000 1766.000
\hypertarget{c01:cue:818}{}不要讓它繼續公開。
% cue 819 1766.000 1768.610
\hypertarget{c01:cue:819}{}這種對個人資訊。
% cue 820 1768.610 1771.170
\hypertarget{c01:cue:820}{}所以這叫資訊。
% cue 821 1771.170 1772.960
\hypertarget{c01:cue:821}{}英文叫 information\footnote{英文 information privacy 或 informational privacy，即資訊隱私。司法院釋字第 603 號解釋稱之為「資訊隱私權」。\url{https://plato.stanford.edu/entries/privacy/}。}。
% cue 822 1772.960 1773.980
\hypertarget{c01:cue:822}{}也叫 information。
% cue 823 1773.980 1774.980
\hypertarget{c01:cue:823}{}也叫 information。
% cue 824 1774.980 1775.970
\hypertarget{c01:cue:824}{}privacy。
% cue 825 1775.970 1776.960
\hypertarget{c01:cue:825}{}都可以。
% cue 826 1776.960 1778.770
\hypertarget{c01:cue:826}{}剛才有聽懂嗎。
% cue 827 1778.770 1779.770
\hypertarget{c01:cue:827}{}看 22 條。
% cue 828 1779.770 1780.770
\hypertarget{c01:cue:828}{}全公公的資料。
% cue 829 1780.770 1781.860
\hypertarget{c01:cue:829}{}就這麼大部分。
% cue 830 1781.860 1783.150
\hypertarget{c01:cue:830}{}還有其他的。
% cue 831 1783.150 1784.630
\hypertarget{c01:cue:831}{}多了。
% cue 832 1784.630 1786.100
\hypertarget{c01:cue:832}{}要不要講醫學期。

\tc{0:29:46}
% cue 833 1786.100 1786.200
\hypertarget{c01:cue:833}{}老師也開過。
% cue 834 1786.100 1787.100
\hypertarget{c01:cue:834}{}另外一個專題。
% cue 835 1787.100 1788.090
\hypertarget{c01:cue:835}{}言論自由\footnote{中華民國憲法第 11 條：「人民有言論、講學、著作及出版之自由。」司法院釋字第 509 號解釋（2000 年 7 月 7 日）認刑法第 310 條第 1 項及第 2 項的誹謗罪「為防止妨礙他人之自由權利所必要，符合憲法第二十三條規定之意旨」。\url{https://law.moj.gov.tw/LawClass/LawAll.aspx?pcode=A0000001}。}。
% cue 836 1788.090 1790.480
\hypertarget{c01:cue:836}{}言論自由很複雜很複雜很複雜。
% cue 837 1790.480 1791.480
\hypertarget{c01:cue:837}{}我就舉一個例子。
% cue 838 1791.480 1792.600
\hypertarget{c01:cue:838}{}你們同學。
% cue 839 1792.600 1793.780
\hypertarget{c01:cue:839}{}就是在。
% cue 840 1793.780 1794.780
\hypertarget{c01:cue:840}{}常常因為你們。
% cue 841 1794.780 1795.960
\hypertarget{c01:cue:841}{}社群裡面。
% cue 842 1795.960 1797.960
\hypertarget{c01:cue:842}{}沒有 social media。
% cue 843 1797.960 1799.060
\hypertarget{c01:cue:843}{}line 的社群裡面講的話。
% cue 844 1799.060 1801.220
\hypertarget{c01:cue:844}{}覺得它修物理。
% cue 845 1801.220 1802.210
\hypertarget{c01:cue:845}{}它在霸給你。
% cue 846 1802.210 1805.260
\hypertarget{c01:cue:846}{}那你在 line 的族群。
% cue 847 1805.260 1807.290
\hypertarget{c01:cue:847}{}或者是類似這種。
% cue 848 1807.290 1808.220
\hypertarget{c01:cue:848}{}social media 的族群。
% cue 849 1808.220 1809.480
\hypertarget{c01:cue:849}{}你有沒有言論自由。
% cue 850 1809.480 1810.830
\hypertarget{c01:cue:850}{}你的言論自由。
% cue 851 1810.830 1811.830
\hypertarget{c01:cue:851}{}可不可以講假的。
% cue 852 1811.830 1813.340
\hypertarget{c01:cue:852}{}某某某。
% cue 853 1813.340 1814.300
\hypertarget{c01:cue:853}{}跟某某某。
% cue 854 1814.300 1815.310
\hypertarget{c01:cue:854}{}在電話。
% cue 855 1815.310 1816.310
\hypertarget{c01:cue:855}{}他們昨天發生性關係。
% cue 856 1816.310 1818.010
\hypertarget{c01:cue:856}{}覺得你可不可以。
% cue 857 1818.010 1819.010
\hypertarget{c01:cue:857}{}發過這種訊息。
% cue 858 1819.010 1822.400
\hypertarget{c01:cue:858}{}我講話有聽懂嗎。
% cue 859 1822.400 1823.600
\hypertarget{c01:cue:859}{}言論自由有非常多的東西。
% cue 860 1823.600 1825.010
\hypertarget{c01:cue:860}{}可以考。
% cue 861 1825.010 1825.950
\hypertarget{c01:cue:861}{}這一講的也可以學習。
% cue 862 1825.950 1827.650
\hypertarget{c01:cue:862}{}我們這個學期。
% cue 863 1827.650 1828.900
\hypertarget{c01:cue:863}{}權力分立也在上面。

\tc{0:30:28}
% cue 864 1828.900 1830.260
\hypertarget{c01:cue:864}{}總之呢。
% cue 865 1830.260 1831.590
\hypertarget{c01:cue:865}{}你們有興趣。
% cue 866 1831.590 1833.870
\hypertarget{c01:cue:866}{}你們不要那麼被動。
% cue 867 1833.870 1834.810
\hypertarget{c01:cue:867}{}但有興趣。
% cue 868 1834.810 1837.220
\hypertarget{c01:cue:868}{}這個不是我們。
% cue 869 1837.220 1838.220
\hypertarget{c01:cue:869}{}我老師可以完全決定的。
% cue 870 1838.220 1839.320
\hypertarget{c01:cue:870}{}一個人。
% cue 871 1839.320 1840.320
\hypertarget{c01:cue:871}{}對某樣事情有興趣。
% cue 872 1840.320 1842.570
\hypertarget{c01:cue:872}{}大概有幾個原因。
% cue 873 1842.570 1843.570
\hypertarget{c01:cue:873}{}第一個是你個人。
% cue 874 1843.570 1844.520
\hypertarget{c01:cue:874}{}這是主要的。
% cue 875 1844.520 1845.660
\hypertarget{c01:cue:875}{}你的特質。
% cue 876 1845.660 1846.810
\hypertarget{c01:cue:876}{}喜歡什麼東西。
% cue 877 1846.810 1847.910
\hypertarget{c01:cue:877}{}不喜歡什麼東西。
% cue 878 1847.910 1850.090
\hypertarget{c01:cue:878}{}這個是你。
% cue 879 1850.090 1851.090
\hypertarget{c01:cue:879}{}大半是自己天生決定。
% cue 880 1851.090 1853.040
\hypertarget{c01:cue:880}{}後半是你後生的。
% cue 881 1853.040 1855.170
\hypertarget{c01:cue:881}{}後天教育影響環境決定。
% cue 882 1855.170 1856.460
\hypertarget{c01:cue:882}{}這個老師沒辦法。
% cue 883 1856.460 1857.920
\hypertarget{c01:cue:883}{}第二點。
% cue 884 1857.920 1859.000
\hypertarget{c01:cue:884}{}老師可以做的事情是。
% cue 885 1859.000 1861.190
\hypertarget{c01:cue:885}{}如果我的教學是成功的。
% cue 886 1861.190 1863.510
\hypertarget{c01:cue:886}{}我會讓你建立某些事情的興趣。
% cue 887 1863.510 1865.600
\hypertarget{c01:cue:887}{}至少是初步的興趣。
% cue 888 1865.600 1867.080
\hypertarget{c01:cue:888}{}因為我講的太有趣了。
% cue 889 1867.080 1868.210
\hypertarget{c01:cue:889}{}你覺得很好玩。
% cue 890 1868.210 1870.530
\hypertarget{c01:cue:890}{}那有的老師也許很多學問。
% cue 891 1870.530 1871.530
\hypertarget{c01:cue:891}{}但他講不好。
% cue 892 1871.530 1872.530
\hypertarget{c01:cue:892}{}他就不會給你興趣。
% cue 893 1872.530 1874.820
\hypertarget{c01:cue:893}{}如果老師不能帶領。
% cue 894 1874.820 1876.820
\hypertarget{c01:cue:894}{}學生進入他所受校的科目。
% cue 895 1876.820 1879.000
\hypertarget{c01:cue:895}{}那這個老師當然是。
% cue 896 1879.000 1880.000
\hypertarget{c01:cue:896}{}會失敗。

\tc{0:31:20}
% cue 897 1880.000 1882.820
\hypertarget{c01:cue:897}{}那第三個。
% cue 898 1882.820 1884.850
\hypertarget{c01:cue:898}{}就是先遲到你將來的想法。
% cue 899 1884.850 1885.940
\hypertarget{c01:cue:899}{}你將來想說。
% cue 900 1885.940 1886.940
\hypertarget{c01:cue:900}{}是跟這個有關的工作。
% cue 901 1886.940 1888.940
\hypertarget{c01:cue:901}{}你可能會特別有興趣。
% cue 902 1888.940 1891.550
\hypertarget{c01:cue:902}{}比如說你們當中有人想賺錢。
% cue 903 1891.550 1892.760
\hypertarget{c01:cue:903}{}可能對商法。
% cue 904 1892.760 1893.730
\hypertarget{c01:cue:904}{}對民法。
% cue 905 1893.730 1897.250
\hypertarget{c01:cue:905}{}思想也要賺錢。
% cue 906 1897.250 1899.250
\hypertarget{c01:cue:906}{}所以這對你們比較重要。
% cue 907 1899.250 1901.420
\hypertarget{c01:cue:907}{}讓我們知道。

\section{分組與報告}

\tc{0:31:41}
% cue 908 1901.420 1902.390
\hypertarget{c01:cue:908}{}好啦。
% cue 909 1902.390 1903.390
\hypertarget{c01:cue:909}{}先讓一段講到。
% cue 910 1903.390 1904.390
\hypertarget{c01:cue:910}{}教學的目標。
% cue 911 1904.390 1905.390
\hypertarget{c01:cue:911}{}我就是要講這個。
% cue 912 1905.390 1906.900
\hypertarget{c01:cue:912}{}到開始的時候。
% cue 913 1906.900 1907.810
\hypertarget{c01:cue:913}{}老師講一講。
% cue 914 1907.810 1909.620
\hypertarget{c01:cue:914}{}最後呢。
% cue 915 1909.620 1910.740
\hypertarget{c01:cue:915}{}你們就要開始做報告。
% cue 916 1910.740 1912.860
\hypertarget{c01:cue:916}{}那我想你們比較在意的事情。
% cue 917 1912.860 1914.330
\hypertarget{c01:cue:917}{}是做什麼事情。
% cue 918 1914.330 1915.580
\hypertarget{c01:cue:918}{}第一。
% cue 919 1915.580 1917.120
\hypertarget{c01:cue:919}{}我們下面有分組。
% cue 920 1917.120 1918.120
\hypertarget{c01:cue:920}{}每個禮拜要上什麼。
% cue 921 1918.120 1919.690
\hypertarget{c01:cue:921}{}哪些案子。
% cue 922 1919.690 1920.850
\hypertarget{c01:cue:922}{}分到你等一下就要分。
% cue 923 1920.850 1921.850
\hypertarget{c01:cue:923}{}那個是班長。
% cue 924 1921.850 1923.850
\hypertarget{c01:cue:924}{}要跟大家做一起做的事情。

\tc{0:32:03}
% cue 925 1923.850 1925.170
\hypertarget{c01:cue:925}{}他分好。
% cue 926 1925.170 1926.390
\hypertarget{c01:cue:926}{}分組的時候。
% cue 927 1926.390 1927.350
\hypertarget{c01:cue:927}{}聽清楚。
% cue 928 1927.350 1928.600
\hypertarget{c01:cue:928}{}你可以一個人一組。
% cue 929 1928.600 1930.720
\hypertarget{c01:cue:929}{}因為我適應個人化的時代。
% cue 930 1930.720 1932.720
\hypertarget{c01:cue:930}{}我就是一個原子。
% cue 931 1932.720 1933.750
\hypertarget{c01:cue:931}{}不喜歡跟其他人結合。
% cue 932 1933.750 1935.030
\hypertarget{c01:cue:932}{}我這個人。
% cue 933 1935.030 1936.030
\hypertarget{c01:cue:933}{}這個組。
% cue 934 1936.030 1937.030
\hypertarget{c01:cue:934}{}這個案子不認。
% cue 935 1937.030 1938.890
\hypertarget{c01:cue:935}{}我自己來報告。
% cue 936 1938.890 1939.890
\hypertarget{c01:cue:936}{}現在三個和尚教書。
% cue 937 1939.890 1941.400
\hypertarget{c01:cue:937}{}一定能夠說。
% cue 938 1941.400 1942.400
\hypertarget{c01:cue:938}{}兩個和尚教書。
% cue 939 1942.400 1943.330
\hypertarget{c01:cue:939}{}通常也是沒說的。
% cue 940 1943.330 1945.320
\hypertarget{c01:cue:940}{}所以與其跟人家一組。
% cue 941 1945.320 1946.320
\hypertarget{c01:cue:941}{}到說草門都花入山。
% cue 942 1946.320 1947.320
\hypertarget{c01:cue:942}{}還能自己一個人一組。
% cue 943 1947.320 1948.320
\hypertarget{c01:cue:943}{}也可以。
% cue 944 1948.320 1949.440
\hypertarget{c01:cue:944}{}但是我們班上。
% cue 945 1949.440 1950.440
\hypertarget{c01:cue:945}{}沒有那麼多組。
% cue 946 1950.440 1951.560
\hypertarget{c01:cue:946}{}你們四十幾個人。
% cue 947 1951.560 1952.560
\hypertarget{c01:cue:947}{}我拿來幹嘛。
% cue 948 1952.560 1953.560
\hypertarget{c01:cue:948}{}再一次一個人一組。
% cue 949 1953.560 1954.560
\hypertarget{c01:cue:949}{}那上四十幾個禮拜。
% cue 950 1954.560 1955.530
\hypertarget{c01:cue:950}{}你不太可能。
% cue 951 1955.530 1957.690
\hypertarget{c01:cue:951}{}所以一定要分組。
% cue 952 1957.690 1958.690
\hypertarget{c01:cue:952}{}分組就是確保。
% cue 953 1958.690 1960.100
\hypertarget{c01:cue:953}{}到那個時候。
% cue 954 1960.100 1961.100
\hypertarget{c01:cue:954}{}照我們下的那個進度。
% cue 955 1961.100 1962.010
\hypertarget{c01:cue:955}{}那個案子。
% cue 956 1962.010 1963.160
\hypertarget{c01:cue:956}{}有人報告。
% cue 957 1963.160 1965.480
\hypertarget{c01:cue:957}{}我第一眼。

\tc{0:32:45}
% cue 958 1965.480 1966.800
\hypertarget{c01:cue:958}{}第二點報告的格式。
% cue 959 1966.800 1967.800
\hypertarget{c01:cue:959}{}老師會在現場。
% cue 960 1967.800 1968.860
\hypertarget{c01:cue:960}{}可以一道理。
% cue 961 1968.860 1970.240
\hypertarget{c01:cue:961}{}我會請問。
% cue 962 1970.240 1971.370
\hypertarget{c01:cue:962}{}就要回答。
% cue 963 1971.370 1972.330
\hypertarget{c01:cue:963}{}報告的小組。
% cue 964 1972.330 1973.280
\hypertarget{c01:cue:964}{}要來回答。
% cue 965 1973.280 1973.380
\hypertarget{c01:cue:965}{}好。
% cue 966 1973.280 1974.740
\hypertarget{c01:cue:966}{}然後報告小組。
% cue 967 1974.740 1976.740
\hypertarget{c01:cue:967}{}也可以跟同學提問。
% cue 968 1976.740 1977.740
\hypertarget{c01:cue:968}{}跟老師提問。
% cue 969 1977.740 1979.740
\hypertarget{c01:cue:969}{}也可以提出一些有趣的問題。
% cue 970 1979.740 1980.740
\hypertarget{c01:cue:970}{}大家來討論。

\tc{0:33:00}
% cue 971 1980.740 1983.150
\hypertarget{c01:cue:971}{}我會盡量採用。
% cue 972 1983.150 1985.150
\hypertarget{c01:cue:972}{}所謂蘇格拉底的教學法\footnote{英文 Socratic method，美國法學院常用的課堂教學方式，教師以連續提問與被點名的少數學生討論，讓學生自己推出法律概念與原則。芝加哥大學法學院以此培養學生的批判思考能力。\url{https://www.law.uchicago.edu/socratic-method}。}。
% cue 973 1985.150 1986.780
\hypertarget{c01:cue:973}{}但是在臺灣。
% cue 974 1986.780 1987.780
\hypertarget{c01:cue:974}{}這個不有行。
% cue 975 1987.780 1989.070
\hypertarget{c01:cue:975}{}因為有習慣聽老師講。
% cue 976 1989.070 1991.700
\hypertarget{c01:cue:976}{}但我不知道要怎麼辦。
% cue 977 1991.700 1993.020
\hypertarget{c01:cue:977}{}我只能說。
% cue 978 1993.020 1994.080
\hypertarget{c01:cue:978}{}聽老師講。
% cue 979 1994.080 1995.080
\hypertarget{c01:cue:979}{}不是不好。
% cue 980 1995.080 1997.200
\hypertarget{c01:cue:980}{}老師一定是認為。
% cue 981 1997.200 1998.200
\hypertarget{c01:cue:981}{}比較重要。
% cue 982 1998.200 1999.200
\hypertarget{c01:cue:982}{}所以他才會講給你聽。
% cue 983 1999.200 2000.350
\hypertarget{c01:cue:983}{}那其他他就不講。
% cue 984 2000.350 2001.280
\hypertarget{c01:cue:984}{}可是。
% cue 985 2001.280 2003.630
\hypertarget{c01:cue:985}{}這是老師認為重要的。
% cue 986 2003.630 2004.910
\hypertarget{c01:cue:986}{}老師的重要。
% cue 987 2004.910 2005.910
\hypertarget{c01:cue:987}{}不等於你的重要。
% cue 988 2005.910 2008.130
\hypertarget{c01:cue:988}{}那有可能你沒有聽到你要。
% cue 989 2008.130 2009.130
\hypertarget{c01:cue:989}{}這時候怎麼辦呢。
% cue 990 2009.130 2010.510
\hypertarget{c01:cue:990}{}你要提問。
% cue 991 2010.510 2012.540
\hypertarget{c01:cue:991}{}那如果你要很害羞。
% cue 992 2012.540 2013.540
\hypertarget{c01:cue:992}{}又不願意提問。
% cue 993 2013.540 2015.540
\hypertarget{c01:cue:993}{}然後又不願意準備。
% cue 994 2015.540 2016.540
\hypertarget{c01:cue:994}{}又不願意思考。
% cue 995 2016.540 2017.510
\hypertarget{c01:cue:995}{}那你就是。
% cue 996 2017.510 2018.800
\hypertarget{c01:cue:996}{}最多老師講 100 個。
% cue 997 2018.800 2019.800
\hypertarget{c01:cue:997}{}你可能跟思考。
% cue 998 2019.800 2022.340
\hypertarget{c01:cue:998}{}因為你能 catch up 就是很多。
% cue 999 2022.340 2023.340
\hypertarget{c01:cue:999}{}好。

\section{共筆與影音資源}

\tc{0:33:43}
% cue 1000 2023.340 2026.000
\hypertarget{c01:cue:1000}{}那如果你不願意做共筆\footnote{課綱教學目標規定，選課同學「應輪流製作授課共筆，以備期末之開書測驗（Open Book Exam.）」。}是最好。
% cue 1001 2026.000 2028.090
\hypertarget{c01:cue:1001}{}我也不知道從哪一頁開始。
% cue 1002 2028.090 2030.120
\hypertarget{c01:cue:1002}{}以前老師的共筆在老師的出租。
% cue 1003 2030.120 2032.120
\hypertarget{c01:cue:1003}{}以前共筆很快慢。
% cue 1004 2032.120 2034.220
\hypertarget{c01:cue:1004}{}對面那個贏店什麼都要賣。
% cue 1005 2034.220 2036.600
\hypertarget{c01:cue:1005}{}後來老師看到同學太辛苦了。
% cue 1006 2036.600 2037.550
\hypertarget{c01:cue:1006}{}我就說。
% cue 1007 2037.550 2038.560
\hypertarget{c01:cue:1007}{}我就說那無音。
% cue 1008 2038.560 2039.530
\hypertarget{c01:cue:1008}{}所以老師有一段時間。
% cue 1009 2039.530 2042.040
\hypertarget{c01:cue:1009}{}是有對話線。
% cue 1010 2042.040 2045.830
\hypertarget{c01:cue:1010}{}把憲法對話的有聲書\footnote{湯德宗著《對話憲法‧憲法對話：湯德宗憲法有聲書》，天宏出版，上冊（憲法結構）增訂三版 2015 年，下冊（基本人權）增訂三版 2016 年。}。
% cue 1011 2045.830 2046.920
\hypertarget{c01:cue:1011}{}或是有 CD。
% cue 1012 2046.920 2050.300
\hypertarget{c01:cue:1012}{}後來老師就是有 YouTube。
% cue 1013 2050.300 2051.610
\hypertarget{c01:cue:1013}{}這個時代進步了。
% cue 1014 2051.610 2054.570
\hypertarget{c01:cue:1014}{}老師就把單純的語音變成影音了。
% cue 1015 2054.570 2056.810
\hypertarget{c01:cue:1015}{}所以現在網路上演。
% cue 1016 2056.810 2059.120
\hypertarget{c01:cue:1016}{}你們可以在網路平臺。
% cue 1017 2059.120 2061.920
\hypertarget{c01:cue:1017}{}看得到老師在東吳大學。
% cue 1018 2061.920 2064.950
\hypertarget{c01:cue:1018}{}在臺大好幾年的憲法課。
% cue 1019 2064.950 2066.050
\hypertarget{c01:cue:1019}{}從第一講講到最後講。
% cue 1020 2066.050 2068.050
\hypertarget{c01:cue:1020}{}中間有一些碰到的 COVID。
% cue 1021 2068.050 2069.050
\hypertarget{c01:cue:1021}{}有點受影響。
% cue 1022 2069.050 2070.050
\hypertarget{c01:cue:1022}{}在基本上。
% cue 1023 2070.050 2072.270
\hypertarget{c01:cue:1023}{}這講的今年如果沒有。
% cue 1024 2072.270 2074.270
\hypertarget{c01:cue:1024}{}下一年也許不理到。
% cue 1025 2074.270 2075.330
\hypertarget{c01:cue:1025}{}反正你有興趣。
% cue 1026 2075.330 2075.430
\hypertarget{c01:cue:1026}{}你看。

\tc{0:34:35}
% cue 1027 2075.330 2077.460
\hypertarget{c01:cue:1027}{}往來臺大有個開放式課程。
% cue 1028 2077.460 2079.490
\hypertarget{c01:cue:1028}{}但是開放課程。
% cue 1029 2079.490 2080.410
\hypertarget{c01:cue:1029}{}客座比較舊。
% cue 1030 2080.410 2081.460
\hypertarget{c01:cue:1030}{}那是 2016 年。
% cue 1031 2081.460 2082.460
\hypertarget{c01:cue:1031}{}現在 2020 年。
% cue 1032 2082.460 2083.460
\hypertarget{c01:cue:1032}{}有十分進入。
% cue 1033 2083.460 2085.740
\hypertarget{c01:cue:1033}{}大家那時候就有這種想法。
% cue 1034 2085.740 2087.250
\hypertarget{c01:cue:1034}{}他覺得有些客。
% cue 1035 2087.250 2089.270
\hypertarget{c01:cue:1035}{}老師老了會死掉。
% cue 1036 2089.270 2091.490
\hypertarget{c01:cue:1036}{}他把這個經典留起來。
% cue 1037 2091.490 2093.260
\hypertarget{c01:cue:1037}{}他只能拍攝有些。
% cue 1038 2093.260 2094.250
\hypertarget{c01:cue:1038}{}解憶。
% cue 1039 2094.250 2096.750
\hypertarget{c01:cue:1039}{}反正你們想要看。
% cue 1040 2096.750 2097.940
\hypertarget{c01:cue:1040}{}不要跟老師講。
% cue 1041 2097.940 2099.000
\hypertarget{c01:cue:1041}{}老師。
% cue 1042 2099.000 2100.000
\hypertarget{c01:cue:1042}{}我上去。
% cue 1043 2100.000 2101.000
\hypertarget{c01:cue:1043}{}我去年。
% cue 1044 2101.000 2102.000
\hypertarget{c01:cue:1044}{}我老師只有講三小時。
% cue 1045 2102.000 2102.960
\hypertarget{c01:cue:1045}{}你能期待有什麼。
% cue 1046 2102.960 2106.010
\hypertarget{c01:cue:1046}{}這樣子講的有點太。
% cue 1047 2106.010 2107.330
\hypertarget{c01:cue:1047}{}太不像現代人。
% cue 1048 2107.330 2109.480
\hypertarget{c01:cue:1048}{}因為我在現在這個時代。
% cue 1049 2109.480 2110.480
\hypertarget{c01:cue:1049}{}你可以做的事情很多。
% cue 1050 2110.480 2113.150
\hypertarget{c01:cue:1050}{}你可以在網路上。
% cue 1051 2113.150 2114.150
\hypertarget{c01:cue:1051}{}尤其像老師的頻道。
% cue 1052 2114.150 2116.150
\hypertarget{c01:cue:1052}{}你可能很快就找到。
% cue 1053 2116.150 2117.150
\hypertarget{c01:cue:1053}{}能夠讓你感到滿意的東西。
% cue 1054 2117.150 2119.270
\hypertarget{c01:cue:1054}{}就可以自己看。

\section{期末考試與指定教科書}

\tc{0:35:19}
% cue 1055 2119.270 2121.680
\hypertarget{c01:cue:1055}{}我回到剛才要講的事情。
% cue 1056 2121.680 2124.160
\hypertarget{c01:cue:1056}{}那你們最擔心的事情。
% cue 1057 2124.160 2125.160
\hypertarget{c01:cue:1057}{}應該是考試吧。
% cue 1058 2125.160 2129.160
\hypertarget{c01:cue:1058}{}考試呢。
% cue 1059 2129.160 2130.160
\hypertarget{c01:cue:1059}{}事情是這樣子的。
% cue 1060 2130.160 2132.660
\hypertarget{c01:cue:1060}{}如果沒有意外。
% cue 1061 2132.660 2133.660
\hypertarget{c01:cue:1061}{}我希望你們沒有意外。
% cue 1062 2133.660 2135.170
\hypertarget{c01:cue:1062}{}那老師說。
% cue 1063 2135.170 2136.170
\hypertarget{c01:cue:1063}{}在期末考試前。
% cue 1064 2136.170 2137.270
\hypertarget{c01:cue:1064}{}一個月。
% cue 1065 2137.270 2138.390
\hypertarget{c01:cue:1065}{}陸陸續續的。
% cue 1066 2138.390 2140.420
\hypertarget{c01:cue:1066}{}我就會出期末給你們。
% cue 1067 2140.420 2143.090
\hypertarget{c01:cue:1067}{}那通常我都是出申論題。
% cue 1068 2143.090 2144.510
\hypertarget{c01:cue:1068}{}有的時候我母題。
% cue 1069 2144.510 2145.540
\hypertarget{c01:cue:1069}{}有的時候比較多一點。
% cue 1070 2145.540 2146.540
\hypertarget{c01:cue:1070}{}但是我希望不要太多。
% cue 1071 2146.540 2149.150
\hypertarget{c01:cue:1071}{}那考試的時候呢。
% cue 1072 2149.150 2150.150
\hypertarget{c01:cue:1072}{}是 Open Book。
% cue 1073 2150.150 2151.460
\hypertarget{c01:cue:1073}{}所以。
% cue 1074 2151.460 2152.460
\hypertarget{c01:cue:1074}{}權力分立新論的。
% cue 1075 2152.460 2153.460
\hypertarget{c01:cue:1075}{}卷一卷二。
% cue 1076 2153.460 2154.460
\hypertarget{c01:cue:1076}{}我強烈的建議品要買。
% cue 1077 2154.460 2155.390
\hypertarget{c01:cue:1077}{}因為。
% cue 1078 2155.390 2157.600
\hypertarget{c01:cue:1078}{}因為你聽懂得很好。
% cue 1079 2157.600 2158.940
\hypertarget{c01:cue:1079}{}但是寫不出來。
% cue 1080 2158.940 2161.900
\hypertarget{c01:cue:1080}{}那些字裡不會。
% cue 1081 2161.900 2162.900
\hypertarget{c01:cue:1081}{}那些用語也沒有。
% cue 1082 2162.900 2165.280
\hypertarget{c01:cue:1082}{}那個文字要這樣組合。
% cue 1083 2165.280 2167.370
\hypertarget{c01:cue:1083}{}你要從書上去抄。
% cue 1084 2167.370 2168.630
\hypertarget{c01:cue:1084}{}但是你至少有東西可抄。
% cue 1085 2168.630 2170.110
\hypertarget{c01:cue:1085}{}OK。

\tc{0:36:10}
% cue 1086 2170.110 2172.110
\hypertarget{c01:cue:1086}{}所以我會跟你們講。
% cue 1087 2172.110 2173.110
\hypertarget{c01:cue:1087}{}要去讀拿幾遍。
% cue 1088 2173.110 2174.110
\hypertarget{c01:cue:1088}{}拿幾遍。
% cue 1089 2174.110 2175.200
\hypertarget{c01:cue:1089}{}那有問題就要問。
% cue 1090 2175.200 2177.200
\hypertarget{c01:cue:1090}{}這個事情是有很多種。
% cue 1091 2177.200 2178.200
\hypertarget{c01:cue:1091}{}教學的方法。
% cue 1092 2178.200 2179.290
\hypertarget{c01:cue:1092}{}最好在理。
% cue 1093 2179.290 2181.350
\hypertarget{c01:cue:1093}{}老師會跟你們講說。
% cue 1094 2181.350 2182.410
\hypertarget{c01:cue:1094}{}這三本書。
% cue 1095 2182.410 2184.670
\hypertarget{c01:cue:1095}{}是我們這個學期。
% cue 1096 2184.670 2186.620
\hypertarget{c01:cue:1096}{}這個課指定教科書。
% cue 1097 2186.620 2186.720
\hypertarget{c01:cue:1097}{}好。

\tc{0:36:26}
% cue 1098 2186.620 2187.990
\hypertarget{c01:cue:1098}{}那剩下呢。
% cue 1099 2187.990 2188.990
\hypertarget{c01:cue:1099}{}他有可能說。
% cue 1100 2188.990 2189.990
\hypertarget{c01:cue:1100}{}第一個禮拜。
% cue 1101 2189.990 2190.990
\hypertarget{c01:cue:1101}{}請你看第一本書。
% cue 1102 2190.990 2191.990
\hypertarget{c01:cue:1102}{}30 頁到 80 頁。
% cue 1103 2191.990 2194.430
\hypertarget{c01:cue:1103}{}老師上課的時候。
% cue 1104 2194.430 2195.650
\hypertarget{c01:cue:1104}{}就是講兩三分鐘。
% cue 1105 2195.650 2197.770
\hypertarget{c01:cue:1105}{}或者是 20 分鐘。
% cue 1106 2197.770 2198.770
\hypertarget{c01:cue:1106}{}剩下就是討論。
% cue 1107 2198.770 2199.860
\hypertarget{c01:cue:1107}{}他會提一些問題來。
% cue 1108 2199.860 2200.980
\hypertarget{c01:cue:1108}{}當同學就來。
% cue 1109 2200.980 2201.980
\hypertarget{c01:cue:1109}{}對方式討論。
% cue 1110 2201.980 2203.230
\hypertarget{c01:cue:1110}{}開始是兩個人。
% cue 1111 2203.230 2204.230
\hypertarget{c01:cue:1111}{}三個人。
% cue 1112 2204.230 2204.330
\hypertarget{c01:cue:1112}{}五個人。
% cue 1113 2204.230 2205.350
\hypertarget{c01:cue:1113}{}然後後來。
% cue 1114 2205.350 2206.350
\hypertarget{c01:cue:1114}{}越來越少。
% cue 1115 2206.350 2207.350
\hypertarget{c01:cue:1115}{}越來越少。
% cue 1116 2207.350 2209.980
\hypertarget{c01:cue:1116}{}像很難對話的樣。
% cue 1117 2209.980 2211.330
\hypertarget{c01:cue:1117}{}這是上課的好處。
% cue 1118 2211.330 2212.330
\hypertarget{c01:cue:1118}{}但是也是上課的確定。
% cue 1119 2212.330 2215.120
\hypertarget{c01:cue:1119}{}這就是一個教學的方式。
% cue 1120 2215.120 2217.750
\hypertarget{c01:cue:1120}{}最後是期末開書考試。
% cue 1121 2217.750 2219.130
\hypertarget{c01:cue:1121}{}這樣可以了吧。

\section{分週進度與半總統制}

\tc{0:36:59}
% cue 1122 2219.130 2221.190
\hypertarget{c01:cue:1122}{}那我們現在看一下。
% cue 1123 2221.190 2223.190
\hypertarget{c01:cue:1123}{}請這個禮拜的結構。
% cue 1124 2223.190 2224.920
\hypertarget{c01:cue:1124}{}各種安排。
% cue 1125 2224.920 2226.020
\hypertarget{c01:cue:1125}{}今天第一節課。
% cue 1126 2226.020 2229.400
\hypertarget{c01:cue:1126}{}等一下會給你們討論解例跟 YouTube。
% cue 1127 2229.400 2230.940
\hypertarget{c01:cue:1127}{}921。
% cue 1128 2230.940 2231.940
\hypertarget{c01:cue:1128}{}我們上一次課。
% cue 1129 2231.940 2233.320
\hypertarget{c01:cue:1129}{}928 一的方向。
% cue 1130 2233.320 2234.380
\hypertarget{c01:cue:1130}{}921 就是這樣。
% cue 1131 2234.380 2236.020
\hypertarget{c01:cue:1131}{}比較憲法結構\footnote{課綱第 2 週（9 月 21 日）的主題，分總統制、內閣制與半總統制三類，指定閱讀《權力分立新論 卷一》第 1 頁起。第 3 週（9 月 28 日）為教師節放假。}。
% cue 1132 2236.020 2239.470
\hypertarget{c01:cue:1132}{}或者你叫做比較政府構造都可以。
% cue 1133 2239.470 2241.170
\hypertarget{c01:cue:1133}{}總統制。
% cue 1134 2241.170 2242.260
\hypertarget{c01:cue:1134}{}內閣制。
% cue 1135 2242.260 2243.260
\hypertarget{c01:cue:1135}{}跟半總統制\footnote{英文 semi-presidential system。Maurice Duverger 於〈A New Political System Model: Semi-Presidential Government〉（European Journal of Political Research 第 8 卷第 2 期，1980，頁 165–187）以三項要素界定：總統由普選產生、總統享有相當權力、另有對國會負責而掌行政權的總理與內閣。湯德宗稱 1997 年修憲後的現制為「弱勢總統制」，並曾以「貌似半總統制的總統制」稱之（《權力分立新論 卷一》增訂三版，頁 68 註 22）。\url{https://doi.org/10.1111/j.1475-6765.1980.tb00569.x}。}。
% cue 1136 2243.260 2245.000
\hypertarget{c01:cue:1136}{}為什麼放在第一角。
% cue 1137 2245.000 2246.250
\hypertarget{c01:cue:1137}{}因為。
% cue 1138 2246.250 2248.660
\hypertarget{c01:cue:1138}{}這是我們權力分配。
% cue 1139 2248.660 2249.810
\hypertarget{c01:cue:1139}{}比較看得見。
% cue 1140 2249.810 2250.810
\hypertarget{c01:cue:1140}{}不能找的東西。
% cue 1141 2250.810 2252.200
\hypertarget{c01:cue:1141}{}而且。
% cue 1142 2252.200 2253.320
\hypertarget{c01:cue:1142}{}臺灣發生的問題。
% cue 1143 2253.320 2254.290
\hypertarget{c01:cue:1143}{}都在這裡。
% cue 1144 2254.290 2255.950
\hypertarget{c01:cue:1144}{}你認為。
% cue 1145 2255.950 2256.950
\hypertarget{c01:cue:1145}{}其實那日常可以不服屬嗎。
% cue 1146 2256.950 2258.550
\hypertarget{c01:cue:1146}{}那可能要問。
% cue 1147 2258.550 2259.510
\hypertarget{c01:cue:1147}{}是總統制。
% cue 1148 2259.510 2260.510
\hypertarget{c01:cue:1148}{}內閣制。
% cue 1149 2260.510 2261.510
\hypertarget{c01:cue:1149}{}還是半總統制。

\tc{0:37:41}
% cue 1150 2261.510 2262.630
\hypertarget{c01:cue:1150}{}這兩天。
% cue 1151 2262.630 2264.630
\hypertarget{c01:cue:1151}{}許信良\footnote{1941 年生，1991 年與 1996 年兩度當選民主進步黨主席，1997 年與李登輝總統合作推動第四次修憲，自稱以法國第五共和的雙首長制為藍本。2026 年 9 月 12 日在台北政經學院基金會「總統直選 30 年的回顧與前瞻」研討會上表示，行政院院長以副署權對抗國會決議不合憲政原理，建議朝野仿效 1990 年代國是會議與修憲時的合作經驗。\url{https://udn.com/news/story/6656/9780854}。}又跑出來做妖怪。
% cue 1152 2264.630 2265.880
\hypertarget{c01:cue:1152}{}他要來講。
% cue 1153 2265.880 2266.880
\hypertarget{c01:cue:1153}{}我們是半總統制。
% cue 1154 2266.880 2267.880
\hypertarget{c01:cue:1154}{}他講了一大堆。
% cue 1155 2267.880 2270.540
\hypertarget{c01:cue:1155}{}我們先不管他個人的歷程。
% cue 1156 2270.540 2271.830
\hypertarget{c01:cue:1156}{}他心儀。
% cue 1157 2271.830 2273.050
\hypertarget{c01:cue:1157}{}他崇拜。
% cue 1158 2273.050 2275.330
\hypertarget{c01:cue:1158}{}法國的半總統制。
% cue 1159 2275.330 2276.330
\hypertarget{c01:cue:1159}{}當年就是他。
% cue 1160 2276.330 2278.390
\hypertarget{c01:cue:1160}{}他那個死掉的李登輝\footnote{李登輝（1923–2020），1988 年 1 月蔣經國總統逝世後繼任第 7 任總統，1990 年經國民大會選舉為第 8 任總統，1996 年首次總統直接選舉當選第 9 任總統，2000 年 5 月 20 日卸任。\url{https://www.president.gov.tw/Page/86}。}兩個人。
% cue 1161 2278.390 2280.390
\hypertarget{c01:cue:1161}{}聯手修憲\footnote{1997 年 7 月 21 日總統公布第四次修正之憲法增修條文。該次修正的第 3 條第 1 項規定行政院院長由總統任命之，憲法第 55 條停止適用；同條第 2 項第 3 款增設立法院對行政院院長之不信任案，第 2 條第 5 項規定總統於不信任案通過後十日內，經諮詢立法院院長後，得宣告解散立法院。\cite{pres325}。}。
% cue 1162 2280.390 2281.350
\hypertarget{c01:cue:1162}{}修成今天這種模樣。
% cue 1163 2281.350 2283.600
\hypertarget{c01:cue:1163}{}你還好意思跑出來告訴我說。
% cue 1164 2283.600 2285.600
\hypertarget{c01:cue:1164}{}我們現在是半總統制。
% cue 1165 2285.600 2286.940
\hypertarget{c01:cue:1165}{}好。
% cue 1166 2286.940 2287.950
\hypertarget{c01:cue:1166}{}那我問你。
% cue 1167 2287.950 2289.070
\hypertarget{c01:cue:1167}{}你馬上就死。
% cue 1168 2289.070 2290.070
\hypertarget{c01:cue:1168}{}那半總統制。
% cue 1169 2290.070 2292.070
\hypertarget{c01:cue:1169}{}允許這種情況發生嗎。
% cue 1170 2292.070 2293.860
\hypertarget{c01:cue:1170}{}你還覺得是半總統制。

\tc{0:38:13}
% cue 1171 2293.860 2295.860
\hypertarget{c01:cue:1171}{}半總統制有兩個要件。
% cue 1172 2295.860 2296.860
\hypertarget{c01:cue:1172}{}簡單的講。
% cue 1173 2296.860 2298.020
\hypertarget{c01:cue:1173}{}一 總統民選。
% cue 1174 2298.020 2298.960
\hypertarget{c01:cue:1174}{}有自己的權力。
% cue 1175 2298.960 2299.980
\hypertarget{c01:cue:1175}{}不是全部。
% cue 1176 2299.980 2300.980
\hypertarget{c01:cue:1176}{}但是有一部分。
% cue 1177 2300.980 2302.130
\hypertarget{c01:cue:1177}{}第二。
% cue 1178 2302.130 2304.130
\hypertarget{c01:cue:1178}{}行政院長會名叫總理\footnote{法國第五共和憲法第 8 條第 1 項規定：「Le Président de la République nomme le Premier ministre.」（共和國總統任命總理。）總理於政府提出總辭時由總統免職。\cite{const1958}。}。
% cue 1179 2304.130 2305.180
\hypertarget{c01:cue:1179}{}向立法院。
% cue 1180 2305.180 2306.730
\hypertarget{c01:cue:1180}{}向國會負責。
% cue 1181 2306.730 2310.000
\hypertarget{c01:cue:1181}{}現在是國會通過的法律。
% cue 1182 2310.000 2312.380
\hypertarget{c01:cue:1182}{}你向我負責的行政院長。
% cue 1183 2312.380 2314.470
\hypertarget{c01:cue:1183}{}你給誰借了膽子。
% cue 1184 2314.470 2315.620
\hypertarget{c01:cue:1184}{}你幫我沒收啊。
% cue 1185 2315.620 2317.320
\hypertarget{c01:cue:1185}{}啊。
% cue 1186 2317.320 2318.910
\hypertarget{c01:cue:1186}{}告訴我。
% cue 1187 2318.910 2319.890
\hypertarget{c01:cue:1187}{}如果你向我負責的。
% cue 1188 2319.890 2320.890
\hypertarget{c01:cue:1188}{}我會饒過你嗎。
% cue 1189 2320.890 2323.550
\hypertarget{c01:cue:1189}{}我這樣講你就聽懂啊。
% cue 1190 2323.550 2325.550
\hypertarget{c01:cue:1190}{}我們什麼是半總統制。
% cue 1191 2325.550 2326.830
\hypertarget{c01:cue:1191}{}如果是半總統制。
% cue 1192 2326.830 2327.830
\hypertarget{c01:cue:1192}{}會有這種情況發生。
% cue 1193 2327.830 2329.830
\hypertarget{c01:cue:1193}{}這叫死人半總統制。
% cue 1194 2329.830 2330.850
\hypertarget{c01:cue:1194}{}沒有什麼可能 work。
% cue 1195 2330.850 2333.580
\hypertarget{c01:cue:1195}{}你還要聽懂嗎。

\tc{0:38:53}
% cue 1196 2333.580 2335.700
\hypertarget{c01:cue:1196}{}半總統制還有另外一個特徵。
% cue 1197 2335.700 2337.660
\hypertarget{c01:cue:1197}{}就剛才講的兩半。
% cue 1198 2337.660 2338.660
\hypertarget{c01:cue:1198}{}還有這兩個中間的互動。
% cue 1199 2338.660 2340.880
\hypertarget{c01:cue:1199}{}行政不是一個統。
% cue 1200 2340.880 2341.880
\hypertarget{c01:cue:1200}{}是兩個統。
% cue 1201 2341.880 2342.840
\hypertarget{c01:cue:1201}{}一個總統。
% cue 1202 2342.840 2344.260
\hypertarget{c01:cue:1202}{}一個叫總理。
% cue 1203 2344.260 2345.260
\hypertarget{c01:cue:1203}{}這裡面就會發生。
% cue 1204 2345.260 2346.260
\hypertarget{c01:cue:1204}{}一種奇怪的關係。
% cue 1205 2346.260 2347.410
\hypertarget{c01:cue:1205}{}如果總統跟總理。
% cue 1206 2347.410 2348.410
\hypertarget{c01:cue:1206}{}同屬一個黨。
% cue 1207 2348.410 2349.910
\hypertarget{c01:cue:1207}{}還要是蔡英文\footnote{2016 年 5 月 20 日至 2024 年 5 月 20 日任中華民國總統。}。
% cue 1208 2349.910 2350.910
\hypertarget{c01:cue:1208}{}跟他的行政院長。
% cue 1209 2350.910 2351.910
\hypertarget{c01:cue:1209}{}都是民進黨。
% cue 1210 2351.910 2354.360
\hypertarget{c01:cue:1210}{}這是一種情況。

\tc{0:39:14}
% cue 1211 2354.360 2355.450
\hypertarget{c01:cue:1211}{}第二種情況。
% cue 1212 2355.450 2356.790
\hypertarget{c01:cue:1212}{}總統是一個黨的。
% cue 1213 2356.790 2358.970
\hypertarget{c01:cue:1213}{}那行政院長是另外一個黨。
% cue 1214 2358.970 2360.540
\hypertarget{c01:cue:1214}{}那請問。
% cue 1215 2360.540 2363.570
\hypertarget{c01:cue:1215}{}這兩種不同的政治生態。
% cue 1216 2363.570 2364.950
\hypertarget{c01:cue:1216}{}英文叫做。
% cue 1217 2364.950 2365.630
\hypertarget{c01:cue:1217}{}有一個讀文代。
% cue 1218 2365.630 2367.270
\hypertarget{c01:cue:1218}{}它會一樣啊。
% cue 1219 2367.270 2369.560
\hypertarget{c01:cue:1219}{}它的現場操作會一樣啊。
% cue 1220 2369.560 2370.650
\hypertarget{c01:cue:1220}{}總統。
% cue 1221 2370.650 2372.680
\hypertarget{c01:cue:1221}{}我們通常叫做左右共治\footnote{法語 cohabitation，指總統與國會多數分屬不同政黨，由國會多數組閣的情形。法國第五共和曾三度出現，分別始於 1986 年、1993 年與 1997 年；前兩次起於國民議會任期屆滿改選，第三次起於總統解散國民議會後的改選，均由原反對黨取得多數席次。}。
% cue 1222 2372.680 2373.680
\hypertarget{c01:cue:1222}{}一個向左一個向右。
% cue 1223 2373.680 2375.680
\hypertarget{c01:cue:1223}{}如果不是同黨。
% cue 1224 2375.680 2376.680
\hypertarget{c01:cue:1224}{}同黨當然沒有話講。
% cue 1225 2376.680 2377.900
\hypertarget{c01:cue:1225}{}當然是總統。
% cue 1226 2377.900 2380.020
\hypertarget{c01:cue:1226}{}你行政院長不聽。
% cue 1227 2380.020 2381.020
\hypertarget{c01:cue:1227}{}我也當不了行政院。
% cue 1228 2381.020 2383.560
\hypertarget{c01:cue:1228}{}但是我們兩不同黨的時候。
% cue 1229 2383.560 2385.560
\hypertarget{c01:cue:1229}{}行政院是向立法院的負責。
% cue 1230 2385.560 2387.590
\hypertarget{c01:cue:1230}{}最正常的情況兩種。
% cue 1231 2387.590 2388.620
\hypertarget{c01:cue:1231}{}第一。
% cue 1232 2388.620 2390.680
\hypertarget{c01:cue:1232}{}行政院要經立法院同意\footnote{中華民國憲法第 55 條第 1 項規定：「行政院院長由總統提名，經立法院同意任命之。」1997 年修正的憲法增修條文第 3 條第 1 項明定該條停止適用。\url{https://law.moj.gov.tw/LawClass/LawAll.aspx?pcode=A0000001}。}。
% cue 1233 2390.680 2391.640
\hypertarget{c01:cue:1233}{}然後才能上任。

\tc{0:39:51}
% cue 1234 2391.640 2394.140
\hypertarget{c01:cue:1234}{}比如說有一個菲傭。
% cue 1235 2394.140 2395.450
\hypertarget{c01:cue:1235}{}他到你家來。
% cue 1236 2395.450 2396.710
\hypertarget{c01:cue:1236}{}他說。
% cue 1237 2396.710 2397.710
\hypertarget{c01:cue:1237}{}老爺。
% cue 1238 2397.710 2398.800
\hypertarget{c01:cue:1238}{}我很會做菜。
% cue 1239 2398.800 2399.800
\hypertarget{c01:cue:1239}{}家事料理的非常好。
% cue 1240 2399.800 2400.800
\hypertarget{c01:cue:1240}{}而且我手腳很乾淨。
% cue 1241 2400.800 2401.800
\hypertarget{c01:cue:1241}{}都不會脫節的。
% cue 1242 2401.800 2404.430
\hypertarget{c01:cue:1242}{}我們明天要來上班。
% cue 1243 2404.430 2407.670
\hypertarget{c01:cue:1243}{}你說。
% cue 1244 2407.670 2408.690
\hypertarget{c01:cue:1244}{}我沒有叫你來啊。
% cue 1245 2408.690 2409.690
\hypertarget{c01:cue:1245}{}你可以來嗎。
% cue 1246 2409.690 2411.330
\hypertarget{c01:cue:1246}{}你放心啊。
% cue 1247 2411.330 2413.360
\hypertarget{c01:cue:1247}{}我菲傭向老闆負責。
% cue 1248 2413.360 2414.300
\hypertarget{c01:cue:1248}{}我們明天就來。
% cue 1249 2414.300 2417.300
\hypertarget{c01:cue:1249}{}沒有人。
% cue 1250 2417.300 2418.300
\hypertarget{c01:cue:1250}{}你沒有同意。
% cue 1251 2418.300 2419.300
\hypertarget{c01:cue:1251}{}有人說。
% cue 1252 2419.300 2420.300
\hypertarget{c01:cue:1252}{}他向你負責。
% cue 1253 2420.300 2422.240
\hypertarget{c01:cue:1253}{}他可以上任嗎。
% cue 1254 2422.240 2423.490
\hypertarget{c01:cue:1254}{}這是不是笑話。
% cue 1255 2423.490 2425.080
\hypertarget{c01:cue:1255}{}好。

\tc{0:40:25}
% cue 1256 2425.080 2426.620
\hypertarget{c01:cue:1256}{}第二種情況。
% cue 1257 2426.620 2427.620
\hypertarget{c01:cue:1257}{}最好是兩者都有。
% cue 1258 2427.620 2428.620
\hypertarget{c01:cue:1258}{}但是起碼要有一個。
% cue 1259 2428.620 2430.590
\hypertarget{c01:cue:1259}{}不然你絕對不會是犯同。
% cue 1260 2430.590 2431.970
\hypertarget{c01:cue:1260}{}第二種情況是。
% cue 1261 2431.970 2433.120
\hypertarget{c01:cue:1261}{}你在這邊做了三年。
% cue 1262 2433.120 2434.120
\hypertarget{c01:cue:1262}{}非優做了三年。
% cue 1263 2434.120 2436.630
\hypertarget{c01:cue:1263}{}現在的膽子越來越大。
% cue 1264 2436.630 2438.100
\hypertarget{c01:cue:1264}{}跟你講說每天要有菜。
% cue 1265 2438.100 2439.100
\hypertarget{c01:cue:1265}{}魚肉。
% cue 1266 2439.100 2440.100
\hypertarget{c01:cue:1266}{}你現在菜裡沒有。
% cue 1267 2440.100 2441.050
\hypertarget{c01:cue:1267}{}只剩肉。
% cue 1268 2441.050 2442.320
\hypertarget{c01:cue:1268}{}我已經這麼胖。
% cue 1269 2442.320 2443.420
\hypertarget{c01:cue:1269}{}我不要再吃肉了。
% cue 1270 2443.420 2444.550
\hypertarget{c01:cue:1270}{}我已經警告你了。
% cue 1271 2444.550 2445.660
\hypertarget{c01:cue:1271}{}三十次了。
% cue 1272 2445.660 2446.920
\hypertarget{c01:cue:1272}{}明天你再跟我做肉。
% cue 1273 2446.920 2447.830
\hypertarget{c01:cue:1273}{}沒有菜。
% cue 1274 2447.830 2449.060
\hypertarget{c01:cue:1274}{}You're fired。
% cue 1275 2449.060 2450.100
\hypertarget{c01:cue:1275}{}我學 Trump\footnote{Donald J. Trump，美國第 45 任與第 47 任總統。\url{https://www.whitehouse.gov/administration/donald-j-trump/}。} 講話。

\tc{0:40:50}
% cue 1276 2450.100 2452.180
\hypertarget{c01:cue:1276}{}好。
% cue 1277 2452.180 2452.280
\hypertarget{c01:cue:1277}{}你一看他。
% cue 1278 2452.180 2453.180
\hypertarget{c01:cue:1278}{}他真正就沒有菜。
% cue 1279 2453.180 2454.180
\hypertarget{c01:cue:1279}{}只有肉。
% cue 1280 2454.180 2455.210
\hypertarget{c01:cue:1280}{}那你就看他。
% cue 1281 2455.210 2456.650
\hypertarget{c01:cue:1281}{}You're fired。
% cue 1282 2456.650 2457.630
\hypertarget{c01:cue:1282}{}我講話。
% cue 1283 2457.630 2458.650
\hypertarget{c01:cue:1283}{}我講話。
% cue 1284 2458.650 2459.930
\hypertarget{c01:cue:1284}{}我講話。
% cue 1285 2459.930 2462.750
\hypertarget{c01:cue:1285}{}我講話。
% cue 1286 2462.750 2464.070
\hypertarget{c01:cue:1286}{}有聽懂嗎。
% cue 1287 2464.070 2465.060
\hypertarget{c01:cue:1287}{}好。

\tc{0:41:05}
% cue 1288 2465.060 2466.900
\hypertarget{c01:cue:1288}{}那臺灣的情況很像。
% cue 1289 2466.900 2467.900
\hypertarget{c01:cue:1289}{}臺灣表面上說。
% cue 1290 2467.900 2468.890
\hypertarget{c01:cue:1290}{}立法院隨時可以通過不信任案\footnote{中華民國憲法增修條文第 3 條第 2 項第 3 款：「立法院得經全體立法委員三分之一以上連署，對行政院院長提出不信任案。不信任案提出七十二小時後，應於四十八小時內以記名投票表決之。如經全體立法委員二分之一以上贊成，行政院院長應於十日內提出辭職，並得同時呈請總統解散立法院；不信任案如未獲通過，一年內不得對同一行政院院長再提不信任案。」\url{https://law.moj.gov.tw/LawClass/LawAll.aspx?pcode=A0000002}。}。
% cue 1291 2468.890 2471.120
\hypertarget{c01:cue:1291}{}可以倒閣。
% cue 1292 2471.120 2472.120
\hypertarget{c01:cue:1292}{}好啊。
% cue 1293 2472.120 2473.240
\hypertarget{c01:cue:1293}{}你通過啊。
% cue 1294 2473.240 2474.520
\hypertarget{c01:cue:1294}{}通過的現在問題很多。
% cue 1295 2474.520 2475.490
\hypertarget{c01:cue:1295}{}我們慢慢講。

\tc{0:41:15}
% cue 1296 2475.490 2476.740
\hypertarget{c01:cue:1296}{}第一動。
% cue 1297 2476.740 2477.740
\hypertarget{c01:cue:1297}{}誰通過立法院通過。
% cue 1298 2477.740 2480.300
\hypertarget{c01:cue:1298}{}接下來誰可以做反應。
% cue 1299 2480.300 2481.300
\hypertarget{c01:cue:1299}{}立法院通過。
% cue 1300 2481.300 2483.580
\hypertarget{c01:cue:1300}{}不信任案是對誰不信任。
% cue 1301 2483.580 2484.580
\hypertarget{c01:cue:1301}{}是對行政院長不信任。
% cue 1302 2484.580 2485.830
\hypertarget{c01:cue:1302}{}所以接下來。
% cue 1303 2485.830 2486.830
\hypertarget{c01:cue:1303}{}憲法的增修條文會說。
% cue 1304 2486.830 2490.170
\hypertarget{c01:cue:1304}{}行政院長他必須在通過以後。
% cue 1305 2490.170 2491.240
\hypertarget{c01:cue:1305}{}10 天內辭職。
% cue 1306 2491.240 2492.140
\hypertarget{c01:cue:1306}{}他沒有別的選項。

\tc{0:41:32}
% cue 1307 2492.140 2495.090
\hypertarget{c01:cue:1307}{}但是辭職的同時。
% cue 1308 2495.090 2496.090
\hypertarget{c01:cue:1308}{}他可以呈請總統。
% cue 1309 2496.090 2498.060
\hypertarget{c01:cue:1309}{}賺 100 元。
% cue 1310 2498.060 2499.250
\hypertarget{c01:cue:1310}{}這個像一隻回馬槍。
% cue 1311 2499.250 2500.500
\hypertarget{c01:cue:1311}{}你可以殺我。
% cue 1312 2500.500 2501.500
\hypertarget{c01:cue:1312}{}我就不能報復嗎。
% cue 1313 2501.500 2502.500
\hypertarget{c01:cue:1313}{}忘記了。
% cue 1314 2502.500 2503.790
\hypertarget{c01:cue:1314}{}我有回馬槍。
% cue 1315 2503.790 2505.810
\hypertarget{c01:cue:1315}{}我根本的老闆講。
% cue 1316 2505.810 2506.810
\hypertarget{c01:cue:1316}{}老闆。
% cue 1317 2506.810 2508.220
\hypertarget{c01:cue:1317}{}我是你的分身。
% cue 1318 2508.220 2509.220
\hypertarget{c01:cue:1318}{}你的分身給人幹掉了。
% cue 1319 2509.220 2510.280
\hypertarget{c01:cue:1319}{}你替我報出去。

\tc{0:41:50}
% cue 1320 2510.280 2511.690
\hypertarget{c01:cue:1320}{}那這時候。
% cue 1321 2511.690 2512.690
\hypertarget{c01:cue:1321}{}要看老闆怎麼辦。
% cue 1322 2512.690 2513.780
\hypertarget{c01:cue:1322}{}老闆有兩種選擇。
% cue 1323 2513.780 2514.820
\hypertarget{c01:cue:1323}{}因為憲法並沒有說。
% cue 1324 2514.820 2518.140
\hypertarget{c01:cue:1324}{}這種情況發生。
% cue 1325 2518.140 2519.140
\hypertarget{c01:cue:1325}{}總統一定要解散\footnote{中華民國憲法增修條文第 2 條第 5 項規定，總統於立法院通過對行政院院長之不信任案後十日內，經諮詢立法院院長後，「得宣告解散立法院」；戒嚴或緊急命令生效期間不得解散。立法院解散後，應於六十日內舉行立法委員選舉。\url{https://law.moj.gov.tw/LawClass/LawAll.aspx?pcode=A0000002}。}。
% cue 1326 2519.140 2520.480
\hypertarget{c01:cue:1326}{}不會。
% cue 1327 2520.480 2524.100
\hypertarget{c01:cue:1327}{}沒有。
% cue 1328 2524.100 2525.100
\hypertarget{c01:cue:1328}{}總統可以忍諒一下。
% cue 1329 2525.100 2526.640
\hypertarget{c01:cue:1329}{}認為這個時候。
% cue 1330 2526.640 2528.520
\hypertarget{c01:cue:1330}{}講話這個時候。
% cue 1331 2528.520 2529.630
\hypertarget{c01:cue:1331}{}如果。
% cue 1332 2529.630 2530.760
\hypertarget{c01:cue:1332}{}如果。
% cue 1333 2530.760 2531.760
\hypertarget{c01:cue:1333}{}立法院真的倒閣。
% cue 1334 2531.760 2532.760
\hypertarget{c01:cue:1334}{}如果。
% cue 1335 2532.760 2534.300
\hypertarget{c01:cue:1335}{}然後你認為賴清德\footnote{中華民國第 16 任總統，2024 年 5 月 20 日就任。\url{https://www.president.gov.tw/Page/694}。}會解散。
% cue 1336 2534.300 2535.300
\hypertarget{c01:cue:1336}{}會解散。
% cue 1337 2535.300 2536.780
\hypertarget{c01:cue:1337}{}會解散。
% cue 1338 2536.780 2537.780
\hypertarget{c01:cue:1338}{}國會重選嗎。

\tc{0:42:17}
% cue 1339 2537.780 2539.840
\hypertarget{c01:cue:1339}{}我猜想賴清德的考慮只有一個。
% cue 1340 2539.840 2540.930
\hypertarget{c01:cue:1340}{}如果該算。
% cue 1341 2540.930 2543.280
\hypertarget{c01:cue:1341}{}我的席次會比較多還是比較少。
% cue 1342 2543.280 2545.630
\hypertarget{c01:cue:1342}{}如果從目前這個情況看。
% cue 1343 2545.630 2547.300
\hypertarget{c01:cue:1343}{}不一定會多。
% cue 1344 2547.300 2548.300
\hypertarget{c01:cue:1344}{}不少就不算。
% cue 1345 2548.300 2549.580
\hypertarget{c01:cue:1345}{}對不對。
% cue 1346 2549.580 2550.680
\hypertarget{c01:cue:1346}{}所以呢。
% cue 1347 2550.680 2551.800
\hypertarget{c01:cue:1347}{}拒絕。
% cue 1348 2551.800 2552.800
\hypertarget{c01:cue:1348}{}可以。
% cue 1349 2552.800 2553.860
\hypertarget{c01:cue:1349}{}可不可以拒絕。
% cue 1350 2553.860 2555.120
\hypertarget{c01:cue:1350}{}可以嗎。
% cue 1351 2555.120 2556.110
\hypertarget{c01:cue:1351}{}拒絕不拒絕。
% cue 1352 2556.110 2557.230
\hypertarget{c01:cue:1352}{}是一件事。
% cue 1353 2557.230 2558.290
\hypertarget{c01:cue:1353}{}不少。
% cue 1354 2558.290 2559.290
\hypertarget{c01:cue:1354}{}不少就不算。
% cue 1355 2559.290 2560.410
\hypertarget{c01:cue:1355}{}所以呢。
% cue 1356 2560.410 2562.140
\hypertarget{c01:cue:1356}{}拒絕。
% cue 1357 2562.140 2563.140
\hypertarget{c01:cue:1357}{}可以。
% cue 1358 2563.140 2564.360
\hypertarget{c01:cue:1358}{}拒絕不拒絕。
% cue 1359 2564.360 2565.360
\hypertarget{c01:cue:1359}{}拒絕是一件事。

\tc{0:42:45}
% cue 1360 2565.360 2566.350
\hypertarget{c01:cue:1360}{}原來行政院長。
% cue 1361 2566.350 2567.350
\hypertarget{c01:cue:1361}{}搞不搞。
% cue 1362 2567.350 2567.450
\hypertarget{c01:cue:1362}{}搞啦。
% cue 1363 2567.350 2569.240
\hypertarget{c01:cue:1363}{}如果他拒絕。
% cue 1364 2569.240 2570.240
\hypertarget{c01:cue:1364}{}解散。
% cue 1365 2570.240 2571.200
\hypertarget{c01:cue:1365}{}行政院長。
% cue 1366 2571.200 2572.200
\hypertarget{c01:cue:1366}{}10 天內還是要。
% cue 1367 2572.200 2573.390
\hypertarget{c01:cue:1367}{}下臺。
% cue 1368 2573.390 2574.390
\hypertarget{c01:cue:1368}{}下臺以後。
% cue 1369 2574.390 2575.290
\hypertarget{c01:cue:1369}{}誰任命行政院長。

\tc{0:42:55}
% cue 1370 2575.290 2576.520
\hypertarget{c01:cue:1370}{}看照我們現在的憲法。
% cue 1371 2576.520 2578.430
\hypertarget{c01:cue:1371}{}這就臺灣的一面遊戲去看。
% cue 1372 2578.430 2579.520
\hypertarget{c01:cue:1372}{}第三條。
% cue 1373 2579.520 2580.520
\hypertarget{c01:cue:1373}{}第一項。
% cue 1374 2580.520 2581.520
\hypertarget{c01:cue:1374}{}第一句。
% cue 1375 2581.520 2582.520
\hypertarget{c01:cue:1375}{}行政院長由總統任命之\footnote{中華民國憲法增修條文第 3 條第 1 項：「行政院院長由總統任命之。行政院院長辭職或出缺時，在總統未任命行政院院長前，由行政院副院長暫行代理。憲法第五十五條之規定，停止適用。」\url{https://law.moj.gov.tw/LawClass/LawAll.aspx?pcode=A0000002}。}。
% cue 1376 2582.520 2584.860
\hypertarget{c01:cue:1376}{}句話。
% cue 1377 2584.860 2585.860
\hypertarget{c01:cue:1377}{}沒有了。
% cue 1378 2585.860 2585.960
\hypertarget{c01:cue:1378}{}沒有什麼意思。
% cue 1379 2585.860 2587.850
\hypertarget{c01:cue:1379}{}就這麼短。
% cue 1380 2587.850 2588.850
\hypertarget{c01:cue:1380}{}誰任命行政院長。
% cue 1381 2588.850 2590.490
\hypertarget{c01:cue:1381}{}他有義務要反映。
% cue 1382 2590.490 2591.490
\hypertarget{c01:cue:1382}{}立法院的生態嗎。
% cue 1383 2591.490 2593.550
\hypertarget{c01:cue:1383}{}他有要向誰來爭取。
% cue 1384 2593.550 2594.550
\hypertarget{c01:cue:1384}{}商量的義務嗎。
% cue 1385 2594.550 2595.520
\hypertarget{c01:cue:1385}{}都沒有。

\tc{0:43:15}
% cue 1386 2595.520 2596.830
\hypertarget{c01:cue:1386}{}所以卓榮泰\footnote{2024 年 5 月 20 日就任行政院院長，曾任立法委員、行政院秘書長，2019 年至 2020 年任民主進步黨主席。\url{https://www.ey.gov.tw/Page/3C927BC7FA83C969}。}倒了。
% cue 1387 2596.830 2597.830
\hypertarget{c01:cue:1387}{}換卓榮泰。
% cue 1388 2597.830 2600.110
\hypertarget{c01:cue:1388}{}你們不喜歡他。
% cue 1389 2600.110 2601.110
\hypertarget{c01:cue:1389}{}再倒啊。
% cue 1390 2601.110 2602.230
\hypertarget{c01:cue:1390}{}你就換卓榮泰。
% cue 1391 2602.230 2603.260
\hypertarget{c01:cue:1391}{}二。
% cue 1392 2603.260 2605.250
\hypertarget{c01:cue:1392}{}行政院長這麼多。
% cue 1393 2605.250 2606.250
\hypertarget{c01:cue:1393}{}沒有人做嗎。
% cue 1394 2606.250 2608.180
\hypertarget{c01:cue:1394}{}小朋友當行政院長。
% cue 1395 2608.180 2609.180
\hypertarget{c01:cue:1395}{}跟我講。
% cue 1396 2609.180 2611.810
\hypertarget{c01:cue:1396}{}爽。
% cue 1397 2611.810 2612.810
\hypertarget{c01:cue:1397}{}約四十幾萬。
% cue 1398 2612.810 2615.730
\hypertarget{c01:cue:1398}{}爽不爽。
% cue 1399 2615.730 2617.460
\hypertarget{c01:cue:1399}{}擔心沒有人做。
% cue 1400 2617.460 2618.460
\hypertarget{c01:cue:1400}{}你只要擔心的是。
% cue 1401 2618.460 2619.460
\hypertarget{c01:cue:1401}{}卓榮泰一。
% cue 1402 2619.460 2620.460
\hypertarget{c01:cue:1402}{}也許還有一點長處。
% cue 1403 2620.460 2622.800
\hypertarget{c01:cue:1403}{}到了卓榮泰二。
% cue 1404 2622.800 2623.800
\hypertarget{c01:cue:1404}{}可能零長處。
% cue 1405 2623.800 2625.150
\hypertarget{c01:cue:1405}{}到了卓榮泰三。
% cue 1406 2625.150 2627.110
\hypertarget{c01:cue:1406}{}那簡直別賴清楚。
% cue 1407 2627.110 2630.110
\hypertarget{c01:cue:1407}{}還要賴清楚。

\tc{0:43:50}
% cue 1408 2630.110 2633.270
\hypertarget{c01:cue:1408}{}那如果現在他說。
% cue 1409 2633.270 2633.370
\hypertarget{c01:cue:1409}{}好。
% cue 1410 2633.270 2634.460
\hypertarget{c01:cue:1410}{}我認為賭一把。
% cue 1411 2634.460 2635.740
\hypertarget{c01:cue:1411}{}那就賭了。
% cue 1412 2635.740 2637.630
\hypertarget{c01:cue:1412}{}解散不貴。
% cue 1413 2637.630 2639.010
\hypertarget{c01:cue:1413}{}解散完畢後重選。
% cue 1414 2639.010 2640.260
\hypertarget{c01:cue:1414}{}重選先不管說。
% cue 1415 2640.260 2643.290
\hypertarget{c01:cue:1415}{}重選中間幾個月的空檔。
% cue 1416 2643.290 2644.290
\hypertarget{c01:cue:1416}{}這會發生什麼事情。
% cue 1417 2644.290 2646.280
\hypertarget{c01:cue:1417}{}他去告誡也怎麼辦。
% cue 1418 2646.280 2647.340
\hypertarget{c01:cue:1418}{}等等等等。
% cue 1419 2647.340 2649.370
\hypertarget{c01:cue:1419}{}他事實上是不能去告誡。

\tc{0:44:09}
% cue 1420 2649.370 2650.790
\hypertarget{c01:cue:1420}{}如果這麼假設。
% cue 1421 2650.790 2653.350
\hypertarget{c01:cue:1421}{}那改完。
% cue 1422 2653.350 2654.350
\hypertarget{c01:cue:1422}{}改選。
% cue 1423 2654.350 2655.350
\hypertarget{c01:cue:1423}{}改選以後。
% cue 1424 2655.350 2656.570
\hypertarget{c01:cue:1424}{}國民黨加民眾黨\footnote{2024 年 1 月 13 日第 11 屆立法委員選舉，113 席中，國民黨 52 席、民進黨 51 席、民眾黨 8 席，另有無黨籍 2 席。\url{https://www.cna.com.tw/news/aipl/202401130361.aspx}。}。
% cue 1425 2656.570 2657.570
\hypertarget{c01:cue:1425}{}比目前的情況。
% cue 1426 2657.570 2658.790
\hypertarget{c01:cue:1426}{}照樣過半。
% cue 1427 2658.790 2660.790
\hypertarget{c01:cue:1427}{}而且其實更多。
% cue 1428 2660.790 2663.580
\hypertarget{c01:cue:1428}{}現在是五十二加八。
% cue 1429 2663.580 2664.730
\hypertarget{c01:cue:1429}{}最多六十。
% cue 1430 2664.730 2666.980
\hypertarget{c01:cue:1430}{}他現在改選變七十。
% cue 1431 2666.980 2667.980
\hypertarget{c01:cue:1431}{}反正五百一十三。
% cue 1432 2667.980 2668.080
\hypertarget{c01:cue:1432}{}多得很。
% cue 1433 2667.980 2668.980
\hypertarget{c01:cue:1433}{}變七十。

\tc{0:44:28}
% cue 1434 2668.980 2671.570
\hypertarget{c01:cue:1434}{}這兩黨強烈要求說。
% cue 1435 2671.570 2673.670
\hypertarget{c01:cue:1435}{}其正院校可以百院負責。
% cue 1436 2673.670 2674.960
\hypertarget{c01:cue:1436}{}所以我們要出來組閣。
% cue 1437 2674.960 2677.110
\hypertarget{c01:cue:1437}{}你這裡滑。
% cue 1438 2677.110 2681.230
\hypertarget{c01:cue:1438}{}看你喊到喉嚨都破。
% cue 1439 2681.230 2682.520
\hypertarget{c01:cue:1439}{}有沒有用。

\tc{0:44:42}
% cue 1440 2682.520 2683.520
\hypertarget{c01:cue:1440}{}先把條文寫在這裡。
% cue 1441 2683.520 2685.770
\hypertarget{c01:cue:1441}{}行政院長由總統任命之。
% cue 1442 2685.770 2687.240
\hypertarget{c01:cue:1442}{}聽不懂嗎。
% cue 1443 2687.240 2688.620
\hypertarget{c01:cue:1443}{}吵什麼吵。
% cue 1444 2688.620 2689.620
\hypertarget{c01:cue:1444}{}阿別當年就講過。
% cue 1445 2689.620 2690.650
\hypertarget{c01:cue:1445}{}這怪我嗎。
% cue 1446 2690.650 2691.940
\hypertarget{c01:cue:1446}{}怪我阿別的道理嗎。
% cue 1447 2691.940 2692.940
\hypertarget{c01:cue:1447}{}這怪李登輝啊。
% cue 1448 2692.940 2695.380
\hypertarget{c01:cue:1448}{}他其實還講漏了一句。
% cue 1449 2695.380 2697.530
\hypertarget{c01:cue:1449}{}怪李登輝跟許信良。
% cue 1450 2697.530 2698.530
\hypertarget{c01:cue:1450}{}他還看出來講話。
% cue 1451 2698.530 2700.420
\hypertarget{c01:cue:1451}{}休憲休止。
% cue 1452 2700.420 2705.410
\hypertarget{c01:cue:1452}{}這樣的人很見話。

\tc{0:45:05}
% cue 1453 2705.410 2706.520
\hypertarget{c01:cue:1453}{}如果今天老師跟妳講。
% cue 1454 2706.520 2708.540
\hypertarget{c01:cue:1454}{}大概根本就不知道許信良。
% cue 1455 2708.540 2711.660
\hypertarget{c01:cue:1455}{}許信良是國民黨。
% cue 1456 2711.660 2713.390
\hypertarget{c01:cue:1456}{}後來變成民進黨。
% cue 1457 2713.390 2715.830
\hypertarget{c01:cue:1457}{}阿別以前是國民。
% cue 1458 2715.830 2717.660
\hypertarget{c01:cue:1458}{}後來也變成民進黨。
% cue 1459 2717.660 2719.480
\hypertarget{c01:cue:1459}{}很多民進黨人。
% cue 1460 2719.480 2720.480
\hypertarget{c01:cue:1460}{}都是國民黨出色的。
% cue 1461 2720.480 2721.480
\hypertarget{c01:cue:1461}{}被什麼大體驗壞。
% cue 1462 2721.480 2724.460
\hypertarget{c01:cue:1462}{}好啦 這個不重要。

\tc{0:45:24}
% cue 1463 2724.460 2725.490
\hypertarget{c01:cue:1463}{}講到這裡。
% cue 1464 2725.490 2727.660
\hypertarget{c01:cue:1464}{}那可見現在叫。
% cue 1465 2727.660 2728.660
\hypertarget{c01:cue:1465}{}虛有其他人。
% cue 1466 2728.660 2730.170
\hypertarget{c01:cue:1466}{}所以我們不是犯人。
% cue 1467 2730.170 2731.430
\hypertarget{c01:cue:1467}{}好啦 講到這樣說法。

\section{分週進度與分組建議}

\tc{0:45:31}
% cue 1468 2731.430 2733.730
\hypertarget{c01:cue:1468}{}下面就看到每個禮拜。
% cue 1469 2733.730 2735.730
\hypertarget{c01:cue:1469}{}老師要講教請。
% cue 1470 2735.730 2736.640
\hypertarget{c01:cue:1470}{}就繼續講吧。
% cue 1471 2736.640 2738.040
\hypertarget{c01:cue:1471}{}我只是要跟你們講說。
% cue 1472 2738.040 2740.290
\hypertarget{c01:cue:1472}{}給老師要講的。
% cue 1473 2740.290 2741.290
\hypertarget{c01:cue:1473}{}老師畫號。
% cue 1474 2741.290 2742.250
\hypertarget{c01:cue:1474}{}畫完畢了以後。
% cue 1475 2742.250 2743.250
\hypertarget{c01:cue:1475}{}就要改定了。
% cue 1476 2743.250 2744.640
\hypertarget{c01:cue:1476}{}就要去做做的事情。
% cue 1477 2744.640 2752.860
\hypertarget{c01:cue:1477}{}一直到 10 月 12 號開始。
% cue 1478 2752.860 2753.860
\hypertarget{c01:cue:1478}{}就改定了。
% cue 1479 2753.860 2754.860
\hypertarget{c01:cue:1479}{}有沒有看到。
% cue 1480 2754.860 2755.860
\hypertarget{c01:cue:1480}{}講義上看到了嗎。
% cue 1481 2755.860 2758.750
\hypertarget{c01:cue:1481}{}31 號解釋\footnote{司法院釋字第 31 號解釋（1954 年 1 月 29 日）：國家發生重大變故，事實上不能依法辦理次屆選舉時，「在第二屆委員未能依法選出集會與召集以前，自應仍由第一屆立法委員、監察委員繼續行使其職權」。\url{https://cons.judicial.gov.tw/docdata.aspx?fid=100&id=310212}。}。
% cue 1482 2758.750 2759.750
\hypertarget{c01:cue:1482}{}261 號解釋\footnote{司法院釋字第 261 號解釋（1990 年 6 月 21 日）：第一屆未定期改選之中央民意代表，除事實上已不能行使職權或經常不行使職權者應即查明解職外，「其餘應於中華民國八十年十二月三十一日以前終止行使職權」，並由中央政府適時辦理全國性之次屆中央民意代表選舉。\cite{jy261}。}。
% cue 1483 2759.750 2762.000
\hypertarget{c01:cue:1483}{}還有 114 憲判 1\footnote{中華民國憲法法庭 114 年憲判字第 1 號判決（2025 年 12 月 19 日）宣告 2025 年 1 月 23 日修正公布的憲法訴訟法第 4 條第 3 項、第 30 條第 2 項至第 6 項及第 95 條「立法程序有明顯重大瑕疵，違背憲法正當立法程序，且違反憲法權力分立原則，均牴觸憲法」，自判決公告之日起失其效力。課綱列為第 5 週（10 月 12 日）「憲政危機」的案件。\cite{ccj114-1}。}。

\tc{0:46:02}
% cue 1484 2762.000 2763.000
\hypertarget{c01:cue:1484}{}所以這裡呢。
% cue 1485 2763.000 2764.410
\hypertarget{c01:cue:1485}{}老師的建議就是。
% cue 1486 2764.410 2766.720
\hypertarget{c01:cue:1486}{}31 和 261 算 1 組。
% cue 1487 2766.720 2769.260
\hypertarget{c01:cue:1487}{}114 憲判 1 算第二組。
% cue 1488 2769.260 2771.700
\hypertarget{c01:cue:1488}{}到這裡我們不能再過了。

\tc{0:46:11}
% cue 1489 2771.700 2772.700
\hypertarget{c01:cue:1489}{}需要有一個班長。
% cue 1490 2772.700 2773.800
\hypertarget{c01:cue:1490}{}再問一次。
% cue 1491 2773.800 2776.360
\hypertarget{c01:cue:1491}{}哪一個人自願當班長。
% cue 1492 2776.360 2777.300
\hypertarget{c01:cue:1492}{}好 現在講。
% cue 1493 2777.300 2778.650
\hypertarget{c01:cue:1493}{}要做的事情。
% cue 1494 2778.650 2781.080
\hypertarget{c01:cue:1494}{}我現在是全力一路相對。
% cue 1495 2781.080 2782.360
\hypertarget{c01:cue:1495}{}班長什麼好處呢。
% cue 1496 2782.360 2784.540
\hypertarget{c01:cue:1496}{}學期總成績加兩分。
% cue 1497 2784.540 2787.880
\hypertarget{c01:cue:1497}{}你看這是不是跟大家一樣的。
% cue 1498 2787.880 2789.970
\hypertarget{c01:cue:1498}{}但是班長有服務。
% cue 1499 2789.970 2790.960
\hypertarget{c01:cue:1499}{}所以學期服務。
% cue 1500 2790.960 2793.380
\hypertarget{c01:cue:1500}{}學期總成績加兩分。
% cue 1501 2793.380 2797.210
\hypertarget{c01:cue:1501}{}好 還有沒有人自願要當班長。
% cue 1502 2797.210 2799.870
\hypertarget{c01:cue:1502}{}夠嗎。
% cue 1503 2799.870 2800.870
\hypertarget{c01:cue:1503}{}要加 20 分嗎。
% cue 1504 2800.870 2803.430
\hypertarget{c01:cue:1504}{}好 那你看看。
% cue 1505 2803.430 2804.430
\hypertarget{c01:cue:1505}{}記得怎麼開這個東西。
% cue 1506 2804.430 2805.430
\hypertarget{c01:cue:1506}{}好吧。

\tc{0:46:45}
% cue 1507 2805.430 2807.430
\hypertarget{c01:cue:1507}{}你看 10 月 11。
% cue 1508 2807.430 2808.430
\hypertarget{c01:cue:1508}{}10 月 12 號。
% cue 1509 2808.430 2809.430
\hypertarget{c01:cue:1509}{}班會老師走。
% cue 1510 2809.430 2811.460
\hypertarget{c01:cue:1510}{}你們就必須要有人來認領。
% cue 1511 2811.460 2812.650
\hypertarget{c01:cue:1511}{}這東西是用這樣子。
% cue 1512 2812.650 2813.650
\hypertarget{c01:cue:1512}{}教練式的。
% cue 1513 2813.650 2814.770
\hypertarget{c01:cue:1513}{}先講先語。
% cue 1514 2814.770 2815.770
\hypertarget{c01:cue:1514}{}明天就可以想。
% cue 1515 2815.770 2816.860
\hypertarget{c01:cue:1515}{}那個解釋。
% cue 1516 2816.860 2817.860
\hypertarget{c01:cue:1516}{}你有興趣。
% cue 1517 2817.860 2818.860
\hypertarget{c01:cue:1517}{}就趕快登記就對了。
% cue 1518 2818.860 2820.010
\hypertarget{c01:cue:1518}{}誰先拿。
% cue 1519 2820.010 2821.010
\hypertarget{c01:cue:1519}{}誰先贏。

\tc{0:47:01}
% cue 1520 2821.010 2823.480
\hypertarget{c01:cue:1520}{}然後在我們下。
% cue 1521 2823.480 2824.480
\hypertarget{c01:cue:1521}{}10 月 16 號。
% cue 1522 2824.480 2825.540
\hypertarget{c01:cue:1522}{}這裡看起來。
% cue 1523 2825.540 2826.540
\hypertarget{c01:cue:1523}{}就不只是兩組。
% cue 1524 2826.540 2828.140
\hypertarget{c01:cue:1524}{}我會建議呢。
% cue 1525 2828.140 2832.510
\hypertarget{c01:cue:1525}{}3 號\footnote{司法院釋字第 3 號解釋（1952 年 5 月 21 日）：憲法第 87 條明定考試院得向立法院提出法律案，對監察院則無明文，解釋基於「五權分治，平等相維之體制」，認「監察院關於所掌事項，得向立法院提出法律案，實與憲法之精神相符」。\url{https://cons.judicial.gov.tw/docdata.aspx?fid=100&id=310184}。} 175 號\footnote{司法院釋字第 175 號解釋（1982 年 5 月 25 日）：「司法院為國家最高司法機關，基於五權分治彼此相維之憲政體制，就其所掌有關司法機關之組織及司法權行使之事項，得向立法院提出法律案。」\url{https://cons.judicial.gov.tw/docdata.aspx?fid=100&id=310356}。}一組。
% cue 1526 2832.510 2835.310
\hypertarget{c01:cue:1526}{}14 號\footnote{司法院釋字第 14 號解釋（1953 年 3 月 21 日）依制憲經過認定，「在制憲者之意，當以立監委員為直接或間接之民意代表，均不認其為監察權行使之對象」，國民大會代表與省縣議會議員亦同。\url{https://cons.judicial.gov.tw/docdata.aspx?fid=100&id=310195}。}跟 319 號\footnote{司法院釋字第 319 號解釋（1993 年 6 月 4 日）：考試機關之閱卷委員於試卷彌封時評定成績，彌封開拆後，除依形式觀察即可發見成績有顯然錯誤者外，「不應循應考人之要求任意再行評閱」；應考人申請複查考試成績處理辦法第 8 條與憲法尚無牴觸。\url{https://cons.judicial.gov.tw/docdata.aspx?fid=100&id=310500}。}一組。
% cue 1527 2835.310 2837.720
\hypertarget{c01:cue:1527}{}然後 461\footnote{司法院釋字第 461 號解釋（1998 年 7 月 24 日）：參謀總長非憲法所定之部會首長，無須至立法院院會備質詢；但屬憲法第 67 條第 2 項所指政府人員，「除非因執行關係國家安全之軍事業務而有正當理由外，不得拒絕應邀到會備詢」。\url{https://cons.judicial.gov.tw/docdata.aspx?fid=100&id=310642}。}、632\footnote{司法院釋字第 632 號解釋（2007 年 8 月 15 日）：「總統如消極不為提名，或立法院消極不行使同意權，致監察院無從行使職權、發揮功能，國家憲政制度之完整因而遭受破壞，自為憲法所不許。」\url{https://cons.judicial.gov.tw/docdata.aspx?fid=100&id=310813}。} 一組。
% cue 1528 2837.720 2839.290
\hypertarget{c01:cue:1528}{}基本上我放置。
% cue 1529 2839.290 2840.220
\hypertarget{c01:cue:1529}{}橫條的。
% cue 1530 2840.220 2841.420
\hypertarget{c01:cue:1530}{}大概都是有原因。
% cue 1531 2841.420 2843.410
\hypertarget{c01:cue:1531}{}但你要把上面四個放置一組。
% cue 1532 2843.410 2844.370
\hypertarget{c01:cue:1532}{}我也不反對。

\tc{0:47:24}
% cue 1533 2844.370 2846.650
\hypertarget{c01:cue:1533}{}那下面。
% cue 1534 2846.650 2847.620
\hypertarget{c01:cue:1534}{}11 月 2 號。
% cue 1535 2847.620 2848.620
\hypertarget{c01:cue:1535}{}387\footnote{司法院釋字第 387 號解釋（1995 年 10 月 13 日）基於民意政治與責任政治之原理，認「立法委員任期屆滿改選後第一次集會前，行政院院長自應向總統提出辭職」。\url{https://cons.judicial.gov.tw/docdata.aspx?fid=100&id=310568}。} 1 組。
% cue 1536 2848.620 2849.610
\hypertarget{c01:cue:1536}{}419\footnote{司法院釋字第 419 號解釋（1996 年 12 月 31 日）：副總統兼任行政院院長「與憲法設置副總統及行政院院長職位分由不同之人擔任之本旨未盡相符」；行政院院長於新任總統就職時提出總辭，係禮貌性辭職，「並非其憲法上之義務」。湯德宗曾為該案提出鑑定意見書，收於《權力分立新論 卷一》。\url{https://cons.judicial.gov.tw/docdata.aspx?fid=100&id=310600}。} 1 組。
% cue 1537 2849.610 2850.860
\hypertarget{c01:cue:1537}{}11 月 9 號。
% cue 1538 2850.860 2851.860
\hypertarget{c01:cue:1538}{}520\footnote{司法院釋字第 520 號解釋（2001 年 1 月 15 日）審查行政院停止執行核四預算一案，認因重要政策變更涉及法定預算之停止執行時，應「本行政院對立法院負責之憲法意旨暨尊重立法院對國家重要事項之參與決策權」，由行政院院長或有關部會首長適時向立法院提出報告並備質詢。《權力分立新論 卷一》所收專文為〈憲法結構與違憲審查：司法院大法官釋字第五二〇號解釋評釋〉（增訂三版頁 109 起）。\url{https://cons.judicial.gov.tw/docdata.aspx?fid=100&id=310701}。} 很複雜。
% cue 1539 2851.860 2852.860
\hypertarget{c01:cue:1539}{}單獨有一組。
% cue 1540 2852.860 2854.990
\hypertarget{c01:cue:1540}{}老師在視線區。
% cue 1541 2854.990 2855.990
\hypertarget{c01:cue:1541}{}在這個。
% cue 1542 2855.990 2858.020
\hypertarget{c01:cue:1542}{}權力分立新論的卷一吧。
% cue 1543 2858.020 2859.080
\hypertarget{c01:cue:1543}{}有一篇專門的文章。
% cue 1544 2859.080 2860.080
\hypertarget{c01:cue:1544}{}就是在評估而已。
% cue 1545 2860.080 2862.240
\hypertarget{c01:cue:1545}{}跟權力分立的關係。

\tc{0:47:42}
% cue 1546 2862.240 2864.490
\hypertarget{c01:cue:1546}{}然後 11 月 16 號。
% cue 1547 2864.490 2865.490
\hypertarget{c01:cue:1547}{}585\footnote{司法院釋字第 585 號解釋（2004 年 12 月 15 日）審查 2004 年 9 月 24 日公布施行的三一九槍擊事件真相調查特別委員會條例（真調會條例），宣告禁止被調查人出境等多項規定有違憲法意旨，自解釋公布之日起失其效力。\url{https://cons.judicial.gov.tw/docdata.aspx?fid=100&id=310766}。}。
% cue 1548 2865.490 2866.490
\hypertarget{c01:cue:1548}{}633\footnote{司法院釋字第 633 號解釋（2007 年 9 月 28 日）審查 2006 年 5 月 1 日修正公布的真調會條例，宣告第 8 條之 2 第 3 項關於罰鍰部分及第 4 項與釋字第 585 號解釋意旨不符，第 11 條第 3 項與權力分立制衡原則不符，均自解釋公布之日起失其效力。\url{https://cons.judicial.gov.tw/docdata.aspx?fid=100&id=310814}。}。
% cue 1549 2866.490 2868.550
\hypertarget{c01:cue:1549}{}585 是真調會第一個案子。
% cue 1550 2868.550 2870.770
\hypertarget{c01:cue:1550}{}宣告違憲以後。
% cue 1551 2870.770 2871.800
\hypertarget{c01:cue:1551}{}他們又修法。
% cue 1552 2871.800 2873.800
\hypertarget{c01:cue:1552}{}修法又被大法官打了第二槍。
% cue 1553 2873.800 2874.830
\hypertarget{c01:cue:1553}{}所以 633。
% cue 1554 2874.830 2876.260
\hypertarget{c01:cue:1554}{}然後 627\footnote{司法院釋字第 627 號解釋（2007 年 6 月 15 日）認憲法第 52 條所定總統之刑事豁免權「為一種暫時性之程序障礙，而非總統就其犯罪行為享有實體之免責權」，並就總統之國家機密特權及其範圍作成解釋。\url{https://cons.judicial.gov.tw/docdata.aspx?fid=100&id=310808}。}。
% cue 1555 2876.260 2878.400
\hypertarget{c01:cue:1555}{}613\footnote{司法院釋字第 613 號解釋（2006 年 7 月 21 日）認國家通訊傳播委員會組織法第 4 條關於委員依政黨（團）在立法院席次比例推薦的選任程序規定，「實質上幾近完全剝奪行政院之人事決定權」，違反責任政治暨權力分立原則。\url{https://cons.judicial.gov.tw/docdata.aspx?fid=100&id=310794}。} 一個人。
% cue 1556 2878.400 2879.400
\hypertarget{c01:cue:1556}{}645\footnote{司法院釋字第 645 號解釋（2008 年 7 月 11 日）認公民投票法第 16 條第 1 項賦予立法院公民投票提案權與憲法尚無牴觸；同法第 35 條第 1 項由各政黨依立法院席次比例推荐行政院公民投票審議委員會委員，「實質上完全剝奪行政院依憲法應享有之人事任命決定權」，牴觸權力分立原則。\url{https://cons.judicial.gov.tw/docdata.aspx?fid=100&id=310826}。}。
% cue 1557 2879.400 2881.940
\hypertarget{c01:cue:1557}{}499\footnote{司法院釋字第 499 號解釋（2000 年 3 月 24 日）審查 1999 年 9 月 15 日修正公布的第五次憲法增修條文，認國民大會以無記名投票三讀通過，「此項修憲行為有明顯重大瑕疵，已違反修憲條文發生效力之基本規範」。\cite{jy499}。} 跟 721\footnote{司法院釋字第 721 號解釋（2014 年 6 月 6 日）：憲法增修條文第 4 條第 1 項及第 2 項關於單一選區兩票制之並立制、政黨比例代表席次及政黨門檻規定部分，「並未違反現行憲法賴以存立之自由民主憲政秩序」。\url{https://cons.judicial.gov.tw/docdata.aspx?fid=100&id=310902}。} 最好是同一組。
% cue 1558 2881.940 2882.940
\hypertarget{c01:cue:1558}{}這兩組。
% cue 1559 2882.940 2886.210
\hypertarget{c01:cue:1559}{}721 跟 190 是密切相關的。
% cue 1560 2886.210 2887.350
\hypertarget{c01:cue:1560}{}325\footnote{司法院釋字第 325 號解釋（1993 年 7 月 23 日）：1992 年憲法增修後監察院已非中央民意機構，原屬監察院職權中之彈劾、糾舉、糾正權及為行使此等職權之調查權，「仍應專由監察院行使」；立法院為行使職權，得經院會或委員會之決議，要求有關機關提供參考資料，必要時並得調閱文件原本。\url{https://cons.judicial.gov.tw/docdata.aspx?fid=100&id=310506}。} 跟 729\footnote{司法院釋字第 729 號解釋（2015 年 5 月 1 日）：「基於權力分立與制衡原則，且為保障檢察機關獨立行使職權，對於偵查中之案件，立法院自不得向其調閱相關卷證。」已偵查終結之案件，則須基於目的與範圍均屬明確之特定議案等要件始得調閱。\url{https://cons.judicial.gov.tw/docdata.aspx?fid=100&id=310910}。} 是兩個。

\tc{0:48:07}
% cue 1561 2887.350 2889.350
\hypertarget{c01:cue:1561}{}然後 113 憲判 9\footnote{中華民國憲法法庭 113 年憲判字第 9 號判決（2024 年 10 月 25 日）審查 2024 年 6 月 24 日修正公布的立法院職權行使法修正條文及刑法第 141 條之 1，主文分就聽取總統國情報告、質詢、人事同意權、調查權之行使、聽證會之舉行各部分作成宣告，並宣告刑法第 141 條之 1 違憲，自判決宣示之日起失其效力。\cite{ccj113-9}。}。
% cue 1562 2889.350 2891.380
\hypertarget{c01:cue:1562}{}老師 youtube 上面有一集。
% cue 1563 2891.380 2892.660
\hypertarget{c01:cue:1563}{}有看過這集的舉手。
% cue 1564 2892.660 2894.810
\hypertarget{c01:cue:1564}{}這集舉手有十萬人。
% cue 1565 2894.810 2895.840
\hypertarget{c01:cue:1565}{}十幾萬。
% cue 1566 2895.840 2896.840
\hypertarget{c01:cue:1566}{}十二萬。
% cue 1567 2896.840 2899.410
\hypertarget{c01:cue:1567}{}沒有看到這集。
% cue 1568 2899.410 2900.590
\hypertarget{c01:cue:1568}{}113 憲判 9叫什麼。
% cue 1569 2900.590 2903.660
\hypertarget{c01:cue:1569}{}13 憲判 9。
% cue 1570 2903.660 2905.210
\hypertarget{c01:cue:1570}{}13 憲判 9。
% cue 1571 2905.210 2906.270
\hypertarget{c01:cue:1571}{}連攻聲。
% cue 1572 2906.270 2908.070
\hypertarget{c01:cue:1572}{}有些人。
% cue 1573 2908.070 2909.840
\hypertarget{c01:cue:1573}{}有些人。
% cue 1574 2909.840 2910.810
\hypertarget{c01:cue:1574}{}有些人。
% cue 1575 2910.810 2912.350
\hypertarget{c01:cue:1575}{}有些人。
% cue 1576 2912.350 2913.920
\hypertarget{c01:cue:1576}{}有些人。
% cue 1577 2913.920 2914.920
\hypertarget{c01:cue:1577}{}沒有人。
% cue 1578 2914.920 2915.920
\hypertarget{c01:cue:1578}{}沒有人。

\tc{0:48:35}
\speaker{學生}
% cue 1579 2915.920 2921.620
\hypertarget{c01:cue:1579}{}立法院。

\tc{0:48:41}
\speaker{老師}
% cue 1580 2921.620 2923.670
\hypertarget{c01:cue:1580}{}什麼。
% cue 1581 2923.670 2926.680
\hypertarget{c01:cue:1581}{}只有這麼簡單嗎。
% cue 1582 2926.680 2927.680
\hypertarget{c01:cue:1582}{}太有複雜。
% cue 1583 2927.680 2928.680
\hypertarget{c01:cue:1583}{}所以你比較好的講法。
% cue 1584 2928.680 2929.680
\hypertarget{c01:cue:1584}{}再給你一次機會。
% cue 1585 2929.680 2930.800
\hypertarget{c01:cue:1585}{}你要講什麼。
% cue 1586 2930.800 2931.800
\hypertarget{c01:cue:1586}{}這是有關於什麼東西。
% cue 1587 2931.800 2933.760
\hypertarget{c01:cue:1587}{}是否和諧的安全。

\tc{0:48:53}
\speaker{學生}
% cue 1588 2933.760 2937.670
\hypertarget{c01:cue:1588}{}立法院。
% cue 1589 2937.670 2938.670
\hypertarget{c01:cue:1589}{}職權行使。

\tc{0:48:58}
\speaker{老師}
% cue 1590 2938.670 2941.550
\hypertarget{c01:cue:1590}{}懂嗎。

\tc{0:49:01}
% cue 1591 2941.550 2946.660
\hypertarget{c01:cue:1591}{}自從 2024 年年底。
% cue 1592 2946.660 2950.950
\hypertarget{c01:cue:1592}{}是嗎。
% cue 1593 2950.950 2952.430
\hypertarget{c01:cue:1593}{}改選嗎。
% cue 1594 2952.430 2953.430
\hypertarget{c01:cue:1594}{}國民黨加民眾黨。
% cue 1595 2953.430 2954.430
\hypertarget{c01:cue:1595}{}在立法院。
% cue 1596 2954.430 2955.780
\hypertarget{c01:cue:1596}{}都說這時候。
% cue 1597 2955.780 2956.780
\hypertarget{c01:cue:1597}{}你一定會再有不能。
% cue 1598 2956.780 2958.000
\hypertarget{c01:cue:1598}{}我要了。
% cue 1599 2958.000 2958.910
\hypertarget{c01:cue:1599}{}誰敢打。
% cue 1600 2958.910 2960.060
\hypertarget{c01:cue:1600}{}我不要了。
% cue 1601 2960.060 2961.250
\hypertarget{c01:cue:1601}{}誰敢過。
% cue 1602 2961.250 2962.350
\hypertarget{c01:cue:1602}{}以前過去罷免。
% cue 1603 2962.350 2964.940
\hypertarget{c01:cue:1603}{}蔡英文實在是這樣的。
% cue 1604 2964.940 2966.000
\hypertarget{c01:cue:1604}{}但是呢。
% cue 1605 2966.000 2968.530
\hypertarget{c01:cue:1605}{}賴清德命不好。
% cue 1606 2968.530 2969.530
\hypertarget{c01:cue:1606}{}賴清安定隊的就是。
% cue 1607 2969.530 2970.530
\hypertarget{c01:cue:1607}{}強大的後衛。
% cue 1608 2970.530 2971.530
\hypertarget{c01:cue:1608}{}他的頭痛。
% cue 1609 2971.530 2972.530
\hypertarget{c01:cue:1609}{}一切的頭痛。
% cue 1610 2972.530 2973.530
\hypertarget{c01:cue:1610}{}都是來自於這個後衛。

\tc{0:49:33}
% cue 1611 2973.530 2974.810
\hypertarget{c01:cue:1611}{}所以近不近的。
% cue 1612 2974.810 2975.810
\hypertarget{c01:cue:1612}{}去年一整年。
% cue 1613 2975.810 2976.810
\hypertarget{c01:cue:1613}{}要幹嘛。
% cue 1614 2976.810 2978.600
\hypertarget{c01:cue:1614}{}要搞大罷免\footnote{2025 年 7 月 26 日舉行 24 名國民黨籍立法委員與新竹市長高虹安的罷免投票，8 月 23 日舉行另 7 名國民黨籍立法委員的罷免投票，依中央選舉委員會開票結果均未通過。\url{https://news.pts.org.tw/hotTopic/524}。}。
% cue 1615 2978.600 2979.600
\hypertarget{c01:cue:1615}{}怎麼辦呢。
% cue 1616 2979.600 2980.600
\hypertarget{c01:cue:1616}{}這個人太討厭啊。
% cue 1617 2980.600 2981.940
\hypertarget{c01:cue:1617}{}他有沒有幹掉。
% cue 1618 2981.940 2983.250
\hypertarget{c01:cue:1618}{}罷免掉。
% cue 1619 2983.250 2984.250
\hypertarget{c01:cue:1619}{}如果罷免幾席。
% cue 1620 2984.250 2985.250
\hypertarget{c01:cue:1620}{}那我就會增加。
% cue 1621 2985.250 2986.250
\hypertarget{c01:cue:1621}{}我會幾席。
% cue 1622 2986.250 2987.660
\hypertarget{c01:cue:1622}{}對不對。
% cue 1623 2987.660 2987.760
\hypertarget{c01:cue:1623}{}好。
% cue 1624 2987.660 2988.660
\hypertarget{c01:cue:1624}{}那沒有之勾。

\tc{0:49:48}
% cue 1625 2988.660 2989.750
\hypertarget{c01:cue:1625}{}那這個從這裡開始。
% cue 1626 2989.750 2990.750
\hypertarget{c01:cue:1626}{}就是因為。
% cue 1627 2990.750 2993.180
\hypertarget{c01:cue:1627}{}國會醒來了。
% cue 1628 2993.180 2994.690
\hypertarget{c01:cue:1628}{}國會發現。
% cue 1629 2994.690 2995.880
\hypertarget{c01:cue:1629}{}自己的權力。
% cue 1630 2995.880 2996.880
\hypertarget{c01:cue:1630}{}被閹割了。
% cue 1631 2996.880 2997.880
\hypertarget{c01:cue:1631}{}差不多了。
% cue 1632 2997.880 2998.880
\hypertarget{c01:cue:1632}{}所以他們就把。
% cue 1633 2998.880 2999.880
\hypertarget{c01:cue:1633}{}那個法律改了。
% cue 1634 2999.880 3000.880
\hypertarget{c01:cue:1634}{}原來就有。
% cue 1635 3000.880 3001.880
\hypertarget{c01:cue:1635}{}叫做。
% cue 1636 3001.880 3003.880
\hypertarget{c01:cue:1636}{}國立法院職權行使法。
% cue 1637 3003.880 3005.130
\hypertarget{c01:cue:1637}{}立法院怎麼行使。
% cue 1638 3005.130 3007.150
\hypertarget{c01:cue:1638}{}簡單的講。
% cue 1639 3007.150 3008.150
\hypertarget{c01:cue:1639}{}很多職權都被閹割掉。
% cue 1640 3008.150 3010.150
\hypertarget{c01:cue:1640}{}沒有辦法行使。

\tc{0:50:10}
% cue 1641 3010.150 3011.660
\hypertarget{c01:cue:1641}{}所以他們就修法。
% cue 1642 3011.660 3012.720
\hypertarget{c01:cue:1642}{}把立法院職權行使法。
% cue 1643 3012.720 3013.720
\hypertarget{c01:cue:1643}{}做了修正。
% cue 1644 3013.720 3017.300
\hypertarget{c01:cue:1644}{}那你覺得民進黨是數的嗎。
% cue 1645 3017.300 3018.970
\hypertarget{c01:cue:1645}{}你修。
% cue 1646 3018.970 3020.000
\hypertarget{c01:cue:1646}{}你儘管修。
% cue 1647 3020.000 3021.090
\hypertarget{c01:cue:1647}{}我聲請我們的哥們。
% cue 1648 3021.090 3023.090
\hypertarget{c01:cue:1648}{}大法官來解釋一下。
% cue 1649 3023.090 3025.110
\hypertarget{c01:cue:1649}{}快。
% cue 1650 3025.110 3026.180
\hypertarget{c01:cue:1650}{}大法官馬上就做成解釋。
% cue 1651 3026.180 3027.780
\hypertarget{c01:cue:1651}{}而這個解釋。
% cue 1652 3027.780 3028.780
\hypertarget{c01:cue:1652}{}依依三線。
% cue 1653 3028.780 3031.030
\hypertarget{c01:cue:1653}{}態級幾乎幾乎一句查。
% cue 1654 3031.030 3035.350
\hypertarget{c01:cue:1654}{}幾乎所有的修改條文都會寫。
% cue 1655 3035.350 3039.520
\hypertarget{c01:cue:1655}{}你剛剛講說立法院的調查權\footnote{司法院釋字第 585 號解釋：「立法院為有效行使憲法所賦予之立法職權，本其固有之權能自得享有一定之調查權」；此項調查權「乃立法院行使其憲法職權所必要之輔助性權力」，所得調查之對象或事項並非毫無限制。\url{https://cons.judicial.gov.tw/docdata.aspx?fid=100&id=310766}。}。
% cue 1656 3039.520 3041.520
\hypertarget{c01:cue:1656}{}調查權只是立法院行使的。
% cue 1657 3041.520 3043.520
\hypertarget{c01:cue:1657}{}諸多職權中的一種。
% cue 1658 3043.520 3044.960
\hypertarget{c01:cue:1658}{}就是這個。

\tc{0:50:44}
% cue 1659 3044.960 3047.620
\hypertarget{c01:cue:1659}{}所以我剛才跟你講。
% cue 1660 3047.620 3048.650
\hypertarget{c01:cue:1660}{}再想一下。
% cue 1661 3048.650 3048.750
\hypertarget{c01:cue:1661}{}好了。

\tc{0:50:48}
% cue 1662 3048.650 3051.260
\hypertarget{c01:cue:1662}{}然後中央地方權限劃分\footnote{中華民國憲法第十章「中央與地方之權限」，第 107 條至第 111 條。課綱第 16 週（12 月 28 日）以此為題，列釋字第 498、738、769 號解釋及 111 年憲判字第 6 號判決。\url{https://law.moj.gov.tw/LawClass/LawAll.aspx?pcode=A0000001}。}。
% cue 1663 3051.260 3053.480
\hypertarget{c01:cue:1663}{}這裡有四個案子。
% cue 1664 3053.480 3054.510
\hypertarget{c01:cue:1664}{}你們看要分析。
% cue 1665 3054.510 3055.890
\hypertarget{c01:cue:1665}{}最後呢。
% cue 1666 3055.890 3057.890
\hypertarget{c01:cue:1666}{}一月四號就是期末考。
% cue 1667 3057.890 3060.150
\hypertarget{c01:cue:1667}{}12 月 28 號的那個。
% cue 1668 3060.150 3061.150
\hypertarget{c01:cue:1668}{}是一定會考慮。
% cue 1669 3061.150 3062.240
\hypertarget{c01:cue:1669}{}但是我們抱歉。
% cue 1670 3062.240 3063.240
\hypertarget{c01:cue:1670}{}是講到這次。

\tc{0:51:03}
% cue 1671 3063.240 3065.490
\hypertarget{c01:cue:1671}{}然後你發現旁邊有出記。
% cue 1672 3065.490 3066.490
\hypertarget{c01:cue:1672}{}有沒有。
% cue 1673 3066.490 3067.490
\hypertarget{c01:cue:1673}{}心路 627。
% cue 1674 3067.490 3068.490
\hypertarget{c01:cue:1674}{}心路 1019。
% cue 1675 3068.490 3069.550
\hypertarget{c01:cue:1675}{}就是在老師。
% cue 1676 3069.550 3071.550
\hypertarget{c01:cue:1676}{}釋憲學思心路當中。
% cue 1677 3071.550 3072.550
\hypertarget{c01:cue:1677}{}第幾頁。
% cue 1678 3072.550 3073.550
\hypertarget{c01:cue:1678}{}你會看到這個東西。
% cue 1679 3073.550 3075.180
\hypertarget{c01:cue:1679}{}我都盡量幫你。
% cue 1680 3075.180 3076.170
\hypertarget{c01:cue:1680}{}都做好了。
% cue 1681 3076.170 3078.230
\hypertarget{c01:cue:1681}{}113 憲判 9的旁邊。
% cue 1682 3078.230 3079.230
\hypertarget{c01:cue:1682}{}是不是有寫 UYT。
% cue 1683 3079.230 3081.230
\hypertarget{c01:cue:1683}{}YT 就是 Youtube。
% cue 1684 3081.230 3082.230
\hypertarget{c01:cue:1684}{}你在 Youtube 上面。
% cue 1685 3082.230 3083.230
\hypertarget{c01:cue:1685}{}也看到這個情況。
% cue 1686 3083.230 3084.320
\hypertarget{c01:cue:1686}{}所以你們分到這一組的人。
% cue 1687 3084.320 3085.320
\hypertarget{c01:cue:1687}{}要報告。
% cue 1688 3085.320 3086.280
\hypertarget{c01:cue:1688}{}沒有那麼難。
% cue 1689 3086.280 3087.820
\hypertarget{c01:cue:1689}{}老師在不同地方。
% cue 1690 3087.820 3088.820
\hypertarget{c01:cue:1690}{}都已經做了提示了。

\tc{0:51:28}
% cue 1691 3088.820 3091.330
\hypertarget{c01:cue:1691}{}那借 1 借幾借幾。
% cue 1692 3091.330 3092.330
\hypertarget{c01:cue:1692}{}現在這些號碼。
% cue 1693 3092.330 3093.330
\hypertarget{c01:cue:1693}{}也許又有一點修改。
% cue 1694 3093.330 3095.390
\hypertarget{c01:cue:1694}{}因為這是舊版的。
% cue 1695 3095.390 3096.450
\hypertarget{c01:cue:1695}{}四版\footnote{課綱所列指定教科書為增訂五版（2026），先前各版頁碼不同。《權力分立新論 卷一》的〈新世紀憲改工程〉一文，增訂三版（元照，2005）為頁 59 起，增訂四版（2014）為頁 53 起；課綱第 4 週所列「卷一／53」與後者相同。}。
% cue 1696 3096.450 3097.450
\hypertarget{c01:cue:1696}{}現在要出一個是五版。
% cue 1697 3097.450 3098.670
\hypertarget{c01:cue:1697}{}那五版跟四版呢。
% cue 1698 3098.670 3100.980
\hypertarget{c01:cue:1698}{}老師每一次出書。
% cue 1699 3100.980 3102.520
\hypertarget{c01:cue:1699}{}一定負責任。
% cue 1700 3102.520 3103.520
\hypertarget{c01:cue:1700}{}從讀改到我。
% cue 1701 3103.520 3105.930
\hypertarget{c01:cue:1701}{}沒有說我不看過可以出去了。
% cue 1702 3105.930 3108.050
\hypertarget{c01:cue:1702}{}那這次我做了很多修正。
% cue 1703 3108.050 3110.140
\hypertarget{c01:cue:1703}{}因為時間不一樣。
% cue 1704 3110.140 3112.190
\hypertarget{c01:cue:1704}{}四版的時候是 2014 年。
% cue 1705 3112.190 3113.880
\hypertarget{c01:cue:1705}{}很久很久以前。
% cue 1706 3113.880 3115.740
\hypertarget{c01:cue:1706}{}現在是 2026 年。
% cue 1707 3115.740 3118.150
\hypertarget{c01:cue:1707}{}有很多新案子。
% cue 1708 3118.150 3119.150
\hypertarget{c01:cue:1708}{}進來順便講一下。

\tc{0:51:59}
% cue 1709 3119.150 3121.610
\hypertarget{c01:cue:1709}{}好。
% cue 1710 3121.610 3122.600
\hypertarget{c01:cue:1710}{}那等一下班長。
% cue 1711 3122.600 3124.920
\hypertarget{c01:cue:1711}{}你現在都知道要怎麼分了。
% cue 1712 3124.920 3126.070
\hypertarget{c01:cue:1712}{}就問大家。
% cue 1713 3126.070 3127.070
\hypertarget{c01:cue:1713}{}就來重新登記。
% cue 1714 3127.070 3128.360
\hypertarget{c01:cue:1714}{}你要登記了。
% cue 1715 3128.360 3129.360
\hypertarget{c01:cue:1715}{}不然到了。
% cue 1716 3129.360 3131.330
\hypertarget{c01:cue:1716}{}那時候就沒辦法開始。
% cue 1717 3131.330 3132.330
\hypertarget{c01:cue:1717}{}你這次如果沒有做完。
% cue 1718 3132.330 3133.330
\hypertarget{c01:cue:1718}{}也沒關係。
% cue 1719 3133.330 3134.330
\hypertarget{c01:cue:1719}{}下禮拜老師在上次。
% cue 1720 3134.330 3135.330
\hypertarget{c01:cue:1720}{}你還有一次機會。
% cue 1721 3135.330 3137.710
\hypertarget{c01:cue:1721}{}但是就要定下來。

\tc{0:52:17}
% cue 1722 3137.710 3140.930
\hypertarget{c01:cue:1722}{}那如果有人要參加我們這個班。
% cue 1723 3140.930 3141.850
\hypertarget{c01:cue:1723}{}表示你要加選。
% cue 1724 3141.850 3142.890
\hypertarget{c01:cue:1724}{}非常歡迎。
% cue 1725 3142.890 3143.860
\hypertarget{c01:cue:1725}{}你要去拉人。
% cue 1726 3143.860 3144.860
\hypertarget{c01:cue:1726}{}也非常歡迎。
% cue 1727 3144.860 3146.950
\hypertarget{c01:cue:1727}{}但是你要記得。
% cue 1728 3146.950 3149.010
\hypertarget{c01:cue:1728}{}如果你沒有參加報告。
% cue 1729 3149.010 3150.070
\hypertarget{c01:cue:1729}{}這個報告欄位。
% cue 1730 3150.070 3152.070
\hypertarget{c01:cue:1730}{}你這個分數就沒有了。

\section{成績配分、出缺席與分組報告}

\tc{0:52:32}
% cue 1731 3152.070 3153.810
\hypertarget{c01:cue:1731}{}那先讓大家看一下。
% cue 1732 3153.810 3154.810
\hypertarget{c01:cue:1732}{}怎樣的評分。
% cue 1733 3154.810 3157.640
\hypertarget{c01:cue:1733}{}看一下下頁。
% cue 1734 3157.640 3159.020
\hypertarget{c01:cue:1734}{}考評的方式。
% cue 1735 3159.020 3160.480
\hypertarget{c01:cue:1735}{}考試的時候。
% cue 1736 3160.480 3162.620
\hypertarget{c01:cue:1736}{}老師期望考分數是 100 分。
% cue 1737 3162.620 3163.620
\hypertarget{c01:cue:1737}{}但是呢。
% cue 1738 3163.620 3167.240
\hypertarget{c01:cue:1738}{}出席跟課堂表現。
% cue 1739 3167.240 3168.240
\hypertarget{c01:cue:1739}{}還各有 10 分。
% cue 1740 3168.240 3170.660
\hypertarget{c01:cue:1740}{}所以你不要來跟老師爭。
% cue 1741 3170.660 3171.610
\hypertarget{c01:cue:1741}{}說老師。
% cue 1742 3171.610 3173.630
\hypertarget{c01:cue:1742}{}我怎麼沒有平時成績。
% cue 1743 3173.630 3175.630
\hypertarget{c01:cue:1743}{}我加上去怎麼會怎樣怎樣。
% cue 1744 3175.630 3177.770
\hypertarget{c01:cue:1744}{}我事實上不能寫 120。
% cue 1745 3177.770 3178.770
\hypertarget{c01:cue:1745}{}因為學校的老師。
% cue 1746 3178.770 3179.770
\hypertarget{c01:cue:1746}{}別的老師不來找我孫長。
% cue 1747 3179.770 3180.730
\hypertarget{c01:cue:1747}{}說沒有。
% cue 1748 3180.730 3183.030
\hypertarget{c01:cue:1748}{}他老師在會弱學生嗎。
% cue 1749 3183.030 3184.090
\hypertarget{c01:cue:1749}{}你為什麼 100。
% cue 1750 3184.090 3185.090
\hypertarget{c01:cue:1750}{}我們都 120 分。
% cue 1751 3185.090 3187.090
\hypertarget{c01:cue:1751}{}所以大家都要上一課。
% cue 1752 3187.090 3188.540
\hypertarget{c01:cue:1752}{}我也是這個意思。

\tc{0:53:08}
% cue 1753 3188.540 3190.830
\hypertarget{c01:cue:1753}{}所以我在講裡面。
% cue 1754 3190.830 3191.830
\hypertarget{c01:cue:1754}{}期末考試。
% cue 1755 3191.830 3193.860
\hypertarget{c01:cue:1755}{}他現在配分 80\footnote{課綱考評配分為出席百分之十、期末考試百分之八十、課堂討論百分之十。}。
% cue 1756 3193.860 3196.300
\hypertarget{c01:cue:1756}{}這是大個率。
% cue 1757 3196.300 3197.300
\hypertarget{c01:cue:1757}{}我的意思是說。
% cue 1758 3197.300 3199.330
\hypertarget{c01:cue:1758}{}考試成績仍然是 100。
% cue 1759 3199.330 3202.160
\hypertarget{c01:cue:1759}{}那在這個 100 做清楚。
% cue 1760 3202.160 3204.220
\hypertarget{c01:cue:1760}{}你去加出席。
% cue 1761 3204.220 3205.310
\hypertarget{c01:cue:1761}{}去加課堂。
% cue 1762 3205.310 3206.850
\hypertarget{c01:cue:1762}{}課堂表現當中。
% cue 1763 3206.850 3208.850
\hypertarget{c01:cue:1763}{}一定會有一個成績。
% cue 1764 3208.850 3210.780
\hypertarget{c01:cue:1764}{}是你們的報告成績。
% cue 1765 3210.780 3212.000
\hypertarget{c01:cue:1765}{}就是我們分數。
% cue 1766 3212.000 3213.190
\hypertarget{c01:cue:1766}{}所以如果你沒有分在組裡面。
% cue 1767 3213.190 3215.190
\hypertarget{c01:cue:1767}{}你沒有這個分數報告。
% cue 1768 3215.190 3217.190
\hypertarget{c01:cue:1768}{}你的課堂成績的 10 分。
% cue 1769 3217.190 3218.190
\hypertarget{c01:cue:1769}{}你是加不到的。

\tc{0:53:38}
% cue 1770 3218.190 3221.150
\hypertarget{c01:cue:1770}{}通常我不會扣分。
% cue 1771 3221.150 3222.430
\hypertarget{c01:cue:1771}{}但是有一種情況。
% cue 1772 3222.430 3223.330
\hypertarget{c01:cue:1772}{}我不懷恥。
% cue 1773 3223.330 3224.480
\hypertarget{c01:cue:1773}{}比方說這一組。
% cue 1774 3224.480 3225.760
\hypertarget{c01:cue:1774}{}分到你。
% cue 1775 3225.760 3226.760
\hypertarget{c01:cue:1775}{}如果這一組今天給我們放鴿子。
% cue 1776 3226.760 3227.760
\hypertarget{c01:cue:1776}{}開對床。
% cue 1777 3227.760 3228.940
\hypertarget{c01:cue:1777}{}那只要扣分。
% cue 1778 3228.940 3230.640
\hypertarget{c01:cue:1778}{}太平負責任。
% cue 1779 3230.640 3231.870
\hypertarget{c01:cue:1779}{}全班等你。
% cue 1780 3231.870 3233.780
\hypertarget{c01:cue:1780}{}不來。
% cue 1781 3233.780 3235.320
\hypertarget{c01:cue:1781}{}只要要半過分。
% cue 1782 3235.320 3236.220
\hypertarget{c01:cue:1782}{}這是不是形容詐欺。
% cue 1783 3236.220 3238.950
\hypertarget{c01:cue:1783}{}對嗎。
% cue 1784 3238.950 3239.950
\hypertarget{c01:cue:1784}{}大家講好的分工。
% cue 1785 3239.950 3240.950
\hypertarget{c01:cue:1785}{}所以我貢獻你的。
% cue 1786 3240.950 3241.910
\hypertarget{c01:cue:1786}{}我貢獻你的。
% cue 1787 3241.910 3242.940
\hypertarget{c01:cue:1787}{}你貢獻你的。
% cue 1788 3242.940 3243.940
\hypertarget{c01:cue:1788}{}到了我貢獻我貢獻完了。
% cue 1789 3243.940 3244.940
\hypertarget{c01:cue:1789}{}該你貢獻你。
% cue 1790 3244.940 3245.940
\hypertarget{c01:cue:1790}{}我貢獻了。
% cue 1791 3245.940 3247.410
\hypertarget{c01:cue:1791}{}這不要過分。

\tc{0:54:07}
% cue 1792 3247.410 3248.420
\hypertarget{c01:cue:1792}{}那出席的話呢。
% cue 1793 3248.420 3249.440
\hypertarget{c01:cue:1793}{}我貢獻還沒有決定。
% cue 1794 3249.440 3250.500
\hypertarget{c01:cue:1794}{}基本上呢。
% cue 1795 3250.500 3251.520
\hypertarget{c01:cue:1795}{}就是全部到的人。
% cue 1796 3251.520 3253.480
\hypertarget{c01:cue:1796}{}一定會取得幾多加分。
% cue 1797 3253.480 3255.480
\hypertarget{c01:cue:1797}{}但是方法很簡單。
% cue 1798 3255.480 3257.480
\hypertarget{c01:cue:1798}{}我講一下一個。
% cue 1799 3257.480 3258.510
\hypertarget{c01:cue:1799}{}我只是沒有決定。
% cue 1800 3258.510 3259.510
\hypertarget{c01:cue:1800}{}要加幾分。
% cue 1801 3259.510 3260.480
\hypertarget{c01:cue:1801}{}我現在寫的。
% cue 1802 3260.480 3261.480
\hypertarget{c01:cue:1802}{}左邊的這個字。
% cue 1803 3261.480 3262.480
\hypertarget{c01:cue:1803}{}你又幾次沒有完。
% cue 1804 3262.480 3264.410
\hypertarget{c01:cue:1804}{}缺席的次數。
% cue 1805 3264.410 3267.390
\hypertarget{c01:cue:1805}{}缺席零次。
% cue 1806 3267.390 3268.390
\hypertarget{c01:cue:1806}{}是不是全勤。
% cue 1807 3268.390 3269.840
\hypertarget{c01:cue:1807}{}那我這裡一定加分。
% cue 1808 3269.840 3270.880
\hypertarget{c01:cue:1808}{}我不知道加幾。
% cue 1809 3270.880 3273.110
\hypertarget{c01:cue:1809}{}我們估計寫加四。

\tc{0:54:33}
% cue 1810 3273.110 3275.330
\hypertarget{c01:cue:1810}{}為什麼我們說要估計。
% cue 1811 3275.330 3277.450
\hypertarget{c01:cue:1811}{}因為我給我看卷子改完。
% cue 1812 3277.450 3279.960
\hypertarget{c01:cue:1812}{}卷子最高分 94。
% cue 1813 3279.960 3281.020
\hypertarget{c01:cue:1813}{}我還加它四分。
% cue 1814 3281.020 3282.020
\hypertarget{c01:cue:1814}{}這個卷分 100。
% cue 1815 3282.020 3282.980
\hypertarget{c01:cue:1815}{}這只通嗎。
% cue 1816 3282.980 3283.980
\hypertarget{c01:cue:1816}{}那這時候。
% cue 1817 3283.980 3285.550
\hypertarget{c01:cue:1817}{}我們可能就要下降。
% cue 1818 3285.550 3287.870
\hypertarget{c01:cue:1818}{}那如果現在大家都不及格。
% cue 1819 3287.870 3288.870
\hypertarget{c01:cue:1819}{}那這分數就要多高一點。
% cue 1820 3288.870 3290.090
\hypertarget{c01:cue:1820}{}這樣聽懂嗎。

\tc{0:54:50}
% cue 1821 3290.090 3292.090
\hypertarget{c01:cue:1821}{}你缺席一次。
% cue 1822 3292.090 3294.090
\hypertarget{c01:cue:1822}{}那我們就加這個一半。
% cue 1823 3294.090 3295.660
\hypertarget{c01:cue:1823}{}不管它是多少。
% cue 1824 3295.660 3296.680
\hypertarget{c01:cue:1824}{}那一半。
% cue 1825 3296.680 3298.200
\hypertarget{c01:cue:1825}{}如果你缺席兩次。
% cue 1826 3298.200 3300.320
\hypertarget{c01:cue:1826}{}那就是不加減加零。
% cue 1827 3300.320 3301.920
\hypertarget{c01:cue:1827}{}如果你缺席到。
% cue 1828 3301.920 3302.880
\hypertarget{c01:cue:1828}{}這兩次。
% cue 1829 3302.880 3304.610
\hypertarget{c01:cue:1829}{}你缺席了三次。
% cue 1830 3304.610 3306.610
\hypertarget{c01:cue:1830}{}那就要開始扣分了。
% cue 1831 3306.610 3307.960
\hypertarget{c01:cue:1831}{}扣一分。
% cue 1832 3307.960 3309.150
\hypertarget{c01:cue:1832}{}四次呢。
% cue 1833 3309.150 3310.280
\hypertarget{c01:cue:1833}{}扣兩分。
% cue 1834 3310.280 3311.440
\hypertarget{c01:cue:1834}{}五次呢。
% cue 1835 3311.440 3313.430
\hypertarget{c01:cue:1835}{}扣四分。
% cue 1836 3313.430 3314.620
\hypertarget{c01:cue:1836}{}六次的話呢。
% cue 1837 3314.620 3315.720
\hypertarget{c01:cue:1837}{}扣八分。
% cue 1838 3315.720 3317.940
\hypertarget{c01:cue:1838}{}因為它總共只有十分。
% cue 1839 3317.940 3318.890
\hypertarget{c01:cue:1839}{}扣到八分。
% cue 1840 3318.890 3319.900
\hypertarget{c01:cue:1840}{}我們再扣。
% cue 1841 3319.900 3322.270
\hypertarget{c01:cue:1841}{}這樣有過欣甫嗎。

\tc{0:55:22}
% cue 1842 3322.270 3323.680
\hypertarget{c01:cue:1842}{}好。
% cue 1843 3323.680 3324.680
\hypertarget{c01:cue:1843}{}那班長要做另外一件事。
% cue 1844 3324.680 3327.220
\hypertarget{c01:cue:1844}{}以前上個學期。
% cue 1845 3327.220 3328.220
\hypertarget{c01:cue:1845}{}我看的同學。
% cue 1846 3328.220 3329.220
\hypertarget{c01:cue:1846}{}都是這個班長。
% cue 1847 3329.220 3330.600
\hypertarget{c01:cue:1847}{}都這麼點。
% cue 1848 3330.600 3331.660
\hypertarget{c01:cue:1848}{}同學。
% cue 1849 3331.660 3332.660
\hypertarget{c01:cue:1849}{}我們今天點。
% cue 1850 3332.660 3334.200
\hypertarget{c01:cue:1850}{}我給你十次點。
% cue 1851 3334.200 3335.670
\hypertarget{c01:cue:1851}{}我們這一點上課十。
% cue 1852 3335.670 3337.250
\hypertarget{c01:cue:1852}{}你只要點十次。
% cue 1853 3337.250 3338.380
\hypertarget{c01:cue:1853}{}你自己決定什麼時候要點。
% cue 1854 3338.380 3339.420
\hypertarget{c01:cue:1854}{}好。
% cue 1855 3339.420 3340.480
\hypertarget{c01:cue:1855}{}那要點的時候。
% cue 1856 3340.480 3341.480
\hypertarget{c01:cue:1856}{}下課了。
% cue 1857 3341.480 3342.480
\hypertarget{c01:cue:1857}{}現在下課了。
% cue 1858 3342.480 3343.480
\hypertarget{c01:cue:1858}{}我們現在在座的人。
% cue 1859 3343.480 3344.480
\hypertarget{c01:cue:1859}{}我們要點名。
% cue 1860 3344.480 3345.510
\hypertarget{c01:cue:1860}{}然後你登到公司。
% cue 1861 3345.510 3346.510
\hypertarget{c01:cue:1861}{}再點一個什麼。
% cue 1862 3346.510 3347.570
\hypertarget{c01:cue:1862}{}Broadcast 還是什麼。
% cue 1863 3347.570 3348.530
\hypertarget{c01:cue:1863}{}那打打勾。
% cue 1864 3348.530 3349.530
\hypertarget{c01:cue:1864}{}打勾。
% cue 1865 3349.530 3350.750
\hypertarget{c01:cue:1865}{}所以有來就勾。
% cue 1866 3350.750 3352.130
\hypertarget{c01:cue:1866}{}沒有來就沒有勾。
% cue 1867 3352.130 3355.310
\hypertarget{c01:cue:1867}{}自己決定。

\tc{0:55:55}
% cue 1868 3355.310 3355.410
\hypertarget{c01:cue:1868}{}OK。
% cue 1869 3355.310 3356.280
\hypertarget{c01:cue:1869}{}那這個是可做公平的。
% cue 1870 3356.280 3358.280
\hypertarget{c01:cue:1870}{}這個東西最後要公開。
% cue 1871 3358.280 3360.440
\hypertarget{c01:cue:1871}{}老師不會去管這個同學的事情。
% cue 1872 3360.440 3363.900
\hypertarget{c01:cue:1872}{}大家對於出缺席跟同學。
% cue 1873 3363.900 3364.900
\hypertarget{c01:cue:1873}{}因為計算到分數。
% cue 1874 3364.900 3367.060
\hypertarget{c01:cue:1874}{}你可能會覺得你有來。
% cue 1875 3367.060 3368.150
\hypertarget{c01:cue:1875}{}但是你沒有被點到。
% cue 1876 3368.150 3370.150
\hypertarget{c01:cue:1876}{}或者你明明沒有來就要點到。
% cue 1877 3370.150 3372.150
\hypertarget{c01:cue:1877}{}那你大概不會講。
% cue 1878 3372.150 3373.240
\hypertarget{c01:cue:1878}{}沒有關係。
% cue 1879 3373.240 3374.240
\hypertarget{c01:cue:1879}{}這跟班長講就好。

\tc{0:56:14}
% cue 1880 3374.240 3376.430
\hypertarget{c01:cue:1880}{}然後分組要告訴我。
% cue 1881 3376.430 3377.590
\hypertarget{c01:cue:1881}{}兩個禮拜以後。
% cue 1882 3377.590 3378.590
\hypertarget{c01:cue:1882}{}要告訴我分組是誰。
% cue 1883 3378.590 3380.770
\hypertarget{c01:cue:1883}{}每一組的一個負責人。
% cue 1884 3380.770 3381.730
\hypertarget{c01:cue:1884}{}然後下面是誰。
% cue 1885 3381.730 3382.760
\hypertarget{c01:cue:1885}{}好不好。
% cue 1886 3382.760 3383.820
\hypertarget{c01:cue:1886}{}那到時時間到了。
% cue 1887 3383.820 3384.820
\hypertarget{c01:cue:1887}{}就要進行他們。
% cue 1888 3384.820 3385.820
\hypertarget{c01:cue:1888}{}摩摩。
% cue 1889 3385.820 3386.820
\hypertarget{c01:cue:1889}{}下禮拜該你喔。
% cue 1890 3386.820 3386.920
\hypertarget{c01:cue:1890}{}等一下。
% cue 1891 3386.820 3388.100
\hypertarget{c01:cue:1891}{}不要說都不知道。

\tc{0:56:28}
% cue 1892 3388.100 3390.100
\hypertarget{c01:cue:1892}{}那你們報告要怎麼報告。
% cue 1893 3390.100 3392.100
\hypertarget{c01:cue:1893}{}下禮拜我會給一個格式。
% cue 1894 3392.100 3393.160
\hypertarget{c01:cue:1894}{}供參考。
% cue 1895 3393.160 3396.720
\hypertarget{c01:cue:1895}{}總之老師是用對話的方式給你們講的。
% cue 1896 3396.720 3398.890
\hypertarget{c01:cue:1896}{}所以不一定是完全照一個格式。
% cue 1897 3398.890 3400.850
\hypertarget{c01:cue:1897}{}那個格式有些過沒什麼意思。
% cue 1898 3400.850 3402.170
\hypertarget{c01:cue:1898}{}有的時候很有意思。
% cue 1899 3402.170 3404.840
\hypertarget{c01:cue:1899}{}這個我整整我講的。

\tc{0:56:44}
% cue 1900 3404.840 3405.830
\hypertarget{c01:cue:1900}{}好到這裡講完了。
% cue 1901 3405.830 3409.010
\hypertarget{c01:cue:1901}{}請問有什麼問題請舉手。

\tc{0:56:49}
\speaker{學生}
% cue 1902 3409.010 3410.870
\hypertarget{c01:cue:1902}{}答問。
% cue 1903 3410.870 3413.660
\hypertarget{c01:cue:1903}{}一年了。
% cue 1904 3413.660 3416.900
\hypertarget{c01:cue:1904}{}但居然是有相關。

\tc{0:56:56}
\speaker{老師}
% cue 1905 3416.900 3417.900
\hypertarget{c01:cue:1905}{}最少一人。
% cue 1906 3417.900 3418.900
\hypertarget{c01:cue:1906}{}不然自己就沒人了。
% cue 1907 3418.900 3419.870
\hypertarget{c01:cue:1907}{}不能超過五人。
% cue 1908 3419.870 3422.110
\hypertarget{c01:cue:1908}{}為什麼不能超過五人。
% cue 1909 3422.110 3423.310
\hypertarget{c01:cue:1909}{}你們如果偷來。
% cue 1910 3423.310 3424.310
\hypertarget{c01:cue:1910}{}我就要五個人無人。
% cue 1911 3424.310 3425.310
\hypertarget{c01:cue:1911}{}讓我們這班到底很不夠分啊。
% cue 1912 3425.310 3427.910
\hypertarget{c01:cue:1912}{}所以這樣五組五六組都完了。
% cue 1913 3427.910 3430.710
\hypertarget{c01:cue:1913}{}太多人是沒辦法做事。
% cue 1914 3430.710 3432.990
\hypertarget{c01:cue:1914}{}因為知道我們同學都很忙。
% cue 1915 3432.990 3434.090
\hypertarget{c01:cue:1915}{}約在一起商量一下。
% cue 1916 3434.090 3436.500
\hypertarget{c01:cue:1916}{}然後到最後誰報告。
% cue 1917 3436.500 3437.500
\hypertarget{c01:cue:1917}{}就被老師罵。
% cue 1918 3437.500 3438.500
\hypertarget{c01:cue:1918}{}就電電。
% cue 1919 3438.500 3439.500
\hypertarget{c01:cue:1919}{}電到反正都是報告。
% cue 1920 3439.500 3440.690
\hypertarget{c01:cue:1920}{}但如果不講話。
% cue 1921 3440.690 3440.790
\hypertarget{c01:cue:1921}{}人都沒事。
% cue 1922 3440.690 3441.690
\hypertarget{c01:cue:1922}{}這樣子都沒。
% cue 1923 3441.690 3445.240
\hypertarget{c01:cue:1923}{}那我不希望看到這樣。

\tc{0:57:25}
% cue 1924 3445.240 3446.770
\hypertarget{c01:cue:1924}{}以前有同學。
% cue 1925 3446.770 3447.770
\hypertarget{c01:cue:1925}{}會寫信給老師說。
% cue 1926 3447.770 3449.110
\hypertarget{c01:cue:1926}{}老師我要跟你抱怨。
% cue 1927 3449.110 3451.210
\hypertarget{c01:cue:1927}{}我們今天這組四個人。
% cue 1928 3451.210 3453.370
\hypertarget{c01:cue:1928}{}其他三百人都沒有在做事。
% cue 1929 3453.370 3455.200
\hypertarget{c01:cue:1929}{}然後我今天講的。
% cue 1930 3455.200 3456.450
\hypertarget{c01:cue:1930}{}都被疑問。
% cue 1931 3456.450 3457.450
\hypertarget{c01:cue:1931}{}我都被問到。
% cue 1932 3457.450 3458.480
\hypertarget{c01:cue:1932}{}太多很尷尬。
% cue 1933 3458.480 3460.600
\hypertarget{c01:cue:1933}{}那你不管給我交際。
% cue 1934 3460.600 3463.730
\hypertarget{c01:cue:1934}{}請你不能給另外三個人加任何人。
% cue 1935 3463.730 3466.170
\hypertarget{c01:cue:1935}{}所以好不要這樣子。
% cue 1936 3466.170 3467.510
\hypertarget{c01:cue:1936}{}那我沒辦法。
% cue 1937 3467.510 3468.510
\hypertarget{c01:cue:1937}{}那你們要講。
% cue 1938 3468.510 3468.610
\hypertarget{c01:cue:1938}{}當然也是可以講。
% cue 1939 3468.510 3469.510
\hypertarget{c01:cue:1939}{}老師來上帝吧。

\tc{0:57:49}
% cue 1940 3469.510 3471.080
\hypertarget{c01:cue:1940}{}好了。
% cue 1941 3471.080 3472.080
\hypertarget{c01:cue:1941}{}還有什麼問題。
% cue 1942 3472.080 3474.930
\hypertarget{c01:cue:1942}{}走。
% cue 1943 3474.930 3475.930
\hypertarget{c01:cue:1943}{}有什麼問題。
% cue 1944 3475.930 3479.970
\hypertarget{c01:cue:1944}{}沒有。
% cue 1945 3479.970 3481.190
\hypertarget{c01:cue:1945}{}走。

\tc{0:58:01}
% cue 1946 3481.190 3481.290
\hypertarget{c01:cue:1946}{}沒有我們現在做。
% cue 1947 3481.190 3482.190
\hypertarget{c01:cue:1947}{}剛才我已經問過你們了。
% cue 1948 3482.190 3483.190
\hypertarget{c01:cue:1948}{}11 月 18 號。
% cue 1949 3483.190 3485.190
\hypertarget{c01:cue:1949}{}請你們把時間登錄起來。
% cue 1950 3485.190 3486.390
\hypertarget{c01:cue:1950}{}那第二件事情。
% cue 1951 3486.390 3490.040
\hypertarget{c01:cue:1951}{}請你們現在把老師的簡歷打開。
% cue 1952 3490.040 3492.230
\hypertarget{c01:cue:1952}{}在這個老師的解釋你們現在觀看。
% cue 1953 3492.230 3494.230
\hypertarget{c01:cue:1953}{}與其思考問題的同時。
% cue 1954 3494.230 3496.230
\hypertarget{c01:cue:1954}{}我現在要來聽你們的意見。

\section{全班點名與學生發言}

\tc{0:58:16}
% cue 1955 3496.230 3497.190
\hypertarget{c01:cue:1955}{}那現在呢。
% cue 1956 3497.190 3498.740
\hypertarget{c01:cue:1956}{}這可能是我這一次。
% cue 1957 3498.740 3500.960
\hypertarget{c01:cue:1957}{}唯一的這個學期有機會。
% cue 1958 3500.960 3501.960
\hypertarget{c01:cue:1958}{}全班點名一次。
% cue 1959 3501.960 3502.960
\hypertarget{c01:cue:1959}{}那我現在點。
% cue 1960 3502.960 3504.020
\hypertarget{c01:cue:1960}{}如果沒有點到。
% cue 1961 3504.020 3506.020
\hypertarget{c01:cue:1961}{}就表示你沒有在名單上。
% cue 1962 3506.020 3508.020
\hypertarget{c01:cue:1962}{}你自始要知道自己不在名單上。
% cue 1963 3508.020 3509.020
\hypertarget{c01:cue:1963}{}你要去加選擇。
% cue 1964 3509.020 3510.020
\hypertarget{c01:cue:1964}{}你過來。
% cue 1965 3510.020 3511.020
\hypertarget{c01:cue:1965}{}我現在要來聽。

\tc{0:58:31}
% cue 1966 3511.020 3512.070
\hypertarget{c01:cue:1966}{}法五第一。
% cue 1967 3512.070 3513.140
\hypertarget{c01:cue:1967}{}唐聖涵。
% cue 1968 3513.140 3514.080
\hypertarget{c01:cue:1968}{}請舉手。
% cue 1969 3514.080 3515.300
\hypertarget{c01:cue:1969}{}請問一下。
% cue 1970 3515.300 3516.300
\hypertarget{c01:cue:1970}{}有什麼話來跟老師講。
% cue 1971 3516.300 3518.450
\hypertarget{c01:cue:1971}{}或者你要跟學妹講。
% cue 1972 3518.450 3519.540
\hypertarget{c01:cue:1972}{}你現在五年級了。
% cue 1973 3519.540 3520.570
\hypertarget{c01:cue:1973}{}為什麼這麼勇敢。
% cue 1974 3520.570 3521.570
\hypertarget{c01:cue:1974}{}還要學湯老師\footnote{湯德宗，1956 年生，2011 年 10 月 1 日至 2019 年 9 月 30 日任中華民國司法院大法官，由馬英九總統提名；卸任後受東吳大學禮聘為王紹堉先生講座教授。\url{https://cons.judicial.gov.tw/judge.aspx?fid=9&id=10190&class_per_id=40024}。}的歌。
% cue 1975 3521.570 3527.270
\hypertarget{c01:cue:1975}{}真言。
% cue 1976 3527.270 3528.260
\hypertarget{c01:cue:1976}{}現在馬學生不動。
% cue 1977 3528.260 3529.260
\hypertarget{c01:cue:1977}{}還要把他讀起來。
% cue 1978 3529.260 3531.350
\hypertarget{c01:cue:1978}{}變得離開東吳有遺憾。
% cue 1979 3531.350 3533.370
\hypertarget{c01:cue:1979}{}一定要記得。
% cue 1980 3533.370 3534.340
\hypertarget{c01:cue:1980}{}在東吳沒有傷口的湯老師。
% cue 1981 3534.340 3536.340
\hypertarget{c01:cue:1981}{}可不要出來過東吳。
% cue 1982 3536.340 3537.720
\hypertarget{c01:cue:1982}{}好嗎。

\tc{0:58:57}
% cue 1983 3537.720 3538.880
\hypertarget{c01:cue:1983}{}法五第一。
% cue 1984 3538.880 3541.230
\hypertarget{c01:cue:1984}{}陽光會。
% cue 1985 3541.230 3546.550
\hypertarget{c01:cue:1985}{}或者你想要通過講師。
% cue 1986 3546.550 3547.960
\hypertarget{c01:cue:1986}{}這是你。
% cue 1987 3547.960 3549.920
\hypertarget{c01:cue:1987}{}讓老師了解你的時候。
% cue 1988 3549.920 3573.890
\hypertarget{c01:cue:1988}{}那你要稍微等一等。
% cue 1989 3573.890 3575.990
\hypertarget{c01:cue:1989}{}像我剛才講那個專題。
% cue 1990 3575.990 3576.990
\hypertarget{c01:cue:1990}{}就只好碰到。
% cue 1991 3576.990 3578.270
\hypertarget{c01:cue:1991}{}那我就會講。

\tc{0:59:38}
\speaker{學生}
% cue 1992 3578.270 3579.270
\hypertarget{c01:cue:1992}{}那現在就是。
% cue 1993 3579.270 3580.260
\hypertarget{c01:cue:1993}{}我不知道你們。
% cue 1994 3580.260 3582.260
\hypertarget{c01:cue:1994}{}我不知道我可能的觀眾在哪裡。
% cue 1995 3582.260 3585.310
\hypertarget{c01:cue:1995}{}我就隨便跳跳跳跳這樣子。

\tc{0:59:45}
\speaker{老師}
% cue 1996 3585.310 3586.660
\hypertarget{c01:cue:1996}{}好 謝謝你。
% cue 1997 3586.660 3587.630
\hypertarget{c01:cue:1997}{}我聽懂了。
% cue 1998 3587.630 3588.630
\hypertarget{c01:cue:1998}{}我會一部把它列入考慮。
% cue 1999 3588.630 3589.790
\hypertarget{c01:cue:1999}{}可以嗎。
% cue 2000 3589.790 3591.040
\hypertarget{c01:cue:2000}{}好。

\tc{0:59:51}
% cue 2001 3591.040 3592.040
\hypertarget{c01:cue:2001}{}你有看過老師的影片嗎。
% cue 2002 3592.040 3594.090
\hypertarget{c01:cue:2002}{}沒有關係。
% cue 2003 3594.090 3596.090
\hypertarget{c01:cue:2003}{}從今天開始趕快去加入。
% cue 2004 3596.090 3597.790
\hypertarget{c01:cue:2004}{}好不好。
% cue 2005 3597.790 3598.790
\hypertarget{c01:cue:2005}{}反正的影片都在那裡。
% cue 2006 3598.790 3599.760
\hypertarget{c01:cue:2006}{}自己慢慢看。

\tc{0:59:59}
% cue 2007 3599.760 3600.760
\hypertarget{c01:cue:2007}{}法五第一。
% cue 2008 3600.760 3601.760
\hypertarget{c01:cue:2008}{}劉其為。
% cue 2009 3601.760 3603.200
\hypertarget{c01:cue:2009}{}好。
% cue 2010 3603.200 3604.840
\hypertarget{c01:cue:2010}{}來請說。

\tc{1:00:04}
\speaker{學生}
% cue 2011 3604.840 3607.150
\hypertarget{c01:cue:2011}{}我們學老師想辦法很多。
% cue 2012 3607.150 3608.150
\hypertarget{c01:cue:2012}{}然後想。
% cue 2013 3608.150 3609.150
\hypertarget{c01:cue:2013}{}上課的影片才會。

\tc{1:00:09}
\speaker{老師}
% cue 2014 3609.150 3610.150
\hypertarget{c01:cue:2014}{}好。
% cue 2015 3610.150 3611.340
\hypertarget{c01:cue:2015}{}那給你一個請求好不好。
% cue 2016 3611.340 3613.340
\hypertarget{c01:cue:2016}{}老師我上課沒想辦法。
% cue 2017 3613.340 3614.340
\hypertarget{c01:cue:2017}{}下課來跟我說。
% cue 2018 3614.340 3617.010
\hypertarget{c01:cue:2018}{}那如果我上課想不夠詳細。
% cue 2019 3617.010 3618.050
\hypertarget{c01:cue:2019}{}就舉手了。
% cue 2020 3618.050 3619.640
\hypertarget{c01:cue:2020}{}好不好。
% cue 2021 3619.640 3619.740
\hypertarget{c01:cue:2021}{}可以嗎。
% cue 2022 3619.640 3620.640
\hypertarget{c01:cue:2022}{}年級了吧。
% cue 2023 3620.640 3621.700
\hypertarget{c01:cue:2023}{}你要珍惜你再多一個機會。
% cue 2024 3621.700 3622.650
\hypertarget{c01:cue:2024}{}太多了。
% cue 2025 3622.650 3624.250
\hypertarget{c01:cue:2025}{}真的。

\tc{1:00:24}
% cue 2026 3624.250 3625.630
\hypertarget{c01:cue:2026}{}法五新。
% cue 2027 3625.630 3626.880
\hypertarget{c01:cue:2027}{}張明星。
% cue 2028 3626.880 3629.950
\hypertarget{c01:cue:2028}{}說。
% cue 2029 3629.950 3633.260
\hypertarget{c01:cue:2029}{}看過我的影片。
% cue 2030 3633.260 3634.260
\hypertarget{c01:cue:2030}{}看哪一個。

\tc{1:00:34}
\speaker{學生}
% cue 2031 3634.260 3636.130
\hypertarget{c01:cue:2031}{}多一遍。
% cue 2032 3636.130 3638.400
\hypertarget{c01:cue:2032}{}我也改革了。

\tc{1:00:38}
\speaker{老師}
% cue 2033 3638.400 3640.690
\hypertarget{c01:cue:2033}{}你覺得老師是不是很大膽。
% cue 2034 3640.690 3642.740
\hypertarget{c01:cue:2034}{}你腿之後很近。
% cue 2035 3642.740 3643.740
\hypertarget{c01:cue:2035}{}蛤。
% cue 2036 3643.740 3645.340
\hypertarget{c01:cue:2036}{}蠢蠢蠢。
% cue 2037 3645.340 3646.690
\hypertarget{c01:cue:2037}{}我跟你講。
% cue 2038 3646.690 3647.690
\hypertarget{c01:cue:2038}{}我會改革內視劇。
% cue 2039 3647.690 3649.870
\hypertarget{c01:cue:2039}{}因為他問題非常多。
% cue 2040 3649.870 3652.060
\hypertarget{c01:cue:2040}{}你看他的判決數真是我的。
% cue 2041 3652.060 3654.240
\hypertarget{c01:cue:2041}{}就樂樂等了。
% cue 2042 3654.240 3655.660
\hypertarget{c01:cue:2042}{}讀都讀不完。
% cue 2043 3655.660 3656.750
\hypertarget{c01:cue:2043}{}讀了前段。
% cue 2044 3656.750 3657.750
\hypertarget{c01:cue:2044}{}看不懂後段。
% cue 2045 3657.750 3658.750
\hypertarget{c01:cue:2045}{}讀了後段。
% cue 2046 3658.750 3659.750
\hypertarget{c01:cue:2046}{}忘記前段。
% cue 2047 3659.750 3660.700
\hypertarget{c01:cue:2047}{}就是這麼亂的東西。
% cue 2048 3660.700 3663.000
\hypertarget{c01:cue:2048}{}現在要一個人在很短的時間內。
% cue 2049 3663.000 3665.000
\hypertarget{c01:cue:2049}{}一針借曲告訴我們這個判決。
% cue 2050 3665.000 3666.000
\hypertarget{c01:cue:2050}{}根本錯在哪裡。
% cue 2051 3666.000 3667.450
\hypertarget{c01:cue:2051}{}很難吧。

\tc{1:01:07}
% cue 2052 3667.450 3668.480
\hypertarget{c01:cue:2052}{}下集。
% cue 2053 3668.480 3669.480
\hypertarget{c01:cue:2053}{}是在講體力。
% cue 2054 3669.480 3671.050
\hypertarget{c01:cue:2054}{}這個案子開始。
% cue 2055 3671.050 3673.010
\hypertarget{c01:cue:2055}{}他的體力跟你講的不一樣。
% cue 2056 3673.010 3675.460
\hypertarget{c01:cue:2056}{}一般人有感覺嗎。
% cue 2057 3675.460 3676.460
\hypertarget{c01:cue:2057}{}一般人沒有那麼大法。
% cue 2058 3676.460 3677.460
\hypertarget{c01:cue:2058}{}外面不會寫判決。
% cue 2059 3677.460 3678.460
\hypertarget{c01:cue:2059}{}他不會有感覺。
% cue 2060 3678.460 3681.600
\hypertarget{c01:cue:2060}{}所以我現在是想要這樣講。

\tc{1:01:21}
% cue 2061 3681.600 3682.820
\hypertarget{c01:cue:2061}{}好 謝謝你。
% cue 2062 3682.820 3683.820
\hypertarget{c01:cue:2062}{}有什麼對老師。
% cue 2063 3683.820 3684.820
\hypertarget{c01:cue:2063}{}有什麼欺許。
% cue 2064 3684.820 3687.840
\hypertarget{c01:cue:2064}{}再講更多的八卦。
% cue 2065 3687.840 3690.090
\hypertarget{c01:cue:2065}{}更多的八卦請來。
% cue 2066 3690.090 3692.090
\hypertarget{c01:cue:2066}{}新書發表會題。
% cue 2067 3692.090 3693.630
\hypertarget{c01:cue:2067}{}嘿嘿嘿。
% cue 2068 3693.630 3693.730
\hypertarget{c01:cue:2068}{}好不好。
% cue 2069 3693.630 3695.450
\hypertarget{c01:cue:2069}{}你要說什麼。
% cue 2070 3695.450 3696.450
\hypertarget{c01:cue:2070}{}對不起。
% cue 2071 3696.450 3699.810
\hypertarget{c01:cue:2071}{}好。

\tc{1:01:39}
% cue 2072 3699.810 3700.810
\hypertarget{c01:cue:2072}{}法五 C。
% cue 2073 3700.810 3701.970
\hypertarget{c01:cue:2073}{}劉以樂。
% cue 2074 3701.970 3710.080
\hypertarget{c01:cue:2074}{}好喔。
% cue 2075 3710.080 3714.570
\hypertarget{c01:cue:2075}{}那你喜歡那些練馬。
% cue 2076 3714.570 3724.680
\hypertarget{c01:cue:2076}{}你現在這位學習。

\tc{1:02:04}
\speaker{學生}
% cue 2077 3724.680 3725.680
\hypertarget{c01:cue:2077}{}我沒有錄可以參加。
% cue 2078 3725.680 3728.500
\hypertarget{c01:cue:2078}{}只是我沒有錄影機。
% cue 2079 3728.500 3729.630
\hypertarget{c01:cue:2079}{}其實也請不錄。

\tc{1:02:09}
\speaker{老師}
% cue 2080 3729.630 3730.850
\hypertarget{c01:cue:2080}{}你還是可以。
% cue 2081 3730.850 3731.850
\hypertarget{c01:cue:2081}{}隨時都可以提問。
% cue 2082 3731.850 3734.030
\hypertarget{c01:cue:2082}{}那如果你覺得老師的問題。
% cue 2083 3734.030 3735.030
\hypertarget{c01:cue:2083}{}問得不夠好。
% cue 2084 3735.030 3736.030
\hypertarget{c01:cue:2084}{}你還可以說。
% cue 2085 3736.030 3738.030
\hypertarget{c01:cue:2085}{}老師我可不可以再補充一個問題。
% cue 2086 3738.030 3739.340
\hypertarget{c01:cue:2086}{}讓你來當老師。
% cue 2087 3739.340 3741.360
\hypertarget{c01:cue:2087}{}這樣好不好。
% cue 2088 3741.360 3742.360
\hypertarget{c01:cue:2088}{}試試看。
% cue 2089 3742.360 3743.350
\hypertarget{c01:cue:2089}{}好。

\tc{1:02:23}
% cue 2090 3743.350 3744.350
\hypertarget{c01:cue:2090}{}解育停。
% cue 2091 3744.350 3746.300
\hypertarget{c01:cue:2091}{}解育停。
% cue 2092 3746.300 3748.420
\hypertarget{c01:cue:2092}{}請說。

\tc{1:02:28}
\speaker{學生}
% cue 2093 3748.420 3749.350
\hypertarget{c01:cue:2093}{}沒有你可以 S。
% cue 2094 3749.350 3751.140
\hypertarget{c01:cue:2094}{}我就。

\tc{1:02:31}
\speaker{老師}
% cue 2095 3751.140 3752.140
\hypertarget{c01:cue:2095}{}好。
% cue 2096 3752.140 3754.510
\hypertarget{c01:cue:2096}{}然後你又成這樣。
% cue 2097 3754.510 3760.560
\hypertarget{c01:cue:2097}{}那你上次至少上網的課次。
% cue 2098 3760.560 3761.560
\hypertarget{c01:cue:2098}{}好一點。
% cue 2099 3761.560 3762.560
\hypertarget{c01:cue:2099}{}你要跟人家講說。
% cue 2100 3762.560 3764.560
\hypertarget{c01:cue:2100}{}那種難的先把我都上得這麼好。
% cue 2101 3764.560 3765.620
\hypertarget{c01:cue:2101}{}可見我不是先把。
% cue 2102 3765.620 3767.160
\hypertarget{c01:cue:2102}{}對嗎。
% cue 2103 3767.160 3767.260
\hypertarget{c01:cue:2103}{}對嗎。
% cue 2104 3767.160 3768.150
\hypertarget{c01:cue:2104}{}我這樣子。
% cue 2105 3768.150 3770.370
\hypertarget{c01:cue:2105}{}我不相信有人上得比我難。
% cue 2106 3770.370 3772.880
\hypertarget{c01:cue:2106}{}那講內容豐富的程度。
% cue 2107 3772.880 3774.100
\hypertarget{c01:cue:2107}{}應該沒有。
% cue 2108 3774.100 3775.140
\hypertarget{c01:cue:2108}{}考試的困難度。
% cue 2109 3775.140 3778.100
\hypertarget{c01:cue:2109}{}那就更讓你們透透。

\tc{1:02:58}
% cue 2110 3778.100 3779.100
\hypertarget{c01:cue:2110}{}好。
% cue 2111 3779.100 3780.320
\hypertarget{c01:cue:2111}{}神樂哲。
% cue 2112 3780.320 3781.370
\hypertarget{c01:cue:2112}{}法三 B。
% cue 2113 3781.370 3782.440
\hypertarget{c01:cue:2113}{}欸。
% cue 2114 3782.440 3784.250
\hypertarget{c01:cue:2114}{}教育老師。
% cue 2115 3784.250 3785.480
\hypertarget{c01:cue:2115}{}教育老師。
% cue 2116 3785.480 3786.480
\hypertarget{c01:cue:2116}{}教育老師。
% cue 2117 3786.480 3788.150
\hypertarget{c01:cue:2117}{}那你高育前位置。
% cue 2118 3788.150 3790.150
\hypertarget{c01:cue:2118}{}你覺得老師是不是跟人家很不一樣。

\tc{1:03:10}
\speaker{學生}
% cue 2119 3790.150 3793.300
\hypertarget{c01:cue:2119}{}其實我自己的主要。
% cue 2120 3793.300 3794.580
\hypertarget{c01:cue:2120}{}對呀。

\tc{1:03:14}
\speaker{老師}
% cue 2121 3794.580 3799.480
\hypertarget{c01:cue:2121}{}用聽的是你們聽到老師想告訴你的。
% cue 2122 3799.480 3803.720
\hypertarget{c01:cue:2122}{}但是你不能讓老師知道你想要知道的事情。
% cue 2123 3803.720 3804.720
\hypertarget{c01:cue:2123}{}因為我不知道你要想知道。
% cue 2124 3804.720 3805.910
\hypertarget{c01:cue:2124}{}所以我必須要用對話。

\tc{1:03:25}
% cue 2125 3805.910 3807.200
\hypertarget{c01:cue:2125}{}欸。
% cue 2126 3807.200 3810.260
\hypertarget{c01:cue:2126}{}最好的方法是你照你的方法去學習。
% cue 2127 3810.260 3812.450
\hypertarget{c01:cue:2127}{}老師是在從旁幫忙的。
% cue 2128 3812.450 3813.470
\hypertarget{c01:cue:2128}{}像老師拚肌。
% cue 2129 3813.470 3814.560
\hypertarget{c01:cue:2129}{}他自己本來就會拚。
% cue 2130 3814.560 3815.690
\hypertarget{c01:cue:2130}{}他這個拚不好。
% cue 2131 3815.690 3816.690
\hypertarget{c01:cue:2131}{}你幫他一下。
% cue 2132 3816.690 3819.070
\hypertarget{c01:cue:2132}{}那這樣子的好處是呢。
% cue 2133 3819.070 3821.980
\hypertarget{c01:cue:2133}{}小孩子從小就會很有 ambition。
% cue 2134 3821.980 3823.360
\hypertarget{c01:cue:2134}{}他有雄心。
% cue 2135 3823.360 3825.360
\hypertarget{c01:cue:2135}{}他知道自己要聽什麼。
% cue 2136 3825.360 3827.450
\hypertarget{c01:cue:2136}{}而不是老師說這裡加這裡。
% cue 2137 3827.450 3828.450
\hypertarget{c01:cue:2137}{}那教完下一步。
% cue 2138 3828.450 3829.450
\hypertarget{c01:cue:2138}{}老師如果不講他沒有。
% cue 2139 3829.450 3830.890
\hypertarget{c01:cue:2139}{}就停在這裡。
% cue 2140 3830.890 3831.840
\hypertarget{c01:cue:2140}{}他完全沒有概念。

\tc{1:03:51}
% cue 2141 3831.840 3833.050
\hypertarget{c01:cue:2141}{}這是我們教學上的問題。
% cue 2142 3833.050 3834.050
\hypertarget{c01:cue:2142}{}但我只能說。
% cue 2143 3834.050 3835.620
\hypertarget{c01:cue:2143}{}唉。
% cue 2144 3835.620 3836.620
\hypertarget{c01:cue:2144}{}恭喜你能透透我們。
% cue 2145 3836.620 3837.620
\hypertarget{c01:cue:2145}{}但是有點太慢了。
% cue 2146 3837.620 3840.070
\hypertarget{c01:cue:2146}{}因為我們臺灣整個就是這種風力士。
% cue 2147 3840.070 3842.210
\hypertarget{c01:cue:2147}{}就是這樣大。
% cue 2148 3842.210 3843.210
\hypertarget{c01:cue:2148}{}希望被養大。
% cue 2149 3843.210 3845.460
\hypertarget{c01:cue:2149}{}這個生涯是要改成什麼版。
% cue 2150 3845.460 3846.460
\hypertarget{c01:cue:2150}{}這很難改變。

\tc{1:04:06}
% cue 2151 3846.460 3849.200
\hypertarget{c01:cue:2151}{}那你剛開始跟我上課。
% cue 2152 3849.200 3850.200
\hypertarget{c01:cue:2152}{}因為我們不喜歡老師。
% cue 2153 3850.200 3851.200
\hypertarget{c01:cue:2153}{}跳來跳去。
% cue 2154 3851.200 3852.610
\hypertarget{c01:cue:2154}{}跳來跳去。
% cue 2155 3852.610 3853.770
\hypertarget{c01:cue:2155}{}讓我更多人。
% cue 2156 3853.770 3854.770
\hypertarget{c01:cue:2156}{}你厲害。
% cue 2157 3854.770 3855.920
\hypertarget{c01:cue:2157}{}我就跟你講得更多。
% cue 2158 3855.920 3856.920
\hypertarget{c01:cue:2158}{}你比較弱。
% cue 2159 3856.920 3857.920
\hypertarget{c01:cue:2159}{}我就跟你講得更多。
% cue 2160 3857.920 3859.330
\hypertarget{c01:cue:2160}{}但是我一定要在這個時間。
% cue 2161 3859.330 3860.330
\hypertarget{c01:cue:2161}{}把這樣子講回。
% cue 2162 3860.330 3862.860
\hypertarget{c01:cue:2162}{}讓我這樣子就交代不過去。
% cue 2163 3862.860 3864.200
\hypertarget{c01:cue:2163}{}但是你很強。
% cue 2164 3864.200 3865.200
\hypertarget{c01:cue:2164}{}我可以對你很強。
% cue 2165 3865.200 3867.450
\hypertarget{c01:cue:2165}{}你可以問我難度問題。
% cue 2166 3867.450 3869.920
\hypertarget{c01:cue:2166}{}我來說是看你。
% cue 2167 3869.920 3871.460
\hypertarget{c01:cue:2167}{}愈強則強。
% cue 2168 3871.460 3872.490
\hypertarget{c01:cue:2168}{}愈弱則弱。

\tc{1:04:32}
% cue 2169 3872.490 3872.590
\hypertarget{c01:cue:2169}{}好。
% cue 2170 3872.490 3873.490
\hypertarget{c01:cue:2170}{}法三 A。
% cue 2171 3873.490 3874.490
\hypertarget{c01:cue:2171}{}會我好。
% cue 2172 3874.490 3875.490
\hypertarget{c01:cue:2172}{}在哪裡。
% cue 2173 3875.490 3880.040
\hypertarget{c01:cue:2173}{}請教。
% cue 2174 3880.040 3880.140
\hypertarget{c01:cue:2174}{}說什麼。
% cue 2175 3880.040 3881.810
\hypertarget{c01:cue:2175}{}噢謝謝你。
% cue 2176 3881.810 3883.770
\hypertarget{c01:cue:2176}{}沒有用完沒錢。
% cue 2177 3883.770 3886.380
\hypertarget{c01:cue:2177}{}不要跟老師講。
% cue 2178 3886.380 3887.380
\hypertarget{c01:cue:2178}{}老師賺錢包在我身上。
% cue 2179 3887.380 3888.660
\hypertarget{c01:cue:2179}{}我會說好。
% cue 2180 3888.660 3890.980
\hypertarget{c01:cue:2180}{}我們開始打打課。
% cue 2181 3890.980 3891.980
\hypertarget{c01:cue:2181}{}好來可以開箱。
% cue 2182 3891.980 3892.980
\hypertarget{c01:cue:2182}{}謝謝你。

\tc{1:04:52}
% cue 2183 3892.980 3904.840
\hypertarget{c01:cue:2183}{}法三 A。
% cue 2184 3904.840 3905.840
\hypertarget{c01:cue:2184}{}未重電。
% cue 2185 3905.840 3906.840
\hypertarget{c01:cue:2185}{}速重電。
% cue 2186 3906.840 3907.840
\hypertarget{c01:cue:2186}{}對不起。

\tc{1:05:07}
\speaker{學生}
% cue 2187 3907.840 3912.650
\hypertarget{c01:cue:2187}{}轉學來。

\tc{1:05:12}
\speaker{老師}
% cue 2188 3912.650 3914.120
\hypertarget{c01:cue:2188}{}從哪裡轉來。

\tc{1:05:14}
\speaker{學生}
% cue 2189 3914.120 3915.210
\hypertarget{c01:cue:2189}{}跟亞洲。

\tc{1:05:15}
\speaker{老師}
% cue 2190 3915.210 3916.240
\hypertarget{c01:cue:2190}{}喔亞洲。
% cue 2191 3916.240 3917.560
\hypertarget{c01:cue:2191}{}我以為去過你們學校耶。
% cue 2192 3917.560 3924.130
\hypertarget{c01:cue:2192}{}我沒有散播像樣的學校。
% cue 2193 3924.130 3926.440
\hypertarget{c01:cue:2193}{}真的嗎。
% cue 2194 3926.440 3929.070
\hypertarget{c01:cue:2194}{}希望有一個老師是我的學生。
% cue 2195 3929.070 3930.070
\hypertarget{c01:cue:2195}{}跑步腳。
% cue 2196 3930.070 3931.100
\hypertarget{c01:cue:2196}{}看到以前了。

\tc{1:05:31}
% cue 2197 3931.100 3932.320
\hypertarget{c01:cue:2197}{}好。
% cue 2198 3932.320 3933.350
\hypertarget{c01:cue:2198}{}那這個字典是什麼。

\tc{1:05:33}
\speaker{學生}
% cue 2199 3933.350 3935.350
\hypertarget{c01:cue:2199}{}波品。
% cue 2200 3935.350 3937.710
\hypertarget{c01:cue:2200}{}波品琴。

\tc{1:05:37}
\speaker{老師}
% cue 2201 3937.710 3938.710
\hypertarget{c01:cue:2201}{}好來。
% cue 2202 3938.710 3939.700
\hypertarget{c01:cue:2202}{}你說波品琴。
% cue 2203 3939.700 3945.400
\hypertarget{c01:cue:2203}{}去試試看。
% cue 2204 3945.400 3946.650
\hypertarget{c01:cue:2204}{}尤其是我們這個課。
% cue 2205 3946.650 3947.650
\hypertarget{c01:cue:2205}{}如果我們開始課服。
% cue 2206 3947.650 3948.840
\hypertarget{c01:cue:2206}{}那就會比較有信。
% cue 2207 3948.840 3950.970
\hypertarget{c01:cue:2207}{}我們現在這個上課過程中。
% cue 2208 3950.970 3951.970
\hypertarget{c01:cue:2208}{}也都是在幫助大家。
% cue 2209 3951.970 3953.970
\hypertarget{c01:cue:2209}{}正確的輪動提示。
% cue 2210 3953.970 3956.480
\hypertarget{c01:cue:2210}{}但是而且都是屬於權力分立。
% cue 2211 3956.480 3959.730
\hypertarget{c01:cue:2211}{}但是不一定是疑難雜症。
% cue 2212 3959.730 3963.340
\hypertarget{c01:cue:2212}{}那老師寫在世界學時興路易的。
% cue 2213 3963.340 3965.400
\hypertarget{c01:cue:2213}{}那就比較有些學問。
% cue 2214 3965.400 3967.400
\hypertarget{c01:cue:2214}{}他的線是這一刻。
% cue 2215 3967.400 3968.400
\hypertarget{c01:cue:2215}{}串這刻。
% cue 2216 3968.400 3969.680
\hypertarget{c01:cue:2216}{}串這刻。
% cue 2217 3969.680 3971.120
\hypertarget{c01:cue:2217}{}他有一串。
% cue 2218 3971.120 3973.120
\hypertarget{c01:cue:2218}{}那這樣一串的圖跟一個個圖。
% cue 2219 3973.120 3974.120
\hypertarget{c01:cue:2219}{}就有點不一樣。
% cue 2220 3974.120 3975.300
\hypertarget{c01:cue:2220}{}那再說吧。

\tc{1:06:15}
% cue 2221 3975.300 3976.360
\hypertarget{c01:cue:2221}{}好。
% cue 2222 3976.360 3978.450
\hypertarget{c01:cue:2222}{}這個是有子欠嗎。
% cue 2223 3978.450 3980.350
\hypertarget{c01:cue:2223}{}在哪裡。
% cue 2224 3980.350 3981.350
\hypertarget{c01:cue:2224}{}還是有子音。
% cue 2225 3981.350 3982.470
\hypertarget{c01:cue:2225}{}這個是音還是欠。

\tc{1:06:22}
\speaker{學生}
% cue 2226 3982.470 3983.820
\hypertarget{c01:cue:2226}{}音。

\tc{1:06:23}
\speaker{老師}
% cue 2227 3983.820 3984.920
\hypertarget{c01:cue:2227}{}有子音。
% cue 2228 3984.920 3985.920
\hypertarget{c01:cue:2228}{}來請說。

\tc{1:06:25}
% cue 2229 3985.920 3987.590
\hypertarget{c01:cue:2229}{}斯卡三 D。

\tc{1:06:27}
\speaker{學生}
% cue 2230 3987.590 4014.560
\hypertarget{c01:cue:2230}{}現是現在叫不同。
% cue 2231 4014.560 4017.240
\hypertarget{c01:cue:2231}{}也有一些是以前不同。
% cue 2232 4017.240 4019.230
\hypertarget{c01:cue:2232}{}現在變成這樣子。

\tc{1:06:59}
\speaker{老師}
% cue 2233 4019.230 4020.570
\hypertarget{c01:cue:2233}{}來。
% cue 2234 4020.570 4021.600
\hypertarget{c01:cue:2234}{}書祥語。
% cue 2235 4021.600 4026.300
\hypertarget{c01:cue:2235}{}對呀。
% cue 2236 4026.300 4027.690
\hypertarget{c01:cue:2236}{}你要說什麼。
% cue 2237 4027.690 4028.750
\hypertarget{c01:cue:2237}{}對。
% cue 2238 4028.750 4029.660
\hypertarget{c01:cue:2238}{}再看著。
% cue 2239 4029.660 4031.860
\hypertarget{c01:cue:2239}{}好。
% cue 2240 4031.860 4031.960
\hypertarget{c01:cue:2240}{}謝謝。
% cue 2241 4031.860 4032.860
\hypertarget{c01:cue:2241}{}繼續保持。
% cue 2242 4032.860 4034.860
\hypertarget{c01:cue:2242}{}這樣一個好奇心是好的。

\tc{1:07:14}
% cue 2243 4034.860 4036.110
\hypertarget{c01:cue:2243}{}法二 A 郭明彥。
% cue 2244 4036.110 4049.610
\hypertarget{c01:cue:2244}{}你又看過我自己。
% cue 2245 4049.610 4051.870
\hypertarget{c01:cue:2245}{}哪裡真是嗎。
% cue 2246 4051.870 4053.210
\hypertarget{c01:cue:2246}{}因為我沒交過你。
% cue 2247 4053.210 4065.150
\hypertarget{c01:cue:2247}{}聽說是不準的。
% cue 2248 4065.150 4069.090
\hypertarget{c01:cue:2248}{}謝謝你。
% cue 2249 4069.090 4070.210
\hypertarget{c01:cue:2249}{}郭採琪。
% cue 2250 4070.210 4071.200
\hypertarget{c01:cue:2250}{}好。
% cue 2251 4071.200 4075.370
\hypertarget{c01:cue:2251}{}謝謝。
% cue 2252 4075.370 4076.370
\hypertarget{c01:cue:2252}{}我也希望。

\tc{1:07:56}
% cue 2253 4076.370 4077.560
\hypertarget{c01:cue:2253}{}來。
% cue 2254 4077.560 4079.320
\hypertarget{c01:cue:2254}{}書繪傑。
% cue 2255 4079.320 4081.190
\hypertarget{c01:cue:2255}{}法二 A。
% cue 2256 4081.190 4083.270
\hypertarget{c01:cue:2256}{}用這個人喔。
% cue 2257 4083.270 4086.890
\hypertarget{c01:cue:2257}{}高調囉。
% cue 2258 4086.890 4087.950
\hypertarget{c01:cue:2258}{}法二 A。
% cue 2259 4087.950 4089.590
\hypertarget{c01:cue:2259}{}徐香瑜。
% cue 2260 4089.590 4093.720
\hypertarget{c01:cue:2260}{}S。
% cue 2261 4093.720 4094.720
\hypertarget{c01:cue:2261}{}法二 A。
% cue 2262 4094.720 4096.360
\hypertarget{c01:cue:2262}{}傅如琴。
% cue 2263 4096.360 4101.100
\hypertarget{c01:cue:2263}{}對。
% cue 2264 4101.100 4102.100
\hypertarget{c01:cue:2264}{}我就說過了。
% cue 2265 4102.100 4103.100
\hypertarget{c01:cue:2265}{}沒有上過我的課。
% cue 2266 4103.100 4104.100
\hypertarget{c01:cue:2266}{}不要說來歐東。

\tc{1:08:24}
% cue 2267 4104.100 4105.860
\hypertarget{c01:cue:2267}{}法二 A。
% cue 2268 4105.860 4107.080
\hypertarget{c01:cue:2268}{}張亞橋。
% cue 2269 4107.080 4110.260
\hypertarget{c01:cue:2269}{}法二 A。
% cue 2270 4110.260 4111.390
\hypertarget{c01:cue:2270}{}湯鬼。
% cue 2271 4111.390 4112.530
\hypertarget{c01:cue:2271}{}傑嗎。
% cue 2272 4112.530 4113.530
\hypertarget{c01:cue:2272}{}這個是傑嗎。

\tc{1:08:33}
\speaker{學生}
% cue 2273 4113.530 4115.550
\hypertarget{c01:cue:2273}{}傻。
% cue 2274 4115.550 4115.650
\hypertarget{c01:cue:2274}{}傻。

\tc{1:08:35}
\speaker{老師}
% cue 2275 4115.550 4116.550
\hypertarget{c01:cue:2275}{}這個是傻啊。
% cue 2276 4116.550 4117.770
\hypertarget{c01:cue:2276}{}有捲傻的傻。
% cue 2277 4117.770 4119.770
\hypertarget{c01:cue:2277}{}還是沒有捲傻的傻。
% cue 2278 4119.770 4121.020
\hypertarget{c01:cue:2278}{}蛤。

\tc{1:08:41}
\speaker{學生}
% cue 2279 4121.020 4122.020
\hypertarget{c01:cue:2279}{}有捲傻的傻。

\tc{1:08:42}
\speaker{老師}
% cue 2280 4122.020 4123.210
\hypertarget{c01:cue:2280}{}你也心疼啊。
% cue 2281 4123.210 4125.300
\hypertarget{c01:cue:2281}{}你是哪裡人啊。

\tc{1:08:45}
\speaker{學生}
% cue 2282 4125.300 4127.450
\hypertarget{c01:cue:2282}{}高雄人。

\tc{1:08:47}
\speaker{老師}
% cue 2283 4127.450 4128.450
\hypertarget{c01:cue:2283}{}我知道。
% cue 2284 4128.450 4129.450
\hypertarget{c01:cue:2284}{}你是臺灣人。
% cue 2285 4129.450 4130.420
\hypertarget{c01:cue:2285}{}還是客家人。
% cue 2286 4130.420 4131.420
\hypertarget{c01:cue:2286}{}還是什麼人。
% cue 2287 4131.420 4136.090
\hypertarget{c01:cue:2287}{}對。
% cue 2288 4136.090 4137.330
\hypertarget{c01:cue:2288}{}好真的喔。
% cue 2289 4137.330 4138.630
\hypertarget{c01:cue:2289}{}性疼的人不多啊。
% cue 2290 4138.630 4140.840
\hypertarget{c01:cue:2290}{}那個在高雄。
% cue 2291 4140.840 4142.840
\hypertarget{c01:cue:2291}{}在臺海有個新他們沒有。

\tc{1:09:02}
\speaker{學生}
% cue 2292 4142.840 4144.250
\hypertarget{c01:cue:2292}{}那個是。
% cue 2293 4144.250 4145.790
\hypertarget{c01:cue:2293}{}抗日人群。

\tc{1:09:05}
\speaker{老師}
% cue 2294 4145.790 4149.690
\hypertarget{c01:cue:2294}{}欸。

\tc{1:09:09}
% cue 2295 4149.690 4152.230
\hypertarget{c01:cue:2295}{}那有什麼話要講嗎。
% cue 2296 4152.230 4153.320
\hypertarget{c01:cue:2296}{}回哲。
% cue 2297 4153.320 4160.400
\hypertarget{c01:cue:2297}{}中原啊。
% cue 2298 4160.400 4161.360
\hypertarget{c01:cue:2298}{}那恭喜你啊。
% cue 2299 4161.360 4163.090
\hypertarget{c01:cue:2299}{}今天來之前就有。
% cue 2300 4163.090 4167.090
\hypertarget{c01:cue:2300}{}你叫什麼。
% cue 2301 4167.090 4168.090
\hypertarget{c01:cue:2301}{}我。
% cue 2302 4168.090 4168.190
\hypertarget{c01:cue:2302}{}讀我。
% cue 2303 4168.090 4170.820
\hypertarget{c01:cue:2303}{}真的。
% cue 2304 4170.820 4170.920
\hypertarget{c01:cue:2304}{}為什麼。
% cue 2305 4170.820 4173.140
\hypertarget{c01:cue:2305}{}我在中原有聽過你嗎。
% cue 2306 4173.140 4175.160
\hypertarget{c01:cue:2306}{}對啊。
% cue 2307 4175.160 4176.600
\hypertarget{c01:cue:2307}{}哦。
% cue 2308 4176.600 4177.600
\hypertarget{c01:cue:2308}{}對啊。
% cue 2309 4177.600 4179.580
\hypertarget{c01:cue:2309}{}中原還沒有請我去新書發表。

\tc{1:09:39}
% cue 2310 4179.580 4182.140
\hypertarget{c01:cue:2310}{}我去了很多地方。
% cue 2311 4182.140 4183.140
\hypertarget{c01:cue:2311}{}在 YouTube 上。
% cue 2312 4183.140 4185.140
\hypertarget{c01:cue:2312}{}老師有官網\footnote{湯德宗的個人網站，網址 https://dennis-tang.site ，課綱列為參考資源與教師個人網址。站上設履歷、大法官大事記、近期活動、著作、授課、影音天地等欄。}知道嗎。
% cue 2313 4185.140 4186.170
\hypertarget{c01:cue:2313}{}這個上面有官網。
% cue 2314 4186.170 4188.170
\hypertarget{c01:cue:2314}{}官網你會看到我的活動。
% cue 2315 4188.170 4190.490
\hypertarget{c01:cue:2315}{}我有去我就會播一張照片啊。
% cue 2316 4190.490 4192.030
\hypertarget{c01:cue:2316}{}講講這樣。
% cue 2317 4192.030 4193.030
\hypertarget{c01:cue:2317}{}其實中原也有學生。
% cue 2318 4193.030 4194.030
\hypertarget{c01:cue:2318}{}但他沒有找我。
% cue 2319 4194.030 4195.230
\hypertarget{c01:cue:2319}{}就知道為什麼。
% cue 2320 4195.230 4196.440
\hypertarget{c01:cue:2320}{}可能現在太難糾集。
% cue 2321 4196.440 4198.720
\hypertarget{c01:cue:2321}{}因為找湯老師來。
% cue 2322 4198.720 4199.780
\hypertarget{c01:cue:2322}{}如果來人不多。
% cue 2323 4199.780 4200.780
\hypertarget{c01:cue:2323}{}不好意思啊。
% cue 2324 4200.780 4202.160
\hypertarget{c01:cue:2324}{}那要去找那麼多人。
% cue 2325 4202.160 4203.160
\hypertarget{c01:cue:2325}{}這很累啊。
% cue 2326 4203.160 4204.610
\hypertarget{c01:cue:2326}{}對不對。
% cue 2327 4204.610 4206.350
\hypertarget{c01:cue:2327}{}沒有關係。
% cue 2328 4206.350 4208.390
\hypertarget{c01:cue:2328}{}希望回到中原去過敵。
% cue 2329 4208.390 4209.960
\hypertarget{c01:cue:2329}{}嘿嘿嘿。
% cue 2330 4209.960 4211.160
\hypertarget{c01:cue:2330}{}可以看完這樣子。
% cue 2331 4211.160 4213.470
\hypertarget{c01:cue:2331}{}你可以在包括在中原拜託。

\tc{1:10:13}
% cue 2332 4213.470 4214.980
\hypertarget{c01:cue:2332}{}好來。
% cue 2333 4214.980 4215.890
\hypertarget{c01:cue:2333}{}然後。
% cue 2334 4215.890 4217.040
\hypertarget{c01:cue:2334}{}詹欣寧。
% cue 2335 4217.040 4217.140
\hypertarget{c01:cue:2335}{}來。
% cue 2336 4217.040 4218.620
\hypertarget{c01:cue:2336}{}請說。
% cue 2337 4218.620 4224.550
\hypertarget{c01:cue:2337}{}上課跟女兒有什麼異樣。
% cue 2338 4224.550 4225.550
\hypertarget{c01:cue:2338}{}沒什麼關係啊。
% cue 2339 4225.550 4226.550
\hypertarget{c01:cue:2339}{}看我們突然跟教授。
% cue 2340 4226.550 4227.580
\hypertarget{c01:cue:2340}{}沒果必來的管理。
% cue 2341 4227.580 4229.090
\hypertarget{c01:cue:2341}{}好來下來。

\tc{1:10:29}
% cue 2342 4229.090 4230.090
\hypertarget{c01:cue:2342}{}法二 B。
% cue 2343 4230.090 4231.280
\hypertarget{c01:cue:2343}{}減江雲。
% cue 2344 4231.280 4247.740
\hypertarget{c01:cue:2344}{}我也不敢說我豐富。
% cue 2345 4247.740 4248.700
\hypertarget{c01:cue:2345}{}我只當八年而已。
% cue 2346 4248.700 4251.080
\hypertarget{c01:cue:2346}{}法二 B。
% cue 2347 4251.080 4252.080
\hypertarget{c01:cue:2347}{}王一柔。
% cue 2348 4252.080 4253.910
\hypertarget{c01:cue:2348}{}請說。
% cue 2349 4253.910 4258.680
\hypertarget{c01:cue:2349}{}阿爾冰。
% cue 2350 4258.680 4259.680
\hypertarget{c01:cue:2350}{}王再興。
% cue 2351 4259.680 4261.510
\hypertarget{c01:cue:2351}{}王再興。
% cue 2352 4261.510 4262.510
\hypertarget{c01:cue:2352}{}沒有人。
% cue 2353 4262.510 4264.850
\hypertarget{c01:cue:2353}{}不掉喔。

\tc{1:11:04}
% cue 2354 4264.850 4267.290
\hypertarget{c01:cue:2354}{}通法二 B。
% cue 2355 4267.290 4268.280
\hypertarget{c01:cue:2355}{}張貴明。

\tc{1:11:08}
\speaker{學生}
% cue 2356 4268.280 4272.220
\hypertarget{c01:cue:2356}{}上學挺好。
% cue 2357 4272.220 4273.220
\hypertarget{c01:cue:2357}{}不推薦的哈佛\footnote{Harvard Law School，哈佛大學法學院，在美國麻薩諸塞州劍橋。湯德宗 1978 年畢業於國立臺灣大學法律學系，1981 年取得臺大法律研究所碩士，1984 年取得哈佛法學院 LL.M.，1989 年取得美國杜蘭大學（Tulane University）法學院 S.J.D.。\href{https://dennis-tang.site/\%E5\%B1\%A5\%E6\%AD\%B7/\%EF\%BC\%88\%E5\%B1\%A5\%E6\%AD\%B7\%E9\%A0\%81\%E5\%9C\%96\%E6\%AA\%94\%E3\%80\%8C\%E5\%AD\%B8\%E6\%AD\%B7\%E3\%80\%8D\%EF\%BC\%9ANTU\%20\%E6\%B3\%95\%E5\%BE\%8B\%E7\%B3\%BB\%E5\%AD\%B8\%E5\%A3\%AB\%201978\%E3\%80\%81\%E6\%B3\%95\%E5\%BE\%8B\%E7\%A0\%94\%E7\%A9\%B6\%E6\%89\%80\%E7\%A2\%A9\%E5\%A3\%AB\%201981\%E3\%80\%81Harvard\%20LLM\%201984\%E3\%80\%81Tulane\%20SJD\%201989\%EF\%BC\%89}{https://dennis-tang.site/履歷/（履歷頁圖檔「學歷」：NTU 法律系學士 1978、法律研究所碩士 1981、Harvard LLM 1984、Tulane SJD 1989）}。}新鮮。
% cue 2358 4273.220 4274.220
\hypertarget{c01:cue:2358}{}然後老師。
% cue 2359 4274.220 4276.210
\hypertarget{c01:cue:2359}{}就在哈佛高度。
% cue 2360 4276.210 4277.390
\hypertarget{c01:cue:2360}{}然後。
% cue 2361 4277.390 4278.650
\hypertarget{c01:cue:2361}{}變很激。
% cue 2362 4278.650 4280.020
\hypertarget{c01:cue:2362}{}所以我跟學姐一樣。
% cue 2363 4280.020 4281.020
\hypertarget{c01:cue:2363}{}很期待。
% cue 2364 4281.020 4283.610
\hypertarget{c01:cue:2364}{}蘇格拉底。

\tc{1:11:23}
\speaker{老師}
% cue 2365 4283.610 4284.610
\hypertarget{c01:cue:2365}{}蘇格拉底是的教學。
% cue 2366 4284.610 4285.610
\hypertarget{c01:cue:2366}{}基本上在很多都是。
% cue 2367 4285.610 4287.090
\hypertarget{c01:cue:2367}{}不是只有在哈佛。
% cue 2368 4287.090 4289.310
\hypertarget{c01:cue:2368}{}但是給我興趣聽哈佛。
% cue 2369 4289.310 4290.310
\hypertarget{c01:cue:2369}{}我都實在可以講一些。

\tc{1:11:30}
% cue 2370 4290.310 4291.910
\hypertarget{c01:cue:2370}{}但是我這個人是你的隨性。
% cue 2371 4291.910 4294.450
\hypertarget{c01:cue:2371}{}你有問題我就比較會回答。
% cue 2372 4294.450 4295.510
\hypertarget{c01:cue:2372}{}你這樣泛泛問我。
% cue 2373 4295.510 4296.510
\hypertarget{c01:cue:2373}{}那我就要看。
% cue 2374 4296.510 4300.020
\hypertarget{c01:cue:2374}{}假如今天秋天落葉很漂亮。
% cue 2375 4300.020 4303.010
\hypertarget{c01:cue:2375}{}我會忽然會想起我在畫裡面。
% cue 2376 4303.010 4304.800
\hypertarget{c01:cue:2376}{}我會跟你講。
% cue 2377 4304.800 4306.800
\hypertarget{c01:cue:2377}{}踩在你的落葉上的感覺。
% cue 2378 4306.800 4309.630
\hypertarget{c01:cue:2378}{}你們有過這種經驗嗎。
% cue 2379 4309.630 4310.630
\hypertarget{c01:cue:2379}{}這一踩下去。
% cue 2380 4310.630 4312.630
\hypertarget{c01:cue:2380}{}叫陷下去的彈起來一點點。
% cue 2381 4312.630 4314.750
\hypertarget{c01:cue:2381}{}就是秋天的落葉。
% cue 2382 4314.750 4317.960
\hypertarget{c01:cue:2382}{}現在還太早了。
% cue 2383 4317.960 4318.960
\hypertarget{c01:cue:2383}{}大概到九月裡。
% cue 2384 4318.960 4320.400
\hypertarget{c01:cue:2384}{}九月初。
% cue 2385 4320.400 4322.030
\hypertarget{c01:cue:2385}{}來月子都是在。
% cue 2386 4322.030 4324.440
\hypertarget{c01:cue:2386}{}那樹葉都到什麼程度。
% cue 2387 4324.440 4328.070
\hypertarget{c01:cue:2387}{}要用電的機器這樣吹。
% cue 2388 4328.070 4329.420
\hypertarget{c01:cue:2388}{}然後用很大的袋子。
% cue 2389 4329.420 4330.420
\hypertarget{c01:cue:2389}{}很多袋子。
% cue 2390 4330.420 4331.420
\hypertarget{c01:cue:2390}{}很多袋子做起來。
% cue 2391 4331.420 4333.960
\hypertarget{c01:cue:2391}{}也是有什麼。

\tc{1:12:13}
% cue 2392 4333.960 4335.540
\hypertarget{c01:cue:2392}{}好那張維銘講。
% cue 2393 4335.540 4336.630
\hypertarget{c01:cue:2393}{}RB 林品銀。

\tc{1:12:16}
\speaker{學生}
% cue 2394 4336.630 4337.630
\hypertarget{c01:cue:2394}{}我是。

\tc{1:12:17}
\speaker{老師}
% cue 2395 4337.630 4340.670
\hypertarget{c01:cue:2395}{}你剛剛不是想到班長。
% cue 2396 4340.670 4341.670
\hypertarget{c01:cue:2396}{}什麼退。

\tc{1:12:21}
\speaker{學生}
% cue 2397 4341.670 4342.670
\hypertarget{c01:cue:2397}{}沒有。
% cue 2398 4342.670 4343.670
\hypertarget{c01:cue:2398}{}沒有。
% cue 2399 4343.670 4344.670
\hypertarget{c01:cue:2399}{}沒有。
% cue 2400 4344.670 4345.950
\hypertarget{c01:cue:2400}{}因為我。
% cue 2401 4345.950 4346.950
\hypertarget{c01:cue:2401}{}可以說比較好戰一點。
% cue 2402 4346.950 4349.330
\hypertarget{c01:cue:2402}{}然後聽說老師很比較難。

\tc{1:12:29}
\speaker{老師}
% cue 2403 4349.330 4351.990
\hypertarget{c01:cue:2403}{}好很好。
% cue 2404 4351.990 4353.150
\hypertarget{c01:cue:2404}{}我自己對。
% cue 2405 4353.150 4357.090
\hypertarget{c01:cue:2405}{}也有些想。
% cue 2406 4357.090 4358.050
\hypertarget{c01:cue:2406}{}真的嗎。
% cue 2407 4358.050 4359.050
\hypertarget{c01:cue:2407}{}真的嗎。

\tc{1:12:39}
% cue 2408 4359.050 4361.710
\hypertarget{c01:cue:2408}{}你喜歡去外遊嗎。
% cue 2409 4361.710 4362.710
\hypertarget{c01:cue:2409}{}為什麼。

\tc{1:12:42}
\speaker{學生}
% cue 2410 4362.710 4362.810
\hypertarget{c01:cue:2410}{}沒。
% cue 2411 4362.710 4366.270
\hypertarget{c01:cue:2411}{}沒。

\tc{1:12:46}
\speaker{老師}
% cue 2412 4366.270 4367.620
\hypertarget{c01:cue:2412}{}真的嗎。

\tc{1:12:47}
\speaker{學生}
% cue 2413 4367.620 4369.990
\hypertarget{c01:cue:2413}{}然後在屋下的時候也有去。

\tc{1:12:49}
\speaker{老師}
% cue 2414 4369.990 4373.940
\hypertarget{c01:cue:2414}{}我怪你為什麼不想去 Harvard。

\tc{1:12:53}
\speaker{學生}
% cue 2415 4373.940 4377.440
\hypertarget{c01:cue:2415}{}跟外遊比較小。
% cue 2416 4377.440 4379.070
\hypertarget{c01:cue:2416}{}然後在。
% cue 2417 4379.070 4380.070
\hypertarget{c01:cue:2417}{}比較像。
% cue 2418 4380.070 4383.940
\hypertarget{c01:cue:2418}{}喜歡都市。

\tc{1:13:03}
\speaker{老師}
% cue 2419 4383.940 4384.940
\hypertarget{c01:cue:2419}{}那當然了。
% cue 2420 4384.940 4386.090
\hypertarget{c01:cue:2420}{}那不過。
% cue 2421 4386.090 4388.080
\hypertarget{c01:cue:2421}{}Harvard 離都市也很近。
% cue 2422 4388.080 4390.080
\hypertarget{c01:cue:2422}{}我也和查理師傅。
% cue 2423 4390.080 4391.140
\hypertarget{c01:cue:2423}{}這樣還沒到班校。
% cue 2424 4391.140 4392.580
\hypertarget{c01:cue:2424}{}結果就可以了。
% cue 2425 4392.580 4394.580
\hypertarget{c01:cue:2425}{}而且他有機率可以坐過去。
% cue 2426 4394.580 4397.080
\hypertarget{c01:cue:2426}{}這講太多了嗎。

\tc{1:13:17}
% cue 2427 4397.080 4399.070
\hypertarget{c01:cue:2427}{}侯瑜雲。
% cue 2428 4399.070 4400.070
\hypertarget{c01:cue:2428}{}法二 D。
% cue 2429 4400.070 4402.860
\hypertarget{c01:cue:2429}{}請說。
% cue 2430 4402.860 4408.210
\hypertarget{c01:cue:2430}{}上一次上誰的簽。
% cue 2431 4408.210 4411.840
\hypertarget{c01:cue:2431}{}上一次誰的簽。

\tc{1:13:31}
\speaker{學生}
% cue 2432 4411.840 4414.470
\hypertarget{c01:cue:2432}{}孟基老師的名字。

\tc{1:13:34}
\speaker{老師}
% cue 2433 4414.470 4416.420
\hypertarget{c01:cue:2433}{}他會哭喔。
% cue 2434 4416.420 4418.840
\hypertarget{c01:cue:2434}{}都會講那些。

\tc{1:13:38}
% cue 2435 4418.840 4421.080
\hypertarget{c01:cue:2435}{}法二 D。
% cue 2436 4421.080 4422.080
\hypertarget{c01:cue:2436}{}劉玉安。
% cue 2437 4422.080 4426.180
\hypertarget{c01:cue:2437}{}你講完了嗎。
% cue 2438 4426.180 4427.270
\hypertarget{c01:cue:2438}{}劉安。
% cue 2439 4427.270 4428.270
\hypertarget{c01:cue:2439}{}劉安。

\tc{1:13:48}
% cue 2440 4428.270 4429.260
\hypertarget{c01:cue:2440}{}法二 D。

\tc{1:13:49}
\speaker{學生}
% cue 2441 4429.260 4450.880
\hypertarget{c01:cue:2441}{}其實每一個法院的判決。
% cue 2442 4450.880 4451.880
\hypertarget{c01:cue:2442}{}他都有時空背景。
% cue 2443 4451.880 4454.170
\hypertarget{c01:cue:2443}{}但是他基本上是個案。
% cue 2444 4454.170 4457.150
\hypertarget{c01:cue:2444}{}他的影響源是出贏而已。
% cue 2445 4457.150 4459.520
\hypertarget{c01:cue:2445}{}那大法官人其實是影響大。
% cue 2446 4459.520 4461.670
\hypertarget{c01:cue:2446}{}而且大法官人其實是比較具有政治性的。
% cue 2447 4461.670 4463.830
\hypertarget{c01:cue:2447}{}所以那個時空背景。
% cue 2448 4463.830 4466.100
\hypertarget{c01:cue:2448}{}就更重要。

\tc{1:14:26}
\speaker{老師}
% cue 2449 4466.100 4466.200
\hypertarget{c01:cue:2449}{}好。

\tc{1:14:26}
% cue 2450 4466.100 4467.100
\hypertarget{c01:cue:2450}{}然後歐玉婷。
% cue 2451 4467.100 4468.000
\hypertarget{c01:cue:2451}{}法二 D。
% cue 2452 4468.000 4468.100
\hypertarget{c01:cue:2452}{}你。
% cue 2453 4468.000 4474.270
\hypertarget{c01:cue:2453}{}知道以前的嗎。
% cue 2454 4474.270 4477.130
\hypertarget{c01:cue:2454}{}那你覺得。
% cue 2455 4477.130 4478.130
\hypertarget{c01:cue:2455}{}你的影片給你什麼感覺。

\tc{1:14:38}
\speaker{學生}
% cue 2456 4478.130 4480.700
\hypertarget{c01:cue:2456}{}這老師上課。
% cue 2457 4480.700 4483.070
\hypertarget{c01:cue:2457}{}有點壓迫喔。

\tc{1:14:43}
\speaker{老師}
% cue 2458 4483.070 4484.070
\hypertarget{c01:cue:2458}{}好。
% cue 2459 4484.070 4485.070
\hypertarget{c01:cue:2459}{}可不可以跟我講一下。
% cue 2460 4485.070 4486.070
\hypertarget{c01:cue:2460}{}為什麼覺得有點壓迫。
% cue 2461 4486.070 4487.040
\hypertarget{c01:cue:2461}{}真的很好嗎。

\tc{1:14:47}
\speaker{學生}
% cue 2462 4487.040 4493.610
\hypertarget{c01:cue:2462}{}學生才說。
% cue 2463 4493.610 4494.620
\hypertarget{c01:cue:2463}{}這樣我才有參與。

\tc{1:14:54}
\speaker{老師}
% cue 2464 4494.620 4499.530
\hypertarget{c01:cue:2464}{}那所以你應該跟我說。
% cue 2465 4499.530 4500.530
\hypertarget{c01:cue:2465}{}老師我覺得很用調。
% cue 2466 4500.530 4504.000
\hypertarget{c01:cue:2466}{}而且壓迫好像是老師很壞。
% cue 2467 4504.000 4506.000
\hypertarget{c01:cue:2467}{}我去壓迫你做了什麼事情。
% cue 2468 4506.000 4508.000
\hypertarget{c01:cue:2468}{}或者老師給你們很大的壓力這樣子。
% cue 2469 4508.000 4509.410
\hypertarget{c01:cue:2469}{}好。
% cue 2470 4509.410 4510.410
\hypertarget{c01:cue:2470}{}好。

\tc{1:15:10}
% cue 2471 4510.410 4510.510
\hypertarget{c01:cue:2471}{}烏秋彭。
% cue 2472 4510.410 4512.480
\hypertarget{c01:cue:2472}{}我說。
% cue 2473 4512.480 4521.200
\hypertarget{c01:cue:2473}{}謝謝。
% cue 2474 4521.200 4522.810
\hypertarget{c01:cue:2474}{}謝利宇。
% cue 2475 4522.810 4524.670
\hypertarget{c01:cue:2475}{}謝利宇。
% cue 2476 4524.670 4525.770
\hypertarget{c01:cue:2476}{}法二 D。
% cue 2477 4525.770 4526.770
\hypertarget{c01:cue:2477}{}沒有喔。
% cue 2478 4526.770 4528.750
\hypertarget{c01:cue:2478}{}不掉喔。

\tc{1:15:28}
% cue 2479 4528.750 4529.750
\hypertarget{c01:cue:2479}{}然後法二 D。
% cue 2480 4529.750 4531.380
\hypertarget{c01:cue:2480}{}曾心瑜。
% cue 2481 4531.380 4532.870
\hypertarget{c01:cue:2481}{}那句是你。
% cue 2482 4532.870 4542.540
\hypertarget{c01:cue:2482}{}好請說。
% cue 2483 4542.540 4544.170
\hypertarget{c01:cue:2483}{}喔。
% cue 2484 4544.170 4545.170
\hypertarget{c01:cue:2484}{}可以告訴我為什麼開心嗎。

\tc{1:15:45}
% cue 2485 4545.170 4565.240
\hypertarget{c01:cue:2485}{}我跟你講。
% cue 2486 4565.240 4566.240
\hypertarget{c01:cue:2486}{}非常簡單的講是。
% cue 2487 4566.240 4568.490
\hypertarget{c01:cue:2487}{}民法刑法都比憲法。
% cue 2488 4568.490 4571.220
\hypertarget{c01:cue:2488}{}那句話。
% cue 2489 4571.220 4573.790
\hypertarget{c01:cue:2489}{}憲法相對比較抽象。
% cue 2490 4573.790 4576.290
\hypertarget{c01:cue:2490}{}那我跟一義上來講。
% cue 2491 4576.290 4579.410
\hypertarget{c01:cue:2491}{}抽象的東西本來就能學。
% cue 2492 4579.410 4582.740
\hypertarget{c01:cue:2492}{}越抽象越難學。
% cue 2493 4582.740 4583.680
\hypertarget{c01:cue:2493}{}好。

\tc{1:16:23}
% cue 2494 4583.680 4585.300
\hypertarget{c01:cue:2494}{}法二 D。
% cue 2495 4585.300 4586.390
\hypertarget{c01:cue:2495}{}陳朱如。
% cue 2496 4586.390 4589.120
\hypertarget{c01:cue:2496}{}朱如喔。
% cue 2497 4589.120 4590.080
\hypertarget{c01:cue:2497}{}是。
% cue 2498 4590.080 4600.880
\hypertarget{c01:cue:2498}{}請問一下你上學期的憲法。
% cue 2499 4600.880 4602.880
\hypertarget{c01:cue:2499}{}是講人權還是講政府體制。

\tc{1:16:42}
\speaker{學生}
% cue 2500 4602.880 4608.170
\hypertarget{c01:cue:2500}{}直接講人權。
% cue 2501 4608.170 4609.230
\hypertarget{c01:cue:2501}{}都有。
% cue 2502 4609.230 4610.230
\hypertarget{c01:cue:2502}{}都有上到。

\tc{1:16:50}
\speaker{老師}
% cue 2503 4610.230 4611.230
\hypertarget{c01:cue:2503}{}喔喔那太厲害了。
% cue 2504 4611.230 4613.230
\hypertarget{c01:cue:2504}{}三小時可以上這麼多。
% cue 2505 4613.230 4615.840
\hypertarget{c01:cue:2505}{}那太厲害了。

\tc{1:16:55}
\speaker{學生}
% cue 2506 4615.840 4616.840
\hypertarget{c01:cue:2506}{}也是公文講的。
% cue 2507 4616.840 4618.220
\hypertarget{c01:cue:2507}{}也是公文講的。

\tc{1:16:58}
\speaker{老師}
% cue 2508 4618.220 4619.690
\hypertarget{c01:cue:2508}{}法二 D。
% cue 2509 4619.690 4620.790
\hypertarget{c01:cue:2509}{}這個字是怎麼念。
% cue 2510 4620.790 4621.690
\hypertarget{c01:cue:2510}{}黃凱。
% cue 2511 4621.690 4623.050
\hypertarget{c01:cue:2511}{}黃凱。
% cue 2512 4623.050 4627.810
\hypertarget{c01:cue:2512}{}凱。
% cue 2513 4627.810 4629.350
\hypertarget{c01:cue:2513}{}在哪裡。
% cue 2514 4629.350 4630.350
\hypertarget{c01:cue:2514}{}喔我沒有看到。

\tc{1:17:10}
\speaker{學生}
% cue 2515 4630.350 4633.240
\hypertarget{c01:cue:2515}{}也是上公文講的。
% cue 2516 4633.240 4634.460
\hypertarget{c01:cue:2516}{}課。
% cue 2517 4634.460 4635.450
\hypertarget{c01:cue:2517}{}然後就是。
% cue 2518 4635.450 4636.450
\hypertarget{c01:cue:2518}{}我確實覺得這個。
% cue 2519 4636.450 4638.690
\hypertarget{c01:cue:2519}{}憲法時間非常多。
% cue 2520 4638.690 4639.690
\hypertarget{c01:cue:2520}{}也不抽象。
% cue 2521 4639.690 4640.780
\hypertarget{c01:cue:2521}{}就是。
% cue 2522 4640.780 4642.540
\hypertarget{c01:cue:2522}{}他那個。
% cue 2523 4642.540 4643.540
\hypertarget{c01:cue:2523}{}很去正大。
% cue 2524 4643.540 4644.540
\hypertarget{c01:cue:2524}{}他就丟一些那個。
% cue 2525 4644.540 4645.880
\hypertarget{c01:cue:2525}{}超過。
% cue 2526 4645.880 4648.400
\hypertarget{c01:cue:2526}{}十幾。
% cue 2527 4648.400 4651.160
\hypertarget{c01:cue:2527}{}反正就是。
% cue 2528 4651.160 4652.160
\hypertarget{c01:cue:2528}{}丟了這篇。
% cue 2529 4652.160 4653.990
\hypertarget{c01:cue:2529}{}對。
% cue 2530 4653.990 4654.990
\hypertarget{c01:cue:2530}{}就是講了一些美古法。
% cue 2531 4654.990 4655.990
\hypertarget{c01:cue:2531}{}美的一些東西。
% cue 2532 4655.990 4656.990
\hypertarget{c01:cue:2532}{}然後去做。
% cue 2533 4656.990 4658.180
\hypertarget{c01:cue:2533}{}然後我就。
% cue 2534 4658.180 4660.070
\hypertarget{c01:cue:2534}{}時間上的關係。

\tc{1:17:40}
\speaker{老師}
% cue 2535 4660.070 4661.070
\hypertarget{c01:cue:2535}{}好。

\section{憲法第八條法定程序與釋字第 384 號}

\tc{1:17:41}
% cue 2536 4661.070 4662.070
\hypertarget{c01:cue:2536}{}我問你一個很簡單的事。
% cue 2537 4662.070 4664.070
\hypertarget{c01:cue:2537}{}其他人們都可以聽著回答。
% cue 2538 4664.070 4665.510
\hypertarget{c01:cue:2538}{}當老師跟你講說。
% cue 2539 4665.510 4667.510
\hypertarget{c01:cue:2539}{}有正當程序的時候。
% cue 2540 4667.510 4668.860
\hypertarget{c01:cue:2540}{}老師有沒有交代。
% cue 2541 4668.860 4671.020
\hypertarget{c01:cue:2541}{}在中華民國憲法。
% cue 2542 4671.020 4673.610
\hypertarget{c01:cue:2542}{}正當程序的根基。
% cue 2543 4673.610 4675.760
\hypertarget{c01:cue:2543}{}他的基礎出於憲法的地方。
% cue 2544 4675.760 4679.590
\hypertarget{c01:cue:2544}{}因為我們憲法裡面沒有正當程序。
% cue 2545 4679.590 4680.540
\hypertarget{c01:cue:2545}{}這是不是。

\tc{1:18:00}
\speaker{學生}
% cue 2546 4680.540 4686.350
\hypertarget{c01:cue:2546}{}有有有。
% cue 2547 4686.350 4687.350
\hypertarget{c01:cue:2547}{}他說他是從某一條。
% cue 2548 4687.350 4688.350
\hypertarget{c01:cue:2548}{}有稱之間。

\tc{1:18:08}
\speaker{老師}
% cue 2549 4688.350 4689.660
\hypertarget{c01:cue:2549}{}哪一條。

\tc{1:18:09}
\speaker{學生}
% cue 2550 4689.660 4692.030
\hypertarget{c01:cue:2550}{}下次有機會我們再講。
% cue 2551 4692.030 4693.030
\hypertarget{c01:cue:2551}{}第八條\footnote{中華民國憲法第 8 條第 1 項規定：「人民身體之自由應予保障。除現行犯之逮捕由法律另定外，非經司法或警察機關依法定程序，不得逮捕拘禁。非由法院依法定程序，不得審問處罰。非依法定程序之逮捕、拘禁、審問、處罰，得拒絕之。」條文所用的詞是「依法定程序」，未見「正當程序」字樣。\url{https://law.moj.gov.tw/LawClass/LawAll.aspx?pcode=A0000001}。}。

\tc{1:18:13}
\speaker{老師}
% cue 2552 4693.030 4695.340
\hypertarget{c01:cue:2552}{}第八條。

\tc{1:18:15}
% cue 2553 4695.340 4695.440
\hypertarget{c01:cue:2553}{}好。
% cue 2554 4695.340 4696.340
\hypertarget{c01:cue:2554}{}那第八條。
% cue 2555 4696.340 4697.340
\hypertarget{c01:cue:2555}{}為什麼。
% cue 2556 4697.340 4698.630
\hypertarget{c01:cue:2556}{}他要跟你講。
% cue 2557 4698.630 4699.630
\hypertarget{c01:cue:2557}{}為什麼可以從那裡。
% cue 2558 4699.630 4700.630
\hypertarget{c01:cue:2558}{}顏色出來。
% cue 2559 4700.870 4703.890
\hypertarget{c01:cue:2559}{}這不只是大法官自己解釋的問題而已。
% cue 2560 4703.890 4704.890
\hypertarget{c01:cue:2560}{}而是有根據的。
% cue 2561 4704.890 4708.210
\hypertarget{c01:cue:2561}{}你想想看第八條裡面最常出現的姿勢說。
% cue 2562 4708.210 4711.170
\hypertarget{c01:cue:2562}{}非依法律正當程序不得逮捕拘禁。
% cue 2563 4711.170 4713.170
\hypertarget{c01:cue:2563}{}非依法律正當程序。
% cue 2564 4713.170 4716.550
\hypertarget{c01:cue:2564}{}非依法定程序不得逮捕拘禁。
% cue 2565 4716.550 4718.600
\hypertarget{c01:cue:2565}{}非依法定程序。

\tc{1:18:38}
% cue 2566 4718.600 4721.780
\hypertarget{c01:cue:2566}{}比如法定程序。
% cue 2567 4721.780 4723.780
\hypertarget{c01:cue:2567}{}大法官在釋字三百八十四號\footnote{司法院釋字第 384 號解釋（1995 年 7 月 28 日）。解釋文稱憲法第 8 條第 1 項所稱「依法定程序」，「係指凡限制人民身體自由之處置，不問其是否屬於刑事被告之身分，國家機關所依據之程序，須以法律規定，其內容更須實質正當，並符合憲法第二十三條所定相關之條件」；檢肅流氓條例第 5 條、第 6 條、第 7 條、第 12 條、第 21 條各規定，經宣告至遲於 1996 年 12 月 31 日失其效力。\url{https://cons.judicial.gov.tw/docdata.aspx?fid=100&id=310565}。}。
% cue 2568 4723.780 4726.850
\hypertarget{c01:cue:2568}{}法定程序意義往上提升。
% cue 2569 4726.850 4727.880
\hypertarget{c01:cue:2569}{}法定程序。
% cue 2570 4727.880 4732.030
\hypertarget{c01:cue:2570}{}如果你除了在照字面前做文意解釋。
% cue 2571 4732.030 4734.060
\hypertarget{c01:cue:2571}{}是不是法律所規定的程序。
% cue 2572 4734.060 4736.380
\hypertarget{c01:cue:2572}{}就叫法定程序。
% cue 2573 4736.380 4738.370
\hypertarget{c01:cue:2573}{}那如果是這樣。
% cue 2574 4738.370 4739.370
\hypertarget{c01:cue:2574}{}那法律規定的亂七八糟。
% cue 2575 4739.370 4741.370
\hypertarget{c01:cue:2575}{}那也是個程序啊。
% cue 2576 4741.370 4742.520
\hypertarget{c01:cue:2576}{}大法官說不可以。
% cue 2577 4742.520 4744.580
\hypertarget{c01:cue:2577}{}因為這個條法很重要。
% cue 2578 4744.580 4746.580
\hypertarget{c01:cue:2578}{}它代表的是一種憲法上的精神。
% cue 2579 4746.580 4750.580
\hypertarget{c01:cue:2579}{}所以它實際上是必須符合憲法要求的。
% cue 2580 4750.580 4752.790
\hypertarget{c01:cue:2580}{}這種才叫做法定程序。
% cue 2581 4752.790 4753.770
\hypertarget{c01:cue:2581}{}如果不符合。
% cue 2582 4753.770 4755.890
\hypertarget{c01:cue:2582}{}這個法律會必須告個憲。
% cue 2583 4755.890 4758.620
\hypertarget{c01:cue:2583}{}它沒有符合正當法律程序\footnote{英文 due process of law。美國聯邦憲法增修條文第 5 條與第 14 條第 1 項均規定，非經正當法律程序，不得剝奪任何人之生命、自由或財產；第 14 條以各州為規範對象。\url{https://www.archives.gov/founding-docs/bill-of-rights-transcript}。}。
% cue 2584 4758.620 4762.370
\hypertarget{c01:cue:2584}{}不是單獨的法律程序規定。
% cue 2585 4762.370 4763.560
\hypertarget{c01:cue:2585}{}就只有 due。
% cue 2586 4763.560 4764.870
\hypertarget{c01:cue:2586}{}所以才會有問題。
% cue 2587 4764.870 4765.870
\hypertarget{c01:cue:2587}{}How much is due?\footnote{美國聯邦最高法院 Morrissey v. Brewer, 408 U.S. 471 (1972)（1972 年 6 月 29 日判決）第 481 頁謂：「Once it is determined that due process applies, the question remains what process is due.」（確定應適用正當程序之後，仍須決定應給予何種程序。）\url{https://www.law.cornell.edu/supremecourt/text/408/471}。}
% cue 2588 4765.870 4767.150
\hypertarget{c01:cue:2588}{}How much process?
% cue 2589 4767.150 4770.220
\hypertarget{c01:cue:2589}{}How much procedure is due？

\section{全班點名（續）}

\tc{1:19:30}
% cue 2590 4770.220 4771.790
\hypertarget{c01:cue:2590}{}好，下一位。
% cue 2591 4771.790 4774.170
\hypertarget{c01:cue:2591}{}法二 C，李冠儀。
% cue 2592 4774.170 4776.100
\hypertarget{c01:cue:2592}{}李冠儀。
% cue 2593 4776.100 4777.100
\hypertarget{c01:cue:2593}{}有沒有，XX。
% cue 2594 4777.100 4779.670
\hypertarget{c01:cue:2594}{}好，好，志平。
% cue 2595 4779.670 4780.640
\hypertarget{c01:cue:2595}{}所以指示班長。
% cue 2596 4780.640 4783.470
\hypertarget{c01:cue:2596}{}好，等一下下課來告訴我你的電話。
% cue 2597 4783.470 4784.470
\hypertarget{c01:cue:2597}{}讓我找不到你。
% cue 2598 4784.470 4785.620
\hypertarget{c01:cue:2598}{}等一下加個 9。
% cue 2599 4785.620 4788.710
\hypertarget{c01:cue:2599}{}法二 C，李冠儀請說。

\tc{1:19:48}
\speaker{學生}
% cue 2600 4788.710 4808.550
\hypertarget{c01:cue:2600}{}Dail。

\tc{1:20:08}
\speaker{老師}
% cue 2601 4808.550 4810.310
\hypertarget{c01:cue:2601}{}那你很勇敢。
% cue 2602 4810.310 4812.310
\hypertarget{c01:cue:2602}{}等你跑來這裡想要補充一下。
% cue 2603 4812.310 4813.430
\hypertarget{c01:cue:2603}{}希望不要讓你太失望。

\tc{1:20:13}
% cue 2604 4813.430 4816.940
\hypertarget{c01:cue:2604}{}好，法二 C，歐乘佑。
% cue 2605 4816.940 4821.970
\hypertarget{c01:cue:2605}{}好，也是開放式課程。

\tc{1:20:21}
\speaker{學生}
% cue 2606 4821.970 4822.970
\hypertarget{c01:cue:2606}{}很早咧。
% cue 2607 4822.970 4826.020
\hypertarget{c01:cue:2607}{}就是因為老師的開放式課程。
% cue 2608 4826.020 4828.120
\hypertarget{c01:cue:2608}{}老師的課程，幫我來講反式在講。
% cue 2609 4828.120 4835.840
\hypertarget{c01:cue:2609}{}簡單跟唯一。

\tc{1:20:35}
\speaker{老師}
% cue 2610 4835.840 4836.840
\hypertarget{c01:cue:2610}{}我們先講一半一半。
% cue 2611 4836.840 4838.990
\hypertarget{c01:cue:2611}{}有一半是在講權力分立。
% cue 2612 4838.990 4841.200
\hypertarget{c01:cue:2612}{}有一半是在講基本人權。
% cue 2613 4841.200 4843.250
\hypertarget{c01:cue:2613}{}雖然基本人權的條文很少。
% cue 2614 4843.250 4845.770
\hypertarget{c01:cue:2614}{}但它的重要性很高。
% cue 2615 4845.770 4846.770
\hypertarget{c01:cue:2615}{}權力分立的條文很多。
% cue 2616 4846.770 4848.770
\hypertarget{c01:cue:2616}{}因為權力要分立需要有很多規定。
% cue 2617 4848.770 4850.960
\hypertarget{c01:cue:2617}{}但是它就差不多一半一半。
% cue 2618 4850.960 4852.950
\hypertarget{c01:cue:2618}{}所以老師一半講權力分立。
% cue 2619 4852.950 4855.680
\hypertarget{c01:cue:2619}{}一半講基本人權。
% cue 2620 4855.680 4857.030
\hypertarget{c01:cue:2620}{}大概是這樣子。

\tc{1:20:57}
% cue 2621 4857.030 4859.310
\hypertarget{c01:cue:2621}{}好，然後下面這個詞我也不太會名。
% cue 2622 4859.310 4861.310
\hypertarget{c01:cue:2622}{}這叫歐將律。
% cue 2623 4861.310 4862.330
\hypertarget{c01:cue:2623}{}歐將路。
% cue 2624 4862.330 4880.000
\hypertarget{c01:cue:2624}{}江 無什麼？ 5-5-5 無 江 入？ 無什麼？ 無 佑 露 佑 佑 佑 佑 佑 佑 佑 佑 佑 佑 佑 佑 佑 佑 佑 佑 佑 佑 佑 佑 佑 佑 佑 佑 佑 佑 佑 佑 佑 佑 佑 佑 佑 佑 佑 好 謝謝 佑 佑 佑 佑 佑 佑 佑。
% cue 2625 4880.000 4889.710
\hypertarget{c01:cue:2625}{}有字是什麼呢？ 回家去查一查因為你的證明字很特別。
% cue 2626 4889.710 4894.870
\hypertarget{c01:cue:2626}{}啥歌曲 要不來是爸爸曲 不知道啊 回去問一下為什麼。
% cue 2627 4894.870 4899.460
\hypertarget{c01:cue:2627}{}幫其查查開這個證明字是因為很多人取名字是有原因的。
% cue 2628 4899.460 4901.070
\hypertarget{c01:cue:2628}{}也許他們會對你講話。

\tc{1:21:41}
% cue 2629 4901.070 4902.350
\hypertarget{c01:cue:2629}{}好來 請說。

\tc{1:21:42}
\speaker{學生}
% cue 2630 4902.350 4903.650
\hypertarget{c01:cue:2630}{}PASS。

\tc{1:21:43}
\speaker{老師}
% cue 2631 4903.650 4904.540
\hypertarget{c01:cue:2631}{}PASS 好。
% cue 2632 4904.540 4916.410
\hypertarget{c01:cue:2632}{}PASS 林玉環限迫於十字。
% cue 2633 4916.410 4917.200
\hypertarget{c01:cue:2633}{}限迫。
% cue 2634 4917.200 4919.020
\hypertarget{c01:cue:2634}{}好 限迫之四。
% cue 2635 4919.020 4922.930
\hypertarget{c01:cue:2635}{}然後你收過蘇格拉底的三個。
% cue 2636 4922.930 4925.100
\hypertarget{c01:cue:2636}{}然後就有話。

\tc{1:22:05}
% cue 2637 4925.100 4927.400
\hypertarget{c01:cue:2637}{}好 那你要幫我一個忙。
% cue 2638 4927.400 4929.080
\hypertarget{c01:cue:2638}{}你要訓練自己說話。
% cue 2639 4929.080 4931.820
\hypertarget{c01:cue:2639}{}你要講到說 找到講話的機會。
% cue 2640 4931.820 4933.740
\hypertarget{c01:cue:2640}{}所以勇敢舉手講話。
% cue 2641 4933.740 4936.470
\hypertarget{c01:cue:2641}{}我已經講過了 你講什麼都會被扣分。
% cue 2642 4936.470 4938.440
\hypertarget{c01:cue:2642}{}但只有可能會被加分。
% cue 2643 4938.440 4940.650
\hypertarget{c01:cue:2643}{}那你必須要勇敢表示一切。
% cue 2644 4940.650 4941.870
\hypertarget{c01:cue:2644}{}你怎麼有可能會被加分。
% cue 2645 4941.870 4943.500
\hypertarget{c01:cue:2645}{}你都不講 你要什麼加分。
% cue 2646 4943.500 4945.170
\hypertarget{c01:cue:2646}{}我剛剛也是這麼小孩的朋友。
% cue 2647 4945.170 4947.410
\hypertarget{c01:cue:2647}{}要講出來事情。
% cue 2648 4947.410 4949.520
\hypertarget{c01:cue:2648}{}我給你講。

\tc{1:22:29}
\speaker{學生}
% cue 2649 4949.520 4951.440
\hypertarget{c01:cue:2649}{}師 庭 您。
% cue 2650 4951.440 4958.800
\hypertarget{c01:cue:2650}{}覺得比較好。
% cue 2651 4958.800 4960.660
\hypertarget{c01:cue:2651}{}模糊。
% cue 2652 4960.660 4963.430
\hypertarget{c01:cue:2652}{}幫我聽懂 幫我聽懂。

\tc{1:22:43}
\speaker{老師}
% cue 2653 4963.430 4965.680
\hypertarget{c01:cue:2653}{}那現在老師只能幫你小壞。
% cue 2654 4965.680 4966.650
\hypertarget{c01:cue:2654}{}就是有關。
% cue 2655 4966.650 4968.770
\hypertarget{c01:cue:2655}{}而且我們不是把現場的。
% cue 2656 4968.770 4970.960
\hypertarget{c01:cue:2656}{}有關於權力分立的各章都講。
% cue 2657 4970.960 4973.390
\hypertarget{c01:cue:2657}{}我只是挑一些具有代表性的案子。
% cue 2658 4973.390 4976.360
\hypertarget{c01:cue:2658}{}當然我希望這些具有代表性的案子。
% cue 2659 4976.360 4979.250
\hypertarget{c01:cue:2659}{}慢慢的幫你勾勒出一個完整的。
% cue 2660 4979.250 4982.860
\hypertarget{c01:cue:2660}{}所謂的五權分治 平等相維\footnote{司法院釋字第 175 號解釋（1982 年 5 月 25 日）解釋文：「司法院為國家最高司法機關，基於五權分治彼此相維之憲政體制，就其所掌有關司法機關之組織及司法權行使之事項，得向立法院提出法律案。」湯德宗《權力分立新論 卷一》以「五權分治、平等相維」稱中華民國憲法的五權分立制衡；課綱第 4 週的主題為中華民國憲法結構，副題「五權分治、平等相維」。\url{https://cons.judicial.gov.tw/docdata.aspx?fid=100&id=310356}。}。
% cue 2661 4982.860 4984.400
\hypertarget{c01:cue:2661}{}他終於長的什麼樣子。
% cue 2662 4984.400 4986.480
\hypertarget{c01:cue:2662}{}他跟其他的三權分立。
% cue 2663 4986.480 4988.530
\hypertarget{c01:cue:2663}{}有些地方高度相似。
% cue 2664 4988.530 4990.400
\hypertarget{c01:cue:2664}{}有些地方完全不一樣。
% cue 2665 4990.400 4991.340
\hypertarget{c01:cue:2665}{}好不好。

\tc{1:23:11}
% cue 2666 4991.340 4994.250
\hypertarget{c01:cue:2666}{}許花華說。

\tc{1:23:14}
\speaker{學生}
% cue 2667 4994.250 5003.590
\hypertarget{c01:cue:2667}{}哪一個老師叫我跟他講。
% cue 2668 5003.590 5004.230
\hypertarget{c01:cue:2668}{}對不對。

\tc{1:23:24}
\speaker{老師}
% cue 2669 5004.230 5006.920
\hypertarget{c01:cue:2669}{}哪一個老師叫 張志中嗎。
% cue 2670 5006.920 5009.210
\hypertarget{c01:cue:2670}{}張志中嗎。
% cue 2671 5009.210 5009.310
\hypertarget{c01:cue:2671}{}張志中。
% cue 2672 5009.210 5011.190
\hypertarget{c01:cue:2672}{}哦 因為他留得嘛。
% cue 2673 5011.190 5013.280
\hypertarget{c01:cue:2673}{}但是實際上。
% cue 2674 5013.280 5015.690
\hypertarget{c01:cue:2674}{}這次我聽懂你的意思。
% cue 2675 5015.690 5016.970
\hypertarget{c01:cue:2675}{}應該要回到本土來。

\tc{1:23:36}
\speaker{學生}
% cue 2676 5016.970 5017.930
\hypertarget{c01:cue:2676}{}對啊 廢國。
% cue 2677 5017.930 5020.320
\hypertarget{c01:cue:2677}{}對。

\tc{1:23:40}
\speaker{老師}
% cue 2678 5020.320 5024.980
\hypertarget{c01:cue:2678}{}那你要問我才有辦法回應。
% cue 2679 5024.980 5025.540
\hypertarget{c01:cue:2679}{}你問我。
% cue 2680 5025.540 5026.590
\hypertarget{c01:cue:2680}{}我照理一下。

\tc{1:23:46}
% cue 2681 5026.590 5028.760
\hypertarget{c01:cue:2681}{}現在這個名單。
% cue 2682 5028.760 5030.040
\hypertarget{c01:cue:2682}{}我都全部念完了。
% cue 2683 5030.040 5032.290
\hypertarget{c01:cue:2683}{}沒有在這個名單上的同學。
% cue 2684 5032.290 5033.090
\hypertarget{c01:cue:2684}{}舉手我看一下。
% cue 2685 5033.090 5034.520
\hypertarget{c01:cue:2685}{}應該有。
% cue 2686 5034.520 5036.540
\hypertarget{c01:cue:2686}{}那我把這張紙傳下去。
% cue 2687 5036.540 5037.660
\hypertarget{c01:cue:2687}{}你在上面填給你。
% cue 2688 5037.660 5039.970
\hypertarget{c01:cue:2688}{}不然就永遠不會出現在我的名字。
% cue 2689 5039.970 5041.170
\hypertarget{c01:cue:2689}{}我就沒辦法叫你。
% cue 2690 5041.170 5043.310
\hypertarget{c01:cue:2690}{}一位在這裡。
% cue 2691 5043.310 5044.190
\hypertarget{c01:cue:2691}{}第二位剛這邊。
% cue 2692 5044.190 5044.670
\hypertarget{c01:cue:2692}{}好。
% cue 2693 5044.670 5046.020
\hypertarget{c01:cue:2693}{}你之後傳給他。
% cue 2694 5046.020 5047.690
\hypertarget{c01:cue:2694}{}然後你看到他在下一層。
% cue 2695 5047.690 5049.220
\hypertarget{c01:cue:2695}{}然後你之後在這邊。
% cue 2696 5049.220 5050.320
\hypertarget{c01:cue:2696}{}這邊有兩個人。
% cue 2697 5050.320 5051.630
\hypertarget{c01:cue:2697}{}這樣可以嗎。
% cue 2698 5051.630 5052.270
\hypertarget{c01:cue:2698}{}五個人。
% cue 2699 5052.270 5053.450
\hypertarget{c01:cue:2699}{}我這樣搶沒有錯。
% cue 2700 5053.450 5053.930
\hypertarget{c01:cue:2700}{}好。

\tc{1:24:13}
% cue 2701 5053.930 5054.730
\hypertarget{c01:cue:2701}{}其他同學。
% cue 2702 5054.730 5056.570
\hypertarget{c01:cue:2702}{}現在請看一下老師的簡歷。
% cue 2703 5056.570 5057.900
\hypertarget{c01:cue:2703}{}現在開放題目。
% cue 2704 5057.900 5059.780
\hypertarget{c01:cue:2704}{}你們暗謀什麼。
% cue 2705 5059.780 5060.300
\hypertarget{c01:cue:2705}{}為什麼。
% cue 2706 5060.300 5060.930
\hypertarget{c01:cue:2706}{}但是。
% cue 2707 5060.930 5062.050
\hypertarget{c01:cue:2707}{}今天老師的回答。
% cue 2708 5062.050 5063.250
\hypertarget{c01:cue:2708}{}緊線在課堂上。
% cue 2709 5063.250 5064.400
\hypertarget{c01:cue:2709}{}除了真的教室以外。
% cue 2710 5064.400 5065.100
\hypertarget{c01:cue:2710}{}我們該不會。
% cue 2711 5065.100 5069.630
\hypertarget{c01:cue:2711}{}你不是舉手嗎。
% cue 2712 5069.630 5070.990
\hypertarget{c01:cue:2712}{}沒有人跟你搶。
% cue 2713 5070.990 5071.900
\hypertarget{c01:cue:2713}{}我就不看。

\tc{1:24:31}
% cue 2714 5071.900 5082.380
\hypertarget{c01:cue:2714}{}我當時去的行星。
% cue 2715 5082.380 5085.070
\hypertarget{c01:cue:2715}{}現在我相信應該還是。
% cue 2716 5085.070 5086.440
\hypertarget{c01:cue:2716}{}他們的學校。
% cue 2717 5086.440 5087.160
\hypertarget{c01:cue:2717}{}Law School。
% cue 2718 5087.160 5088.760
\hypertarget{c01:cue:2718}{}叫做 Professional School。
% cue 2719 5088.760 5089.560
\hypertarget{c01:cue:2719}{}因為我們單子的話。
% cue 2720 5089.560 5090.760
\hypertarget{c01:cue:2720}{}就是職業學校。
% cue 2721 5090.760 5092.790
\hypertarget{c01:cue:2721}{}他們非要職業就業。
% cue 2722 5092.790 5093.350
\hypertarget{c01:cue:2722}{}當律師。
% cue 2723 5093.350 5094.350
\hypertarget{c01:cue:2723}{}當法官的證明。
% cue 2724 5094.350 5096.110
\hypertarget{c01:cue:2724}{}他們的法官實際上。
% cue 2725 5096.110 5097.550
\hypertarget{c01:cue:2725}{}一定是律師轉任。
% cue 2726 5097.550 5099.240
\hypertarget{c01:cue:2726}{}所以就是培養律師。
% cue 2727 5099.240 5100.710
\hypertarget{c01:cue:2727}{}就是法官司的人。
% cue 2728 5100.710 5101.490
\hypertarget{c01:cue:2728}{}好。
% cue 2729 5101.490 5103.520
\hypertarget{c01:cue:2729}{}所以嚴格的來講。
% cue 2730 5103.520 5104.680
\hypertarget{c01:cue:2730}{}沒有為了。
% cue 2731 5104.680 5106.580
\hypertarget{c01:cue:2731}{}像我們這些高生。
% cue 2732 5106.580 5108.420
\hypertarget{c01:cue:2732}{}追求高深學問的人開始。
% cue 2733 5108.420 5109.530
\hypertarget{c01:cue:2733}{}不是為了。
% cue 2734 5109.530 5110.520
\hypertarget{c01:cue:2734}{}都是。
% cue 2735 5110.520 5112.050
\hypertarget{c01:cue:2735}{}他們就是三年。
% cue 2736 5112.050 5113.170
\hypertarget{c01:cue:2736}{}三年就被，
% cue 2737 5113.170 5114.750
\hypertarget{c01:cue:2737}{}像我們私隱。
% cue 2738 5114.750 5114.850
\hypertarget{c01:cue:2738}{}三年。
% cue 2739 5114.820 5116.870
\hypertarget{c01:cue:2739}{}那因為，
% cue 2740 5116.870 5118.060
\hypertarget{c01:cue:2740}{}哈爾格爾的學生。
% cue 2741 5118.060 5119.280
\hypertarget{c01:cue:2741}{}這樣講。
% cue 2742 5119.280 5119.920
\hypertarget{c01:cue:2742}{}應該。
% cue 2743 5119.920 5121.010
\hypertarget{c01:cue:2743}{}我們可以同意吧。
% cue 2744 5121.010 5122.740
\hypertarget{c01:cue:2744}{}全世界。
% cue 2745 5122.740 5123.460
\hypertarget{c01:cue:2745}{}包含。
% cue 2746 5123.460 5125.070
\hypertarget{c01:cue:2746}{}你要進到哈爾格爾。
% cue 2747 5125.070 5127.990
\hypertarget{c01:cue:2747}{}你不是領域。
% cue 2748 5127.990 5129.220
\hypertarget{c01:cue:2748}{}你不到哈爾格爾。
% cue 2749 5129.220 5130.830
\hypertarget{c01:cue:2749}{}每一年都有跳舞。
% cue 2750 5130.830 5134.380
\hypertarget{c01:cue:2750}{}而且都是放流。
% cue 2751 5134.380 5135.310
\hypertarget{c01:cue:2751}{}就到哈爾格爾。
% cue 2752 5135.310 5137.200
\hypertarget{c01:cue:2752}{}他不可能再維持 A+。
% cue 2753 5137.200 5138.510
\hypertarget{c01:cue:2753}{}這個可能。
% cue 2754 5138.510 5139.560
\hypertarget{c01:cue:2754}{}不出去會跳舞。
% cue 2755 5139.560 5140.460
\hypertarget{c01:cue:2755}{}每年都都是 A。
% cue 2756 5140.460 5141.670
\hypertarget{c01:cue:2756}{}我們不知道到底。
% cue 2757 5141.670 5142.510
\hypertarget{c01:cue:2757}{}什麼想不開。
% cue 2758 5142.510 5143.530
\hypertarget{c01:cue:2758}{}我們不知道。

\tc{1:25:43}
% cue 2759 5143.530 5144.370
\hypertarget{c01:cue:2759}{}那，
% cue 2760 5144.370 5145.320
\hypertarget{c01:cue:2760}{}他的課。
% cue 2761 5145.320 5146.520
\hypertarget{c01:cue:2761}{}就是開給。
% cue 2762 5146.520 5148.680
\hypertarget{c01:cue:2762}{}他們叫 JD。
% cue 2763 5148.680 5150.730
\hypertarget{c01:cue:2763}{}為什麼叫 JD。
% cue 2764 5150.730 5154.120
\hypertarget{c01:cue:2764}{}JD 就是 Juris Doctor 的簡稱。
% cue 2765 5154.120 5154.980
\hypertarget{c01:cue:2765}{}就是，
% cue 2766 5154.980 5159.300
\hypertarget{c01:cue:2766}{}因為你已經拿到了一個大學的學士學位。
% cue 2767 5159.300 5162.480
\hypertarget{c01:cue:2767}{}才能來念 Professional School。
% cue 2768 5162.480 5164.150
\hypertarget{c01:cue:2768}{}那你是想要做律師。
% cue 2769 5164.150 5164.970
\hypertarget{c01:cue:2769}{}所以你來念什麼。
% cue 2770 5164.970 5166.970
\hypertarget{c01:cue:2770}{}你可以去其他的 Professional School。
% cue 2771 5166.970 5168.690
\hypertarget{c01:cue:2771}{}你要從你的 Design School。
% cue 2772 5168.690 5170.130
\hypertarget{c01:cue:2772}{}你去學設計。
% cue 2773 5170.130 5171.410
\hypertarget{c01:cue:2773}{}出來叫設計師。
% cue 2774 5171.410 5172.810
\hypertarget{c01:cue:2774}{}等等。
% cue 2775 5172.810 5175.820
\hypertarget{c01:cue:2775}{}那因為是大學之後的學位。
% cue 2776 5175.820 5177.450
\hypertarget{c01:cue:2776}{}所以如果沒有給你 D。
% cue 2777 5177.450 5178.380
\hypertarget{c01:cue:2777}{}會來嗎？
% cue 2778 5178.380 5179.850
\hypertarget{c01:cue:2778}{}用 JD。
% cue 2779 5179.850 5181.270
\hypertarget{c01:cue:2779}{}但是大家都知道。
% cue 2780 5181.270 5182.090
\hypertarget{c01:cue:2780}{}JD。
% cue 2781 5182.090 5184.050
\hypertarget{c01:cue:2781}{}學制上是第一個法律學位。

\tc{1:26:24}
% cue 2782 5184.050 5186.410
\hypertarget{c01:cue:2782}{}那在我們這裡。
% cue 2783 5186.410 5188.270
\hypertarget{c01:cue:2783}{}我們的大學部的時候。
% cue 2784 5188.270 5190.410
\hypertarget{c01:cue:2784}{}法律系是放在大學部。
% cue 2785 5190.410 5192.200
\hypertarget{c01:cue:2785}{}所以我們通常這種的書學。
% cue 2786 5192.200 5193.520
\hypertarget{c01:cue:2786}{}叫 LLB。
% cue 2787 5193.520 5195.600
\hypertarget{c01:cue:2787}{}就是 Bachelor of Laws。
% cue 2788 5195.600 5197.970
\hypertarget{c01:cue:2788}{}那如果我們要再念。
% cue 2789 5197.970 5199.530
\hypertarget{c01:cue:2789}{}在臺灣也是 LL。
% cue 2790 5199.530 5201.620
\hypertarget{c01:cue:2790}{}在美國你 JD 念完。
% cue 2791 5201.620 5202.340
\hypertarget{c01:cue:2791}{}如果，
% cue 2792 5202.340 5203.340
\hypertarget{c01:cue:2792}{}還要再念。
% cue 2793 5203.340 5204.690
\hypertarget{c01:cue:2793}{}那也是 LL。
% cue 2794 5204.690 5206.240
\hypertarget{c01:cue:2794}{}它的差別。
% cue 2795 5206.240 5208.110
\hypertarget{c01:cue:2795}{}是在 LL 開始。
% cue 2796 5208.110 5209.900
\hypertarget{c01:cue:2796}{}就是要做學術了。
% cue 2797 5209.900 5211.850
\hypertarget{c01:cue:2797}{}所以你會有一個主題。
% cue 2798 5211.410 5213.410
\hypertarget{c01:cue:2798}{}你會有一個範圍。
% cue 2799 5213.410 5214.410
\hypertarget{c01:cue:2799}{}你不可能什麼都學。
% cue 2800 5214.410 5215.570
\hypertarget{c01:cue:2800}{}像你現在大學裡面。
% cue 2801 5215.570 5216.570
\hypertarget{c01:cue:2801}{}是不是什麼都要學。
% cue 2802 5216.570 5217.570
\hypertarget{c01:cue:2802}{}剛剛人家憲法刑法。
% cue 2803 5217.570 5218.570
\hypertarget{c01:cue:2803}{}民法刑法。
% cue 2804 5218.570 5220.630
\hypertarget{c01:cue:2804}{}訴訟法什麼都要學。
% cue 2805 5220.630 5222.010
\hypertarget{c01:cue:2805}{}上法。
% cue 2806 5222.010 5223.160
\hypertarget{c01:cue:2806}{}美國也是這樣。
% cue 2807 5223.160 5225.900
\hypertarget{c01:cue:2807}{}那你將來到了 LL.M.\footnote{Master of Laws，法學碩士。哈佛法學院的 LL.M. 為一年制學位，每年約有來自 65 國左右的 180 名學生。\url{https://hls.harvard.edu/graduate-program/ll-m-program/}。} 以後。
% cue 2808 5225.900 5227.310
\hypertarget{c01:cue:2808}{}會有點專業的。
% cue 2809 5227.310 5228.530
\hypertarget{c01:cue:2809}{}不可能什麼都不用。

\tc{1:27:08}
% cue 2810 5228.530 5230.840
\hypertarget{c01:cue:2810}{}那你如果要進 SJD。
% cue 2811 5230.840 5234.320
\hypertarget{c01:cue:2811}{}他是 Doctor of Juridical Science。
% cue 2812 5234.320 5236.410
\hypertarget{c01:cue:2812}{}那這個就是一定要有論文。
% cue 2813 5236.410 5238.660
\hypertarget{c01:cue:2813}{}那 LL.M. 不一定要有論文。
% cue 2814 5238.660 5239.620
\hypertarget{c01:cue:2814}{}你可以自己選。
% cue 2815 5239.620 5240.740
\hypertarget{c01:cue:2815}{}那我當然是學有論文。
% cue 2816 5240.740 5242.990
\hypertarget{c01:cue:2816}{}在臺大經驗的都說的事情。
% cue 2817 5242.990 5244.490
\hypertarget{c01:cue:2817}{}開玩笑。
% cue 2818 5244.490 5245.490
\hypertarget{c01:cue:2818}{}臺大的說是考慮了三年。
% cue 2819 5245.490 5246.490
\hypertarget{c01:cue:2819}{}又要念德文。
% cue 2820 5246.490 5247.490
\hypertarget{c01:cue:2820}{}又要念法文。
% cue 2821 5247.490 5247.590
\hypertarget{c01:cue:2821}{}又跟英文。
% cue 2822 5247.490 5249.320
\hypertarget{c01:cue:2822}{}類似的。
% cue 2823 5249.320 5250.570
\hypertarget{c01:cue:2823}{}他們要寫論文。
% cue 2824 5250.570 5251.520
\hypertarget{c01:cue:2824}{}我的擅長。
% cue 2825 5251.520 5252.760
\hypertarget{c01:cue:2825}{}我的專長是寫論文。
% cue 2826 5252.760 5253.820
\hypertarget{c01:cue:2826}{}我的專長是寫論文。
% cue 2827 5253.820 5254.820
\hypertarget{c01:cue:2827}{}我的專長不見得是考試。
% cue 2828 5254.820 5255.820
\hypertarget{c01:cue:2828}{}尤其是英文考試。
% cue 2829 5255.820 5257.260
\hypertarget{c01:cue:2829}{}對不對？
% cue 2830 5257.260 5258.260
\hypertarget{c01:cue:2830}{}但是很多學校的 LL.M.。
% cue 2831 5258.260 5259.260
\hypertarget{c01:cue:2831}{}是沒有寫論文。
% cue 2832 5259.260 5261.180
\hypertarget{c01:cue:2832}{}寫論文。
% cue 2833 5261.180 5262.720
\hypertarget{c01:cue:2833}{}不容不如一些人家人知道。
% cue 2834 5262.720 5264.740
\hypertarget{c01:cue:2834}{}這是另外一件事。

\tc{1:27:44}
% cue 2835 5264.740 5267.930
\hypertarget{c01:cue:2835}{}那有沒有為我們專門開的課。
% cue 2836 5267.930 5268.030
\hypertarget{c01:cue:2836}{}沒有。
% cue 2837 5267.930 5268.930
\hypertarget{c01:cue:2837}{}所有人各都是 JD。
% cue 2838 5268.930 5270.120
\hypertarget{c01:cue:2838}{}所有來的學生。
% cue 2839 5270.120 5271.570
\hypertarget{c01:cue:2839}{}都是 JD 的學生。
% cue 2840 5271.570 5272.690
\hypertarget{c01:cue:2840}{}但是你會看。
% cue 2841 5272.690 5274.720
\hypertarget{c01:cue:2841}{}有一些課程。
% cue 2842 5274.720 5275.780
\hypertarget{c01:cue:2842}{}它就是比較難的。
% cue 2843 5275.780 5278.340
\hypertarget{c01:cue:2843}{}那這個時候呢，
% cue 2844 5278.340 5280.220
\hypertarget{c01:cue:2844}{}就會有比較多的 LL.M.。
% cue 2845 5280.220 5282.270
\hypertarget{c01:cue:2845}{}和 SJD 的學生。
% cue 2846 5282.270 5284.840
\hypertarget{c01:cue:2846}{}SJD 可能很少很少。
% cue 2847 5284.840 5287.280
\hypertarget{c01:cue:2847}{}少年學校。
% cue 2848 5287.280 5288.560
\hypertarget{c01:cue:2848}{}少年學校。
% cue 2849 5288.560 5295.350
\hypertarget{c01:cue:2849}{}少年學校。
% cue 2850 5295.350 5298.030
\hypertarget{c01:cue:2850}{}我的什麼。

\tc{1:28:18}
% cue 2851 5298.030 5300.900
\hypertarget{c01:cue:2851}{}每個人的情況不一樣。
% cue 2852 5300.900 5302.990
\hypertarget{c01:cue:2852}{}我們不是叫做 Personal Statement。
% cue 2853 5302.990 5305.020
\hypertarget{c01:cue:2853}{}我們只有兩種東西。
% cue 2854 5305.020 5306.210
\hypertarget{c01:cue:2854}{}有一種東西呢，
% cue 2855 5306.210 5308.200
\hypertarget{c01:cue:2855}{}比較像是你的 CV。
% cue 2856 5308.200 5310.550
\hypertarget{c01:cue:2856}{}那 CV 的這種東西就是，
% cue 2857 5310.550 5314.080
\hypertarget{c01:cue:2857}{}我在過去大學的成績。
% cue 2858 5314.080 5315.240
\hypertarget{c01:cue:2858}{}我參加了社團。
% cue 2859 5315.240 5317.210
\hypertarget{c01:cue:2859}{}我得到過的獎。
% cue 2860 5317.210 5318.330
\hypertarget{c01:cue:2860}{}比方說，
% cue 2861 5318.330 5320.550
\hypertarget{c01:cue:2861}{}老師是民國 66 年。
% cue 2862 5320.550 5324.790
\hypertarget{c01:cue:2862}{}今年節全國金融大會的主席。
% cue 2863 5324.790 5326.490
\hypertarget{c01:cue:2863}{}現在還沒有這種智慧。
% cue 2864 5326.490 5327.490
\hypertarget{c01:cue:2864}{}我們不知道。
% cue 2865 5327.490 5330.100
\hypertarget{c01:cue:2865}{}反正我有很多獎。
% cue 2866 5330.100 5331.280
\hypertarget{c01:cue:2866}{}什麼什麼比賽。
% cue 2867 5331.280 5331.380
\hypertarget{c01:cue:2867}{}辯論賽。
% cue 2868 5331.280 5332.280
\hypertarget{c01:cue:2868}{}什麼獎。
% cue 2869 5332.280 5333.270
\hypertarget{c01:cue:2869}{}什麼獎。
% cue 2870 5333.270 5334.240
\hypertarget{c01:cue:2870}{}什麼獎。
% cue 2871 5334.240 5335.560
\hypertarget{c01:cue:2871}{}論文賽。
% cue 2872 5335.560 5335.660
\hypertarget{c01:cue:2872}{}什麼獎。
% cue 2873 5335.560 5336.560
\hypertarget{c01:cue:2873}{}我學生社會。
% cue 2874 5336.560 5337.810
\hypertarget{c01:cue:2874}{}我幹過什麼什麼。
% cue 2875 5337.810 5338.750
\hypertarget{c01:cue:2875}{}這就是我的 CV。

\tc{1:28:58}
% cue 2876 5338.750 5343.070
\hypertarget{c01:cue:2876}{}那我們是跟請 LL.M. 的人。
% cue 2877 5343.070 5345.200
\hypertarget{c01:cue:2877}{}也要寫一個。
% cue 2878 5345.200 5347.900
\hypertarget{c01:cue:2878}{}Standing Man。
% cue 2879 5347.900 5350.600
\hypertarget{c01:cue:2879}{}那我在那個裡面會講。
% cue 2880 5350.600 5353.720
\hypertarget{c01:cue:2880}{}我過去我的學術研究。
% cue 2881 5353.720 5354.720
\hypertarget{c01:cue:2881}{}是做什麼的。
% cue 2882 5354.720 5357.260
\hypertarget{c01:cue:2882}{}比方說我在臺大的碩士論文。
% cue 2883 5357.260 5359.320
\hypertarget{c01:cue:2883}{}去用哪些資料。
% cue 2884 5359.320 5360.320
\hypertarget{c01:cue:2884}{}寫什麼題目。
% cue 2885 5360.320 5361.320
\hypertarget{c01:cue:2885}{}我有什麼想法。
% cue 2886 5361.320 5363.630
\hypertarget{c01:cue:2886}{}像我現在來 Harvard。
% cue 2887 5363.630 5366.630
\hypertarget{c01:cue:2887}{}我打算還要繼續在哪個領域。
% cue 2888 5366.630 5367.850
\hypertarget{c01:cue:2888}{}我跟其他人一樣。
% cue 2889 5367.850 5368.850
\hypertarget{c01:cue:2889}{}我是公法類的。
% cue 2890 5368.850 5371.070
\hypertarget{c01:cue:2890}{}我是教育公費生出去的。
% cue 2891 5371.070 5373.070
\hypertarget{c01:cue:2891}{}所以我只修。
% cue 2892 5373.070 5373.170
\hypertarget{c01:cue:2892}{}我。
% cue 2893 5373.070 5374.070
\hypertarget{c01:cue:2893}{}我不是學校。
% cue 2894 5374.070 5375.070
\hypertarget{c01:cue:2894}{}不是國家的國際。
% cue 2895 5375.070 5376.070
\hypertarget{c01:cue:2895}{}我基本上。
% cue 2896 5376.070 5378.030
\hypertarget{c01:cue:2896}{}我只管公法類的東西。
% cue 2897 5378.030 5379.160
\hypertarget{c01:cue:2897}{}而且我只上沒過。
% cue 2898 5379.160 5381.700
\hypertarget{c01:cue:2898}{}上那些其他。
% cue 2899 5381.700 5382.700
\hypertarget{c01:cue:2899}{}都沒有什麼。
% cue 2900 5382.700 5383.980
\hypertarget{c01:cue:2900}{}中國法國的國。
% cue 2901 5383.980 5385.650
\hypertarget{c01:cue:2901}{}我是純的。
% cue 2902 5385.650 5387.350
\hypertarget{c01:cue:2902}{}為學術區。
% cue 2903 5387.350 5389.670
\hypertarget{c01:cue:2903}{}所以這是我的時代第一代。
% cue 2904 5389.670 5390.670
\hypertarget{c01:cue:2904}{}這個是幫助。
% cue 2905 5390.670 5391.670
\hypertarget{c01:cue:2905}{}要不要募取的人。
% cue 2906 5391.670 5392.670
\hypertarget{c01:cue:2906}{}他來說決定。
% cue 2907 5392.670 5394.210
\hypertarget{c01:cue:2907}{}這個人。
% cue 2908 5394.210 5395.180
\hypertarget{c01:cue:2908}{}還要不要。
% cue 2909 5395.180 5397.010
\hypertarget{c01:cue:2909}{}我是很清楚的。
% cue 2910 5397.010 5398.130
\hypertarget{c01:cue:2910}{}那至於我申請博士的時候。

\tc{1:29:58}
% cue 2911 5398.130 5399.130
\hypertarget{c01:cue:2911}{}會更複雜。
% cue 2912 5399.130 5402.310
\hypertarget{c01:cue:2912}{}這個時候我已經有一個很清楚的論文集合。
% cue 2913 5402.310 5403.820
\hypertarget{c01:cue:2913}{}我要針對這個論文集合。
% cue 2914 5403.820 5405.850
\hypertarget{c01:cue:2914}{}我要寫一下我的動機。
% cue 2915 5405.850 5407.040
\hypertarget{c01:cue:2915}{}我為什麼要做這個。
% cue 2916 5407.040 5409.130
\hypertarget{c01:cue:2916}{}從我的研究目的到我的研究方法。
% cue 2917 5409.130 5411.130
\hypertarget{c01:cue:2917}{}以及各個章節。
% cue 2918 5411.130 5412.130
\hypertarget{c01:cue:2918}{}我打算要做些什麼。
% cue 2919 5412.130 5413.100
\hypertarget{c01:cue:2919}{}我事實上已經有想法。
% cue 2920 5413.100 5414.100
\hypertarget{c01:cue:2920}{}我去完成了。
% cue 2921 5414.100 5417.760
\hypertarget{c01:cue:2921}{}那這個時候他審查的。
% cue 2922 5417.760 5418.940
\hypertarget{c01:cue:2922}{}主要是。
% cue 2923 5418.940 5421.310
\hypertarget{c01:cue:2923}{}他們叫 SJD Proposal。
% cue 2924 5421.310 5422.420
\hypertarget{c01:cue:2924}{}論文大。
% cue 2925 5422.420 5423.770
\hypertarget{c01:cue:2925}{}論文大。
% cue 2926 5423.770 5425.440
\hypertarget{c01:cue:2926}{}所以是不同的行星。
% cue 2927 5425.440 5426.370
\hypertarget{c01:cue:2927}{}不同的市場。

\tc{1:30:26}
% cue 2928 5426.370 5428.730
\hypertarget{c01:cue:2928}{}但如果你在美國要念 JD。
% cue 2929 5428.730 5429.730
\hypertarget{c01:cue:2929}{}他就不考慮這個。
% cue 2930 5429.730 5431.730
\hypertarget{c01:cue:2930}{}他可能會合成 GRE\footnote{Graduate Record Examination，美國研究所入學考試。哈佛法學院於 2017 年 3 月 8 日宣布試行方案，自 2017 年秋季起，申請三年制 J.D. 者得以 GRE 或 LSAT（Law School Admission Test）成績擇一提出。\url{https://hls.harvard.edu/today/gre/}。} 的成績。
% cue 2931 5431.730 5434.330
\hypertarget{c01:cue:2931}{}JD 本身可能要搞 GRE。
% cue 2932 5434.330 5435.330
\hypertarget{c01:cue:2932}{}但他有一個特別的。
% cue 2933 5435.330 5437.450
\hypertarget{c01:cue:2933}{}叫做 LSAT。
% cue 2934 5437.450 5438.540
\hypertarget{c01:cue:2934}{}Law School Admission Test。
% cue 2935 5438.540 5439.830
\hypertarget{c01:cue:2935}{}他要看一個付出高度。
% cue 2936 5439.830 5441.980
\hypertarget{c01:cue:2936}{}他用這種方法來判斷。
% cue 2937 5441.980 5442.980
\hypertarget{c01:cue:2937}{}你的個性。
% cue 2938 5442.980 5443.980
\hypertarget{c01:cue:2938}{}你的性向。
% cue 2939 5443.980 5445.980
\hypertarget{c01:cue:2939}{}適不適合念 JD。
% cue 2940 5445.980 5459.230
\hypertarget{c01:cue:2940}{}所以我要怎麼說呢，
% cue 2941 5459.230 5461.230
\hypertarget{c01:cue:2941}{}我就說美國人就是這樣。
% cue 2942 5461.230 5469.680
\hypertarget{c01:cue:2942}{}看得出來什麼。
% cue 2943 5469.680 5470.740
\hypertarget{c01:cue:2943}{}你說很難看的。
% cue 2944 5470.740 5471.740
\hypertarget{c01:cue:2944}{}所以每天主張的才。

\tc{1:31:11}
% cue 2945 5471.740 5474.820
\hypertarget{c01:cue:2945}{}都要這麼頂的學生。
% cue 2946 5474.820 5475.810
\hypertarget{c01:cue:2946}{}那我要怎麼。
% cue 2947 5475.810 5477.320
\hypertarget{c01:cue:2947}{}主張怎麼覺得。
% cue 2948 5477.320 5478.920
\hypertarget{c01:cue:2948}{}或者我怎麼。
% cue 2949 5478.920 5479.980
\hypertarget{c01:cue:2949}{}學生。
% cue 2950 5479.980 5482.610
\hypertarget{c01:cue:2950}{}會不會讓你解憋。
% cue 2951 5482.610 5482.710
\hypertarget{c01:cue:2951}{}第一個。
% cue 2952 5482.610 5483.610
\hypertarget{c01:cue:2952}{}我們外國人。
% cue 2953 5483.610 5485.540
\hypertarget{c01:cue:2953}{}大概他不會讓你解憋。
% cue 2954 5485.540 5488.210
\hypertarget{c01:cue:2954}{}因為他不要去跟他搶聲音。
% cue 2955 5488.210 5489.880
\hypertarget{c01:cue:2955}{}所以外國人。
% cue 2956 5489.880 5490.880
\hypertarget{c01:cue:2956}{}大概沒有辦法面解憋。

\tc{1:31:30}
% cue 2957 5490.880 5494.610
\hypertarget{c01:cue:2957}{}那我那一年的 LL.M.。
% cue 2958 5494.610 5495.630
\hypertarget{c01:cue:2958}{}Harvard 來講。
% cue 2959 5495.630 5497.830
\hypertarget{c01:cue:2959}{}我紀錄是 80 的。
% cue 2960 5497.830 5501.070
\hypertarget{c01:cue:2960}{}但是通常都會超級。
% cue 2961 5501.070 5504.670
\hypertarget{c01:cue:2961}{}那臺灣一年大概就五課。
% cue 2962 5504.670 5505.950
\hypertarget{c01:cue:2962}{}我簡單的講就是這樣。
% cue 2963 5505.950 5509.260
\hypertarget{c01:cue:2963}{}那你就可以看到你的反應。
% cue 2964 5509.260 5511.260
\hypertarget{c01:cue:2964}{}你在臺灣也如果是 TOW。
% cue 2965 5511.260 5512.730
\hypertarget{c01:cue:2965}{}就是希望試試看。
% cue 2966 5512.730 5514.790
\hypertarget{c01:cue:2966}{}我那一年是有 Harvard。
% cue 2967 5514.790 5515.790
\hypertarget{c01:cue:2967}{}跟 EOA。
% cue 2968 5515.790 5516.790
\hypertarget{c01:cue:2968}{}不斷那種的學識。
% cue 2969 5516.790 5517.790
\hypertarget{c01:cue:2969}{}所以我選擇。
% cue 2970 5517.790 5524.580
\hypertarget{c01:cue:2970}{}他都不願。

\tc{1:32:04}
% cue 2971 5524.580 5527.430
\hypertarget{c01:cue:2971}{}所以沒有特別為我。
% cue 2972 5527.430 5528.430
\hypertarget{c01:cue:2972}{}不拜國人。
% cue 2973 5528.430 5529.430
\hypertarget{c01:cue:2973}{}或者是為 LL.M.。
% cue 2974 5529.430 5530.400
\hypertarget{c01:cue:2974}{}上一課。
% cue 2975 5530.400 5532.480
\hypertarget{c01:cue:2975}{}他開的是一般JD的課。
% cue 2976 5532.480 5533.670
\hypertarget{c01:cue:2976}{}那你自己去選。
% cue 2977 5533.670 5534.930
\hypertarget{c01:cue:2977}{}那所以我會選的。
% cue 2978 5534.930 5535.930
\hypertarget{c01:cue:2978}{}都是憲法。
% cue 2979 5535.930 5536.930
\hypertarget{c01:cue:2979}{}行政法。
% cue 2980 5536.930 5537.930
\hypertarget{c01:cue:2980}{}憲法專題。
% cue 2981 5537.930 5538.930
\hypertarget{c01:cue:2981}{}行政法專題。
% cue 2982 5538.930 5540.630
\hypertarget{c01:cue:2982}{}像我以前我選過。
% cue 2983 5540.630 5541.630
\hypertarget{c01:cue:2983}{}一些憲法專題。
% cue 2984 5541.630 5543.680
\hypertarget{c01:cue:2984}{}有一些行政法的教育。
% cue 2985 5543.680 5544.830
\hypertarget{c01:cue:2985}{}像我比方說。

\tc{1:32:24}
% cue 2986 5544.830 5546.130
\hypertarget{c01:cue:2986}{}關於福利國家的。
% cue 2987 5546.130 5548.510
\hypertarget{c01:cue:2987}{}但是你要看。
% cue 2988 5548.510 5550.000
\hypertarget{c01:cue:2988}{}第一個。
% cue 2989 5550.000 5551.250
\hypertarget{c01:cue:2989}{}老師要夠多。
% cue 2990 5551.250 5552.250
\hypertarget{c01:cue:2990}{}他的課長會多。
% cue 2991 5552.250 5553.230
\hypertarget{c01:cue:2991}{}第二個。
% cue 2992 5553.230 5555.440
\hypertarget{c01:cue:2992}{}那正好你家的老師。
% cue 2993 5555.440 5556.440
\hypertarget{c01:cue:2993}{}他們也沒有休假。
% cue 2994 5556.440 5558.100
\hypertarget{c01:cue:2994}{}他有開合。
% cue 2995 5558.100 5558.200
\hypertarget{c01:cue:2995}{}他有什麼戰會。
% cue 2996 5558.100 5559.100
\hypertarget{c01:cue:2996}{}開合。
% cue 2997 5559.100 5562.700
\hypertarget{c01:cue:2997}{}所以這是非常欲起的。
% cue 2998 5562.700 5563.700
\hypertarget{c01:cue:2998}{}所以真要問到好多老師。
% cue 2999 5563.700 5565.980
\hypertarget{c01:cue:2999}{}人就像星星。
% cue 3000 5565.980 5568.360
\hypertarget{c01:cue:3000}{}現在就這個星星在這裡。
% cue 3001 5568.360 5570.330
\hypertarget{c01:cue:3001}{}那明年下約他移動到這裡。
% cue 3002 5570.330 5572.540
\hypertarget{c01:cue:3002}{}那你去的正法就成掉了。
% cue 3003 5572.540 5575.980
\hypertarget{c01:cue:3003}{}所以這個是你要自己通到你才會。
% cue 3004 5575.980 5576.970
\hypertarget{c01:cue:3004}{}你才會有感慨。
% cue 3005 5576.970 5579.000
\hypertarget{c01:cue:3005}{}反正我現在，
% cue 3006 5579.000 5580.000
\hypertarget{c01:cue:3006}{}對我講就是講的大概。
% cue 3007 5580.000 5582.640
\hypertarget{c01:cue:3007}{}慢慢清楚。

\section{大法官提名的經過}

\tc{1:33:02}
% cue 3008 5582.640 5583.820
\hypertarget{c01:cue:3008}{}還有其他問題提出。

\tc{1:33:03}
\speaker{學生}
% cue 3009 5583.820 5585.810
\hypertarget{c01:cue:3009}{}想當大法官。
% cue 3010 5585.810 5586.810
\hypertarget{c01:cue:3010}{}對這種被提。
% cue 3011 5586.810 5587.810
\hypertarget{c01:cue:3011}{}你們的過去。

\tc{1:33:07}
\speaker{老師}
% cue 3012 5587.810 5592.110
\hypertarget{c01:cue:3012}{}你現在講的這個事情。
% cue 3013 5592.110 5593.400
\hypertarget{c01:cue:3013}{}還沒有公開。
% cue 3014 5593.400 5594.400
\hypertarget{c01:cue:3014}{}因為老師的書。
% cue 3015 5594.400 5596.400
\hypertarget{c01:cue:3015}{}只是講我在大法官。
% cue 3016 5596.400 5597.400
\hypertarget{c01:cue:3016}{}任內做的事情。
% cue 3017 5597.400 5599.810
\hypertarget{c01:cue:3017}{}這算是我對於全國人民。
% cue 3018 5599.810 5602.060
\hypertarget{c01:cue:3018}{}我領了大家八年的薪水。
% cue 3019 5602.060 5604.340
\hypertarget{c01:cue:3019}{}做了一個工作報告。
% cue 3020 5604.340 5607.890
\hypertarget{c01:cue:3020}{}看的地方。
% cue 3021 5607.890 5608.860
\hypertarget{c01:cue:3021}{}我也是沒看。
% cue 3022 5608.860 5611.030
\hypertarget{c01:cue:3022}{}很重要。

\tc{1:33:31}
% cue 3023 5611.030 5612.030
\hypertarget{c01:cue:3023}{}那你現在講的這個事情。
% cue 3024 5612.030 5613.850
\hypertarget{c01:cue:3024}{}第一。
% cue 3025 5613.850 5616.240
\hypertarget{c01:cue:3025}{}有每個人的機緣。
% cue 3026 5616.240 5617.460
\hypertarget{c01:cue:3026}{}第二。
% cue 3027 5617.460 5619.450
\hypertarget{c01:cue:3027}{}這裡面一定有一些政治的理由。
% cue 3028 5619.450 5620.990
\hypertarget{c01:cue:3028}{}因為這麼多人。
% cue 3029 5620.990 5621.990
\hypertarget{c01:cue:3029}{}要告訴我。
% cue 3030 5621.990 5623.690
\hypertarget{c01:cue:3030}{}因為朋友會大會。
% cue 3031 5623.690 5624.690
\hypertarget{c01:cue:3031}{}那對我來講。
% cue 3032 5624.690 5626.810
\hypertarget{c01:cue:3032}{}我只能簡單的唯一。
% cue 3033 5626.810 5629.570
\hypertarget{c01:cue:3033}{}就是說我將來。
% cue 3034 5629.570 5630.720
\hypertarget{c01:cue:3034}{}我知道我將來。
% cue 3035 5630.720 5633.900
\hypertarget{c01:cue:3035}{}現在有點猶豫。
% cue 3036 5633.900 5636.470
\hypertarget{c01:cue:3036}{}如果我還活著久一點。
% cue 3037 5636.470 5638.690
\hypertarget{c01:cue:3037}{}而且我認為。
% cue 3038 5638.690 5640.680
\hypertarget{c01:cue:3038}{}算重要人物的話，
% cue 3039 5640.680 5642.680
\hypertarget{c01:cue:3039}{}如果我認為我不是重要人物。
% cue 3040 5642.680 5643.640
\hypertarget{c01:cue:3040}{}我就不會。
% cue 3041 5643.640 5645.860
\hypertarget{c01:cue:3041}{}如果是活著久一點。
% cue 3042 5645.860 5647.810
\hypertarget{c01:cue:3042}{}而且自認為是重要人物的話，
% cue 3043 5647.810 5650.040
\hypertarget{c01:cue:3043}{}我也許可能會出一個回憶路。
% cue 3044 5650.040 5651.550
\hypertarget{c01:cue:3044}{}那如果我出回憶路。
% cue 3045 5651.550 5653.700
\hypertarget{c01:cue:3045}{}你現在講的這個。
% cue 3046 5653.700 5654.700
\hypertarget{c01:cue:3046}{}我就會提。

\tc{1:34:14}
% cue 3047 5654.700 5656.620
\hypertarget{c01:cue:3047}{}那簡單的回憶路。
% cue 3048 5656.620 5657.620
\hypertarget{c01:cue:3048}{}是這樣子的。
% cue 3049 5657.620 5660.950
\hypertarget{c01:cue:3049}{}我的經營大部分。
% cue 3050 5660.950 5663.230
\hypertarget{c01:cue:3050}{}是很曲折的。
% cue 3051 5663.230 5665.730
\hypertarget{c01:cue:3051}{}而且呢，
% cue 3052 5665.730 5667.890
\hypertarget{c01:cue:3052}{}是有一點從天上掉下來的。
% cue 3053 5667.890 5670.770
\hypertarget{c01:cue:3053}{}首先呢，
% cue 3054 5670.770 5674.560
\hypertarget{c01:cue:3054}{}這個是我第三次被提名\footnote{2011 年馬英九總統提名邵燕玲、陳碧玉、黃璽君、羅昌發為司法院大法官，邵燕玲婉辭後，總統府改提名湯德宗。湯德宗、陳碧玉、黃璽君、羅昌發四人經立法院同意，任期自 2011 年 10 月 1 日至 2019 年 9 月 30 日。\url{https://newtalk.tw/news/view/2011-04-01/13238}。}。
% cue 3055 5674.560 5676.290
\hypertarget{c01:cue:3055}{}第一次被提名的時候。
% cue 3056 5676.290 5677.290
\hypertarget{c01:cue:3056}{}是司法院推薦。
% cue 3057 5677.290 5678.290
\hypertarget{c01:cue:3057}{}全個司法院。
% cue 3058 5678.290 5679.290
\hypertarget{c01:cue:3058}{}只有推薦。
% cue 3059 5679.290 5682.080
\hypertarget{c01:cue:3059}{}這個名單放到阿扁\footnote{陳水扁的暱稱。陳水扁於 2000 年 5 月 20 日至 2008 年 5 月 20 日任中華民國總統。}的桌上。
% cue 3060 5682.080 5684.030
\hypertarget{c01:cue:3060}{}阿扁說他們的桌是誰？
% cue 3061 5684.030 5685.000
\hypertarget{c01:cue:3061}{}認識。
% cue 3062 5685.000 5685.100
\hypertarget{c01:cue:3062}{}劃掉。
% cue 3063 5685.000 5686.770
\hypertarget{c01:cue:3063}{}沒了。
% cue 3064 5686.770 5690.780
\hypertarget{c01:cue:3064}{}第二次的時候。
% cue 3065 5690.780 5693.100
\hypertarget{c01:cue:3065}{}我在中央研究院籌備。

\tc{1:34:53}
% cue 3066 5693.100 5694.440
\hypertarget{c01:cue:3066}{}是籌備處\footnote{中央研究院法律學研究所籌備處於 2004 年 7 月 1 日成立，由湯德宗任籌備處主任；2011 年 7 月 1 日正式成所，湯德宗為首任所長，同年 10 月轉任司法院大法官。籌備處成立之前，中研院的法學研究以中山人文社會科學研究所法律組為據點。\url{https://publication.iias.sinica.edu.tw/2021ch00.pdf}。}。
% cue 3067 5694.440 5695.540
\hypertarget{c01:cue:3067}{}是籌備處的處。
% cue 3068 5695.540 5696.540
\hypertarget{c01:cue:3068}{}人後來才當所長。
% cue 3069 5696.540 5697.830
\hypertarget{c01:cue:3069}{}我們中央院比較特別。
% cue 3070 5697.830 5703.650
\hypertarget{c01:cue:3070}{}籌備處要經過不定的相當期間。
% cue 3071 5703.650 5705.830
\hypertarget{c01:cue:3071}{}不確定的相當期間都會提。
% cue 3072 5705.830 5706.980
\hypertarget{c01:cue:3072}{}我籌備了七年。
% cue 3073 5706.980 5709.200
\hypertarget{c01:cue:3073}{}算是很快很快很快的。
% cue 3074 5709.200 5710.260
\hypertarget{c01:cue:3074}{}在最近這二十年來。
% cue 3075 5710.260 5712.250
\hypertarget{c01:cue:3075}{}沒有人被我關注七年。
% cue 3076 5712.250 5715.440
\hypertarget{c01:cue:3076}{}那那時候呢，
% cue 3077 5715.440 5717.200
\hypertarget{c01:cue:3077}{}有一個機會。
% cue 3078 5717.200 5718.610
\hypertarget{c01:cue:3078}{}第二次。
% cue 3079 5718.610 5719.610
\hypertarget{c01:cue:3079}{}大家不要再多。
% cue 3080 5719.610 5720.580
\hypertarget{c01:cue:3080}{}是我第二次被提名。
% cue 3081 5720.580 5721.960
\hypertarget{c01:cue:3081}{}但是那一次呢，
% cue 3082 5721.960 5726.570
\hypertarget{c01:cue:3082}{}話就不好說了。
% cue 3083 5726.570 5727.570
\hypertarget{c01:cue:3083}{}好。
% cue 3084 5727.570 5728.860
\hypertarget{c01:cue:3084}{}算被做掉。
% cue 3085 5728.860 5731.430
\hypertarget{c01:cue:3085}{}做掉的理由啊。
% cue 3086 5731.430 5733.190
\hypertarget{c01:cue:3086}{}講白一點就是，
% cue 3087 5733.190 5734.570
\hypertarget{c01:cue:3087}{}這個人比較不會交際。
% cue 3088 5734.570 5736.570
\hypertarget{c01:cue:3088}{}你沒有期望會過很久。
% cue 3089 5736.570 5738.300
\hypertarget{c01:cue:3089}{}不知道什麼。
% cue 3090 5738.300 5739.260
\hypertarget{c01:cue:3090}{}這是什麼邏輯。
% cue 3091 5739.260 5740.430
\hypertarget{c01:cue:3091}{}那就是關係不夠難。
% cue 3092 5740.430 5741.780
\hypertarget{c01:cue:3092}{}掉掉。
% cue 3093 5741.780 5742.690
\hypertarget{c01:cue:3093}{}這是第二次。

\tc{1:35:42}
% cue 3094 5742.690 5744.190
\hypertarget{c01:cue:3094}{}所以這一次是第三次。
% cue 3095 5744.190 5745.700
\hypertarget{c01:cue:3095}{}而且呢，
% cue 3096 5745.700 5747.700
\hypertarget{c01:cue:3096}{}因為我前面有兩次的經驗。
% cue 3097 5747.700 5750.210
\hypertarget{c01:cue:3097}{}所以我這次是被大家公認。
% cue 3098 5750.210 5751.210
\hypertarget{c01:cue:3098}{}還在最底面。
% cue 3099 5751.210 5753.170
\hypertarget{c01:cue:3099}{}但是有沒有用呢，
% cue 3100 5753.170 5754.170
\hypertarget{c01:cue:3100}{}沒有用。
% cue 3101 5754.170 5755.490
\hypertarget{c01:cue:3101}{}進去了以後。
% cue 3102 5755.490 5756.710
\hypertarget{c01:cue:3102}{}那個時候是火力島。
% cue 3103 5756.710 5758.150
\hypertarget{c01:cue:3103}{}射的死案。
% cue 3104 5758.150 5760.150
\hypertarget{c01:cue:3104}{}都知道火力島誰很生。
% cue 3105 5760.150 5762.380
\hypertarget{c01:cue:3105}{}火力島。
% cue 3106 5762.380 5764.040
\hypertarget{c01:cue:3106}{}因為我們講的。
% cue 3107 5764.040 5765.040
\hypertarget{c01:cue:3107}{}降噹文化也好。
% cue 3108 5765.040 5766.040
\hypertarget{c01:cue:3108}{}觀點政治也好。
% cue 3109 5766.040 5767.010
\hypertarget{c01:cue:3109}{}反正你進去裡面。
% cue 3110 5767.010 5768.010
\hypertarget{c01:cue:3110}{}你不知道他。
% cue 3111 5768.010 5769.010
\hypertarget{c01:cue:3111}{}你不是在講故事啊。
% cue 3112 5769.010 5771.010
\hypertarget{c01:cue:3112}{}我是最擺在最前面的時候進去。
% cue 3113 5771.010 5774.450
\hypertarget{c01:cue:3113}{}但是總統府過來發布了。
% cue 3114 5774.450 5775.450
\hypertarget{c01:cue:3114}{}新聞記者會上面。
% cue 3115 5775.450 5776.400
\hypertarget{c01:cue:3115}{}沒有人有什麼。
% cue 3116 5776.400 5777.470
\hypertarget{c01:cue:3116}{}他也不是記者會。
% cue 3117 5777.470 5781.130
\hypertarget{c01:cue:3117}{}他是一個透過酒肉風聲。
% cue 3118 5781.130 5782.130
\hypertarget{c01:cue:3118}{}給報紙先報寒。
% cue 3119 5782.130 5784.160
\hypertarget{c01:cue:3119}{}所以晚上的晚報。
% cue 3120 5784.160 5785.160
\hypertarget{c01:cue:3120}{}還有第二天的早報。
% cue 3121 5785.160 5786.160
\hypertarget{c01:cue:3121}{}就已經報寒。
% cue 3122 5786.160 5787.960
\hypertarget{c01:cue:3122}{}沒有人有什麼。

\tc{1:36:27}
% cue 3123 5787.960 5789.720
\hypertarget{c01:cue:3123}{}好 那我們講一下。
% cue 3124 5789.720 5790.720
\hypertarget{c01:cue:3124}{}趁死發生的故事。
% cue 3125 5790.720 5792.720
\hypertarget{c01:cue:3125}{}那天我要去中研院路上。
% cue 3126 5792.720 5794.000
\hypertarget{c01:cue:3126}{}在這之前。
% cue 3127 5794.000 5795.060
\hypertarget{c01:cue:3127}{}當然知道。
% cue 3128 5795.060 5796.060
\hypertarget{c01:cue:3128}{}因為很多人會告訴我。
% cue 3129 5796.060 5797.310
\hypertarget{c01:cue:3129}{}說你怎樣怎樣。
% cue 3130 5797.310 5798.440
\hypertarget{c01:cue:3130}{}這次應該。
% cue 3131 5798.440 5799.790
\hypertarget{c01:cue:3131}{}不會委屈你那次了。
% cue 3132 5799.790 5801.810
\hypertarget{c01:cue:3132}{}我想說好啊。
% cue 3133 5801.810 5802.810
\hypertarget{c01:cue:3133}{}應該也是。
% cue 3134 5802.810 5803.810
\hypertarget{c01:cue:3134}{}對不對？
% cue 3135 5803.810 5804.870
\hypertarget{c01:cue:3135}{}那在最前面。
% cue 3136 5804.870 5805.790
\hypertarget{c01:cue:3136}{}開四個。
% cue 3137 5805.790 5806.830
\hypertarget{c01:cue:3137}{}那四個。
% cue 3138 5806.830 5806.930
\hypertarget{c01:cue:3138}{}那四個。
% cue 3139 5806.830 5807.830
\hypertarget{c01:cue:3139}{}我最前面。
% cue 3140 5807.830 5808.830
\hypertarget{c01:cue:3140}{}怎麼也能到我。
% cue 3141 5808.830 5810.080
\hypertarget{c01:cue:3141}{}我這麼想。
% cue 3142 5810.080 5811.330
\hypertarget{c01:cue:3142}{}我也沒有多想。
% cue 3143 5811.330 5812.490
\hypertarget{c01:cue:3143}{}結果呢，
% cue 3144 5812.490 5813.490
\hypertarget{c01:cue:3144}{}報紙出來。
% cue 3145 5813.490 5814.680
\hypertarget{c01:cue:3145}{}不是這麼沒事。
% cue 3146 5814.680 5815.680
\hypertarget{c01:cue:3146}{}在路上開車。
% cue 3147 5815.680 5816.640
\hypertarget{c01:cue:3147}{}一周年院。
% cue 3148 5816.640 5817.640
\hypertarget{c01:cue:3148}{}大概我見到說。
% cue 3149 5817.640 5819.110
\hypertarget{c01:cue:3149}{}沒有半小時。
% cue 3150 5819.110 5820.110
\hypertarget{c01:cue:3150}{}接了兩個電話。
% cue 3151 5820.110 5822.370
\hypertarget{c01:cue:3151}{}我就跟你們描述一下。

\tc{1:37:02}
% cue 3152 5822.370 5823.620
\hypertarget{c01:cue:3152}{}第一個電話說。
% cue 3153 5823.620 5824.610
\hypertarget{c01:cue:3153}{}德宗啊。
% cue 3154 5824.610 5825.930
\hypertarget{c01:cue:3154}{}請你凡事。
% cue 3155 5825.930 5826.030
\hypertarget{c01:cue:3155}{}相信主。
% cue 3156 5825.930 5826.930
\hypertarget{c01:cue:3156}{}凡事。
% cue 3157 5826.930 5828.280
\hypertarget{c01:cue:3157}{}在別人。
% cue 3158 5828.280 5828.380
\hypertarget{c01:cue:3158}{}我要失處。
% cue 3159 5828.280 5830.050
\hypertarget{c01:cue:3159}{}就是第一。
% cue 3160 5830.050 5831.010
\hypertarget{c01:cue:3160}{}電話。
% cue 3161 5831.010 5831.110
\hypertarget{c01:cue:3161}{}第二。
% cue 3162 5831.010 5832.680
\hypertarget{c01:cue:3162}{}第二。
% cue 3163 5832.680 5834.920
\hypertarget{c01:cue:3163}{}因為他們知道我在開車。
% cue 3164 5834.920 5835.980
\hypertarget{c01:cue:3164}{}沒有多講。
% cue 3165 5835.980 5836.980
\hypertarget{c01:cue:3165}{}而且這個話也不能多講。
% cue 3166 5836.980 5838.360
\hypertarget{c01:cue:3166}{}第二。
% cue 3167 5838.360 5838.460
\hypertarget{c01:cue:3167}{}第二。
% cue 3168 5838.360 5838.460
\hypertarget{c01:cue:3168}{}第二。
% cue 3169 5838.360 5839.770
\hypertarget{c01:cue:3169}{}第二。
% cue 3170 5839.770 5841.860
\hypertarget{c01:cue:3170}{}叫故事。
% cue 3171 5841.860 5842.860
\hypertarget{c01:cue:3171}{}德宗啊可以講。
% cue 3172 5842.860 5844.560
\hypertarget{c01:cue:3172}{}很生氣。
% cue 3173 5844.560 5844.660
\hypertarget{c01:cue:3173}{}我已經跟副長。
% cue 3174 5844.560 5845.560
\hypertarget{c01:cue:3174}{}跟打電話。
% cue 3175 5845.560 5848.060
\hypertarget{c01:cue:3175}{}靠幕。
% cue 3176 5848.060 5853.540
\hypertarget{c01:cue:3176}{}說讓我們多聽。
% cue 3177 5853.540 5854.820
\hypertarget{c01:cue:3177}{}這就是兩個電話。

\tc{1:37:34}
% cue 3178 5854.820 5855.960
\hypertarget{c01:cue:3178}{}那我現在，
% cue 3179 5855.960 5857.330
\hypertarget{c01:cue:3179}{}反正把報已經出現了。
% cue 3180 5857.330 5858.330
\hypertarget{c01:cue:3180}{}沒有我的名字。
% cue 3181 5858.330 5860.460
\hypertarget{c01:cue:3181}{}他早報我也沒有看。
% cue 3182 5860.460 5861.560
\hypertarget{c01:cue:3182}{}但是早報。
% cue 3183 5861.560 5862.560
\hypertarget{c01:cue:3183}{}去路上。
% cue 3184 5862.560 5863.560
\hypertarget{c01:cue:3184}{}就看到大家講這些消息。
% cue 3185 5863.560 5865.520
\hypertarget{c01:cue:3185}{}那就一定沒有我嘛。
% cue 3186 5865.520 5867.540
\hypertarget{c01:cue:3186}{}對不對？
% cue 3187 5867.540 5868.920
\hypertarget{c01:cue:3187}{}上來了。
% cue 3188 5868.920 5874.350
\hypertarget{c01:cue:3188}{}大概到 8 點 55 分左右。
% cue 3189 5874.350 5876.350
\hypertarget{c01:cue:3189}{}反正 9 點前後。
% cue 3190 5876.350 5877.830
\hypertarget{c01:cue:3190}{}東府來電。
% cue 3191 5877.830 5881.860
\hypertarget{c01:cue:3191}{}認識。
% cue 3192 5881.860 5882.950
\hypertarget{c01:cue:3192}{}跟我講呢，
% cue 3193 5882.950 5884.430
\hypertarget{c01:cue:3193}{}總統你要見。
% cue 3194 5884.430 5884.530
\hypertarget{c01:cue:3194}{}請你。
% cue 3195 5884.430 5885.490
\hypertarget{c01:cue:3195}{}現在就來。
% cue 3196 5885.490 5887.950
\hypertarget{c01:cue:3196}{}那我幾點幾分。
% cue 3197 5887.950 5890.080
\hypertarget{c01:cue:3197}{}我在第幾號人。
% cue 3198 5890.080 5892.140
\hypertarget{c01:cue:3198}{}從東府進出。
% cue 3199 5892.140 5893.210
\hypertarget{c01:cue:3199}{}就是有一點進出。
% cue 3200 5893.210 5893.310
\hypertarget{c01:cue:3200}{}結果呢，
% cue 3201 5893.210 5897.460
\hypertarget{c01:cue:3201}{}去了。
% cue 3202 5897.460 5899.250
\hypertarget{c01:cue:3202}{}結果以後呢，
% cue 3203 5899.250 5901.860
\hypertarget{c01:cue:3203}{}帶一路走進去。
% cue 3204 5901.860 5902.860
\hypertarget{c01:cue:3204}{}反正總統府。
% cue 3205 5902.860 5904.020
\hypertarget{c01:cue:3205}{}我也很少進去。
% cue 3206 5904.020 5906.900
\hypertarget{c01:cue:3206}{}坐在一個房間。
% cue 3207 5906.900 5907.900
\hypertarget{c01:cue:3207}{}他說你等一下。
% cue 3208 5907.900 5909.310
\hypertarget{c01:cue:3208}{}再過了一路。
% cue 3209 5909.310 5910.370
\hypertarget{c01:cue:3209}{}從那個門。
% cue 3210 5910.370 5911.370
\hypertarget{c01:cue:3210}{}從這個門進來。
% cue 3211 5911.370 5912.910
\hypertarget{c01:cue:3211}{}那個門推開。
% cue 3212 5912.910 5915.340
\hypertarget{c01:cue:3212}{}講話。
% cue 3213 5915.340 5917.200
\hypertarget{c01:cue:3213}{}非常妙的詩詩。
% cue 3214 5917.200 5918.200
\hypertarget{c01:cue:3214}{}從頭到尾。
% cue 3215 5918.200 5920.320
\hypertarget{c01:cue:3215}{}我沒有一句話說。
% cue 3216 5920.320 5922.210
\hypertarget{c01:cue:3216}{}對不起你啊。
% cue 3217 5922.210 5923.210
\hypertarget{c01:cue:3217}{}我怎麼啊。
% cue 3218 5923.210 5924.210
\hypertarget{c01:cue:3218}{}沒有沒有沒有。
% cue 3219 5924.210 5924.310
\hypertarget{c01:cue:3219}{}頭都沒有。
% cue 3220 5924.210 5925.340
\hypertarget{c01:cue:3220}{}哎呀。
% cue 3221 5925.340 5926.340
\hypertarget{c01:cue:3221}{}德宗啊很棒啊。
% cue 3222 5926.340 5927.720
\hypertarget{c01:cue:3222}{}怎樣怎樣怎樣啊。
% cue 3223 5927.720 5930.060
\hypertarget{c01:cue:3223}{}我會麻煩你啊。
% cue 3224 5930.060 5931.060
\hypertarget{c01:cue:3224}{}也沒有提到大火。
% cue 3225 5931.060 5932.820
\hypertarget{c01:cue:3225}{}什麼都沒有。
% cue 3226 5932.820 5933.820
\hypertarget{c01:cue:3226}{}他就走了。
% cue 3227 5933.820 5933.920
\hypertarget{c01:cue:3227}{}講完了。
% cue 3228 5933.820 5934.950
\hypertarget{c01:cue:3228}{}就這樣就講完了。

\tc{1:38:54}
% cue 3229 5934.950 5936.360
\hypertarget{c01:cue:3229}{}好。
% cue 3230 5936.360 5937.520
\hypertarget{c01:cue:3230}{}再過來。
% cue 3231 5937.520 5938.520
\hypertarget{c01:cue:3231}{}剛才帶我進去。
% cue 3232 5938.520 5939.520
\hypertarget{c01:cue:3232}{}那個時候。
% cue 3233 5939.520 5939.620
\hypertarget{c01:cue:3233}{}人物呢，
% cue 3234 5939.520 5940.490
\hypertarget{c01:cue:3234}{}又來帶我啦。
% cue 3235 5940.490 5941.670
\hypertarget{c01:cue:3235}{}說來。
% cue 3236 5941.670 5942.810
\hypertarget{c01:cue:3236}{}待會兒。
% cue 3237 5942.810 5943.860
\hypertarget{c01:cue:3237}{}從這個門進去。
% cue 3238 5943.860 5945.220
\hypertarget{c01:cue:3238}{}那你要進去。
% cue 3239 5945.220 5946.280
\hypertarget{c01:cue:3239}{}路上任何人問你。
% cue 3240 5946.280 5947.280
\hypertarget{c01:cue:3240}{}什麼事情。
% cue 3241 5947.280 5949.600
\hypertarget{c01:cue:3241}{}你都含笑不答。
% cue 3242 5949.600 5951.010
\hypertarget{c01:cue:3242}{}所有的問題。
% cue 3243 5951.010 5952.540
\hypertarget{c01:cue:3243}{}你會來處理。
% cue 3244 5952.540 5953.990
\hypertarget{c01:cue:3244}{}然後呢，
% cue 3245 5953.990 5954.990
\hypertarget{c01:cue:3245}{}你待會兒進去的時候。
% cue 3246 5954.990 5956.120
\hypertarget{c01:cue:3246}{}那個時候。
% cue 3247 5956.120 5957.120
\hypertarget{c01:cue:3247}{}其他幾個大房間。
% cue 3248 5957.120 5958.370
\hypertarget{c01:cue:3248}{}被提名。
% cue 3249 5958.370 5959.690
\hypertarget{c01:cue:3249}{}蓋在裡面了。
% cue 3250 5959.690 5960.690
\hypertarget{c01:cue:3250}{}他們是在裡面。
% cue 3251 5960.690 5961.690
\hypertarget{c01:cue:3251}{}已經坐在裡面。
% cue 3252 5961.690 5961.790
\hypertarget{c01:cue:3252}{}我知道。
% cue 3253 5961.690 5962.690
\hypertarget{c01:cue:3253}{}我是從這個門進來。
% cue 3254 5962.690 5964.100
\hypertarget{c01:cue:3254}{}都跟那個提名。
% cue 3255 5964.100 5965.580
\hypertarget{c01:cue:3255}{}然後接下來。
% cue 3256 5965.580 5966.580
\hypertarget{c01:cue:3256}{}副總統就來了。
% cue 3257 5966.580 5967.550
\hypertarget{c01:cue:3257}{}鏡子路就開了。
% cue 3258 5967.550 5970.020
\hypertarget{c01:cue:3258}{}就這樣上去了。

\tc{1:39:30}
% cue 3259 5970.020 5971.810
\hypertarget{c01:cue:3259}{}那後來很多人。
% cue 3260 5971.810 5972.810
\hypertarget{c01:cue:3260}{}都跟我講說。
% cue 3261 5972.810 5974.870
\hypertarget{c01:cue:3261}{}該是你。
% cue 3262 5974.870 5975.870
\hypertarget{c01:cue:3262}{}怎樣跑都跑掉。
% cue 3263 5975.870 5978.500
\hypertarget{c01:cue:3263}{}那你沒有興趣呢，
% cue 3264 5978.500 5979.650
\hypertarget{c01:cue:3264}{}我們查一個事情。
% cue 3265 5979.650 5981.130
\hypertarget{c01:cue:3265}{}這個事情。
% cue 3266 5981.130 5982.130
\hypertarget{c01:cue:3266}{}為什麼會突然突變。
% cue 3267 5982.130 5983.250
\hypertarget{c01:cue:3267}{}是不是有一個人。
% cue 3268 5983.250 5984.340
\hypertarget{c01:cue:3268}{}那是原來。
% cue 3269 5984.340 5987.630
\hypertarget{c01:cue:3269}{}女的最高法院的法官。
% cue 3270 5987.630 5989.880
\hypertarget{c01:cue:3270}{}他因為涉及到那個案子。
% cue 3271 5989.880 5991.070
\hypertarget{c01:cue:3271}{}性侵。
% cue 3272 5991.070 5992.070
\hypertarget{c01:cue:3272}{}兒童的性侵案。
% cue 3273 5992.070 5992.990
\hypertarget{c01:cue:3273}{}所以大家說。
% cue 3274 5992.990 5994.190
\hypertarget{c01:cue:3274}{}他是。
% cue 3275 5994.190 5996.600
\hypertarget{c01:cue:3275}{}他是一個恐龍法官\footnote{邵燕玲時任最高法院法官，2011 年獲提名為司法院大法官。她擔任審判長的合議庭曾將一件幼童性侵案發回更審，認女童證詞不能證明被告違反其意願，因此被批評為「恐龍法官」；提名公布後社會譁然，邵燕玲婉辭提名。\url{https://newtalk.tw/news/view/2011-04-01/13238}。}。
% cue 3276 5996.600 5998.210
\hypertarget{c01:cue:3276}{}恐龍法官就這樣來的。
% cue 3277 5998.210 5999.530
\hypertarget{c01:cue:3277}{}那個性侵案。
% cue 3278 5999.530 6000.630
\hypertarget{c01:cue:3278}{}然後這個。
% cue 3279 6000.630 6001.630
\hypertarget{c01:cue:3279}{}這個。
% cue 3280 6001.630 6002.630
\hypertarget{c01:cue:3280}{}這個。
% cue 3281 6002.630 6003.600
\hypertarget{c01:cue:3281}{}這個。
% cue 3282 6003.600 6003.700
\hypertarget{c01:cue:3282}{}這個。
% cue 3283 6003.600 6004.600
\hypertarget{c01:cue:3283}{}這個。
% cue 3284 6004.600 6005.600
\hypertarget{c01:cue:3284}{}這個。
% cue 3285 6005.600 6007.550
\hypertarget{c01:cue:3285}{}這個法院的法官。
% cue 3286 6007.550 6009.370
\hypertarget{c01:cue:3286}{}被提名。
% cue 3287 6009.370 6015.300
\hypertarget{c01:cue:3287}{}原來是提名。
% cue 3288 6015.300 6016.300
\hypertarget{c01:cue:3288}{}說跟兒童講。
% cue 3289 6016.300 6017.300
\hypertarget{c01:cue:3289}{}可以去查白明暉。
% cue 3290 6017.300 6019.320
\hypertarget{c01:cue:3290}{}白明暉的案子。
% cue 3291 6019.320 6019.420
\hypertarget{c01:cue:3291}{}性侵。
% cue 3292 6019.320 6021.320
\hypertarget{c01:cue:3292}{}然後就是恐龍法官。
% cue 3293 6021.320 6023.850
\hypertarget{c01:cue:3293}{}就是報紙上都有。

\tc{1:40:23}
% cue 3294 6023.850 6024.850
\hypertarget{c01:cue:3294}{}那我現在回來講。
% cue 3295 6024.850 6025.850
\hypertarget{c01:cue:3295}{}剛才這件事。
% cue 3296 6025.850 6026.850
\hypertarget{c01:cue:3296}{}我剛這是秘密。
% cue 3297 6026.850 6028.330
\hypertarget{c01:cue:3297}{}我都不知道。
% cue 3298 6028.330 6029.330
\hypertarget{c01:cue:3298}{}進去的時候。
% cue 3299 6029.330 6030.940
\hypertarget{c01:cue:3299}{}我白天的經驗。
% cue 3300 6030.940 6031.940
\hypertarget{c01:cue:3300}{}為什麼出來。
% cue 3301 6031.940 6032.940
\hypertarget{c01:cue:3301}{}從後來。
% cue 3302 6032.940 6034.070
\hypertarget{c01:cue:3302}{}洩漏的名單。
% cue 3303 6034.070 6035.160
\hypertarget{c01:cue:3303}{}不是正式公布。
% cue 3304 6035.160 6037.160
\hypertarget{c01:cue:3304}{}正式公布是你那天開機子的那個時候。
% cue 3305 6037.160 6038.080
\hypertarget{c01:cue:3305}{}我又出現人。
% cue 3306 6038.080 6040.760
\hypertarget{c01:cue:3306}{}那天你再去看畫面。
% cue 3307 6040.760 6041.760
\hypertarget{c01:cue:3307}{}你一早到第一次畫面。
% cue 3308 6041.760 6043.070
\hypertarget{c01:cue:3308}{}你應該找得到。
% cue 3309 6043.070 6044.070
\hypertarget{c01:cue:3309}{}我是幾個經驗人裡面。
% cue 3310 6044.070 6046.070
\hypertarget{c01:cue:3310}{}我一定沒有穿整套西裝。
% cue 3311 6046.070 6048.060
\hypertarget{c01:cue:3311}{}因為我不知道這件事情。
% cue 3312 6048.060 6050.350
\hypertarget{c01:cue:3312}{}你可以去證明我講的話是真的。
% cue 3313 6050.350 6052.410
\hypertarget{c01:cue:3313}{}其他人又穿得很整齊。
% cue 3314 6052.410 6053.410
\hypertarget{c01:cue:3314}{}我不知道。
% cue 3315 6053.410 6054.380
\hypertarget{c01:cue:3315}{}我用的時候上班。
% cue 3316 6054.380 6056.660
\hypertarget{c01:cue:3316}{}我去穿一整套西裝上班。
% cue 3317 6056.660 6057.660
\hypertarget{c01:cue:3317}{}整個經驗。
% cue 3318 6057.660 6058.660
\hypertarget{c01:cue:3318}{}做的沒有那麼嚴重。
% cue 3319 6058.660 6059.660
\hypertarget{c01:cue:3319}{}那天也沒有外面。
% cue 3320 6059.660 6060.620
\hypertarget{c01:cue:3320}{}有外面不一樣。
% cue 3321 6060.620 6061.620
\hypertarget{c01:cue:3321}{}有外面不一樣。
% cue 3322 6061.620 6062.940
\hypertarget{c01:cue:3322}{}有外面不一樣。
% cue 3323 6062.940 6064.060
\hypertarget{c01:cue:3323}{}沒有。
% cue 3324 6064.060 6065.060
\hypertarget{c01:cue:3324}{}所以你進一刻。
% cue 3325 6065.060 6066.060
\hypertarget{c01:cue:3325}{}我就穿得很隨便的那種。
% cue 3326 6066.060 6069.200
\hypertarget{c01:cue:3326}{}上下是不一樣的顏色的衣服。
% cue 3327 6069.200 6071.190
\hypertarget{c01:cue:3327}{}那臨時通知我去總統府。
% cue 3328 6071.190 6073.190
\hypertarget{c01:cue:3328}{}那我當然要穿下。
% cue 3329 6073.190 6074.190
\hypertarget{c01:cue:3329}{}我記得我沒有打領帶。
% cue 3330 6074.190 6075.190
\hypertarget{c01:cue:3330}{}是隨便打領帶。
% cue 3331 6075.190 6077.190
\hypertarget{c01:cue:3331}{}還是就是故事成功的形狀。
% cue 3332 6077.190 6078.190
\hypertarget{c01:cue:3332}{}你看你出大不出來。

\tc{1:41:18}
% cue 3333 6078.190 6079.770
\hypertarget{c01:cue:3333}{}那那個事情故事呢，
% cue 3334 6079.770 6083.080
\hypertarget{c01:cue:3334}{}聽說是名單送進了總統府裡頭。
% cue 3335 6083.080 6085.430
\hypertarget{c01:cue:3335}{}有人不講人是誰。
% cue 3336 6085.430 6086.430
\hypertarget{c01:cue:3336}{}我大概知道。
% cue 3337 6086.430 6087.900
\hypertarget{c01:cue:3337}{}跟總統講。
% cue 3338 6087.900 6090.920
\hypertarget{c01:cue:3338}{}先生你要不要競選連任。
% cue 3339 6090.920 6093.270
\hypertarget{c01:cue:3339}{}如果你要競選連任的話呢，
% cue 3340 6093.270 6095.270
\hypertarget{c01:cue:3340}{}你的婦女票現在都對。
% cue 3341 6095.270 6096.900
\hypertarget{c01:cue:3341}{}那怎麼辦呢？
% cue 3342 6096.900 6097.900
\hypertarget{c01:cue:3342}{}都請你女人。
% cue 3343 6097.900 6100.050
\hypertarget{c01:cue:3343}{}重視母性。
% cue 3344 6100.050 6102.330
\hypertarget{c01:cue:3344}{}婚姻問中。
% cue 3345 6102.330 6103.330
\hypertarget{c01:cue:3345}{}你大概會說好。
% cue 3346 6103.330 6104.900
\hypertarget{c01:cue:3346}{}所以進去。
% cue 3347 6104.900 6107.730
\hypertarget{c01:cue:3347}{}本來請你送進去的名單。
% cue 3348 6107.730 6108.730
\hypertarget{c01:cue:3348}{}本來是那樣。
% cue 3349 6108.730 6109.730
\hypertarget{c01:cue:3349}{}送進去的意思是說。
% cue 3350 6109.730 6113.040
\hypertarget{c01:cue:3350}{}當時去查報紙上有一個大法官。
% cue 3351 6113.040 6114.040
\hypertarget{c01:cue:3351}{}審薦小組\footnote{大法官提名審薦小組，在總統提名司法院大法官之前審議人選。2011 年提名案中，副總統蕭萬長承認審薦過程「不夠周延」；民間監督司法院大法官人選聯盟批評審薦小組組成不當、提名程序不夠嚴謹，主張小組應納入各界人士。\url{https://newtalk.tw/news/view/2011-04-01/13238}。}。
% cue 3352 6114.040 6115.100
\hypertarget{c01:cue:3352}{}規見神明小組。
% cue 3353 6115.100 6116.100
\hypertarget{c01:cue:3353}{}審薦小組。
% cue 3354 6116.100 6117.280
\hypertarget{c01:cue:3354}{}是建議的名單。
% cue 3355 6117.280 6118.280
\hypertarget{c01:cue:3355}{}但是對外公布的。
% cue 3356 6118.280 6120.920
\hypertarget{c01:cue:3356}{}但是有人都不知道。
% cue 3357 6120.920 6123.230
\hypertarget{c01:cue:3357}{}送進去的時候是這樣。
% cue 3358 6123.230 6125.150
\hypertarget{c01:cue:3358}{}但是因為有人去講啊。
% cue 3359 6125.150 6127.450
\hypertarget{c01:cue:3359}{}所以出來的名單幾乎都是女。
% cue 3360 6127.450 6128.760
\hypertarget{c01:cue:3360}{}一共四個大法官。
% cue 3361 6128.760 6129.760
\hypertarget{c01:cue:3361}{}只有一個是男的。
% cue 3362 6129.760 6131.910
\hypertarget{c01:cue:3362}{}這個男的是我們父中。
% cue 3363 6131.910 6133.100
\hypertarget{c01:cue:3363}{}替他求情人。
% cue 3364 6133.100 6135.660
\hypertarget{c01:cue:3364}{}三個月。

\tc{1:42:15}
% cue 3365 6135.660 6138.700
\hypertarget{c01:cue:3365}{}那找到這三個出現了大問題。
% cue 3366 6138.700 6142.080
\hypertarget{c01:cue:3366}{}這三個人各有一樣。
% cue 3367 6142.080 6143.080
\hypertarget{c01:cue:3367}{}剛才我講的這個。
% cue 3368 6143.080 6144.100
\hypertarget{c01:cue:3368}{}就是我們法官。
% cue 3369 6144.100 6145.090
\hypertarget{c01:cue:3369}{}女的。
% cue 3370 6145.090 6146.570
\hypertarget{c01:cue:3370}{}第二個女的。
% cue 3371 6146.570 6149.790
\hypertarget{c01:cue:3371}{}一個最高行政法院的法官。
% cue 3372 6149.790 6151.750
\hypertarget{c01:cue:3372}{}連續七年考基因的。
% cue 3373 6151.750 6153.870
\hypertarget{c01:cue:3373}{}法官要連續考基因的。
% cue 3374 6153.870 6154.870
\hypertarget{c01:cue:3374}{}這非常難度。
% cue 3375 6154.870 6156.730
\hypertarget{c01:cue:3375}{}連續七年。
% cue 3376 6156.730 6158.690
\hypertarget{c01:cue:3376}{}十個偉大。
% cue 3377 6158.690 6159.690
\hypertarget{c01:cue:3377}{}所以你看這個總統是。
% cue 3378 6159.690 6160.680
\hypertarget{c01:cue:3378}{}長期忽丟。
% cue 3379 6160.680 6162.870
\hypertarget{c01:cue:3379}{}滿久不知道自己在做什麼。
% cue 3380 6162.870 6165.140
\hypertarget{c01:cue:3380}{}第三個女的。
% cue 3381 6165.140 6166.490
\hypertarget{c01:cue:3381}{}是由美國綠卡\footnote{綠卡、楓葉卡分別是美國與加拿大永久居留證的俗稱。2011 年的大法官被提名人中，陳碧玉被質疑曾擁有美國國籍與永久居留權，羅昌發曾擁有加拿大永久居留權，民間監督司法院大法官人選聯盟因此呼籲被提名人婉拒提名。\url{https://newtalk.tw/news/view/2011-04-14/13544}。}。
% cue 3382 6166.490 6169.320
\hypertarget{c01:cue:3382}{}第四個那個男的。
% cue 3383 6169.320 6170.320
\hypertarget{c01:cue:3383}{}是由加拿大楓葉卡。

\tc{1:42:50}
% cue 3384 6170.320 6172.990
\hypertarget{c01:cue:3384}{}所以這四個人就報了。
% cue 3385 6172.990 6175.110
\hypertarget{c01:cue:3385}{}你看那個報是要全部在講。
% cue 3386 6175.110 6177.020
\hypertarget{c01:cue:3386}{}罵罵罵罵犯天的。
% cue 3387 6177.020 6179.320
\hypertarget{c01:cue:3387}{}所以馬一條要衝提名衝提名。
% cue 3388 6179.320 6180.320
\hypertarget{c01:cue:3388}{}不要忘記當時在那裡。
% cue 3389 6180.320 6183.710
\hypertarget{c01:cue:3389}{}那些人。
% cue 3390 6183.710 6185.090
\hypertarget{c01:cue:3390}{}是很會被作亂的。
% cue 3391 6185.090 6187.240
\hypertarget{c01:cue:3391}{}這個題目太好做了。
% cue 3392 6187.240 6190.630
\hypertarget{c01:cue:3392}{}結果總統府顯然。
% cue 3393 6190.630 6192.630
\hypertarget{c01:cue:3393}{}我不知道他們怎麼決定。
% cue 3394 6192.630 6193.750
\hypertarget{c01:cue:3394}{}但顯然是我。
% cue 3395 6193.750 6195.750
\hypertarget{c01:cue:3395}{}拿我更多中華當黨進來。
% cue 3396 6195.750 6198.900
\hypertarget{c01:cue:3396}{}接下來就沒事。
% cue 3397 6198.900 6200.050
\hypertarget{c01:cue:3397}{}這真的沒事。
% cue 3398 6200.050 6201.050
\hypertarget{c01:cue:3398}{}但是他們三個。
% cue 3399 6201.050 6203.050
\hypertarget{c01:cue:3399}{}他們總統副總統。
% cue 3400 6203.050 6204.050
\hypertarget{c01:cue:3400}{}加上司法院院長\footnote{2011 年 3 月 31 日晚間，馬英九總統、蕭萬長副總統與司法院院長賴浩敏就邵燕玲的提名發布新聞稿致歉，4 月 1 日中午三人召開記者會再度致歉。賴浩敏於 2010 年 10 月 13 日至 2016 年 10 月 31 日任司法院院長並為大法官。\url{https://newtalk.tw/news/view/2011-04-01/13238}。}三個。
% cue 3401 6204.050 6205.050
\hypertarget{c01:cue:3401}{}一起出來。
% cue 3402 6205.050 6207.050
\hypertarget{c01:cue:3402}{}連一個鞠躬報欺三次。
% cue 3403 6207.050 6208.810
\hypertarget{c01:cue:3403}{}那報紙都知道。
% cue 3404 6208.810 6210.220
\hypertarget{c01:cue:3404}{}就為了他的提名。

\tc{1:43:30}
% cue 3405 6210.220 6212.220
\hypertarget{c01:cue:3405}{}但是後來問題都解決了。
% cue 3406 6212.220 6215.440
\hypertarget{c01:cue:3406}{}那個綠卡的放棄。
% cue 3407 6215.440 6216.440
\hypertarget{c01:cue:3407}{}楓葉卡的放棄。
% cue 3408 6216.440 6218.650
\hypertarget{c01:cue:3408}{}那個女的說是不嚴重。
% cue 3409 6218.650 6219.950
\hypertarget{c01:cue:3409}{}很太不設定。
% cue 3410 6219.950 6220.950
\hypertarget{c01:cue:3410}{}她只是表現不好。
% cue 3411 6220.950 6221.950
\hypertarget{c01:cue:3411}{}也不嚴重。
% cue 3412 6221.950 6223.200
\hypertarget{c01:cue:3412}{}那對重就是，
% cue 3413 6223.200 6225.530
\hypertarget{c01:cue:3413}{}這個實在太差了。
% cue 3414 6225.530 6227.860
\hypertarget{c01:cue:3414}{}尤其是民間團體。
% cue 3415 6227.860 6228.860
\hypertarget{c01:cue:3415}{}貴起來抗凌。
% cue 3416 6228.860 6230.400
\hypertarget{c01:cue:3416}{}所以這個壞掉。
% cue 3417 6230.400 6232.650
\hypertarget{c01:cue:3417}{}我是原來在上面做個武器群。

\tc{1:43:52}
% cue 3418 6232.650 6234.890
\hypertarget{c01:cue:3418}{}所以講到這個。
% cue 3419 6234.890 6236.240
\hypertarget{c01:cue:3419}{}誰叫你問的。
% cue 3420 6236.240 6239.090
\hypertarget{c01:cue:3420}{}我以為是女士充。
% cue 3421 6239.090 6241.850
\hypertarget{c01:cue:3421}{}不過我沒有講的。
% cue 3422 6241.850 6243.140
\hypertarget{c01:cue:3422}{}我跟妳講說。
% cue 3423 6243.140 6246.110
\hypertarget{c01:cue:3423}{}老師這一生從來很傳奇。
% cue 3424 6246.110 6248.480
\hypertarget{c01:cue:3424}{}看老師是不是沒有個性的。
% cue 3425 6248.480 6250.600
\hypertarget{c01:cue:3425}{}我也不鳥人。
% cue 3426 6250.600 6251.600
\hypertarget{c01:cue:3426}{}我也不買誰的帳。
% cue 3427 6251.600 6252.660
\hypertarget{c01:cue:3427}{}我也不跟你買之心。
% cue 3428 6252.660 6253.660
\hypertarget{c01:cue:3428}{}但是呢，
% cue 3429 6253.660 6255.120
\hypertarget{c01:cue:3429}{}就是覺得。
% cue 3430 6255.120 6256.420
\hypertarget{c01:cue:3430}{}這些事情。
% cue 3431 6256.420 6257.420
\hypertarget{c01:cue:3431}{}有什麼好真。

\tc{1:44:17}
% cue 3432 6257.420 6258.570
\hypertarget{c01:cue:3432}{}如果妳要問我。
% cue 3433 6258.570 6259.720
\hypertarget{c01:cue:3433}{}應該問我說。
% cue 3434 6259.720 6260.720
\hypertarget{c01:cue:3434}{}老師那妳後面。
% cue 3435 6260.720 6263.170
\hypertarget{c01:cue:3435}{}我現在最兄弟那裡很慘。
% cue 3436 6263.170 6264.340
\hypertarget{c01:cue:3436}{}被誰卡被我的同學。
% cue 3437 6264.340 6265.340
\hypertarget{c01:cue:3437}{}誰。
% cue 3438 6265.340 6266.440
\hypertarget{c01:cue:3438}{}蔡英文。
% cue 3439 6266.440 6267.530
\hypertarget{c01:cue:3439}{}我是同屆同班的。
% cue 3440 6267.530 6269.770
\hypertarget{c01:cue:3440}{}才到法律系的。
% cue 3441 6269.770 6272.340
\hypertarget{c01:cue:3441}{}那我不去到大法官。
% cue 3442 6272.340 6273.340
\hypertarget{c01:cue:3442}{}我帶你給妳先當。
% cue 3443 6273.340 6274.340
\hypertarget{c01:cue:3443}{}我現在很可憐。
% cue 3444 6274.340 6277.590
\hypertarget{c01:cue:3444}{}但跑到東吳來才有飯吃。
% cue 3445 6277.590 6278.590
\hypertarget{c01:cue:3445}{}所以妳覺得。
% cue 3446 6278.590 6280.170
\hypertarget{c01:cue:3446}{}是要高興。
% cue 3447 6280.170 6281.170
\hypertarget{c01:cue:3447}{}是要退休還是要生氣。
% cue 3448 6281.170 6286.650
\hypertarget{c01:cue:3448}{}人生不算不準。

\tc{1:44:46}
% cue 3449 6286.650 6289.670
\hypertarget{c01:cue:3449}{}我如果還在中研院。
% cue 3450 6289.670 6291.670
\hypertarget{c01:cue:3450}{}我當然知道妳現在為什麼。
% cue 3451 6291.670 6293.270
\hypertarget{c01:cue:3451}{}但是妳覺得。
% cue 3452 6293.270 6295.510
\hypertarget{c01:cue:3452}{}應該不會比曾建仁長。
% cue 3453 6295.510 6297.910
\hypertarget{c01:cue:3453}{}曾建仁是誰？
% cue 3454 6297.910 6298.910
\hypertarget{c01:cue:3454}{}現在中研院的院長。
% cue 3455 6298.910 6303.720
\hypertarget{c01:cue:3455}{}曾建仁是院長是現在。
% cue 3456 6303.720 6304.720
\hypertarget{c01:cue:3456}{}但是在之前。
% cue 3457 6304.720 6306.750
\hypertarget{c01:cue:3457}{}他是中研院的特別研究。
% cue 3458 6306.750 6309.390
\hypertarget{c01:cue:3458}{}老師在離開中研院的時候。
% cue 3459 6309.390 6310.830
\hypertarget{c01:cue:3459}{}我也是特別研究。
% cue 3460 6310.830 6311.800
\hypertarget{c01:cue:3460}{}當然。
% cue 3461 6311.800 6312.800
\hypertarget{c01:cue:3461}{}他是院士。
% cue 3462 6312.800 6313.980
\hypertarget{c01:cue:3462}{}他的心複不叫到。
% cue 3463 6313.980 6314.980
\hypertarget{c01:cue:3463}{}我在那個時候。
% cue 3464 6314.980 6315.950
\hypertarget{c01:cue:3464}{}在我的時候沒有院士。
% cue 3465 6315.950 6318.070
\hypertarget{c01:cue:3465}{}大家都說我被當第一的院士。
% cue 3466 6318.070 6319.070
\hypertarget{c01:cue:3466}{}結果是王生建。
% cue 3467 6319.070 6320.070
\hypertarget{c01:cue:3467}{}當的第一院士。
% cue 3468 6320.070 6321.070
\hypertarget{c01:cue:3468}{}也不太經驗。
% cue 3469 6321.070 6322.130
\hypertarget{c01:cue:3469}{}因為王生建是我們的。
% cue 3470 6322.130 6324.220
\hypertarget{c01:cue:3470}{}中研院法律所籌備處的。
% cue 3471 6324.220 6325.700
\hypertarget{c01:cue:3471}{}他們叫諮詢委員。
% cue 3472 6325.700 6326.620
\hypertarget{c01:cue:3472}{}為了造型。
% cue 3473 6326.620 6331.040
\hypertarget{c01:cue:3473}{}但是我跟妳講說。
% cue 3474 6331.040 6332.080
\hypertarget{c01:cue:3474}{}那些都是沒有錢。
% cue 3475 6332.080 6334.270
\hypertarget{c01:cue:3475}{}有錢的是妳的階級。
% cue 3476 6334.270 6335.270
\hypertarget{c01:cue:3476}{}我們階級是特別的。
% cue 3477 6335.270 6336.430
\hypertarget{c01:cue:3477}{}所以陳建仁\footnote{1998 年當選中央研究院院士，2006 年起任中研院基因體研究中心特聘研究員，2011 年至 2015 年任中研院副院長。2016 年至 2020 年任副總統，卸任後回任中研院特聘研究員並放棄卸任副總統禮遇；2023 年 1 月 31 日至 2024 年 5 月 20 日任行政院院長。總統府於 2026 年 1 月 2 日特任其為中央研究院第 13 任院長，任期自 2026 年 6 月 21 日起。\url{https://udn.com/news/story/7314/9241406}。}是特別的。
% cue 3478 6336.430 6337.430
\hypertarget{c01:cue:3478}{}所以陳建仁是特別的。
% cue 3479 6337.430 6338.520
\hypertarget{c01:cue:3479}{}我也是特別的。

\tc{1:45:38}
% cue 3480 6338.520 6339.520
\hypertarget{c01:cue:3480}{}那我就舉個例子。
% cue 3481 6339.520 6340.580
\hypertarget{c01:cue:3481}{}陳建仁給妳聽。
% cue 3482 6340.580 6341.580
\hypertarget{c01:cue:3482}{}陳建仁不願意放棄特別的。
% cue 3483 6341.580 6343.860
\hypertarget{c01:cue:3483}{}所以他去借掉班。
% cue 3484 6343.860 6344.890
\hypertarget{c01:cue:3484}{}陳建仁業長。
% cue 3485 6344.890 6345.890
\hypertarget{c01:cue:3485}{}借掉班副總。
% cue 3486 6345.890 6350.230
\hypertarget{c01:cue:3486}{}他記得不記得陳建仁。
% cue 3487 6350.230 6351.290
\hypertarget{c01:cue:3487}{}副總統做完。
% cue 3488 6351.290 6352.580
\hypertarget{c01:cue:3488}{}人家說。
% cue 3489 6352.580 6355.020
\hypertarget{c01:cue:3489}{}那妳該領副總統的禮遇\footnote{卸任總統副總統禮遇條例第 3 條第 1 項規定，卸任副總統享有按月致送新臺幣 18 萬元禮遇金等禮遇；同條第 4 項規定，除邀請參加國家大典外，各款禮遇之有效期間與其任職期間相同。第 4 條第 1 項第 1 款規定，卸任後再任有給公職者，停止禮遇。\url{https://law.moj.gov.tw/LawClass/LawAll.aspx?pcode=A0030002}。}啊。
% cue 3490 6355.020 6358.840
\hypertarget{c01:cue:3490}{}後來人家挖出來。
% cue 3491 6358.840 6359.870
\hypertarget{c01:cue:3491}{}他為什麼放棄。
% cue 3492 6359.870 6360.990
\hypertarget{c01:cue:3492}{}因為副總統禮遇。
% cue 3493 6360.990 6361.990
\hypertarget{c01:cue:3493}{}沒有他中研院的。
% cue 3494 6361.990 6363.990
\hypertarget{c01:cue:3494}{}特別嚴訓人欺水來得好。
% cue 3495 6363.990 6365.150
\hypertarget{c01:cue:3495}{}第一。
% cue 3496 6365.150 6365.250
\hypertarget{c01:cue:3496}{}第二。
% cue 3497 6365.150 6366.330
\hypertarget{c01:cue:3497}{}副總統禮遇。
% cue 3498 6366.330 6367.330
\hypertarget{c01:cue:3498}{}你做四年副總統。
% cue 3499 6367.330 6368.330
\hypertarget{c01:cue:3499}{}你只能領四年。
% cue 3500 6368.330 6369.330
\hypertarget{c01:cue:3500}{}他四年以後。
% cue 3501 6369.330 6370.330
\hypertarget{c01:cue:3501}{}你要怎麼樣。
% cue 3502 6370.330 6371.480
\hypertarget{c01:cue:3502}{}會欺負吧。
% cue 3503 6371.480 6372.480
\hypertarget{c01:cue:3503}{}中研院特別嚴訓人。
% cue 3504 6372.480 6373.480
\hypertarget{c01:cue:3504}{}你做到四年。
% cue 3505 6373.480 6374.540
\hypertarget{c01:cue:3505}{}他那一句。
% cue 3506 6374.540 6375.980
\hypertarget{c01:cue:3506}{}他不特別的一個人。
% cue 3507 6375.980 6377.500
\hypertarget{c01:cue:3507}{}所以老四四五四對不大。

\tc{1:46:17}
% cue 3508 6377.500 6381.890
\hypertarget{c01:cue:3508}{}所以你說好。
% cue 3509 6381.890 6382.890
\hypertarget{c01:cue:3509}{}好多人呼吸。
% cue 3510 6382.890 6383.890
\hypertarget{c01:cue:3510}{}哎呀。
% cue 3511 6383.890 6383.990
\hypertarget{c01:cue:3511}{}你當大法官。
% cue 3512 6383.890 6384.890
\hypertarget{c01:cue:3512}{}都要把四四年。
% cue 3513 6384.890 6386.270
\hypertarget{c01:cue:3513}{}都要講一大堆。
% cue 3514 6386.270 6387.270
\hypertarget{c01:cue:3514}{}那怎麼知道。
% cue 3515 6387.270 6387.370
\hypertarget{c01:cue:3515}{}拖到一拆。
% cue 3516 6387.270 6388.270
\hypertarget{c01:cue:3516}{}你拆成的包。
% cue 3517 6388.270 6389.840
\hypertarget{c01:cue:3517}{}退休息。
% cue 3518 6389.840 6391.510
\hypertarget{c01:cue:3518}{}砍了一大堆。
% cue 3519 6391.510 6391.610
\hypertarget{c01:cue:3519}{}對你說人生。
% cue 3520 6391.510 6393.210
\hypertarget{c01:cue:3520}{}怎麼預料的。
% cue 3521 6393.210 6394.260
\hypertarget{c01:cue:3521}{}那你說。
% cue 3522 6394.260 6395.260
\hypertarget{c01:cue:3522}{}這麼被陷害。
% cue 3523 6395.260 6396.260
\hypertarget{c01:cue:3523}{}有害死嗎？
% cue 3524 6396.260 6397.260
\hypertarget{c01:cue:3524}{}也沒有。
% cue 3525 6397.260 6399.190
\hypertarget{c01:cue:3525}{}我都把他整個被弄講話。
% cue 3526 6399.190 6400.190
\hypertarget{c01:cue:3526}{}所以我也算樂觀。
% cue 3527 6400.190 6401.190
\hypertarget{c01:cue:3527}{}而且我也算命。
% cue 3528 6401.190 6403.600
\hypertarget{c01:cue:3528}{}所以你要派死。
% cue 3529 6403.600 6407.790
\hypertarget{c01:cue:3529}{}照理講。
% cue 3530 6407.790 6410.790
\hypertarget{c01:cue:3530}{}你什麼人的薪水都別砍。
% cue 3531 6410.790 6411.730
\hypertarget{c01:cue:3531}{}大法官都不砍。
% cue 3532 6411.730 6413.160
\hypertarget{c01:cue:3532}{}對。
% cue 3533 6413.160 6414.160
\hypertarget{c01:cue:3533}{}你要謊葬他。

\tc{1:46:54}
% cue 3534 6414.160 6415.660
\hypertarget{c01:cue:3534}{}我現在講這個。
% cue 3535 6415.660 6416.660
\hypertarget{c01:cue:3535}{}有興趣去看。
% cue 3536 6416.660 6418.630
\hypertarget{c01:cue:3536}{}五二九號解釋。
% cue 3537 6418.630 6421.520
\hypertarget{c01:cue:3537}{}所以其實我沒有放這個。
% cue 3538 6421.520 6423.800
\hypertarget{c01:cue:3538}{}你進去看五八九\footnote{司法院釋字第 589 號解釋（2005 年 1 月 28 日）。解釋文稱「憲法對特定職位為維護其獨立行使職權而定有任期保障者，其職務之性質與應隨政黨更迭或政策變更而進退之政務人員不同」，其因任期保障所生之信賴利益須充分保護；2004 年 1 月 7 日公布之政務人員退職撫卹條例，對受任期保障之監察委員等人員未設得擇領月退職酬勞金之規定，「顯對其應受保護之信賴利益，並未有合理之保障，與前開憲法意旨有違」。\url{https://cons.judicial.gov.tw/docdata.aspx?fid=100&id=310770}。}。
% cue 3539 6423.800 6425.620
\hypertarget{c01:cue:3539}{}講到這樣。
% cue 3540 6425.620 6426.620
\hypertarget{c01:cue:3540}{}就再多想下去。
% cue 3541 6426.620 6427.620
\hypertarget{c01:cue:3541}{}就不好。
% cue 3542 6427.620 6427.720
\hypertarget{c01:cue:3542}{}好。
% cue 3543 6427.620 6428.970
\hypertarget{c01:cue:3543}{}還有什麼其他問題。
% cue 3544 6428.970 6429.970
\hypertarget{c01:cue:3544}{}行說。
% cue 3545 6429.970 6430.970
\hypertarget{c01:cue:3545}{}現在幾點。
% cue 3546 6430.970 6432.250
\hypertarget{c01:cue:3546}{}我要趕快。
% cue 3547 6432.250 6433.570
\hypertarget{c01:cue:3547}{}鋪立戰場。
% cue 3548 6433.570 6433.670
\hypertarget{c01:cue:3548}{}任務一下。
% cue 3549 6433.570 6435.810
\hypertarget{c01:cue:3549}{}就有四四頭都被你們。
% cue 3550 6435.810 6436.770
\hypertarget{c01:cue:3550}{}混出什麼。
% cue 3551 6436.770 6438.210
\hypertarget{c01:cue:3551}{}關注來。
% cue 3552 6438.210 6439.210
\hypertarget{c01:cue:3552}{}這太危險。
% cue 3553 6439.210 6440.210
\hypertarget{c01:cue:3553}{}我剛好就說過了。
% cue 3554 6440.210 6441.210
\hypertarget{c01:cue:3554}{}第一不準入依。
% cue 3555 6441.210 6442.210
\hypertarget{c01:cue:3555}{}第二出去以後。
% cue 3556 6442.210 6443.210
\hypertarget{c01:cue:3556}{}我就不承認我講話。
% cue 3557 6443.210 6444.210
\hypertarget{c01:cue:3557}{}然後你拿來。
% cue 3558 6444.210 6445.750
\hypertarget{c01:cue:3558}{}你講。
% cue 3559 6445.750 6447.190
\hypertarget{c01:cue:3559}{}是你嗎？
% cue 3560 6447.190 6448.150
\hypertarget{c01:cue:3560}{}怎麼沒有。
% cue 3561 6448.150 6454.200
\hypertarget{c01:cue:3561}{}其他公訣。
% cue 3562 6454.200 6455.170
\hypertarget{c01:cue:3562}{}但是也五十九。
% cue 3563 6455.170 6456.610
\hypertarget{c01:cue:3563}{}對你下課還有二十年。
% cue 3564 6456.610 6458.150
\hypertarget{c01:cue:3564}{}說。
% cue 3565 6458.150 6459.150
\hypertarget{c01:cue:3565}{}等一下你是第二次。
% cue 3566 6459.150 6460.400
\hypertarget{c01:cue:3566}{}先讓別人。
% cue 3567 6460.400 6462.370
\hypertarget{c01:cue:3567}{}其他公訣有多問題。
% cue 3568 6462.370 6463.430
\hypertarget{c01:cue:3568}{}那我才不是有很多人要講話嗎？
% cue 3569 6463.430 6464.430
\hypertarget{c01:cue:3569}{}你不講就是他要一直講。
% cue 3570 6464.430 6466.390
\hypertarget{c01:cue:3570}{}我不是說他講話不對。
% cue 3571 6466.390 6467.390
\hypertarget{c01:cue:3571}{}但是我希望大家可以。
% cue 3572 6467.390 6468.390
\hypertarget{c01:cue:3572}{}繼續講。

\tc{1:47:48}
% cue 3573 6468.390 6469.390
\hypertarget{c01:cue:3573}{}然後我先說。
% cue 3574 6469.390 6475.170
\hypertarget{c01:cue:3574}{}然後我會給你們單子填。
% cue 3575 6475.170 6478.150
\hypertarget{c01:cue:3575}{}我會給你們在課堂上填單子。
% cue 3576 6478.150 6481.150
\hypertarget{c01:cue:3576}{}那因為你們是老師的學生。
% cue 3577 6481.150 6484.110
\hypertarget{c01:cue:3577}{}所以我會用比較優惠的價格給你們。
% cue 3578 6484.110 6487.110
\hypertarget{c01:cue:3578}{}不會是你們替班班的書商買的價格。
% cue 3579 6487.110 6488.590
\hypertarget{c01:cue:3579}{}這是我跟書商的一個例子。
% cue 3580 6488.590 6491.000
\hypertarget{c01:cue:3580}{}我當然是照顧自己的學生。

\tc{1:48:11}
% cue 3581 6491.000 6492.000
\hypertarget{c01:cue:3581}{}有。
% cue 3582 6492.000 6494.300
\hypertarget{c01:cue:3582}{}簡單的問題。
% cue 3583 6494.300 6495.300
\hypertarget{c01:cue:3583}{}還有嗎？
% cue 3584 6495.300 6496.390
\hypertarget{c01:cue:3584}{}有關於你畢業的事情。
% cue 3585 6496.390 6497.390
\hypertarget{c01:cue:3585}{}還有什麼要問嗎？
% cue 3586 6497.390 6498.990
\hypertarget{c01:cue:3586}{}你們都沒有問題嗎？
% cue 3587 6498.990 6501.020
\hypertarget{c01:cue:3587}{}可見你們經驗不夠。
% cue 3588 6501.020 6503.200
\hypertarget{c01:cue:3588}{}應該給我很多奇怪的問題。
% cue 3589 6503.200 6504.200
\hypertarget{c01:cue:3589}{}來 該給我問。
% cue 3590 6504.200 6508.090
\hypertarget{c01:cue:3590}{}好。

\section{答題的要求}

\tc{1:48:28}
% cue 3591 6508.090 6519.250
\hypertarget{c01:cue:3591}{}現在講的這個問題。
% cue 3592 6519.250 6522.080
\hypertarget{c01:cue:3592}{}比較適合我去問。
% cue 3593 6522.080 6523.080
\hypertarget{c01:cue:3593}{}補習班的老師。
% cue 3594 6523.080 6524.300
\hypertarget{c01:cue:3594}{}因為補習班的老師。
% cue 3595 6524.300 6525.560
\hypertarget{c01:cue:3595}{}他們會跟你講。
% cue 3596 6525.560 6528.180
\hypertarget{c01:cue:3596}{}會拆題或解題的方法。
% cue 3597 6528.180 6529.180
\hypertarget{c01:cue:3597}{}不下。
% cue 3598 6529.180 6530.310
\hypertarget{c01:cue:3598}{}是他們的工作。
% cue 3599 6530.310 6532.320
\hypertarget{c01:cue:3599}{}那你問我的話呢，
% cue 3600 6532.320 6533.330
\hypertarget{c01:cue:3600}{}我是一個曾經出題。
% cue 3601 6533.330 6535.490
\hypertarget{c01:cue:3601}{}曾經閱卷改題。
% cue 3602 6535.490 6537.490
\hypertarget{c01:cue:3602}{}或決定法官律師任務局。
% cue 3603 6537.490 6538.490
\hypertarget{c01:cue:3603}{}誰的那種的人。
% cue 3604 6538.490 6539.640
\hypertarget{c01:cue:3604}{}那我的建議呢，
% cue 3605 6539.640 6541.640
\hypertarget{c01:cue:3605}{}其另外一種是這種。
% cue 3606 6541.640 6546.600
\hypertarget{c01:cue:3606}{}官學。
% cue 3607 6546.600 6548.240
\hypertarget{c01:cue:3607}{}他們是私學。
% cue 3608 6548.240 6550.810
\hypertarget{c01:cue:3608}{}我們在國際的觀點不太一樣。
% cue 3609 6550.810 6553.510
\hypertarget{c01:cue:3609}{}那我簡單地回應你的問題。
% cue 3610 6553.510 6554.950
\hypertarget{c01:cue:3610}{}這樣子。

\tc{1:49:14}
% cue 3611 6554.950 6556.360
\hypertarget{c01:cue:3611}{}問題哪來。
% cue 3612 6556.360 6557.580
\hypertarget{c01:cue:3612}{}第一件事情。
% cue 3613 6557.580 6558.830
\hypertarget{c01:cue:3613}{}要看清楚。
% cue 3614 6558.830 6559.830
\hypertarget{c01:cue:3614}{}題目。
% cue 3615 6559.830 6561.240
\hypertarget{c01:cue:3615}{}我不是說。
% cue 3616 6561.240 6562.240
\hypertarget{c01:cue:3616}{}你要去猜。
% cue 3617 6562.240 6563.240
\hypertarget{c01:cue:3617}{}這是誰出的題。
% cue 3618 6563.240 6564.360
\hypertarget{c01:cue:3618}{}因為現在有很多人會。
% cue 3619 6564.360 6565.550
\hypertarget{c01:cue:3619}{}很不幸。
% cue 3620 6565.550 6567.550
\hypertarget{c01:cue:3620}{}我擬他的口氣。
% cue 3621 6567.550 6568.650
\hypertarget{c01:cue:3621}{}我擬他的題目。
% cue 3622 6568.650 6570.650
\hypertarget{c01:cue:3622}{}你們覺得是他要出題。
% cue 3623 6570.650 6571.770
\hypertarget{c01:cue:3623}{}其實。
% cue 3624 6571.770 6572.670
\hypertarget{c01:cue:3624}{}我跟他一直在笑完。
% cue 3625 6572.670 6574.050
\hypertarget{c01:cue:3625}{}我要測試一下。
% cue 3626 6574.050 6575.050
\hypertarget{c01:cue:3626}{}誰是他的實作。
% cue 3627 6575.050 6577.050
\hypertarget{c01:cue:3627}{}你們搭的跟他借錢一樣。
% cue 3628 6577.050 6578.050
\hypertarget{c01:cue:3628}{}我讓你們死一片。
% cue 3629 6578.050 6581.230
\hypertarget{c01:cue:3629}{}那這點我們就不討論。
% cue 3630 6581.230 6582.230
\hypertarget{c01:cue:3630}{}這個很複雜。
% cue 3631 6582.230 6583.670
\hypertarget{c01:cue:3631}{}第一件事情是。
% cue 3632 6583.670 6585.890
\hypertarget{c01:cue:3632}{}你就是要看清楚題目。
% cue 3633 6585.890 6586.890
\hypertarget{c01:cue:3633}{}為什麼？
% cue 3634 6586.890 6589.020
\hypertarget{c01:cue:3634}{}因為你要知道他要問什麼。
% cue 3635 6589.020 6591.560
\hypertarget{c01:cue:3635}{}你要想他要測是什麼？
% cue 3636 6591.560 6592.560
\hypertarget{c01:cue:3636}{}他要測是什麼？
% cue 3637 6592.560 6594.220
\hypertarget{c01:cue:3637}{}他要問你的不同決定分立。
% cue 3638 6594.220 6596.280
\hypertarget{c01:cue:3638}{}還是他問你的不同。
% cue 3639 6596.280 6597.250
\hypertarget{c01:cue:3639}{}現在是什麼制。
% cue 3640 6597.250 6598.660
\hypertarget{c01:cue:3640}{}半總統制。
% cue 3641 6598.660 6599.660
\hypertarget{c01:cue:3641}{}還是老師講。
% cue 3642 6599.660 6601.670
\hypertarget{c01:cue:3642}{}我講的修正式的總統制\footnote{湯德宗《權力分立新論 卷一》（增訂三版，2005）頁 54 圖一〈我國中央政府體制更迭示意圖〉將現制標示為「修正式總統制」；該書稱憲法本文採「修正式內閣制」（頁 62），論 1997 年第四次修憲後的現制，認其已成為總統制，又因覆議表決門檻降低等因素，與美國相較僅屬「弱勢總統制」，並認現制並非半總統制（頁 168–170）。}。
% cue 3643 6601.670 6604.880
\hypertarget{c01:cue:3643}{}你要先知道他想問什麼。

\tc{1:50:04}
% cue 3644 6604.880 6606.530
\hypertarget{c01:cue:3644}{}第二點。
% cue 3645 6606.530 6608.860
\hypertarget{c01:cue:3645}{}你要很清楚地講。
% cue 3646 6608.860 6611.010
\hypertarget{c01:cue:3646}{}現在按照題目的意思。
% cue 3647 6611.010 6613.140
\hypertarget{c01:cue:3647}{}這裡面牽涉到幾個問題。
% cue 3648 6613.140 6615.710
\hypertarget{c01:cue:3648}{}你要先把問題點找到。
% cue 3649 6615.710 6617.440
\hypertarget{c01:cue:3649}{}第三。
% cue 3650 6617.440 6618.660
\hypertarget{c01:cue:3650}{}你的答題。
% cue 3651 6618.660 6619.910
\hypertarget{c01:cue:3651}{}無論你心情如何。
% cue 3652 6619.910 6623.180
\hypertarget{c01:cue:3652}{}或者你怎麼樣手忙腳亂。
% cue 3653 6623.180 6625.400
\hypertarget{c01:cue:3653}{}你要儘可能讓你的字。
% cue 3654 6625.400 6627.780
\hypertarget{c01:cue:3654}{}對考官不要造成危險。
% cue 3655 6627.780 6628.780
\hypertarget{c01:cue:3655}{}現在你們打字。
% cue 3656 6628.780 6631.250
\hypertarget{c01:cue:3656}{}那打字就看你們快很快。
% cue 3657 6631.250 6634.250
\hypertarget{c01:cue:3657}{}打字我不知道他要怎麼處理。
% cue 3658 6634.250 6635.250
\hypertarget{c01:cue:3658}{}因為打字可能電腦被臉拖。
% cue 3659 6635.250 6636.250
\hypertarget{c01:cue:3659}{}啊怎麼樣。
% cue 3660 6636.250 6637.250
\hypertarget{c01:cue:3660}{}不曉得。
% cue 3661 6637.250 6638.630
\hypertarget{c01:cue:3661}{}那以前是要寫的。
% cue 3662 6638.630 6639.720
\hypertarget{c01:cue:3662}{}我們最討厭的一種。
% cue 3663 6639.720 6643.010
\hypertarget{c01:cue:3663}{}那個字寫得密密麻麻。
% cue 3664 6643.010 6644.010
\hypertarget{c01:cue:3664}{}而且黏在一起。
% cue 3665 6644.010 6646.680
\hypertarget{c01:cue:3665}{}這種這種邏輯模。
% cue 3666 6646.680 6647.680
\hypertarget{c01:cue:3666}{}我看到這個。
% cue 3667 6647.680 6648.680
\hypertarget{c01:cue:3667}{}我一定可以看。
% cue 3668 6648.680 6650.810
\hypertarget{c01:cue:3668}{}安扣幾個。

\tc{1:50:50}
% cue 3669 6650.810 6652.120
\hypertarget{c01:cue:3669}{}你知不知道考試的時候。
% cue 3670 6652.120 6653.120
\hypertarget{c01:cue:3670}{}我們閱卷不要求。
% cue 3671 6653.120 6654.120
\hypertarget{c01:cue:3671}{}對老師。
% cue 3672 6654.120 6655.180
\hypertarget{c01:cue:3672}{}要叫你。
% cue 3673 6655.180 6656.850
\hypertarget{c01:cue:3673}{}一進去那些個房間。
% cue 3674 6656.850 6657.850
\hypertarget{c01:cue:3674}{}通通都是這種聲音。
% cue 3675 6657.850 6658.850
\hypertarget{c01:cue:3675}{}痴情疲惑。
% cue 3676 6658.850 6661.390
\hypertarget{c01:cue:3676}{}因為我們要拿我們的豬皮。
% cue 3677 6661.390 6663.510
\hypertarget{c01:cue:3677}{}點在你們的卷子上。
% cue 3678 6663.510 6664.510
\hypertarget{c01:cue:3678}{}跟著你們的劇鬥爭。
% cue 3679 6664.510 6665.760
\hypertarget{c01:cue:3679}{}你黏在一起。
% cue 3680 6665.760 6667.760
\hypertarget{c01:cue:3680}{}我都點不知道該點在哪裡。
% cue 3681 6667.760 6669.980
\hypertarget{c01:cue:3681}{}我看三份的卷子才看一份。
% cue 3682 6669.980 6671.390
\hypertarget{c01:cue:3682}{}這有人考不考驗。
% cue 3683 6671.390 6674.250
\hypertarget{c01:cue:3683}{}而且閱卷錢都很少。
% cue 3684 6674.250 6675.690
\hypertarget{c01:cue:3684}{}以前十二塊。
% cue 3685 6675.690 6677.260
\hypertarget{c01:cue:3685}{}現在調整。
% cue 3686 6677.260 6678.260
\hypertarget{c01:cue:3686}{}已經好多年。
% cue 3687 6678.260 6679.260
\hypertarget{c01:cue:3687}{}我們有全個當大的萬閱卷。
% cue 3688 6679.260 6684.690
\hypertarget{c01:cue:3688}{}所以字體不要太多長。
% cue 3689 6684.690 6687.010
\hypertarget{c01:cue:3689}{}尤其不要不便於閱卷。

\tc{1:51:27}
% cue 3690 6687.010 6688.010
\hypertarget{c01:cue:3690}{}這是找自己。
% cue 3691 6688.010 6689.900
\hypertarget{c01:cue:3691}{}然後你寫作的時候。
% cue 3692 6689.900 6692.860
\hypertarget{c01:cue:3692}{}你要盡可能讓人感覺你有層次。
% cue 3693 6692.860 6695.080
\hypertarget{c01:cue:3693}{}你是層次分明。
% cue 3694 6695.080 6696.080
\hypertarget{c01:cue:3694}{}一二三。
% cue 3695 6696.080 6698.340
\hypertarget{c01:cue:3695}{}你不要統統統混在一起。
% cue 3696 6698.340 6699.340
\hypertarget{c01:cue:3696}{}也許你這一段寫得不好。
% cue 3697 6699.340 6700.340
\hypertarget{c01:cue:3697}{}切一點。
% cue 3698 6700.340 6701.340
\hypertarget{c01:cue:3698}{}下一點再寫。
% cue 3699 6701.340 6702.560
\hypertarget{c01:cue:3699}{}你統統黏在一起。
% cue 3700 6702.560 6703.560
\hypertarget{c01:cue:3700}{}套有五分段。
% cue 3701 6703.560 6704.780
\hypertarget{c01:cue:3701}{}說真的。
% cue 3702 6704.780 6708.290
\hypertarget{c01:cue:3702}{}你要在期待老師在很短的時間裡。
% cue 3703 6708.290 6710.250
\hypertarget{c01:cue:3703}{}完全掌握住你的答題。
% cue 3704 6710.250 6711.470
\hypertarget{c01:cue:3704}{}不容易。
% cue 3705 6711.470 6712.620
\hypertarget{c01:cue:3705}{}比分點。
% cue 3706 6712.620 6713.620
\hypertarget{c01:cue:3706}{}我至少會跟著點看到。
% cue 3707 6713.620 6715.420
\hypertarget{c01:cue:3707}{}我知道。
% cue 3708 6715.420 6716.420
\hypertarget{c01:cue:3708}{}我裡面的三個點都打中。
% cue 3709 6716.420 6718.820
\hypertarget{c01:cue:3709}{}也許你竟然打得非常完整。
% cue 3710 6718.820 6719.820
\hypertarget{c01:cue:3710}{}但是都打中。
% cue 3711 6719.820 6722.030
\hypertarget{c01:cue:3711}{}總是你知道這一大點在行事。
% cue 3712 6722.030 6723.030
\hypertarget{c01:cue:3712}{}要好。

\tc{1:52:03}
% cue 3713 6723.030 6724.570
\hypertarget{c01:cue:3713}{}最後不管你怎麼寫。
% cue 3714 6724.570 6726.110
\hypertarget{c01:cue:3714}{}一定要寫成狀況。
% cue 3715 6726.110 6728.420
\hypertarget{c01:cue:3715}{}最後要下個結論再離開。
% cue 3716 6728.420 6731.850
\hypertarget{c01:cue:3716}{}就算一個老師的人。
% cue 3717 6731.850 6732.850
\hypertarget{c01:cue:3717}{}他前面都看到。
% cue 3718 6732.850 6734.810
\hypertarget{c01:cue:3718}{}他最後結論至少看到。
% cue 3719 6734.810 6736.590
\hypertarget{c01:cue:3719}{}所以你分數也不會太長。
% cue 3720 6736.590 6738.330
\hypertarget{c01:cue:3720}{}我現在講的是官學。
% cue 3721 6738.330 6740.940
\hypertarget{c01:cue:3721}{}官方教人讀書。
% cue 3722 6740.940 6741.940
\hypertarget{c01:cue:3722}{}教人寫作是這麼寫的。
% cue 3723 6741.940 6743.410
\hypertarget{c01:cue:3723}{}我不講私學。
% cue 3724 6743.410 6744.410
\hypertarget{c01:cue:3724}{}私學讓人不知道。
% cue 3725 6744.410 6746.440
\hypertarget{c01:cue:3725}{}你什麼老師要教你什麼。
% cue 3726 6746.440 6748.340
\hypertarget{c01:cue:3726}{}那你需要練習。
% cue 3727 6748.340 6750.950
\hypertarget{c01:cue:3727}{}好事跟平常是一樣的。

\tc{1:52:30}
% cue 3728 6750.950 6753.460
\hypertarget{c01:cue:3728}{}所以你今天跟我答題。
% cue 3729 6753.460 6755.550
\hypertarget{c01:cue:3729}{}老師上課資訊期。
% cue 3730 6755.550 6756.550
\hypertarget{c01:cue:3730}{}開始被你們對答。
% cue 3731 6756.550 6758.860
\hypertarget{c01:cue:3731}{}我問的問題你聽不懂。
% cue 3732 6758.860 6759.860
\hypertarget{c01:cue:3732}{}你卻答不對。
% cue 3733 6759.860 6761.460
\hypertarget{c01:cue:3733}{}所以一定要聽懂。
% cue 3734 6761.460 6762.680
\hypertarget{c01:cue:3734}{}對吧？
% cue 3735 6762.680 6763.640
\hypertarget{c01:cue:3735}{}第二。
% cue 3736 6763.640 6764.640
\hypertarget{c01:cue:3736}{}要回答我之前。
% cue 3737 6764.640 6765.670
\hypertarget{c01:cue:3737}{}一定要稍微組織一下。
% cue 3738 6765.670 6769.010
\hypertarget{c01:cue:3738}{}在旁邊還有一張付卷。
% cue 3739 6769.010 6770.100
\hypertarget{c01:cue:3739}{}可以寫些字。
% cue 3740 6770.100 6771.670
\hypertarget{c01:cue:3740}{}先把自己的要點。
% cue 3741 6771.670 6773.210
\hypertarget{c01:cue:3741}{}一三點下考。
% cue 3742 6773.210 6775.500
\hypertarget{c01:cue:3742}{}想法都把一二三思緒調整。
% cue 3743 6775.500 6776.460
\hypertarget{c01:cue:3743}{}是一二三。
% cue 3744 6776.460 6777.460
\hypertarget{c01:cue:3744}{}不是三二一。
% cue 3745 6777.460 6778.930
\hypertarget{c01:cue:3745}{}還是二三。
% cue 3746 6778.930 6780.000
\hypertarget{c01:cue:3746}{}怎麼讀怎麼寫。
% cue 3747 6780.000 6781.000
\hypertarget{c01:cue:3747}{}讓人家覺得。
% cue 3748 6781.000 6782.000
\hypertarget{c01:cue:3748}{}你這個投稿思緒。
% cue 3749 6782.000 6783.000
\hypertarget{c01:cue:3749}{}是否可以跟你。
% cue 3750 6783.000 6784.570
\hypertarget{c01:cue:3750}{}不然讓七八糟都不理。
% cue 3751 6784.570 6785.520
\hypertarget{c01:cue:3751}{}怎麼。
% cue 3752 6785.520 6786.560
\hypertarget{c01:cue:3752}{}我覺得你很會背熟。
% cue 3753 6786.560 6788.070
\hypertarget{c01:cue:3753}{}但這真正點沒有組織。

\tc{1:53:08}
% cue 3754 6788.070 6791.060
\hypertarget{c01:cue:3754}{}最後要給你機會。
% cue 3755 6791.060 6792.280
\hypertarget{c01:cue:3755}{}寫作的時候。
% cue 3756 6792.280 6795.340
\hypertarget{c01:cue:3756}{}字跡要儘管讓大家看。
% cue 3757 6795.340 6796.780
\hypertarget{c01:cue:3757}{}不要找那個閱卷子。
% cue 3758 6796.780 6798.480
\hypertarget{c01:cue:3758}{}都是老頭子。
% cue 3759 6798.480 6799.600
\hypertarget{c01:cue:3759}{}都很累。
% cue 3760 6799.600 6800.880
\hypertarget{c01:cue:3760}{}其實不是只有你。
% cue 3761 6800.880 6802.940
\hypertarget{c01:cue:3761}{}通常八千個人。
% cue 3762 6802.940 6803.940
\hypertarget{c01:cue:3762}{}一萬元考試。
% cue 3763 6803.940 6806.040
\hypertarget{c01:cue:3763}{}大概兩個禮拜。
% cue 3764 6806.040 6807.260
\hypertarget{c01:cue:3764}{}要閱卷完。
% cue 3765 6807.260 6808.450
\hypertarget{c01:cue:3765}{}結果對不對？
% cue 3766 6808.450 6809.630
\hypertarget{c01:cue:3766}{}我們還要學校上課。
% cue 3767 6809.630 6810.630
\hypertarget{c01:cue:3767}{}還有大堆有個沒有的事情。
% cue 3768 6810.630 6812.850
\hypertarget{c01:cue:3768}{}所以大家都是在課堂中間。
% cue 3769 6812.850 6813.850
\hypertarget{c01:cue:3769}{}跑去裡面。
% cue 3770 6813.850 6814.850
\hypertarget{c01:cue:3770}{}抖抖抖抖抖抖。
% cue 3771 6814.850 6815.810
\hypertarget{c01:cue:3771}{}進到那個講。
% cue 3772 6815.810 6816.800
\hypertarget{c01:cue:3772}{}越迴腿。
% cue 3773 6816.800 6817.800
\hypertarget{c01:cue:3773}{}全部都是這種聲音。
% cue 3774 6817.800 6826.960
\hypertarget{c01:cue:3774}{}你們學到很多東西。

\section{中央研究院歐美研究所與社科所法律組的經歷}

\tc{1:53:46}
% cue 3775 6826.960 6829.180
\hypertarget{c01:cue:3775}{}這個是你們會問的人會到。
% cue 3776 6829.180 6831.470
\hypertarget{c01:cue:3776}{}每學習問的人都不一樣。
% cue 3777 6831.470 6832.720
\hypertarget{c01:cue:3777}{}所以不曉得。
% cue 3778 6832.720 6833.940
\hypertarget{c01:cue:3778}{}你要問什麼。

\tc{1:53:53}
\speaker{學生}
% cue 3779 6833.940 6835.940
\hypertarget{c01:cue:3779}{}是庭安老師可以講一下。
% cue 3780 6835.940 6839.940
\hypertarget{c01:cue:3780}{}在中研院早期的那個美國文化研究所\footnote{中央研究院美國研究中心於 1972 年 3 月 12 日成立，1974 年 7 月改制為美國文化研究所，1991 年 8 月 3 日更名為歐美研究所。\url{https://ir.sinica.edu.tw/cris/ou/ou00023}。}。
% cue 3781 6839.940 6844.260
\hypertarget{c01:cue:3781}{}跟後面歐美研究所的那個經驗。

\tc{1:54:04}
\speaker{老師}
% cue 3782 6844.260 6845.260
\hypertarget{c01:cue:3782}{}這是什麼？
% cue 3783 6845.260 6846.290
\hypertarget{c01:cue:3783}{}經歷啊。
% cue 3784 6846.290 6848.830
\hypertarget{c01:cue:3784}{}我到那邊去的時間其實不長。
% cue 3785 6848.830 6850.410
\hypertarget{c01:cue:3785}{}我是 1985 年。
% cue 3786 6850.410 6852.410
\hypertarget{c01:cue:3786}{}從哈佛拿到碩士以後。
% cue 3787 6852.410 6854.570
\hypertarget{c01:cue:3787}{}進入中研院。
% cue 3788 6854.570 6855.800
\hypertarget{c01:cue:3788}{}1985 年。
% cue 3789 6855.800 6857.140
\hypertarget{c01:cue:3789}{}我還記得我去的時候。
% cue 3790 6857.140 6858.140
\hypertarget{c01:cue:3790}{}是下午熱都掉了。
% cue 3791 6858.140 6859.940
\hypertarget{c01:cue:3791}{}太陽曬了。
% cue 3792 6859.940 6861.080
\hypertarget{c01:cue:3792}{}從日後照。
% cue 3793 6861.080 6863.500
\hypertarget{c01:cue:3793}{}開始報告。
% cue 3794 6863.500 6866.500
\hypertarget{c01:cue:3794}{}那我當時去到那個時候。
% cue 3795 6866.500 6867.420
\hypertarget{c01:cue:3795}{}那個時候。
% cue 3796 6867.420 6869.110
\hypertarget{c01:cue:3796}{}名字叫美國文化也叫做。

\tc{1:54:29}
% cue 3797 6869.110 6871.170
\hypertarget{c01:cue:3797}{}後來我去美國。
% cue 3798 6871.170 6872.170
\hypertarget{c01:cue:3798}{}我做了兩年。
% cue 3799 6872.170 6874.330
\hypertarget{c01:cue:3799}{}然後我就出國就輸了。
% cue 3800 6874.330 6875.330
\hypertarget{c01:cue:3800}{}我借了一筆錢。
% cue 3801 6875.330 6876.290
\hypertarget{c01:cue:3801}{}貸款。
% cue 3802 6876.290 6877.440
\hypertarget{c01:cue:3802}{}因為我們新主人說。
% cue 3803 6877.440 6878.440
\hypertarget{c01:cue:3803}{}乾保貸款。
% cue 3804 6878.440 6879.500
\hypertarget{c01:cue:3804}{}然後我去念。
% cue 3805 6879.500 6881.950
\hypertarget{c01:cue:3805}{}所以我的博士沒有在好。
% cue 3806 6881.950 6882.940
\hypertarget{c01:cue:3806}{}因為我看太貴了。
% cue 3807 6882.940 6885.110
\hypertarget{c01:cue:3807}{}而且我考到我們 TLATE 去。
% cue 3808 6885.110 6887.260
\hypertarget{c01:cue:3808}{}TLATE 對我非常好。
% cue 3809 6887.260 6889.920
\hypertarget{c01:cue:3809}{}那你需要知道這個太多複雜的細節。
% cue 3810 6889.920 6891.920
\hypertarget{c01:cue:3810}{}因為這個人比我還不氣死人。
% cue 3811 6891.920 6893.170
\hypertarget{c01:cue:3811}{}哈老師這裡還是有一個。
% cue 3812 6893.170 6894.610
\hypertarget{c01:cue:3812}{}我總共拿了兩個。
% cue 3813 6894.610 6895.610
\hypertarget{c01:cue:3813}{}我去的人。
% cue 3814 6895.610 6896.610
\hypertarget{c01:cue:3814}{}可能就拿到博士。
% cue 3815 6896.610 6898.640
\hypertarget{c01:cue:3815}{}我的博士論文是有書書的。
% cue 3816 6898.640 6899.640
\hypertarget{c01:cue:3816}{}不是跟蔡英文一樣。
% cue 3817 6899.640 6900.670
\hypertarget{c01:cue:3817}{}550 頁。
% cue 3818 6900.670 6902.760
\hypertarget{c01:cue:3818}{}可以在這個書館找到。
% cue 3819 6902.760 6903.760
\hypertarget{c01:cue:3819}{}中央研究院正式書館。

\tc{1:55:03}
% cue 3820 6903.760 6905.550
\hypertarget{c01:cue:3820}{}所以我就不懂。
% cue 3821 6905.550 6906.550
\hypertarget{c01:cue:3821}{}這些老師。
% cue 3822 6906.550 6909.640
\hypertarget{c01:cue:3822}{}他們沒有出版他們的博士論文。
% cue 3823 6909.640 6910.670
\hypertarget{c01:cue:3823}{}怎麼能當老師。
% cue 3824 6910.670 6912.630
\hypertarget{c01:cue:3824}{}我們在中研院。
% cue 3825 6912.630 6913.720
\hypertarget{c01:cue:3825}{}你要回來。
% cue 3826 6913.720 6914.900
\hypertarget{c01:cue:3826}{}要升等。
% cue 3827 6914.900 6916.470
\hypertarget{c01:cue:3827}{}進來拿路。
% cue 3828 6916.470 6918.170
\hypertarget{c01:cue:3828}{}你沒有路。
% cue 3829 6918.170 6919.170
\hypertarget{c01:cue:3829}{}是不可能升等。
% cue 3830 6919.170 6921.060
\hypertarget{c01:cue:3830}{}進來就進來。
% cue 3831 6921.060 6922.920
\hypertarget{c01:cue:3831}{}我在 Harvard。
% cue 3832 6922.920 6924.070
\hypertarget{c01:cue:3832}{}進 Harvard。

\tc{1:55:24}
% cue 3833 6924.070 6925.070
\hypertarget{c01:cue:3833}{}我進中研院 1985。
% cue 3834 6925.070 6927.070
\hypertarget{c01:cue:3834}{}我之所以能進中研院。
% cue 3835 6927.070 6929.200
\hypertarget{c01:cue:3835}{}那是我還沒有故事。
% cue 3836 6929.200 6930.290
\hypertarget{c01:cue:3836}{}是因為我的碩士論文。
% cue 3837 6930.290 6931.290
\hypertarget{c01:cue:3837}{}在美國的碩士論文。
% cue 3838 6931.290 6933.350
\hypertarget{c01:cue:3838}{}已經登在美國的期刊。
% cue 3839 6933.350 6935.380
\hypertarget{c01:cue:3839}{}有興趣可以檢查。
% cue 3840 6935.380 6937.380
\hypertarget{c01:cue:3840}{}叫 Administrative Law Review\footnote{美國行政法期刊。湯德宗 On the Legal Protection of Civil Service Employees from Arbitrary Dismissal 一文刊於該刊第 37 卷，頁 37–93（1985）。\url{https://dennis-tang.site/scholarly-journals/}。}。
% cue 3841 6937.380 6938.410
\hypertarget{c01:cue:3841}{}在第 37 卷。
% cue 3842 6938.410 6940.760
\hypertarget{c01:cue:3842}{}非常好是這個文字。
% cue 3843 6940.760 6943.920
\hypertarget{c01:cue:3843}{}這在美國。
% cue 3844 6943.920 6945.200
\hypertarget{c01:cue:3844}{}Tick Homer。
% cue 3845 6945.200 6946.200
\hypertarget{c01:cue:3845}{}你去看連學生那麼多。
% cue 3846 6946.200 6948.200
\hypertarget{c01:cue:3846}{}有人有幾個人都在美國。
% cue 3847 6948.200 6949.200
\hypertarget{c01:cue:3847}{}學生呢，
% cue 3848 6949.200 6951.250
\hypertarget{c01:cue:3848}{}是這樣的身份。
% cue 3849 6951.250 6952.280
\hypertarget{c01:cue:3849}{}所以大家。
% cue 3850 6952.280 6953.340
\hypertarget{c01:cue:3850}{}所問的。
% cue 3851 6953.340 6954.360
\hypertarget{c01:cue:3851}{}都做。
% cue 3852 6954.360 6955.680
\hypertarget{c01:cue:3852}{}這樣進來。
% cue 3853 6955.680 6958.430
\hypertarget{c01:cue:3853}{}那因為當時沒有法律所。
% cue 3854 6958.430 6964.510
\hypertarget{c01:cue:3854}{}請問我去哪裡。

\tc{1:56:04}
% cue 3855 6964.510 6965.510
\hypertarget{c01:cue:3855}{}美國所後來。
% cue 3856 6965.510 6966.430
\hypertarget{c01:cue:3856}{}在我出國。
% cue 3857 6966.430 6968.370
\hypertarget{c01:cue:3857}{}出的時候。
% cue 3858 6968.370 6969.330
\hypertarget{c01:cue:3858}{}他改名了。
% cue 3859 6969.330 6971.320
\hypertarget{c01:cue:3859}{}叫歐美。
% cue 3860 6971.320 6973.350
\hypertarget{c01:cue:3860}{}因為他覺得歐洲一樣重視。
% cue 3861 6973.350 6975.020
\hypertarget{c01:cue:3861}{}一樣是重視。
% cue 3862 6975.020 6976.140
\hypertarget{c01:cue:3862}{}所以名字就改。
% cue 3863 6976.140 6977.140
\hypertarget{c01:cue:3863}{}那我回來以後。
% cue 3864 6977.140 6978.140
\hypertarget{c01:cue:3864}{}我就回歐美。
% cue 3865 6978.140 6979.260
\hypertarget{c01:cue:3865}{}出去。
% cue 3866 6979.260 6981.150
\hypertarget{c01:cue:3866}{}我是那邊。
% cue 3867 6981.150 6982.500
\hypertarget{c01:cue:3867}{}留資流行出去。
% cue 3868 6982.500 6983.500
\hypertarget{c01:cue:3868}{}回到那裡去。
% cue 3869 6983.500 6984.480
\hypertarget{c01:cue:3869}{}那那裡去我大概待了兩年。
% cue 3870 6984.480 6986.250
\hypertarget{c01:cue:3870}{}三年以後。

\tc{1:56:26}
% cue 3871 6986.250 6987.250
\hypertarget{c01:cue:3871}{}我就換。
% cue 3872 6987.250 6988.250
\hypertarget{c01:cue:3872}{}我就換到社科所\footnote{中央研究院三民主義研究所於 1981 年 8 月正式成所，1990 年 4 月改稱中山人文社會科學研究所，簡稱社科所；1991 年 5 月社科所法律組正式成立。2004 年 7 月法律學研究所籌備處成立時，社科所轉型為人文社會科學研究中心。\url{https://publication.iias.sinica.edu.tw/40818012.pdf}。}去。
% cue 3873 6988.250 6989.600
\hypertarget{c01:cue:3873}{}到社科所去。
% cue 3874 6989.600 6990.600
\hypertarget{c01:cue:3874}{}社科所就是，
% cue 3875 6990.600 6993.550
\hypertarget{c01:cue:3875}{}以前的三年組。
% cue 3876 6993.550 6994.550
\hypertarget{c01:cue:3876}{}它裡面有很多組。
% cue 3877 6994.550 6996.960
\hypertarget{c01:cue:3877}{}那其中有一個。
% cue 3878 6996.960 6997.950
\hypertarget{c01:cue:3878}{}是叫法律組。
% cue 3879 6997.950 6999.150
\hypertarget{c01:cue:3879}{}那我沒有去以前。
% cue 3880 6999.150 7000.150
\hypertarget{c01:cue:3880}{}只有四個人。
% cue 3881 7000.150 7001.310
\hypertarget{c01:cue:3881}{}他用貴階。
% cue 3882 7001.310 7002.270
\hypertarget{c01:cue:3882}{}五人才能成組。
% cue 3883 7002.270 7003.430
\hypertarget{c01:cue:3883}{}那我去了。
% cue 3884 7003.430 7004.430
\hypertarget{c01:cue:3884}{}所長邀請我去。
% cue 3885 7004.430 7005.620
\hypertarget{c01:cue:3885}{}所長大的一巴掌我去。
% cue 3886 7005.620 7008.420
\hypertarget{c01:cue:3886}{}那個所長是個院士。
% cue 3887 7008.420 7009.610
\hypertarget{c01:cue:3887}{}過世。

\tc{1:56:49}
% cue 3888 7009.610 7011.310
\hypertarget{c01:cue:3888}{}所以老師的這一輩子。
% cue 3889 7011.310 7013.340
\hypertarget{c01:cue:3889}{}也是有人提攜的。
% cue 3890 7013.340 7014.490
\hypertarget{c01:cue:3890}{}就有人看見。
% cue 3891 7014.490 7015.640
\hypertarget{c01:cue:3891}{}我進中研院。
% cue 3892 7015.640 7016.660
\hypertarget{c01:cue:3892}{}歐美所。
% cue 3893 7016.660 7018.210
\hypertarget{c01:cue:3893}{}是一個長輩。
% cue 3894 7018.210 7021.420
\hypertarget{c01:cue:3894}{}我進到社科所。
% cue 3895 7021.420 7022.800
\hypertarget{c01:cue:3895}{}法律組。
% cue 3896 7022.800 7024.020
\hypertarget{c01:cue:3896}{}代理委員組。
% cue 3897 7024.020 7024.120
\hypertarget{c01:cue:3897}{}主任。
% cue 3898 7024.020 7025.020
\hypertarget{c01:cue:3898}{}第四個長輩。
% cue 3899 7025.020 7026.750
\hypertarget{c01:cue:3899}{}那這些人跟我過去。
% cue 3900 7026.750 7027.750
\hypertarget{c01:cue:3900}{}都沒有關係。
% cue 3901 7027.750 7028.750
\hypertarget{c01:cue:3901}{}譬如是說。
% cue 3902 7028.750 7029.750
\hypertarget{c01:cue:3902}{}是我的私交。
% cue 3903 7029.750 7030.750
\hypertarget{c01:cue:3903}{}是我的爸爸。
% cue 3904 7030.750 7033.820
\hypertarget{c01:cue:3904}{}是我的父親。
% cue 3905 7033.820 7034.820
\hypertarget{c01:cue:3905}{}所以有段式律器。
% cue 3906 7034.820 7035.940
\hypertarget{c01:cue:3906}{}反正人的一生當中。
% cue 3907 7035.940 7036.960
\hypertarget{c01:cue:3907}{}一定不碰到貴人。

\tc{1:57:16}
% cue 3908 7036.960 7039.850
\hypertarget{c01:cue:3908}{}像你剛才出去讀書。
% cue 3909 7039.850 7042.070
\hypertarget{c01:cue:3909}{}因為今天上了這個課。
% cue 3910 7042.070 7045.340
\hypertarget{c01:cue:3910}{}然後湯老師用大法官的寫寫給你推薦。
% cue 3911 7045.340 7046.390
\hypertarget{c01:cue:3911}{}結果你進了報了。
% cue 3912 7046.390 7047.510
\hypertarget{c01:cue:3912}{}當然會不會寫寫了。
% cue 3913 7047.510 7048.600
\hypertarget{c01:cue:3913}{}會啊。
% cue 3914 7048.600 7049.600
\hypertarget{c01:cue:3914}{}那老師算不算你恩人。
% cue 3915 7049.600 7051.580
\hypertarget{c01:cue:3915}{}人是肯來說的。
% cue 3916 7051.580 7054.570
\hypertarget{c01:cue:3916}{}那每年都有報復。
% cue 3917 7054.570 7055.570
\hypertarget{c01:cue:3917}{}真寫規見信。
% cue 3918 7055.570 7057.880
\hypertarget{c01:cue:3918}{}不知道是什麼樣的結果。
% cue 3919 7057.880 7059.040
\hypertarget{c01:cue:3919}{}反正每個人情形不懂。
% cue 3920 7059.040 7060.450
\hypertarget{c01:cue:3920}{}最簡單的情形是這樣。
% cue 3921 7060.450 7062.510
\hypertarget{c01:cue:3921}{}沒有特別少的時間。

\tc{1:57:42}
% cue 3922 7062.510 7065.470
\hypertarget{c01:cue:3922}{}後來因為有法律組。
% cue 3923 7065.470 7066.600
\hypertarget{c01:cue:3923}{}不去法律組。
% cue 3924 7066.600 7068.730
\hypertarget{c01:cue:3924}{}後來又待了很久。
% cue 3925 7068.730 7071.300
\hypertarget{c01:cue:3925}{}法律組一直等到一二折。
% cue 3926 7071.300 7073.210
\hypertarget{c01:cue:3926}{}所以我們要成立法律所。
% cue 3927 7073.210 7074.720
\hypertarget{c01:cue:3927}{}請你去當中。
% cue 3928 7074.720 7076.670
\hypertarget{c01:cue:3928}{}做一個復習的復習。
% cue 3929 7076.670 7078.240
\hypertarget{c01:cue:3929}{}就是這樣。
% cue 3930 7078.240 7079.240
\hypertarget{c01:cue:3930}{}在中研院一直這樣轉來轉去。
% cue 3931 7079.240 7080.590
\hypertarget{c01:cue:3931}{}你們可以想到。
% cue 3932 7080.590 7081.590
\hypertarget{c01:cue:3932}{}在很長一段時間。
% cue 3933 7081.590 7082.590
\hypertarget{c01:cue:3933}{}在社科所。
% cue 3934 7082.590 7083.590
\hypertarget{c01:cue:3934}{}很長一段時間。
% cue 3935 7083.590 7085.490
\hypertarget{c01:cue:3935}{}都沒有法律所。

\tc{1:58:05}
% cue 3936 7085.490 7088.160
\hypertarget{c01:cue:3936}{}所以某個意義上。
% cue 3937 7088.160 7089.310
\hypertarget{c01:cue:3937}{}老師跟其他的老師。
% cue 3938 7089.310 7090.210
\hypertarget{c01:cue:3938}{}法律的老師。
% cue 3939 7090.210 7092.750
\hypertarget{c01:cue:3939}{}有一個很大的經驗是不同的。
% cue 3940 7092.750 7095.510
\hypertarget{c01:cue:3940}{}如果你要問。
% cue 3941 7095.510 7097.670
\hypertarget{c01:cue:3941}{}什麼叫寄人籬下。
% cue 3942 7097.670 7098.860
\hypertarget{c01:cue:3942}{}那你要問我這種。
% cue 3943 7098.860 7101.750
\hypertarget{c01:cue:3943}{}寄人籬下你聽懂嗎？
% cue 3944 7101.750 7103.550
\hypertarget{c01:cue:3944}{}中文裡面有。
% cue 3945 7103.550 7106.020
\hypertarget{c01:cue:3945}{}但是你什麼叫寄人籬下。
% cue 3946 7106.020 7108.020
\hypertarget{c01:cue:3946}{}你是引流這種感受的。

\tc{1:58:28}
% cue 3947 7108.020 7110.360
\hypertarget{c01:cue:3947}{}過去一個例子給你聽。
% cue 3948 7110.360 7111.420
\hypertarget{c01:cue:3948}{}我要升等。
% cue 3949 7111.420 7112.420
\hypertarget{c01:cue:3949}{}在社科所。
% cue 3950 7112.420 7114.510
\hypertarget{c01:cue:3950}{}我要升等。
% cue 3951 7114.510 7115.510
\hypertarget{c01:cue:3951}{}審查我的人。
% cue 3952 7115.510 7116.510
\hypertarget{c01:cue:3952}{}我們所裡面的人要投票。
% cue 3953 7116.510 7118.600
\hypertarget{c01:cue:3953}{}絕大部分都不是法律人。
% cue 3954 7118.600 7119.600
\hypertarget{c01:cue:3954}{}只有一個法律人。
% cue 3955 7119.600 7123.580
\hypertarget{c01:cue:3955}{}那你覺得他們審。
% cue 3956 7123.580 7125.830
\hypertarget{c01:cue:3956}{}他們審說可以審不可以。
% cue 3957 7125.830 7128.420
\hypertarget{c01:cue:3957}{}所以就要寄人籬下。

\tc{1:58:48}
% cue 3958 7128.420 7130.740
\hypertarget{c01:cue:3958}{}然後在社科所。
% cue 3959 7130.740 7131.740
\hypertarget{c01:cue:3959}{}這麼多組裡面。
% cue 3960 7131.740 7132.870
\hypertarget{c01:cue:3960}{}律主是最小的。
% cue 3961 7132.870 7134.190
\hypertarget{c01:cue:3961}{}是最後成立的。
% cue 3962 7134.190 7135.280
\hypertarget{c01:cue:3962}{}是最沒用的。
% cue 3963 7135.280 7139.610
\hypertarget{c01:cue:3963}{}那你更不要說。
% cue 3964 7139.610 7141.610
\hypertarget{c01:cue:3964}{}中央研究院最後決定要成立法律所。
% cue 3965 7141.610 7142.830
\hypertarget{c01:cue:3965}{}那大家都看好啊。
% cue 3966 7142.830 7145.660
\hypertarget{c01:cue:3966}{}我們所裡面還有一個人。
% cue 3967 7145.660 7146.880
\hypertarget{c01:cue:3967}{}是我們所裡的同仁。
% cue 3968 7146.880 7149.610
\hypertarget{c01:cue:3968}{}無人的吃裡爬外。
% cue 3969 7149.610 7150.800
\hypertarget{c01:cue:3969}{}跟大家講。
% cue 3970 7150.800 7153.120
\hypertarget{c01:cue:3970}{}說啊我相信這個湯某人啊。
% cue 3971 7153.120 7154.110
\hypertarget{c01:cue:3971}{}當這個所謂。
% cue 3972 7154.110 7155.110
\hypertarget{c01:cue:3972}{}宿主任。
% cue 3973 7155.110 7156.110
\hypertarget{c01:cue:3973}{}不會超過三個月。
% cue 3974 7156.110 7159.020
\hypertarget{c01:cue:3974}{}這會垮掉。
% cue 3975 7159.020 7160.000
\hypertarget{c01:cue:3975}{}沒有垮掉。
% cue 3976 7160.000 7162.360
\hypertarget{c01:cue:3976}{}做成了。
% cue 3977 7162.360 7165.470
\hypertarget{c01:cue:3977}{}就是知道這種環境是什麼。
% cue 3978 7165.470 7167.560
\hypertarget{c01:cue:3978}{}要寄人。
% cue 3979 7167.560 7168.590
\hypertarget{c01:cue:3979}{}就是這個法律。
% cue 3980 7168.590 7169.590
\hypertarget{c01:cue:3980}{}在中央研究院。
% cue 3981 7169.590 7171.810
\hypertarget{c01:cue:3981}{}它不被承認為一個重要的學科。
% cue 3982 7171.810 7173.260
\hypertarget{c01:cue:3982}{}這就是可以說。
% cue 3983 7173.260 7175.260
\hypertarget{c01:cue:3983}{}根本沒有把你看在乾歹的。
% cue 3984 7175.260 7176.420
\hypertarget{c01:cue:3984}{}沒有放在人裡。

\tc{1:59:36}
% cue 3985 7176.420 7178.960
\hypertarget{c01:cue:3985}{}我們不要以為我們這是法律系的。
% cue 3986 7178.960 7179.960
\hypertarget{c01:cue:3986}{}哎唷 我們法律系的這樣。
% cue 3987 7179.960 7180.960
\hypertarget{c01:cue:3987}{}但是你去到那個地方。
% cue 3988 7180.960 7181.960
\hypertarget{c01:cue:3988}{}想發聲音都沒有。
% cue 3989 7181.960 7187.840
\hypertarget{c01:cue:3989}{}要證明你出來的價值。
% cue 3990 7187.840 7190.580
\hypertarget{c01:cue:3990}{}而且要在那邊的偏遠的地方。
% cue 3991 7190.580 7191.490
\hypertarget{c01:cue:3991}{}但在別的大家。
% cue 3992 7191.490 7192.670
\hypertarget{c01:cue:3992}{}別的場合。
% cue 3993 7192.670 7193.670
\hypertarget{c01:cue:3993}{}大家都讓我們看見。

\tc{1:59:53}
% cue 3994 7193.670 7194.730
\hypertarget{c01:cue:3994}{}我算是很幸運的。
% cue 3995 7194.730 7196.730
\hypertarget{c01:cue:3995}{}我從回來 1989 年\footnote{司法院所載大法官簡歷列湯德宗 1989 年 8 月至 2011 年 7 月任國立台灣大學國家發展研究所合聘教授，1989 年 8 月至 2016 年 7 月任私立東吳大學法學院兼任教授。\url{https://cons.judicial.gov.tw/judge.aspx?fid=9&id=10190&class_per_id=40024}。}。
% cue 3996 7196.730 7197.730
\hypertarget{c01:cue:3996}{}我在臺南。
% cue 3997 7197.730 7198.730
\hypertarget{c01:cue:3997}{}雖然不是在法律系。
% cue 3998 7198.730 7199.730
\hypertarget{c01:cue:3998}{}但是至少也是在法律院。
% cue 3999 7199.730 7200.730
\hypertarget{c01:cue:3999}{}不去入客。
% cue 4000 7200.730 7202.110
\hypertarget{c01:cue:4000}{}我一次有上課。
% cue 4001 7202.110 7204.110
\hypertarget{c01:cue:4001}{}東吳也是 1989 年就來了。
% cue 4002 7204.110 7207.120
\hypertarget{c01:cue:4002}{}東吳我剛來的時候。
% cue 4003 7207.120 7208.120
\hypertarget{c01:cue:4003}{}我是叫我還進法。

\section{實務界的升遷與畢業學校}

\tc{2:00:08}
% cue 4004 7208.120 7210.130
\hypertarget{c01:cue:4004}{}什麼問題。

\tc{2:00:10}
\speaker{學生}
% cue 4005 7210.130 7226.460
\hypertarget{c01:cue:4005}{}打壓。
% cue 4006 7226.460 7228.610
\hypertarget{c01:cue:4006}{}為什麼要打壓東吳。
% cue 4007 7228.610 7229.610
\hypertarget{c01:cue:4007}{}或是外國這樣。
% cue 4008 7229.610 7232.880
\hypertarget{c01:cue:4008}{}什麼東西不要有弱勢。
% cue 4009 7232.880 7233.880
\hypertarget{c01:cue:4009}{}就在實務界這樣。

\tc{2:00:33}
\speaker{老師}
% cue 4010 7233.880 7234.910
\hypertarget{c01:cue:4010}{}不好說法。
% cue 4011 7234.910 7235.910
\hypertarget{c01:cue:4011}{}我跟你講。
% cue 4012 7235.910 7237.070
\hypertarget{c01:cue:4012}{}實務界的事情。
% cue 4013 7237.070 7238.570
\hypertarget{c01:cue:4013}{}我不見得懂。
% cue 4014 7238.570 7240.570
\hypertarget{c01:cue:4014}{}因為我也沒有在這裡面碰過。

\tc{2:00:40}
% cue 4015 7240.570 7242.050
\hypertarget{c01:cue:4015}{}但是呢，
% cue 4016 7242.050 7243.330
\hypertarget{c01:cue:4016}{}簡單的講呢，
% cue 4017 7243.330 7245.130
\hypertarget{c01:cue:4017}{}臺灣的實務界。
% cue 4018 7245.130 7247.450
\hypertarget{c01:cue:4018}{}特別是指法官的權限。
% cue 4019 7247.450 7248.990
\hypertarget{c01:cue:4019}{}他們當然有。
% cue 4020 7248.990 7250.750
\hypertarget{c01:cue:4020}{}誰先來誰報告。
% cue 4021 7250.750 7252.940
\hypertarget{c01:cue:4021}{}簡單的講就是問你是。
% cue 4022 7252.940 7254.700
\hypertarget{c01:cue:4022}{}實務界訓。
% cue 4023 7254.700 7255.700
\hypertarget{c01:cue:4023}{}現在叫法官訓練所\footnote{法務部司法官訓練所，其組織條例於 2015 年 6 月 10 日廢止。2013 年 5 月 22 日制定公布、同年 7 月 1 日施行之法務部司法官學院組織法設法務部司法官學院，執行司法官考試錄取人員之培訓；法務部司法官學院司法官訓練規則第 5 條第 1 項規定，司法官學員每期員額定為 90 名。\url{https://mojlaw.moj.gov.tw/LawContent.aspx?LSID=FL009731}。}。
% cue 4024 7255.700 7258.710
\hypertarget{c01:cue:4024}{}打一期的業。
% cue 4025 7258.710 7260.710
\hypertarget{c01:cue:4025}{}那期別對他們兩。
% cue 4026 7260.710 7261.810
\hypertarget{c01:cue:4026}{}為什麼重要。
% cue 4027 7261.810 7262.810
\hypertarget{c01:cue:4027}{}因為我們現在升警長。
% cue 4028 7262.810 7263.870
\hypertarget{c01:cue:4028}{}就升到第 38 期。
% cue 4029 7263.870 7266.410
\hypertarget{c01:cue:4029}{}那你到了 38 期。
% cue 4030 7266.410 7269.100
\hypertarget{c01:cue:4030}{}沒有到。
% cue 4031 7269.100 7270.100
\hypertarget{c01:cue:4031}{}沒有輪不到你。
% cue 4032 7270.100 7271.480
\hypertarget{c01:cue:4032}{}38 期過了。
% cue 4033 7271.480 7272.480
\hypertarget{c01:cue:4033}{}你還沒有升上去。
% cue 4034 7272.480 7273.480
\hypertarget{c01:cue:4034}{}他們表示你不行。
% cue 4035 7273.480 7274.480
\hypertarget{c01:cue:4035}{}就沒有完。
% cue 4036 7274.480 7275.540
\hypertarget{c01:cue:4036}{}所以他們是論資排。
% cue 4037 7275.540 7276.540
\hypertarget{c01:cue:4037}{}被警方去關一下說。
% cue 4038 7276.540 7277.500
\hypertarget{c01:cue:4038}{}我們現在升將軍。
% cue 4039 7277.500 7278.500
\hypertarget{c01:cue:4039}{}要升到第幾期。
% cue 4040 7278.500 7279.590
\hypertarget{c01:cue:4040}{}那個情形是一樣的。
% cue 4041 7279.590 7280.590
\hypertarget{c01:cue:4041}{}這第一點。

\tc{2:01:20}
% cue 4042 7280.590 7281.680
\hypertarget{c01:cue:4042}{}第二點。
% cue 4043 7281.680 7284.410
\hypertarget{c01:cue:4043}{}從純粹的人數上來講。
% cue 4044 7284.410 7286.960
\hypertarget{c01:cue:4044}{}都一定不是最多的。
% cue 4045 7286.960 7288.780
\hypertarget{c01:cue:4045}{}臺大應最多。
% cue 4046 7288.780 7289.970
\hypertarget{c01:cue:4046}{}其實可能我猜想。
% cue 4047 7289.970 7290.970
\hypertarget{c01:cue:4047}{}可能政大跟北大。
% cue 4048 7290.970 7291.900
\hypertarget{c01:cue:4048}{}大概差不多。
% cue 4049 7291.900 7293.130
\hypertarget{c01:cue:4049}{}因為北大會考試。
% cue 4050 7293.130 7294.080
\hypertarget{c01:cue:4050}{}然後再加起來。
% cue 4051 7294.080 7295.600
\hypertarget{c01:cue:4051}{}會是東吳。
% cue 4052 7295.600 7296.600
\hypertarget{c01:cue:4052}{}東吳也不算少。
% cue 4053 7296.600 7297.630
\hypertarget{c01:cue:4053}{}但是不會是多的。
% cue 4054 7297.630 7300.010
\hypertarget{c01:cue:4054}{}所以你要從人數上來講。
% cue 4055 7300.010 7301.010
\hypertarget{c01:cue:4055}{}大家會講。
% cue 4056 7301.010 7303.090
\hypertarget{c01:cue:4056}{}關於是哪個學校畢業的。
% cue 4057 7303.090 7304.090
\hypertarget{c01:cue:4057}{}或者大家有種共同的。
% cue 4058 7304.090 7306.820
\hypertarget{c01:cue:4058}{}所以東吳應該不會佔平。
% cue 4059 7306.820 7308.010
\hypertarget{c01:cue:4059}{}但是我沒聽說。
% cue 4060 7308.010 7309.010
\hypertarget{c01:cue:4060}{}我沒聽說。
% cue 4061 7309.010 7310.010
\hypertarget{c01:cue:4061}{}我知道。
% cue 4062 7310.010 7312.010
\hypertarget{c01:cue:4062}{}有東吳會被打壓的。
% cue 4063 7312.010 7314.910
\hypertarget{c01:cue:4063}{}被特地欺負的。
% cue 4064 7314.910 7316.250
\hypertarget{c01:cue:4064}{}我想不出什麼流。
% cue 4065 7316.250 7317.750
\hypertarget{c01:cue:4065}{}我也不去。
% cue 4066 7317.750 7318.740
\hypertarget{c01:cue:4066}{}我也不去。
% cue 4067 7318.740 7319.770
\hypertarget{c01:cue:4067}{}有什麼流。
% cue 4068 7319.770 7320.770
\hypertarget{c01:cue:4068}{}東吳有哪一點。
% cue 4069 7320.770 7321.770
\hypertarget{c01:cue:4069}{}有人討厭嗎？
% cue 4070 7321.770 7323.830
\hypertarget{c01:cue:4070}{}不幹沒有。
% cue 4071 7323.830 7324.830
\hypertarget{c01:cue:4071}{}所以你只能單純是。
% cue 4072 7324.830 7326.310
\hypertarget{c01:cue:4072}{}其人數的問題。

\tc{2:02:06}
% cue 4073 7326.310 7327.720
\hypertarget{c01:cue:4073}{}至於將來的發展。
% cue 4074 7327.720 7330.060
\hypertarget{c01:cue:4074}{}完全要看個人。
% cue 4075 7330.060 7331.160
\hypertarget{c01:cue:4075}{}那這個就每個人。
% cue 4076 7331.160 7332.690
\hypertarget{c01:cue:4076}{}情況不一樣。
% cue 4077 7332.690 7333.760
\hypertarget{c01:cue:4077}{}有的是臺大的。
% cue 4078 7333.760 7336.170
\hypertarget{c01:cue:4078}{}但是他似乎與順。
% cue 4079 7336.170 7337.170
\hypertarget{c01:cue:4079}{}有的是別的學校。
% cue 4080 7337.170 7338.120
\hypertarget{c01:cue:4080}{}但他似乎與順。
% cue 4081 7338.120 7339.230
\hypertarget{c01:cue:4081}{}你也可以去看一看。
% cue 4082 7339.230 7340.230
\hypertarget{c01:cue:4082}{}現在這些當大官的人。
% cue 4083 7340.230 7341.230
\hypertarget{c01:cue:4083}{}他們背景是什麼？
% cue 4084 7341.230 7342.800
\hypertarget{c01:cue:4084}{}大體上講。
% cue 4085 7342.800 7344.950
\hypertarget{c01:cue:4085}{}我認為臺大還是在大體上。
% cue 4086 7344.950 7346.470
\hypertarget{c01:cue:4086}{}但是這個跟，
% cue 4087 7346.470 7351.320
\hypertarget{c01:cue:4087}{}他是不是學校沒有必然關係。

\tc{2:02:31}
% cue 4088 7351.320 7353.470
\hypertarget{c01:cue:4088}{}好。
% cue 4089 7353.470 7354.820
\hypertarget{c01:cue:4089}{}還有一期。
% cue 4090 7354.820 7358.430
\hypertarget{c01:cue:4090}{}沒有嗎？
% cue 4091 7358.430 7358.530
\hypertarget{c01:cue:4091}{}沒有。
% cue 4092 7358.430 7359.660
\hypertarget{c01:cue:4092}{}就感謝你們。
% cue 4093 7359.660 7360.660
\hypertarget{c01:cue:4093}{}要放老師走。
% cue 4094 7360.660 7363.680
\hypertarget{c01:cue:4094}{}你們都沒有人要問說。
% cue 4095 7363.680 7365.680
\hypertarget{c01:cue:4095}{}老師為什麼去哪裡做客座。
% cue 4096 7365.680 7367.680
\hypertarget{c01:cue:4096}{}那些地方是什麼經驗。
% cue 4097 7367.680 7369.680
\hypertarget{c01:cue:4097}{}都沒有人對這有興趣嗎？
% cue 4098 7369.680 7371.610
\hypertarget{c01:cue:4098}{}還可以講。
% cue 4099 7371.610 7373.050
\hypertarget{c01:cue:4099}{}那有人去走嗎？
% cue 4100 7373.050 7374.050
\hypertarget{c01:cue:4100}{}沒有。
% cue 4101 7374.050 7375.360
\hypertarget{c01:cue:4101}{}還有嗎？
% cue 4102 7375.360 7377.360
\hypertarget{c01:cue:4102}{}有話我就讓你們問最後一集了。
% cue 4103 7377.360 7378.380
\hypertarget{c01:cue:4103}{}沒有話我就。
% cue 4104 7378.380 7379.370
\hypertarget{c01:cue:4104}{}原證到。
% cue 4105 7379.370 7380.320
\hypertarget{c01:cue:4105}{}我就不講。
% cue 4106 7380.320 7382.230
\hypertarget{c01:cue:4106}{}可以。

\tc{2:03:02}
\speaker{學生}
% cue 4107 7382.230 7383.230
\hypertarget{c01:cue:4107}{}那是您說您在大學。
% cue 4108 7383.230 7384.230
\hypertarget{c01:cue:4108}{}就是在臺大的時候。
% cue 4109 7384.230 7385.230
\hypertarget{c01:cue:4109}{}寫說是。
% cue 4110 7385.230 7386.230
\hypertarget{c01:cue:4110}{}是為我們。
% cue 4111 7386.230 7387.290
\hypertarget{c01:cue:4111}{}你會覺得。
% cue 4112 7387.290 7393.150
\hypertarget{c01:cue:4112}{}就是您會覺得。
% cue 4113 7393.150 7394.400
\hypertarget{c01:cue:4113}{}在口頭文。
% cue 4114 7394.400 7395.400
\hypertarget{c01:cue:4114}{}這些語言的方式。
% cue 4115 7395.400 7398.400
\hypertarget{c01:cue:4115}{}就是使用別的國家的語言。

\tc{2:03:18}
\speaker{老師}
% cue 4116 7398.400 7401.000
\hypertarget{c01:cue:4116}{}當然有困難。
% cue 4117 7401.000 7402.670
\hypertarget{c01:cue:4117}{}德文多難。
% cue 4118 7402.670 7404.300
\hypertarget{c01:cue:4118}{}對。
% cue 4119 7404.300 7405.300
\hypertarget{c01:cue:4119}{}但是因為當時。
% cue 4120 7405.300 7406.300
\hypertarget{c01:cue:4120}{}臺大的法律研究所。
% cue 4121 7406.300 7407.360
\hypertarget{c01:cue:4121}{}基本上就是，
% cue 4122 7407.360 7408.360
\hypertarget{c01:cue:4122}{}德文基修班。
% cue 4123 7408.360 7411.440
\hypertarget{c01:cue:4123}{}法律研究所進去。
% cue 4124 7411.440 7412.440
\hypertarget{c01:cue:4124}{}所有人。
% cue 4125 7412.440 7413.440
\hypertarget{c01:cue:4125}{}一個禮拜。
% cue 4126 7413.440 7414.440
\hypertarget{c01:cue:4126}{}八個小時的德文。
% cue 4127 7414.440 7415.440
\hypertarget{c01:cue:4127}{}學分數 0。
% cue 4128 7415.440 7418.870
\hypertarget{c01:cue:4128}{}就這樣子。
% cue 4129 7418.870 7419.970
\hypertarget{c01:cue:4129}{}因為他認為。
% cue 4130 7419.970 7421.410
\hypertarget{c01:cue:4130}{}德國的法律。
% cue 4131 7421.410 7423.410
\hypertarget{c01:cue:4131}{}影響臺灣重大。
% cue 4132 7423.410 7425.790
\hypertarget{c01:cue:4132}{}那很多老師都是德國回海南。
% cue 4133 7425.790 7426.800
\hypertarget{c01:cue:4133}{}所以很希望。
% cue 4134 7426.800 7427.980
\hypertarget{c01:cue:4134}{}來亞地學生。
% cue 4135 7427.980 7428.980
\hypertarget{c01:cue:4135}{}教練去德國。
% cue 4136 7428.980 7430.360
\hypertarget{c01:cue:4136}{}所以他就變成。

\tc{2:03:50}
% cue 4137 7430.360 7432.360
\hypertarget{c01:cue:4137}{}好。
% cue 4138 7432.360 7433.360
\hypertarget{c01:cue:4138}{}好啦。
% cue 4139 7433.360 7434.360
\hypertarget{c01:cue:4139}{}我不能講。
% cue 4140 7434.360 7435.360
\hypertarget{c01:cue:4140}{}那要聽一次。
% cue 4141 7435.360 7436.320
\hypertarget{c01:cue:4141}{}希望他在庭檔面。
% cue 4142 7436.320 7436.420
\hypertarget{c01:cue:4142}{}來吧。
% cue 4143 7436.320 7437.320
\hypertarget{c01:cue:4143}{}請開一個門吧。
% cue 4144 7437.320 7438.320
\hypertarget{c01:cue:4144}{}那麼嚴格。
% cue 4145 7438.320 7439.320
\hypertarget{c01:cue:4145}{}我不是教師。
% cue 4146 7439.320 7440.320
\hypertarget{c01:cue:4146}{}你都在教師。
% cue 4147 7440.320 7441.320
\hypertarget{c01:cue:4147}{}你都在教師。
% cue 4148 7441.320 7442.410
\hypertarget{c01:cue:4148}{}你都在教師。
% cue 4149 7442.410 7443.740
\hypertarget{c01:cue:4149}{}你都在教師。
% cue 4150 7443.740 7456.000
\hypertarget{c01:cue:4150}{}我都在教師。
% cue 4151 7456.000 7456.990
\hypertarget{c01:cue:4151}{}我都在教師。
% cue 4152 7456.990 7458.110
\hypertarget{c01:cue:4152}{}我的天。
% cue 4153 7458.110 7462.140
\hypertarget{c01:cue:4153}{}你從哪裡。
% cue 4154 7462.140 7463.140
\hypertarget{c01:cue:4154}{}來教師。
% cue 4155 7463.140 7463.550
\hypertarget{c01:cue:4155}{}我都在教師。

% 自 Notizen/W08_1102/ 各表由 Skripte/gen_tex.py 產生；改字改表後重產，不手改本檔。
\chapter{釋字第 387 號、第 419 號與行政院院長之總辭（2026-11-02）}

2026 年 11 月 2 日下午課堂錄音（約 125 分鐘）之逐字稿。左邊界為段首時間碼；未標說話人者為湯德宗老師。


\section{購書登記與人數}

\tc{0:00:00}
% cue 1 0.480 14.690
\hypertarget{c08:cue:1}{}慢慢地，按照裡面的案例，一次指定一個、兩個，慢慢看看能不能幫助大家，多看能大法官的起司。但是前提是要有 20 個人，班長，今天我都是二人，登記要買書。
% cue 2 14.690 22.560
\hypertarget{c08:cue:2}{}那我請問問，你那邊登記的看得到嗎？看得出來？

\tc{0:00:22}
\speaker{學生}
% cue 3 22.560 31.380
\hypertarget{c08:cue:3}{}登記的跟學長那邊的表單登記的不太一樣，因為我問他，看他說我真的不會看到，沒有辦法看到表單。

\tc{0:00:31}
\speaker{老師}
% cue 4 31.380 35.140
\hypertarget{c08:cue:4}{}好，那你等一下我問他好了，你等一下我問他好了。
% cue 5 35.140 43.450
\hypertarget{c08:cue:5}{}因為我們班上要有 20 人以上，我才會去開課子，所以你們現在到底有多少人，讓你們自己要，要有了解啊。

\tc{0:00:43}
% cue 6 43.450 73.940
\hypertarget{c08:cue:6}{}或者我們問一個簡單的問題，現在坐在這裡的人，如果你有馬，請你去掃看看，班長幫我數一下，班長站起來數一下，站起來，沒有什麼好科技，數到你就要讓他手放下，這樣子最簡單，請放下，13，所以表示我們班上，如果大家決定要開課子，還插氣梗，意思就是這麼簡單，那這個你們自己去決定。

\tc{0:01:13}
% cue 7 73.940 105.140
\hypertarget{c08:cue:7}{}這件事情也增加，也增加老師的工作量，我不必強求，我不必這麼熱心要上這麼久的課，但是是你們自己要決定的，所以班長有個理解，就是我們班上有 45 人，現在只有 13 人加 7 人，如果湊得到 7 人，當然是在 11 月 18 號以前湊到 7 人，因為 11 月 18 號以後就不會是這個價格了，我就沒有權力在控制這個價格，就得按照出版社他們自己公告，這樣他們去處理。那因為我的新書班，
% cue 8 105.140 137.390
\hypertarget{c08:cue:8}{}所以我把價格當下來，所以是這樣子，那你以後當然買的是可以的，什麼時候湊到 20 人，就要來討論價格的問題，好不好？這樣子發布的兩個公告，有沒有內容不清楚的地方？沒有嗎？好，那我們今天要處理的案件非常多，我是希望我們先從今天的案件先講，然後有剩下的時間再來上禮拜，好不好？那現在，今天的兩個案子非常重要。

\section{釋字第 387 號解釋與行政院院長之總辭}

\tc{0:02:17}
% cue 9 137.390 172.220
\hypertarget{c08:cue:9}{}考試是一定會考，而且，對於了解憲法結構和憲法解釋關聯，中間有關聯，這兩個案子是關鍵。好，那 387 號、解釋內組在哪裡？387 的，誒，班長，為什麼沒有人決定？一個，就你一個人，只有你一個人嗎？那 419 也是你嗎？419 的戰法也是你，你們一邊一個人，好，那開始吧，387。你先跟老師解釋一下 387 號解釋的爭點。

\tc{0:02:52}
% cue 10 172.220 220.530
\hypertarget{c08:cue:10}{}387 號主要的爭執點在哪裡？用你的白話講，不要看文字。大聲一點，你叫什麼名字？法三級，還是法二級？法四 B 啊，法四 B 陳延新，還是王威廷？你是陳延新？好，陳延新請說，用你的白話講給我們聽。387 號解釋在爭什麼？大聲一點。請問第二個問題，這個事情，我們在憲法有沒有明文？如果看完了，你應該能看，
% cue 11 220.530 256.550
\hypertarget{c08:cue:11}{}你應該能回答我啊，有明文規定，除了立法院改選之後，第一次集會以前，行政院長要總辭。有這樣寫嗎？其他同學能回答嗎？他現在不能回答，是他準備的問題。你們能回答是你們比他厲害，或者是你們準備了很好的問題。所以你們如果想答答對了，也是要有加分的。好，沒有嗎？簡單的，好，你請說，有沒有？沒有。好，那所以這個解釋有沒有重要性？

\tc{0:04:16}
% cue 12 256.550 300.280
\hypertarget{c08:cue:12}{}憲法沒有規定。憲法沒有規定。所以才需要解釋。簡單的講是這樣子。那麼完全沒有規定嗎？還是有一個相關的規定。跟本號解釋最相關的規定在哪裡？陳年新，永遠都是你先答。你答不出來，我才讓別人去答。講話怎麼沒有聲音呢？你都沒有看條文呢？好，其他同學，剛才那位同學，你能答對就開始有加分了。答一條跟這個有關的？沒有。來，那好，你說 57 條不太對，
% cue 13 300.280 334.710
\hypertarget{c08:cue:13}{}也不是完全不對，但是不太對。用英文來講呢，It's not the most pertinent。不是最相關條文。最相關條文，請看 55 條第一項。這個表示我們這個班上很危險。如果老師考試教，一定會當掉一半以上。因為已經這個條文講過五遍了。忘記了嗎？55 條第一項說，行政院長是總統提名，經立法院同意。但你認為立法院什麼時候要同意？我們現在立法院以前是 A 黨控制，A 黨佔多數的。

\tc{0:05:34}
% cue 14 334.710 369.840
\hypertarget{c08:cue:14}{}描述民進黨。現在改選以後，國民黨佔多數了。民進黨當時同意的行政院長、國民黨佔多數立法院還要繼續維持嗎？我這樣講得聽懂嗎？每一個解釋一定要從我現在的問題，這種方式開始理解。用背的是沒有用的，因為如果老師考試的時候把前提改了，會不會那個死文字有什麼意義呢？這個解釋是在講什麼東西，那這個為什麼會重要？從這裡開始理解。好，來，剛才執言欣請你大聲念解釋文。

\tc{0:06:09}
% cue 15 369.840 421.130
\hypertarget{c08:cue:15}{}這句話重不重要。如果有一個人要經過立法院同意才能上任。上司機講過很多遍了，從第一次上課講到現在。行政對立法負責的方法第一種。我沒有同意，你不能到我家來工作。你是向我負責的，我是老闆，你是我雇的人，你要來我家，你要我同意。這是行政向立法負責的第一個表現。好，繼續，跳過下一句，跳過下一句，句號是句號是一句，下一句。因為我們今天沒有，
% cue 16 421.130 473.440
\hypertarget{c08:cue:16}{}沒有時間，所以我必須要用這麼快的方法改。憲法第 54 條，跳過下一句。好，出現了，這是剛才念的男同學回答的，規定在 57 條，所以第 57 條規定剩下，有吧？好，下面你要注意了，下面這一句，慢慢念只有一句，大聲一點，聽懂了嗎？你必須經過我同意而任命。好，繼續，又是對我負責，負什麼責，負政治責任，我不要你就滾蛋，沒有理由，繼續。

\tc{0:07:53}
% cue 17 473.440 524.740
\hypertarget{c08:cue:17}{}這八個字圈起來，大法官 387 號解釋的重點，只有民意政治跟責任政治，這是他指更顯示的最基礎。誰代表民意，新改選的立法院代表民意，誰向誰負責，行政院向立法院負責。那我都有了新民意，你行政院長還可以繼續在那裡混嗎？這樣有聽懂嗎？好，繼續。貝老師，這樣一講就打斷了。到這裡裡面的這麼辛苦，到底是不是臺灣人啊？你初小是怎麼長大的啊？我真的不知道該怎麼教你們啊？這些事情是基本功嗎？

\tc{0:08:44}
% cue 18 524.740 560.100
\hypertarget{c08:cue:18}{}你不需要老師講。好，這個句子連合到剛才，連結到剛才老師講的不舒服，要一項。你可以聽懂為什麼有況解釋要解釋嗎？為什麼本屆大法官要做解釋了嗎？你能不能回答？他剛才已經把最關鍵的句子講了。大法官解釋的最基礎就是民意政治跟責任政治。誰代表民意？新改選的立法院。誰向誰負責？行政院向立法院負責。我都有新民意了，那你要到我向我負責。好，所以大法官說，那你當然，
% cue 19 560.100 595.740
\hypertarget{c08:cue:19}{}當然，在改選之後，第一次有新的立委集會以前，行政院長要總辭。那我現在說，這個結論跟剛才 55 條第一項的規定，怎麼產生連結？這個話是大法官應該講，他沒講，老師幫他講完了，今天這個解釋387就結束了。來，所有同學請回答。55 條剛才第一項後已經問過你們了，是什麼啦？那現在這個跟 55 條第一項有什麼關係？大聲一點有回答我，這個對你們有這麼難嗎？

\tc{0:09:57}
% cue 20 597.120 640.620
\hypertarget{c08:cue:20}{}應該不會吧。好，請同學請舉手，舉手的請說，因為所以呢，跟 55 條第一項有什麼關係？你現在要聽清楚老師的問題，大法官說應總辭，老師的問題說應總辭跟 55 條第一項有什麼關係？你應該這麼回答我，如果你不總辭，55 條第一項是不是廢話？你如果沒有總辭，我怎麼同意吧？一個人站在位置上，站在茅坑不拉屎，那個人不走，那我怎麼同意呢？

\tc{0:10:40}
% cue 21 640.620 677.160
\hypertarget{c08:cue:21}{}所以再講白一點，我們這個新科立委白選了，選了什麼？回到原地，我們還這個時候沒有修憲，97 還沒有到，是內閣制，內閣制是不是立法院決定誰能當行政院長？所以老師講過 N 遍了，總統都提名是 gay 的，那是個理性的權利，他必須反應立法院的生態，誰在立法院佔多數，必須同意提名他的黨魁，不然立法院是不會同意的，立法院是笨蛋嗎？

\tc{0:11:17}
% cue 22 677.160 713.850
\hypertarget{c08:cue:22}{}他是烏合之眾嗎？怎麼選贏的呢？是啊，所以我這樣講有聽懂了嗎？我剛才的問題現在可以連結起來了嗎？為什麼 55 條第一項跟本件解釋密切相關？55 條第一項只講了表面，說行政院長要經過立法院同意才能上任，但是原來一個行政院長他就不下來，他不下來，我怎麼逼他下來？現實的問題。好，那我再問，這就是同樣的這項題往下，這個正案子要進鎖，如果回到我們憲法，
% cue 23 713.850 750.210
\hypertarget{c08:cue:23}{}基本上不是沒有規定的。請問你，現在有一個立法院新改選的立委員，一群。這個是這個案子發生在第二屆，就是剛剛全面改選立法院以後，第一次改選的立法院，所以是第二屆立法委員。那我們現在假設到我們當時的原廠記。這個第二屆的立委就通到一個賴平的行政院長，他就不下來，他就不下來。好，那我現在請問你，按照憲法，你怎麼把它逼下來？請回答我，找姚文。上次呢，在這個這一代也是一個女同學，
% cue 24 750.210 785.290
\hypertarget{c08:cue:24}{}在這個區域有回答過我問題，類似這樣。我問你說，內閣制怎麼表現行政向立法負責。你說上任，他經過我同意。然後呢，如果我不喜歡你，可不可以把你趕走。為什麼我們班上對這個問題有這麼難回答？我們已經講來、復去、翻過來、講過去。到現在已經四個禮拜、五個禮拜，中間固然有休息了。但是我不認為你們休息的時候是一個理由，可以把過去講的忘掉。

\tc{0:13:05}
% cue 25 785.290 830.100
\hypertarget{c08:cue:25}{}來，請舉手回答哦，要看第幾條。剛才那個男同學，你說要看第幾條。那你去找五十七，五十七在哪裡。現在我們可不可以自己立法委員發動，把他幹掉。這五十七條都能趕走行政院長，我的問題再白眼。請回答。所有的人都可以設想並且回答。你現在是討厭這個行政院長，你要把他趕走。你怎麼趕走？你如果趕不走，你就不能說他向你負責，因為你趕不走了啊。第二款，第二款講什麼？你大聲一點說，念出來。

\tc{0:13:50}
% cue 26 830.100 866.730
\hypertarget{c08:cue:26}{}你念相關的文字，誰發動？立法院還是行政院？主詞是誰？這個上課都講過了，對我來講一看，我已經跟你們所有分辨了，你們必須要相信我，你們沒有那麼聰明。你們都覺得老師講過了，你就會記住，記住了嗎？我的條文是不是都這樣講？看到條文要先找主詞，現在因為你們已經是我自己條都 Pass 給你們了，我實際上主詞是立法院，對不對？第二款，第二款，當然就是立法院，然後繼續。

\tc{0:14:26}
% cue 27 866.730 902.530
\hypertarget{c08:cue:27}{}那你就隨便找個理由，行政院有很多政策，你不要這麼認真說這條沒有寫，沒有寫說不信任，也沒有寫說我們行政院長上任要怎樣，沒有，但是你是立委員，你不會借路，你不會借斷，你會不會借斷聲勢？好，你先要行政院，好，比如說現在行政院長推動斑斑斑學生，聲聲喝些奶，逃用這個政治，我們就推動把它翻掉，可不可以？就翻掉了，那如果翻掉了，
% cue 28 902.530 936.020
\hypertarget{c08:cue:28}{}那如果翻掉了，你現在看看 57 條第二款，難不難？你要有三分之二的例如，維持原決議，斑斑喝些奶，斑斑不準喝些奶，先決議斑斑不準喝些奶，因為現在的政策是斑斑喝些奶，我們先通過一個決議，說廢止斑斑喝些奶，然後它來覆議，覆議以後，我們有三分之二維持原決議，還是斑斑不準喝些奶，你有二分之一通過原決議，你有沒有辦法有三分之二維持原決議，很緊張。

\tc{0:15:36}
% cue 29 936.020 973.290
\hypertarget{c08:cue:29}{}如果有三分之二維持原決議，你擋高的話立法院永遠到達三分之二維持原決議，目前這個漫漫。所以趕走政院長難不難。這兩個解釋都講，這兩個前提都講完了。所以你們現在看懂，到訪問 387 號解釋都不重要。如果沒有 387 號解釋，我們這種引號修正式的內閣制，到底還算不算內閣制，它真的成大問題。上任的時候，我要同意，他沒有我同意，他就來了。所以他把他趕走，我要三分之二，偏偏沒有三分之二，我有二分之一，多一點。

\tc{0:16:13}
% cue 30 973.290 1006.820
\hypertarget{c08:cue:30}{}我沒有三分之二，又趕不走。那我還叫內閣制嗎？他還叫我負責嗎？那這麼不要臉的人，我怎麼趕不走呢？所以，384 號解釋，不是 384,387 號解釋，非常重要。他確立了，我們作為修正式內閣制，55 條第一項是不容許迴避的。再講一遍，387 號解釋確認，在當時修正式內閣下，修正式內閣制之下，
% cue 31 1006.820 1045.880
\hypertarget{c08:cue:31}{}你絕不可容許迴避 55 條第一項的要件。你上任必須我同意，那我現在就幫你拉一把，舊人還沒走，沒關係，我就告訴他什麼時候該滾蛋。對這個解釋重不重要？來，剛才，陳彥新，本號解釋有沒有反對意見？有沒有不同意見？有還是沒有？誰？好，他的意思是什麼？大意是什麼？他為什麼反對？他反對什麼？你看懂嗎？還是沒有看懂？你如果沒有看懂，是這個騙子，在任何地方，
% cue 32 1045.880 1085.450
\hypertarget{c08:cue:32}{}一個地方都找得到。你到 YouTube 上去花點時間，你都找到了。好，吳庚說什麼？他是政治問題，對不對？所以本案不必解釋，應該解釋。你回來對照一下。如果照吳庚的見解，本案我們不解釋。這些人不走，你也沒辦法。聽懂？沒有聽懂？所以好在多數大法官沒有贊成吳庚的見解。至於吳庚為什麼要說這一件我們不要解釋，不要處理？知道。這個我沒有辦法替他回答。哈哈，這聽得懂嗎？

\section{政治問題理論與 Baker v. Carr}

\tc{0:18:05}
% cue 33 1085.450 1116.790
\hypertarget{c08:cue:33}{}他為什麼這樣講我沒有辦法替他回答為什麼？但是我只要講這個案子，說本件是政治問題，是不通的。好，來，陳妍心，你有沒有辦法回答我什麼叫政治問題？在他的不同意見書裡面他點了一段話，這一段話跟老師的翻譯有點不一樣，但是是可以理解的。他說政治問題在美國的理解是什麼？美國有一個案子叫 Baker vs. Carr，所有的人都找到不同意見，拿出來看，
% cue 34 1116.790 1155.370
\hypertarget{c08:cue:34}{}在小註幾，陳元欣，小註幾，不知道到底是有看還是沒有看，我真的不知道你要怎麼辦。其他人一律看到小註五，Baker vs. Carr，老師說過了，你不需要省錢省到這個樣子，今天是你報告，老師問你的問題，你一向著找，所以你必須印出紙本。上課上到現在第七個禮拜，你們都不聽我的話。我問你，那你有本事你就說，小助五啊，我腦袋好，我過目不忘。

\tc{0:19:15}
% cue 35 1155.370 1189.910
\hypertarget{c08:cue:35}{}問你，那你有辦法做到這樣嗎？那不能，我現在問你，你就不能回答，我腦袋卡住了嗎？小註五，其他同學找到了嗎？這是 Baker vs. Carr，他的原文如下。那老師在世界，在這個權力分立新論，抱歉老師書還沒有改好，所以還沒有出現，我也不想這個到底要怎麼解決，現在我事情越來越多了。那這一段文字，老師有一個精準的翻譯，這個翻譯是不好的翻譯，但是我講意思就好，看到美國聯邦最高法院 Baker vs. Carr 的這個見解，
% cue 36 1189.910 1222.030
\hypertarget{c08:cue:36}{}不是所有人來問我們大法官，來申請我們大法官啟示，我們都可以做啟示。我們需要考慮這個問題的本質，是我們能啟示的嗎？簡單的講，按照 Baker vs. Carr，如果憲法以及明文規定，這件事情應該由某個政治部門，在美國講政治部門，就是民權的，或者總統，就是國會，兩院。如果這個事情說是參議院要同意，那權力已經給了參議院了，才問大法官為什麼呢？我就不能講話了。

\tc{0:20:22}
% cue 37 1222.030 1255.520
\hypertarget{c08:cue:37}{}那說誰誰誰可以做什麼決定，那就不是我們做決定，這個很容易理解嘛。因為我們的白話講，憲法有明文規定，應該要把機關做的，那就不是我們做的，這是第一點。第二件事情，用我的白話講比較清楚，請問這個問題在憲法上有沒有標準答案？如果這個問題沒有標準答案，那我們大法官是憑什麼憲法？大家如果說我有解釋憲法，當然是依據憲法解釋憲法，但是憲法沒給答案，大法官怎麼解釋。

\tc{0:20:55}
% cue 38 1255.520 1292.710
\hypertarget{c08:cue:38}{}那我們舉個例子，舉個例子，好聽清楚，在 1979 年 1 月 1 號以前，美國承認哪一個國家代表全中國，講話中華民國。1979 年以後改了，他承認他說中華民國國為全中國唯一合法的政府，所以他代表中國。好，那麼請問當時就有議員叫高華德，他來到，他說這個不對，這個呢，沒有經過我們條約批准的程序，沒有經過參議院同意，總統不可以跟別的國家建交。

\tc{0:21:32}
% cue 39 1292.710 1326.460
\hypertarget{c08:cue:39}{}這個要經過我們。好，那最高法院的見解，這是政治問題。我們憲法哪裡有規定，要承認哪一個國家，作為某個國家的唯一合法制？有寫嗎？這件事情，憲法上沒有答案，當然是依照當時的政治形勢而為判斷。你要問美國人來後，他說 1979 年以前，我認為中華民國政府，承認中華民國政府最符合美國國家的意義。1979 以後，承認中華人民共和國，
% cue 40 1326.460 1358.660
\hypertarget{c08:cue:40}{}最符合美國國家利益。美國國家利益會隨著時間的經過不同的時空環境而改變。憲法沒有給答案，沒有說一定要承認哪個國家。是不是？所以這種事情，叫我們聯邦、最高法院怎麼決定？難道隨著法院要跳出來說，總統，你做錯了決定？那總統會說，我為什麼錯？憲法哪條說我只能承認誰？不能承認誰？沒有啊！所以最高法院也可以這麼解釋憲法。好啦，回到 Baker vs. Carr，什麼情況下？大法官應該拒絕。

\tc{0:22:38}
% cue 41 1358.660 1389.930
\hypertarget{c08:cue:41}{}第一，這個問題在憲法上已經有明文規定，應該由某個政治部門，不是國會就是總統來做決定。憲法明文寫了，那就我不能再審判，再做決定了。第二件事情就是按照憲法，它沒有一個準據，憲法沒有給答案，那我解釋憲法，憲法又沒有答案，我憑什麼解釋說該怎樣？這種問題就是政治問題，並不是聽清楚，並不是說這個問題本身帶有政治性。

\tc{0:23:09}
% cue 42 1389.930 1425.790
\hypertarget{c08:cue:42}{}你們看 387 有沒有政治性，419 有沒有政治性，哪一個到大法官來的解釋沒有政治性，只要牽涉到權力分立的都有政治性，都有黨派之爭，不可能沒有，只是現在比以前更明顯而已而已而已，這樣就聽懂了嗎？所以不要再用沒有學憲法的事故回答，政治問題，就是這個問題具有高度政治性，所以大法官不能解釋，如果這樣大法官的事故都不能解釋，不是的，按照 Baker vs. Carr原文一部分很差，
% cue 43 1425.790 1457.460
\hypertarget{c08:cue:43}{}很短，你們拿到書再去看老師的解釋，但是我用最簡單的話講，憲法有明文規定，由某個機關做決定，你不能再講話。第二個就是按照憲法，憲法沒有標準答案，憲法沒有提供裁判之準據，那我解釋憲法，憲法都沒有答案，我怎麼解釋啊，我再怎麼解釋就是我的意見，不是憲法的意見，這樣就不行。希望這個標準今天順便學會。好，陳彥新，你有問題嗎？三八七，老師，對於你現在要報告的事情，
% cue 44 1457.460 1500.250
\hypertarget{c08:cue:44}{}老師這麼簡短的說明，有沒有打中，有沒有擊中要害，你自己有沒有問題要問。其他同學，三八七號，聽不懂的，可以有什麼問題請舉手。沒有哦。好，那我問你沒有問題。請問三八七號解釋，今天還有沒有適用。有人搖頭，那搖頭的人告訴我為什麼沒有適用。用你的白話講就可以了。來，請說剛才舉手的是誰？是你嗎？好，你講。誰剛才舉手？好，你說大聲一點。

\tc{0:25:00}
% cue 45 1500.250 1538.300
\hypertarget{c08:cue:45}{}有沒有適用今天立法院改選，以後是不是行政院要總辭？不用？為什麼不用？你這樣子講也可以，但是我建議你講最簡單的方法，你就說老師，因為憲法結構已經變了，已經不是內閣時代，已經不是行政院向立法院負責的時代。現在行政院長是總統任、總統免，誰是老闆，誰向誰負責。這個都聽不懂的話，那就不能當法律系的學生了。你們不要被，
% cue 46 1538.300 1554.320
\hypertarget{c08:cue:46}{}被立法委員騙。尤其不要被李登輝騙了。李登輝是會用漢字騙漢的。因為文字還是寫的。請問我怎麼？他怎麼對我負責？他上任，不要經過我同意。55 條第一項已經停止使用。

\tc{0:25:54}
% cue 47 1554.320 1572.680
\hypertarget{c08:cue:47}{}其次又是你任，又是你免。我立法院人說話的四個都沒有。我怎麼敢怎麼辦？你跟我講啊。我已經講過 N 遍了。這樣清楚了嗎？所以在這種情形下，實際上也不待我講。在之前阿扁政府的時候，
% cue 48 1572.680 1608.890
\hypertarget{c08:cue:48}{}阿扁就有批示，說我們現在 1997 修憲以後，行政院長不需要再隨著立法院改選而總辭。他沒有像我講的這麼白，因為他不是黨憲法。他有他的政治權謀算計，那我不管。好，那我問你，如果你認為還是要逼行政院長總辭，有意義嗎？沒有吧。只要總統不壞，我喜歡卓榮泰可以吧？卓榮泰你辭一下，明天我再認定你，中間差兩小時可不可以？

\tc{0:26:48}
% cue 49 1608.890 1643.750
\hypertarget{c08:cue:49}{}現在是 11.59 分一辭，12.01 都會把你便定回來，這樣可不可以？可以不可以？我認我免都是我說了算，怕什麼？勇敢向前行，別嚇啊！所以這個條文有意義嗎？大法官沒有取消 387，也沒有說 387 號解釋失效，沒有。但是這就是臺灣的問題，我們做總統的人只管他的政治權謀，老百姓是最可憐的笨蛋，這個來騙我一套，那個來又騙我一套，那像江湖某人這樣子講的，幾乎沒有。

\tc{0:27:23}
% cue 50 1643.750 1680.040
\hypertarget{c08:cue:50}{}因為我這樣子左右都不好，你喜歡一定會討厭我，你喜歡陳水扁也討厭我，我終於來只碰到你不喜歡。所以這樣複雜的、頻繁的修憲，而且政府體制從左邊到右邊，覺得我們人民要不要聰明一點，不要不要看得懂現在寫得不瞎的，要不然我們怎麼當主人，我們很容易被封鎖啊，好啦，這是 387，到這裡講完哦，我就不再講，好，來，419，下一位警局所，請問你，419 叫什麼名字？

\section{釋字第 419 號解釋與副總統兼任行政院院長}

\tc{0:28:00}
% cue 51 1680.040 1722.930
\hypertarget{c08:cue:51}{}我看到了，來，現在 387 結束，我們眼睛飄到 419 來，老師今年要感濕接觸，我都要把它講完，好，來，419，現在請問你，在爭什麼，主要爭點，或者你可以簡單的講，副總統可不可以兼任行政院長，這樣子可以吧，大家都記得懂嗎？我們講，憲法課，所以就是永遠危險，那我現在要問，還原級的司令，他現在講副總統兼任行政院長，
% cue 52 1722.930 1735.180
\hypertarget{c08:cue:52}{}得否兼任行政院長，你認為這樣的語句描述是否是精確的，或者你認為更好處理辦法。

\tc{0:28:55}
% cue 53 1735.180 1758.220
\hypertarget{c08:cue:53}{}郭政凡，你五年級囉，你今天上的老師的課算是你性命，以後你這輩子再要碰到像老師這樣的人，他會慢慢找，應該沒有吧，因為你不在東吳，你就碰不到這樣的人。我相信他老師上課不會這樣的人，一定當好我的心思，好好的念完下課，我管你懂不懂，不懂課會什麼關係。來，請回答。

\tc{0:29:18}
% cue 54 1758.220 1805.740
\hypertarget{c08:cue:54}{}這樣子的描述，說副總統可否兼任行政院院長。這樣的描述精確嗎。或者你有更好的，你只是覺得他有沒有符合什麼。我問你，你有沒有想過，這個句子應該倒過來寫。行政院長可不可以兼任副總統。這樣子講比較對。你點頭還是搖頭？點頭。為什麼你點頭？完全答對。老師，老師只要問這個問題，我就可以知道你有沒有看過。來，其他同學可以幫忙補充。

\tc{0:30:05}
% cue 55 1805.740 1844.730
\hypertarget{c08:cue:55}{}為什麼老師講這案子。與其說副總統得否兼任行政院院長。不如說行政院院長得否兼任副總統。請回答為什麼我這樣講。沒有一個人能回答。這表示你們在上課之前沒有一個人看過老師的視頻。也不要緊。沒有一個人認真的去準備今天的報課。你就完全依賴唐勝涵。他如果把你賣掉誤導了。然後湯老師又沒有發現。那你就跟著誤導了。

\tc{0:30:44}
% cue 56 1844.730 1883.570
\hypertarget{c08:cue:56}{}你就跟著背誤導。來，本案的事實，請還原一下。本案裡面牽涉的行政院長是誰？誰被認為你應該要辭去？行政院長。連戰，好。連戰，什麼時候，他有沒有辭掉行政院長？他有沒有經過立法院的同意？在本件的申請書裡面，有沒有提到，你們對於這樣的問題不重視，是因為你們受了大陸法系的毒？毒，你們中毒，太深了。

\tc{0:31:23}
% cue 57 1883.570 1915.670
\hypertarget{c08:cue:57}{}本屆的文句就好啊，講到憲法原之後，不管本案事實是什麼？錯！你如果本案你不了解事實，你永遠不知道為什麼這樣講，甚至於他講錯，你都不知道。所以要還原到事實。本案先開始是行政院長，院長當時是連戰。請問連戰在什麼時候辭去的行政院長嗎？因為我們剛才講 387,387 在 419 之前對不對？387 才說立法院改選完，他要總辭。你認為當時的行政院長該不總辭嗎？

\tc{0:31:55}
% cue 58 1915.670 1948.750
\hypertarget{c08:cue:58}{}那就是公平安違憲哦。聽懂吧？387 在前 419 在後，這個問題已經先前發生過，那所以 387 號解釋的意思是不是連戰應該要辭？請問連戰什麼時候辭呢？這個先生到底會不會描述？剛才連戰第一個 387 號的解釋是在民國 82 年，因為民國 81 年年底選了第二屆，第二屆的立法委員。所以是 82 年初的時候立法委員上任了發現，誒，那我這個立委，
% cue 59 1948.750 1985.880
\hypertarget{c08:cue:59}{}怎麼那麼無路用啊？為什麼我不能同意行政院長？那行政院長說，沒有人叫我走啊？我為什麼要走啊？好，所以就申請大法官提示。好了，那又隔了三年，所以是 82+3,85 年。聽懂了嗎？第三屆立法委員，立法餘任三年當時，又第二，再下一屆立委，再下一屆立委選出來了。那你認為他們會放過連戰嗎？不可能嘛。好，所以連戰什麼時候辭？好，聽懂了嗎？這樣子又聽懂嗎？過去，我已經講過好幾遍。

\tc{0:33:05}
% cue 60 1985.880 2018.550
\hypertarget{c08:cue:60}{}過去我們立法委員是 12 月選，因為立委的任期是 2 月 1 號開始，所以你剛剛講 2 月 20 幾號就這章了。立法院就開議了。那你過辭，連戰辭了沒？辭了。連戰辭了以後，當時國民黨有沒有佔立法院的多數？有。所以呢，連戰有沒有被同意？有。這樣講有聽懂嗎？所以連戰曾經辭職過，立法院也重新同意過。所以就這點來講，
% cue 61 2018.550 2054.790
\hypertarget{c08:cue:61}{}連戰是合憲的行政院長。好，結果中共什麼時候選舉？320，不然就不會有 319 槍擊事件。319 槍擊事件發生在選舉投票的前一天。所以是 320 選舉。什麼時候就職？520 就職。這樣聽懂嗎？好，所以到了 3 月的時候呢，其實這個事情老早就一定定，不是現在定。如果你回去看歷史，是在 84 年的 12 月，連戰李登輝，
% cue 62 2054.790 2088.410
\hypertarget{c08:cue:62}{}李登輝，在國民黨的中常會，他就說了，我們接下來要面對全臺灣第一次的公民直選總統，1996、320。其實不是 320，是 323，公民直選總統。所以呢，我現在請您連戰同志，作為我的搭檔。我選總統，他選副總統。當時 12 月就講了。但是那時候連戰，還是行政院長，就是辭職，然後再變一次。就剛講了 3 月，85 年 3 月已經完成了。那麼到選舉是 3 月 20 選。

\tc{0:34:48}
% cue 63 2088.410 2122.570
\hypertarget{c08:cue:63}{}3 月 23 號選。選了以後，李登輝總統連戰副總統一起當選。選票是在同一張上，好嗎？所以是行政院長兼副總統，不是副總統來兼行政院長。如果你從時間發展的脈絡，你剛剛點頭，你就要講這一段話，誰先誰後，不是我先當副總統再來兼行政院長，是我先做了行政院長，後選上副總統。其實也不是怪我，是人民選，也不是我有辦法，不對嗎？

\tc{0:35:22}
% cue 64 2122.570 2167.240
\hypertarget{c08:cue:64}{}你們人民投票給我啊，這時候我只有一個選擇，我要不要辭職，還是我不辭職，這樣子而已。這樣你有聽懂了嗎？好，以上是第一點。第二點，請問本案是什麼時候做成的？大法官做成解釋的日期在什麼時候？你知道老師為什麼要問你這個？老師為什麼要問你這個？請回答。97 以後，我們的憲政體制就變了。97 以前，剛才問 387，還有沒有適用？97 以前當然適用，97 以後當然不適用。那這個解釋，
% cue 65 2167.240 2201.240
\hypertarget{c08:cue:65}{}在 97 以前跟 97 後關鍵與否，我再講一次。讀死書，如果讀書，如果讀死書，不如不讀書。因為你背的冥頑不靈，你根本就不知道在講什麼。所以這個解釋是什麼時候做的？像這種案件，它的密切就非常重要。這件事情發生，主要是在講從民國 84 年底，或者應該講民國 85 年 2 月，連載以後，被同意以後，繼續當行政院長。到 3 月 5 月，其實是 5 月 20 號，
% cue 66 2201.240 2239.520
\hypertarget{c08:cue:66}{}就職，就職以後才產生兼任的問題。就職以前，他就是行政院長，不是副總統。即使 3 月當選人沒有就任，是 5 月才就任。所以真正的問題是出在 520 就任以後，520 到 531 還有很久。我再問你，97 修憲是哪一些通過人？我們只知道 97 年，但是 97 年的幾月幾號？回答我，其他人可以查，第四次憲法徵修條文國民大會什麼時候通過人都不公布了？

\tc{0:37:19}
% cue 67 2239.520 2275.540
\hypertarget{c08:cue:67}{}現在講 85 年，85 年就是 1996 年，對吧？還沒有 97 嗎？是。那 85 年是+11 嗎？所以是 1996。所以本號解釋是在 96 年做的，還沒有到 97 年，憲法還沒有修改。所以沒有修改以前，我們還是修正式內閣制。你把修正式括起來，暫時不用，因為他不牽涉到他，現在是內閣制。這樣有聽懂了嗎？好，那你在你的問題要這麼回答。在內閣制之下，
% cue 68 2275.540 2329.170
\hypertarget{c08:cue:68}{}行政院向立法院負責。那這個時候的行政院長可不可以兼副總統。我這樣子句子，你應該精準了嗎？那如果你認為不可以，那是什麼理由不可以？你認為可以是怎麼可以？好，說可還是說不可以？怎樣？你有聽懂？大聲一點？最後關鍵的是說太小聲了，咕咕咕咕咕咕咕咕咕。可以接受嗎？從哪裡看到？不能接受，從哪裡看到？不能接受，從哪裡看到？不能接受？經歷法解釋文第一段念完。大聲一點，
% cue 69 2329.170 2374.290
\hypertarget{c08:cue:69}{}大聲一點念完第一段解釋文。好，反過來，如果憲法有明文，大法官大概就沒有解釋了。聽懂嗎？繼續。好，他為什麼講性質相容或不相容？為什麼加這一句話？因為我們先前讀過一個解釋，釋字 30 號，老師講過了嗎？釋字 30 號次度提出來一個標準。A 機關的人要去接 B 機關。本職與兼職，如果顯不相容是不能接的。釋字 30 號就是在講這件事情。他現在這句話，
% cue 70 2374.290 2406.570
\hypertarget{c08:cue:70}{}4-30 號就是在對 3-30 號解釋做回應。我解釋的時候我知道現在是 419。我們的前輩的前輩的前輩的前輩。第 30 號解釋就講了要看本職兼職是不相容。所以我現在要回應這個問題。那麼如果是不相容，那就句話啊不能接啦。除非我把釋字 30 號解釋推翻。我沒有要推翻他的意思。我遵守啦。但是我告訴你釋字 30 號解釋可能無法適用。因為現在是非顯不相容。如果顯不相容。
% cue 71 2406.570 2417.470
\hypertarget{c08:cue:71}{}如果顯不相容。釋字 30 號就決定。決定了。不能解釋。那非顯不相容呢？嗯嗯。問號。不知道怎麼說。好，繼續。唯？

\tc{0:40:17}
\speaker{學生}
% cue 72 2417.470 2454.930
\hypertarget{c08:cue:72}{}我反正是總統，不會任由，不同，會進向。\par\speaker{老師}好，到這裡是句話要暫停下來。現在我講了幾個字，圈起來。如遇。如遇。現在遇到了沒？沒有。現在副總統缺位嗎？總統缺位嗎？總統不能視事嗎？沒有。他說如果碰到。碰到了沒？現在還沒有。如果碰到。
% cue 73 2454.930 2488.610
\hypertarget{c08:cue:73}{}然後要記憲、槍子圈起來。如果碰到。那就會影響到繼任的規定。那繼任的規定。按照憲法是不是總統不在？說副總統、繼任。副總統不在？總統都不在？行政院長、繼任。所以順序是總統、然後副總統、然後行政院長。所以是三個次序。那現在呢？如果副總統跟行政院長是同一人。那是不是就跟憲法原來的事情不太一樣？所以他講未盡相符。未盡相符。圈起來。我有說到完全不符嗎？我有說他顯然。

\tc{0:41:28}
% cue 74 2488.610 2527.960
\hypertarget{c08:cue:74}{}我有說他顯然不符嗎？我只說不完全。好，讀到這裡。聰明的人就要問一個問題。前面你說不是寫不相扶。現在你說跟憲法的本意不盡相符。請問你大法官，究竟是宣告他合憲還是宣告他違憲？還是在跟我打馬虎眼。所以他說我們董好奇事最大問題。他既沒說合憲，也沒說違憲。聽懂？沒有聽懂。郎正涵對於他自己認為，他既不說合憲，亦不說和違憲。
% cue 75 2527.960 2570.930
\hypertarget{c08:cue:75}{}因為他在解釋理由書裡面，可有交代。先回答有還是沒有，你要打屁股。你已經五年級了，這是你在東吳大學最後一次上憲法課。我不相信你離開這個學校以後，還有人會想你這樣子，要你看憲法。有沒有？他在解釋理由書裡面說，答案是有，那你現在要找出來拿一段。因為你現在也不硬，所以我估計你們記得，如果你對證或抵抵抵，因為它是旁枝，我不能花這麼多時間。如果今天只有兩半隻，我會講到。

\tc{0:42:50}
% cue 76 2570.930 2613.290
\hypertarget{c08:cue:76}{}但是我今天用五個案子、六個案子，所以不可能。你去查，解釋理由書。自己，他自己注意到，他怎麼交代。也不是違憲，也不是違憲，合憲，現在我們該怎麼辦。好，那他實際上怎麼辦呢？請你再念本段，最後一句，言佔堅韌的事實。繼續，應依上開解釋意旨，上開解釋意旨清楚嗎？摸摸糊糊糊，不盡相符，也沒有說違憲，也沒有說合憲，更妙的是，應為適當，
% cue 77 2613.290 2655.480
\hypertarget{c08:cue:77}{}什麼叫適當？連戰是辭叫適當？還是不辭叫適當？看到大法官就是這樣子，這不是今天才開始，他對我們的憲法結構，不會像我跟你們講說我們是什麼制，所以應該怎麼解釋？他若是這樣，脈絡也很清楚，他也不知道到底是什麼制，反正從頭到尾，從 387 號解釋就是糊塗炸，是民意政治、責任政治，到 113 憲判9，還在講民意政治跟責任政治，世界都不知道翻幾番，這就是我們的大法官。

\tc{0:44:15}
% cue 78 2655.480 2696.290
\hypertarget{c08:cue:78}{}所以他現在這樣講，你認為他該怎麼處理？請問你最後，最後講事實問題。連戰辭了沒，後來辭了，不是現在辭了，是到了修憲以後辭了，九七的時候。所以這個問題，表面上是解決的，實際上是沒有，因為講了濃濃和谷，也沒有違憲或合憲的宣告，然後為適當之處理，什麼叫適當，也沒講的。所以正在實際上最後解決，是靠政治人物自己下臺，他如果不辭，你也不見得有辦法。那至於連戰為什麼辭，是因為別的政治上的故事。

\tc{0:44:56}
% cue 79 2696.290 2737.230
\hypertarget{c08:cue:79}{}你們有興趣，你去找新聞，我不講，因為這個是八卦，這不是我要講的事情。好，現在我們回到正文來，那我現在請問，修正式的內閣制，修正式打個括弧，內閣制之下，准不准許行政院長由副總統兼任行政院長？請回答，唐聖漢你先，你認為可還是不可？告訴我為什麼？老師上課，點檢一概，沒有廢話。我如果跟沈國雅講，那我可以講四小時後，會知道你們為什麼要寫的事。我們是聽他講話。
% cue 80 2737.230 2782.680
\hypertarget{c08:cue:80}{}我聽不懂。來，我會幹。來，請說。你認為可還是不可？內閣制之下，可不可？聽懂還是聽不懂？講話，不要猶豫。不為他就說，我不知道。跳過。現在因為你們沒有書，沒有辦法跟你們講話。但是我又畫圈圈。行政、立法、司法、考試、監察。這有五個圈圈。這五個圈圈都是雙向、雙向、雙向、雙向。所以他們的社交關係有二十種，上課我們講過一次。在這個五權分治、平等、消費、平等、消費、平等。這個是釋字第幾號。

\tc{0:46:22}
% cue 81 2782.680 2821.160
\hypertarget{c08:cue:81}{}這個圖書裡面沒有包含總統跟國民大會。國民大會展示在這裡。國大、總統在這裡。總統是不是國大選？總統是不是國大罷？所以他們倆之間是他向他負責的？就這樣子。跟這個圖書不問。因為我們的當時是內閣制？內閣主要是講這一塊。行政院、立法院這樣子。其他的跟他沒有關係。這樣有聽懂？好，那我們先來問一頭問題。這個時候，總統，
% cue 82 2821.160 2863.380
\hypertarget{c08:cue:82}{}總統還沒有直選，1996 的直選。對，這時候總統已經直選了。這個時候憲法還沒有修改。你是李登輝，你是權利一票票投出來的。你是有史以來第一個臺灣人選的臺灣人總統。我問你，在你內心世界裡面，你是不是覺得，行政院當然要聽我的，要不要我選假的嗎？你覺得這個圖形在他的心目中還能接受嗎？你選一個總統費這麼大的力，結果這一塊也被你管不到，會同意嗎？所以你就知道，
% cue 83 2863.380 2901.670
\hypertarget{c08:cue:83}{}他在玩遊戲，他在玩個遊戲，我一句話就做不到，教父哥得罪李登輝。李登輝要測試一下，我現在如果不修改憲法，現在還沒有修改憲法。我不修改憲法，但是我已經總統直選了，我想要把我的手伸進來，我要指揮行政院。你認為他要怎麼做？他是一個非常聰明的人，非常聰明。看著等他的人，世界上用屈指可數，沒有基本。他的做法，他不是行政院長叫來罵一個，不是，這太笨了。

\tc{0:48:21}
% cue 84 2901.670 2943.260
\hypertarget{c08:cue:84}{}這個總統跟我一起搭檔當選，是一個副總統。他在沒有當選之前，他究竟是行政院長吧？大選了以後，我叫他不要辭，那麼是不是我的副手就跑去當行政院長呢？那我請問你是連戰，覺得李登輝他是總統，你是他的副總統。在這之前，他提拔你當行政院長。在這之後，他也沒有把你免職。那你認為你要不要跟他好好配合？那再問一個問題，沒有修改之前，我們是內閣子。行政向立法負責，跟總統無關。
% cue 85 2943.260 2978.460
\hypertarget{c08:cue:85}{}總統可以把手伸進行政，用他的副手來接行政院長安。如果這樣子，那行政院是聽總統的，還是行政院是聽立法院的？行政院是向總統負責，還是行政院向立法院負責？我這樣講，有聽懂嗎？這是另外一個有趣的案子。在過去 419 號的解釋，過去，在 Mio 老師講這個之前，大家只看到文字上，什麼未盡相符，適當地處理，咱在說什麼？但是沒有人真正看懂這個解釋。

\tc{0:49:38}
% cue 86 2978.460 3013.720
\hypertarget{c08:cue:86}{}背後適當地處理，這個利害的關係就是我現在講的，已經民選的直選的總統，他還沒有修憲以前，他沒有直接任免行政院長之前，他在想，如果我可以用兼任的方法，讓行政院長和我的副總統去接，那我就可以控制行政啊，我幹嘛修憲？但是沒有想到大法官出來攪局啊，你們這群老徒是老年人，你們搞不懂啊，你們看不懂我，李登輝千秋大的業，你要出來攪局，你也不敢進去，
% cue 87 3013.720 3052.380
\hypertarget{c08:cue:87}{}你明白了腳，然後你腳一半。啊，因為適當處理，於憲法、恩府未盡相。意思嗎？好，不關心，事意句話，其實你放，你做了，我就先拜著。連戰，我沒有準你辭，你就繼續做。到了零年，我修憲修完了，你再辭。修憲修完了，就是總統直接任，直接免，行政院長，這樣叫敬酒不吃，吃罰酒。很好，更好，我現在不管你的婚要做什麼，我只是回到原來的體制。修正式內閣制，還是內閣制，
% cue 88 3052.380 3091.180
\hypertarget{c08:cue:88}{}你不能想像內閣制行政向立法負責，居然總統可以派副手來指揮行政。那這樣子是不是變成是總統向總統負責，或者至少一個公公、一個婆婆，兩個都要聽，這個都是會混亂行政機制。所以大法官即使沒有辦法把它宣告為違憲，它也都出來交易去。其實這個如果按照修正式內閣制，它應該跟 387 號解釋同樣的講法。387 號不是說民意之所在，責任之所在嗎？現在民意在哪？
% cue 89 3091.180 3100.370
\hypertarget{c08:cue:89}{}在立法院啊？責任在哪？行政向立法負責。條文改了嗎？都沒有改啊？你總統跑進來幹什麼？副總統是什麼玩意啊？

\section{行政院院長於新任總統就職時之禮貌性辭職（釋字第 419 號解釋文第二段）}

\tc{0:51:40}
% cue 90 3100.370 3138.930
\hypertarget{c08:cue:90}{}他應該這樣寫，但他不敢寫。好，來，現在看第二大段，請大聲念。這一段是因為前段引申出來的。好，我現在問了兩個問題。第一，他這樣子的描述有錯嗎？行政院院長過去，是不是在總統就任的時候，他就會辭職。他說他這樣的描述有錯嗎？
% cue 91 3138.930 3155.040
\hypertarget{c08:cue:91}{}事實嗎？對呀，這是個事實啊，你就回答，事實就好啦。是。那現在我問你，那為什麼會出現這樣的事實呢？

\tc{0:52:35}
% cue 92 3155.040 3181.280
\hypertarget{c08:cue:92}{}不然你在說什麼，你後面聲音又來越小聲了。你認為為什麼行政院長要總辭？行政對立法負責，你講的對呀，那是行立法院選完你去總辭，那是三八七啊。你現在跟我講是總統就職總辭，你在說 A 還在說 B 啊？
% cue 93 3181.280 3223.710
\hypertarget{c08:cue:93}{}你投昏了，學家同學請回答我，為什麼？你剛剛講的都對，它是威權時期，兩蔣時期。欸，怪了嗎？那為什麼？他要去談，總統改選，他就要去辭呢？為什麼？請回答啊，你們怎麼不回，怎麼不會小說？那老闆，老闆意思我們要猜一猜啊。老闆都已經是新的任期了，我行政院長還要繼續做，難道這麼不要臉要做到老闆來免職嗎？
% cue 94 3223.710 3244.480
\hypertarget{c08:cue:94}{}還是我應該喘我喘意？哎老闆恭喜你啊！新的六年開始啦，拿來我辭職，這樣子你手力彈出來，看你要認你誰，這叫做以退為進，以退為進都不懂你就不能當行任，聽得懂嗎？

\tc{0:54:04}
% cue 95 3244.480 3263.700
\hypertarget{c08:cue:95}{}好啦，那你另外一個比較合理的說法，比較文謅謅的說法是，國家總有新陳代謝啦，哪有行政院長做萬年呢？總統都有任期了，行政院長沒有任期，但就一直做嗎？也不通啊！那我竟然是，
% cue 96 3263.700 3297.550
\hypertarget{c08:cue:96}{}因為總統而當行政院長，這句話要解釋一下。釋字 31 號，是不是說立法委員親自改選？由第一任繼續任？所以剛才才會產生 387。到了第二屆立委產生的時候，前面都沒有例子，因為沒有立法委員改選，並以後行政院長總是沒有。55 條第一項從沒有用過，除了第一次用禮拜，再也沒有用過。民國 37 年那一次是有行政院長，是有經過立法院，以後就沒有出現了。
% cue 97 3297.550 3333.160
\hypertarget{c08:cue:97}{}這是立法院改選，以後要去總辭。沒有這種例子，因為立法院不會改選，第一屆繼續做，所以 55 條就凍住了。那在這種情形下，如果沒有這個慣例，說行政院長在總統就任以後要總辭，是不是無從新陳代謝。那你也可以做得很難看，德國總統把你免職，他當然可以把你免職。有講話有聽懂嗎？好了，所以老師問的問題是有沒有這個事實？有！那我現在進一步的解釋，為什麼有這個事實？

\tc{0:55:33}
% cue 98 3333.160 3385.960
\hypertarget{c08:cue:98}{}請繼續，第二句，這句話，為什麼會出問題？前面已經講了，這種辭職是禮貌性辭職，不是憲法上的義務。憲法上的義務是 387 號講的那種，立法院改選完畢以後，你必須辭，這是義務。總統選舉完跟你行政院沒有關係，總統在這邊，行政院立法院在這邊，他選他的，你選他的，跟他什麼關係？你只需要盯促的是這個，這個是行政院，這個是立法院，你是向他負責人。
% cue 99 3385.960 3432.590
\hypertarget{c08:cue:99}{}你是向他負責人。他如果改，你覺得跟著改，是這一塊。他跟你有什麼關係？沒有關係，你當然能比理他。是這個意思。但是呢，他現在問題來了，就有一個這麼禮貌，不禮貌多的行政院長。他就說總統啊，你都任期到了，你該選了，我也跟你該選了。好啦，這個案子呢，為什麼會出現這一段話？有沒有人能夠回答？520 連戰跟著李登輝當選委政府後，連戰有沒有向總統依照慣例提出總辭？

\tc{0:57:12}
% cue 100 3432.590 3477.260
\hypertarget{c08:cue:100}{}那麼連戰敢不辭嗎？會吧？連戰再怎麼樣也是深諳為官之道。他辭了。辭了以後你都會批了四個字。有誰能答出來是批了四個字的家庭天母？沒有。人會在他的詞子上批了這四個字。這四個字怎麼念？
% cue 101 3477.260 3524.560
\hypertarget{c08:cue:101}{}他們認為，著毋庸議什麼意思？著字有上對下的關係。通常是皇帝對臣子講。簡單的講，連共啊，不必討論啊。我很喜歡你啊，你繼續幹啊。因為他批，著毋庸議，所以連戰才沒有辭的。如果他要叫連戰走，他就可以批 yes。他沒有啊。其實這句話的這四個字就是等於說，我免與未留啦。他沒有讓他辭了。好了，那現在大方會講，這個不是憲法上的義務，但是他辭了。

\tc{0:58:44}
% cue 102 3524.560 3562.670
\hypertarget{c08:cue:102}{}這種事情要怎麼處理呢？大法官說，總統就具有裁量權限。這句話，不管他後面講什麼統治行為，這是包裝率。如果總統沒有裁量權限，他就不是統治。但是問題是，總統為什麼對這種禮貌性辭職有裁量權限呢？回到剛才我講的，在這個時候是不是，這個時候是不是還是內閣制。那麼如果行政院長向國家元首，他是行政連一首，向國家元首，表示禮貌，
% cue 103 3562.670 3602.310
\hypertarget{c08:cue:103}{}不管他是必須什麼考慮，那總統有權利批准，或者有權利慰留嗎？是跟你玩遊戲的，你是總統，不是國家元首，元首，我禮貌性的跟你講，我要辭了，可以準還是可以不準嗎？如果不可以，哪來要裁量權？如果是憲法義務，你沒有批准不批准，他就是一定要這樣做，那你說不是憲法義務，那你有裁量權，那這問題就來了，總統可能等你等很久了，
% cue 104 3602.310 3639.440
\hypertarget{c08:cue:104}{}準，你可以準嗎？你可以準還是你可以不準？我認為這件事情應該沒有裁量權，因為這個是回到原則，行政向立法負責，不是向總統負責，所以我禮貌性我給不，我禮貌很多，我很多理，我向你提出辭職，這不是義務，不是憲法上的義務，那你呢，照說你應該是假性質、假設、假作，只能慰留，你不可以跟我講說準還是不準？

\tc{1:00:39}
% cue 105 3639.440 3691.000
\hypertarget{c08:cue:105}{}說他是統治行為，因為這是很難理解的，然後最重要他講為什麼是統治行為，因為他得到的解釋是非本院應做合憲性審查的事項，這個也不要來問我們沒有權力審，這樣子大法官就躲過一劫，兩邊都不得罪。這個志祺如果從體制上說，是比較容易碎的理由，現在還有第三段，請繼續。他自己講什麼了，還是我講的，對不對？繼續。好，請問你為什麼會加這句話？因為當時有一個附帶的爭議。

\section{立法院決議對總統之拘束力（釋字第 419 號解釋文第三段）}

\tc{1:01:31}
% cue 106 3691.000 3727.190
\hypertarget{c08:cue:106}{}立法院呢，看了總統寫著毋庸議。立法院生氣了，趕不走連戰。立法院通過一個決議，要總統重新提名一個行政院長，對大法官才加這一句話。請問立法院做的決議對總統有拘束力嗎？立法院可以命總統定之行為，請你重新提名院長，或者命他不為一定的行為，你不可以準他怎樣嗎？他說沒有，因為要負責的是行政院向立法院負責。向立法院負責的是行政院，不是，
% cue 107 3727.190 3790.140
\hypertarget{c08:cue:107}{}不是總統，顧語，咨字兩個字圈起來我已經講過了，總統跟五院之間的公文來，叫做咨。立法院咨請總統如何，總統咨立法院如何，算你非常非常棒，非常 bureaucratic，非常官僚，但他就是用什麼？這是他們之間的特權，只有他們可以用什麼？我們其他人？好了，到這裡清楚了嗎？好，那現在 419 號解釋是不是已經說完了？那再來一個問題。

\section{1997 年修憲後副總統或總統兼任行政院院長}

\tc{1:03:10}
% cue 108 3790.140 3834.390
\hypertarget{c08:cue:108}{}419 號解釋與今天還有 419 號。這是在剛才講過了，97 修憲以前。97 修憲以後，副總統能不能接行政院長？97 修憲以後，行政院長誰決定？總統啊！總統任總統免啊！因為我們是總統制啊！那你認為，在這種情況下，總統可不可以指定副總統做行政院長？可不可以？anybody 都可能被他指定啊！他願意就好了啊！但是剩下第二個問題，那麼副總統接任行政院長，他這個人可以？
% cue 109 3834.390 3873.290
\hypertarget{c08:cue:109}{}他這個行政院長，他就必須履行行政院長的職責，他要到立法院去給人家質詢、備詢等等。我講話有聽清楚啊！不是說他可以躲在總統府裡面，不是這個意思啊！我叫你去兼行政院長，你就是行政院長了！所有行政院長的質詢、義務，你要一起承擔。這邊要不必我講吧！不要再去過往文字遊戲啦！臺灣就有這種人，一定會跟你這樣講啊！這樣子啊！不是的！好！再問一個問題，總統自己可不可以兼行政院長，問什麼問題都沒有去啊！

\tc{1:04:33}
% cue 110 3873.290 3910.410
\hypertarget{c08:cue:110}{}我認為這個情況已經快就要回答了。我跟妳講，如果賴清德是一個大有為的總統，看到這個卓榮泰，噗噗噗噗噗的，在立法院都不會答詢，一間到晚害怕損失分，他明年競選連任都會成問題。賴清德在今天，今年我不知道 11 月 28 號投票結果怎麼樣，但是如果投票結果以後，民進黨大敗，六都全部丟光，假定。賴清德都會說，好啦，我也不要再白手，
% cue 111 3910.450 3954.380
\hypertarget{c08:cue:111}{}我自己兼行政院長，我自己立法委員出身，我去立法院成江舌站，絕對比你搶奪了，我去！可不可以？考試考這題會不會有回答？考這題非常具有挑戰性，這個是表示你懂不懂憲法，以及你是不是一個好律師，你會不會有論證？如果那個答案都寫在調問題的案，你幹什麼？你也背不過機器，你問機器人 AI 又被打了嘴，但是我現在問這個問題呢，AI 就不見得被打了，現在，我覺得以後他一定比我，你會被打了？

\tc{1:05:54}
% cue 112 3954.400 3993.370
\hypertarget{c08:cue:112}{}在恩賜言講中，所有講的東西，他都聽進去了，你是換給他，他現在絕對比我厲害了。來，請回答我，總統自己可不可以身兼信任院長，老師在，有沒有在任何地方講過這個？有，你有聽到？在哪裡？很好，你很棒，你在哪裡聽到？我很喜歡你中文，你像是我的粉絲，這樣就對了，你知道你要聽誰的意見，再憑柯文哲的政見來，可以一直講，因為柯文哲的如意，
% cue 113 3993.390 4032.010
\hypertarget{c08:cue:113}{}在講，怎樣怎樣怎樣，這個政府怎樣怎樣，老師就有跟他講，你如果看不完你就去吧，你可不可以，你永遠說你不可以，好，這個問題的答案是這樣子，第一，你要回到憲政體制，第二，回到憲政體制有什麼用，因為總統任免行政院長，我可不可以任自己，你很難說不可以吧，我可以任命一個人去當我的分身，我自己不可以本尊兼分身嗎，我可以任命，你不能任命我自己啊，什麼道理，我知道大多數人會說不可以，老師。

\tc{1:07:12}
% cue 114 4032.010 4070.410
\hypertarget{c08:cue:114}{}這個不可以，我猜想很多理由，第一個是他權力太大，第一個，你真的覺得這樣會權力太大嗎，所有屬於行政院長的權力是不總統賦予的，如果賴清德，如果卓榮泰說的不好，賴清德就會換卓榮泰二，換三，換四，分身，所有的權力都是來自於本尊，本尊跟分身是形式上的差異，權力來源，同一個權力來源，本尊如果沒有了，分身也不可能，句號啊，不是啊，我們講的行政院長是憲法規定的行政院長。

\tc{1:07:50}
% cue 115 4070.410 4106.600
\hypertarget{c08:cue:115}{}他現在的改變，只是總統任免行政院長，他也沒有多一分也沒有少一分，只是來源差別而已。第二，總統去自兼行政院長以後，還是回到剛才那一區，他有兩個身份，他是行政院長，他就必須履行行政院長所有的職責跟義務，他不可以不去備詢，不可以不可以不可以怎樣，他當然有地方可能會有個問題，今天的情況會避免掉，如果現在是因為卓榮泰不副署，然後擔心者說，你沒有東西要，
% cue 116 4106.600 4143.420
\hypertarget{c08:cue:116}{}我副署，因為你沒有副署，那你也有個案子，那所以我就不公布，從而立法院通過的法律就被我們通通沒收了，就不被執行，不存在，沒記啊，這種情況大概就比較難，因為你是行政院長，你又是總統，同一人，你不能再講說，我行政院長沒有副署，所以，是你自己要不要副署，是你自己要不要公布，我把那個兩尊合成一尊，我看見多大白群，倒是有一個意外的好處，他就賴不掉，賴皮總統。

\tc{1:09:03}
% cue 117 4143.420 4176.410
\hypertarget{c08:cue:117}{}他就再不能說這不是我，我的副手沒有副署。事實上大家都知道，以前副署制度是行政院長來對付總統，因為是內閣制，內閣制，他對他的父親，他是國家元首，一起給你，你是誰要公布，因為你是國家元首，所以我很怕你奪權，所以你所有的東西命令都要往副署，表示我任賬，這是用來對付他，現在這個關係不見了，他呢，換了一個位置，他現在從這裡，
% cue 118 4176.410 4191.130
\hypertarget{c08:cue:118}{}從這裡，已經跑到這裡，我已經新的圖形上，我把總人畫到這裡，他都一次改吧，這種時候，還叫他怎樣，你們都是聽我的武士嗎？卓榮泰，你不聽了嗎？你要造反嗎？

\tc{1:09:51}
% cue 119 4191.130 4219.570
\hypertarget{c08:cue:119}{}你聽分別人是，所以我才會講，卓榮泰的副署，他是在跟賴清德演雙簧，唱雙簧，兩個人有泛逸溝通，講好，不要副署，不要害我要公布，那我不副署，你不副署，我公布，就沒有執行的義務，因為東西不存在，這是不是個笑話？

\tc{1:10:19}
% cue 120 4219.570 4264.070
\hypertarget{c08:cue:120}{}那請問立法院通過的東西去哪裡？人間蒸發。他有被撤銷嗎？對。你們問你們自己啊，學法律，你們問你們自己啊，立法院通過的東西，在沒有經過大法官、宣報我一些失效以前，他在不在，笑話說啊，你覺得是不是笑話？他就演給你看啊，好啦，我講太多了，你們自己去想想看，這個解釋419到這裡，全部講完，可以了嗎？還有同學，有問題嗎？可以嗎？有問題，等一下再問好，好不好？可以嗎？我用很短的時間。

\section{釋字第 3 號、第 175 號解釋與五院之法律案提案權}

\tc{1:11:04}
% cue 121 4264.070 4296.870
\hypertarget{c08:cue:121}{}因為你們要啟料，因為現在我們就落後，我不能一直這樣落後，所以我們繼續要追上來，好，現在往上追，我們追到 175 號，釋字第三號解釋，沒有講完，我就有簡單的話講，釋字第三號解釋最重要的是他的解釋方法，釋字第三號解釋的重要是在他的解釋方法，他把所有能用的憲法解釋方法全部用完了，那在所有的方法當中，我幫他加一句話，畫點進一下，
% cue 122 4296.870 4329.960
\hypertarget{c08:cue:122}{}最重要的解釋方法是體系解釋方法，體系解釋方法就是你要看相關的條文，把這個圖案拼起來，不是抓了一個條文，啊用那個文字是這樣這樣這樣，你就跟文是真，那這個條文跟這個條文要不要合，他們是不是同樣在一部憲法裡，彼此是不是要不吐相廢合，所以一定要解釋到這條也通，這條也通，這條也通，每條相關的都通，這樣這個解釋才能被解釋，好，如果從這個觀點講，只要牽涉到憲法結構，
% cue 123 4329.960 4371.950
\hypertarget{c08:cue:123}{}總統是內閣制版總統，最關鍵的解釋方法就是體系解。好了，我覺得釋字第三號，我大體計講完，那我現在只要問一頭問題，如果釋字第三號不做這個解釋，會怎樣？請回答我。三號是哪個同學負責的？你們好，那釋字三號爭什麼？爭監察院怎樣？可不可以對立法院提出法律案。那如果大法官不做這個釋字第三號解釋說可以，那你覺得監察院要怎麼提出法律案？

\tc{1:12:51}
% cue 124 4371.950 4408.150
\hypertarget{c08:cue:124}{}理論上是很多過行政院，然後行政院再送到立法院，那我問，如果行政院要卡死，是不是就被卡死了？如果這樣就卡死了，他有五權分治平等相維，你就挨我一劫來跟我分支什麼？那聽懂嗎？所以大法官五權分治平等相維，這是他對於整個五權憲法的總體理解，一個 overview，他這個 overview 要能成立，他有些地方要做好，第一，兼任的事情不可以，要不然你是我，我是你，大部分在一起，這個不同意，不可以。
% cue 125 4408.150 4449.420
\hypertarget{c08:cue:125}{}所以兼任，我們花很多時間，想，30 號、20 號、15 號等等，不可以兼任。第二件事情是，該有的權力，絕對不能縮損。你可以向立法院，鏈隔向立法院，鏈隔向立法院，大家都可以向立法院，這邊還少了幾條線，20 種。這樣能聽懂嗎？所以這個解釋，對於確定監察院能夠作為平行的一個院，關鍵。好，經過這個解釋以後，我們盤點一下。行政院向立法院負責，行政院永遠會提出法案。立法院本身，立法委員連署員會提案。

\tc{1:14:09}
% cue 126 4449.420 4477.740
\hypertarget{c08:cue:126}{}考試院本來就憲法的規定可以提案。現在剩下的一個問題，那司法院沒有，司法院能不能提案。那就是 175,175 是誰？那 175，大法官說司法院可不可以提案？可以。那我問兩個問題，大法官可以多解釋嗎？有沒有人說司法院就是你，你就是司法院，你怎麼可以自己解釋？
% cue 127 4477.740 4518.930
\hypertarget{c08:cue:127}{}司法院的意思是說，聲請人是監察院？是。 聲請人是監察院，但是他問的問題是你啊？那你自己開始解釋你自己啊？你可不可以這樣解釋？我問的問題有聽懂嗎？司法院自己解釋自己可以向立法院提出法律案，有沒有利益衝突？那司法院也很不耐啊。我們憲法，我們用我能解釋啊。我不解釋，那怎麼辦呢？憲法又沒有規定。第二個可能，請問一下，175 號解釋。

\tc{1:15:18}
% cue 128 4518.930 4560.940
\hypertarget{c08:cue:128}{}175 號還需要解釋嗎？13 號在論理的過程當中，有沒有提到司法院，也可以向立法院提出法律案。杜偉啊，他有講這一句啊。那有一個可能的解釋方法，就是說我們不需要解釋啊。因為釋字第三號其實你講過了。表示說大法官，他做了一個區別。釋字第三號問我們的問題，我剛才問你的，釋字的監察院可不可以。在論述的過程當中，我有談到司法院可不可以。但是這是討論的，dicta。
% cue 129 4560.940 4598.390
\hypertarget{c08:cue:129}{}他不是我的司，阻止，阻止在這裡。這個問題解決了。在阻止當中，旁論附帶提到的東西，沒有舉足。所以，請 175 號針對這個問題，又再做一次解決，有聽懂嗎？好了，到這裡，175 號結束。現在我們講 釋字 14 號，14 號是誰？好，來，請問你 釋字 14 號，上次老師說，這個 14 號其實有問題，什麼意思？我說有問題，我的話是什麼意思？我記得我上次講說有問題。

\section{釋字第 14 號解釋與監察院彈劾對象}

\tc{1:16:38}
% cue 130 4598.390 4644.020
\hypertarget{c08:cue:130}{}有問題嗎？沒有，沒有，好，那你們覺得有沒有問題？先告訴我 14 號，請示說什麼，來，講話，誰可以不可以什麼？好，一句話的意思，說得這麼文謅謅，就是監察院監督的對象，包不包含立委啦，就這樣講一句就敗啦，監察院可不可以監督行政院？可以，監察院主要是來監督行政院，監察院可不可以監督，這個是考試院，可不可以，這個是司法院，可不可以，可以，可以。

\tc{1:17:24}
% cue 131 4644.020 4680.980
\hypertarget{c08:cue:131}{}所以去看憲法相關文字都有，他沒有講到立法院，這是大法官 14 號、解釋說不可以，可以的結果，是不是就造成了，這一條線就斷掉了，立法院有沒有辦法對付監察院，有吧，監察院預算要把立法院通過，監察院的組織法是不是立法院定的，監察院要行使職權是不是要監察法，監察法是不是法，是不是立法院定的，所以立法院永遠可以控制監察院，那監察院監督的對象不包含立委，所以立委，
% cue 132 4680.980 4721.310
\hypertarget{c08:cue:132}{}所以立委、鯊魚、放火、包捕、包裝、包捕疹，都不能對他彈劾，什麼意思嘛，那你這條線就斷掉了，我可以對付你，你不能對付我，他太好了，那你認為這個其實有問題，有用問題，就是我的意思，你要我們全分支平等效率，你要每一個都有毛骨盾，你有毛骨沒有盾，就死定，那猜想一下，為什麼大法官會說立法委員不是監察院所能監督的案，不是監察院所能監督的案，不是監察院所能彈劾，究竟的案，為什麼他要這樣講，\par\speaker{學生}公立大會，好。

\tc{1:18:41}
\speaker{老師}
% cue 133 4721.310 4760.890
\hypertarget{c08:cue:133}{}你現在講的我不知道對不對，但是我可以回應你，國民大會什麼時候管道監督呢，你現在這條線是這樣拉的，我現在是說這條線有沒有這一塊，這一塊被他砍掉了，那你現在跟我說沒有關係有他，那我問你這條怎麼拉到這裡，請問監察委有國大選的嗎，沒有啊，那國大能夠彈劾接軌嗎，沒有啊，所以他跟他沒有關係，他跟他不具有替代行為，好像不是，那還有什麼別的理由可以解釋？

\tc{1:19:20}
% cue 134 4760.890 4797.630
\hypertarget{c08:cue:134}{}他跟他講，他講，他講，他講，他講，他是說，哦，你是民代，間接直接選取才能民意代表，所以我就不能對你行政院長、權利。所以我就不能對你行政院長、權利。為什麼？他講啊，他沒講，他是說，哦，你是民代，間接直接選取才能民意代表，所以我就不能對你行政院長、權利。
% cue 135 4797.630 4834.730
\hypertarget{c08:cue:135}{}行政院長、權利。可是他沒有說，為什麼，為什麼不應該對你行政。我們有沒有白話講，請問，會包賭、包娼、包工程呢？有沒有包含立委？那別人包賭、包賭、包工程，都要被你彈劾。假定他是行政院、法院、假定他是司法院的法官、假定他是考試委員，會不會被你彈劾？那為什麼立委包賭、包賭、包工程，就不答應了什麼道理？你跟我講他是民選，民選就不能彈劾，這又是什麼道理？我現在跟你講清楚，大法官搞錯了一件事情，我說他是搞錯，
% cue 136 4834.730 4868.460
\hypertarget{c08:cue:136}{}你可以不同意說，老師你這樣批評到訪問，太激烈了，你可以說他沒有搞錯，這是什麼問題？在 114 憲判1跟 113 憲判又出現了。簡單的講，責任有兩種，一種的叫政治責任，一種的叫法律責任。現在呢，大法官是把這兩個混為一談了。大法官意思是說，他們是民選的，所以人民選的人民罷免就是負政治責任。剛才 387 也是講，行政向立法負政治責任。

\tc{1:21:08}
% cue 137 4868.460 4903.160
\hypertarget{c08:cue:137}{}有沒有，可以說，還有經濟嗎？這個。但是監察委員要對你提出彈劾糾舉，一定是政治責任嗎？還是他只能是法律責任？你政治責任，你去罷免他，我沒有意見，我沒有，有沒有說可以彈劾的對象，人民就不能罷免？沒有啊，兩個是並存的啊。政治責任是政治責任，誰選你，誰罷免你，向誰負責？總統是不是人民選的，所以人民能不能罷免總統？立委是不是人民選的，所以人民能不能罷免立委？立委是不是人民選的，所以人民能不能罷免立委？立委是不是人民選的，所以人民能不能罷免立委？

\tc{1:21:43}
% cue 138 4903.160 4937.060
\hypertarget{c08:cue:138}{}人民能不能罷免立委。如果不能哪裡來的去年大罷免，我講話已經懂罷免。那為什麼我不可以罷免？你可以啊！我沒有說你這個民意代表，你的民意代表不能被罷免。我現在是講到能不能彈劾。彈劾的要件，在憲法上不太清楚，但是不是完全沒有。你看相關的規定都是違法或失職。所以它的構成要件是違法或失職。所以它的構成要件是違法或失職。失職的情況確實比較難考慮，它有一些政治的考慮。比方說，
% cue 139 4937.060 4971.710
\hypertarget{c08:cue:139}{}我這個政府是擁核的，你呢，就不給我去蓋核四，老師那邊偷偷拉拉拉，那算不算失職。反過來，我是反核的政府，你天哪說要蓋核四，批批你去蓋，那你才失職。但是有沒有違法可能沒有，因為沒有法律的規定，這政策的規定。但是有時候是有違法，剛剛講包賭、包娼、包工程，主要是違法。我沒有講包賭、包娼、包工程，是失職喔，沒有這樣講。好啦，教訓有聽懂了吧？所以大法官呢，是以為，啊。

\tc{1:22:51}
% cue 140 4971.710 5006.140
\hypertarget{c08:cue:140}{}如果立法委員，既然要對人民負政治責任，所以呢，監察院，那對他課法律責任就可以不必了。其實壞的角度，兩者不能並存。對其他的官員不都是並存。請問司法院的大法官，有沒有政治責任，他倒是沒有，他沒有跟誰領責任。所以他只能有法律責任，不然司法院大法官沒法獨立。但其他人，包含這個包含這個，都有。但不管怎樣，你把這條線切開了以後，就變成監察院，沒有辦法對立法，對付立法院。監察院可以對付行政院，
% cue 141 5006.140 5040.710
\hypertarget{c08:cue:141}{}對付司法院、對付考試院。監察院還可以對監察院本身進行監察院。監察院委員有被監察院大學的記錄，好嗎？因為每個人都會違法失職，違法失職就構成我要監督的地方。那為什麼立法委員違法失職不能監督？他的意思就是他把他搞錯了。他認為立法院已經有政治責任了，就不要再有法律責任了。但這就是一件事情，平安心了。有一種立委，他非常會做人，尤其是現在律遊遊的時候，他們是一個黑金組織。我有幫忙，
% cue 142 5040.730 5080.030
\hypertarget{c08:cue:142}{}所以我就分你。所以你們全班每個人都分到我的錢。請問你們會幫忙嗎？我分五百萬，你們每個人都分一萬。反對什麼？我分一千萬，你們都分兩萬。反對什麼？所以我會跟你們都幫你。政治責任，你沒有，你根本不到我。但是法律責任也沒有，那可能，想問我的錢，我才會回來的，我有法律責任。那這就是衛士表的奢侈堅的原因。我想還有聽懂了嗎？好了，十三號解釋，還有什麼問題？有，現在請提出來。那個負責的那一組同學，十三號解釋有問題嗎？

\section{釋字第 319 號與應考人申請複查考試成績}

\tc{1:24:40}
% cue 143 5080.030 5142.870
\hypertarget{c08:cue:143}{}十四號，對不起，十四號。有問題嗎？沒有啊。好了，沒有就結束了。那我們現在繼續往下。我們看一下 319,319 號解釋是哪一組？好，也是你們的。好，請問一下，319 號解釋在講小，我知道，他不要管他，他問的問題是什麼？他的爭執的問題是什麼？在限制裡面引起他本件申請解釋的爭執是什麼？他要求什麼？現在看一下解釋文。解釋文這一段文是，你自己看。
% cue 144 5142.880 5181.860
\hypertarget{c08:cue:144}{}你看應考人申請複查考試成績處理辦法？其中第八條規定內容很多。申請複查考試成績不得要求 1.重新評閱。2.提供參考答案。3.閱覽複印試卷。4.不得要求告知閱卷委員的姓名。5.其他相關資料。他用五件事。你是不是哪一件？他是哪一件人？我熬了他他來的。所以你剛剛講的不太對吧？你不能用猜的，你得要去看他的申請書來回答你。

\tc{1:26:21}
% cue 145 5181.860 5221.180
\hypertarget{c08:cue:145}{}但是就目前的這個情況看起來，可能是他要求要存清書。我們現在所有人都會去考試。這個規定有修改，但是沒有完全廢止。這個規定，老師在當學者的時候，已經多次，而且我還寫了文章，我在考試院的會議裡面我還提一個報告。我認為這個規定是有檢討的空間的。首先，你們現在的考試有兩種試題。一種叫申論題，一種叫測驗題。測驗題是有標準答案的。
% cue 146 5221.180 5259.210
\hypertarget{c08:cue:146}{}它會公布，這個比較沒有問題。那我們講申論題。申論題，不能要求重新評閱。我覺得很合理吧？不然我今天這題得了三分。他得了三十分。重新評閱的結果可能我十五分到零分。那太好。那考試有沒有考不完？每次考有人異議就重新評閱。你們從小到大沒有聽過這樣的考試吧？這等價再說。第二個，提供參考答案。現在你們有看到參考答案嗎？今天在座的，有 75 年級的，有去考試的舉手。
% cue 147 5259.210 5298.240
\hypertarget{c08:cue:147}{}一個都沒有，都沒有，你都沒有去考試。你們去考試試就知道了。現在剛才說有兩類嘛。一類是申論題。申論題是要求考官要提供參考答案。但是參考答案項目公布。尤其是今天我們講的 3192 提示。確認，現在 3192 提示已經很久了，幾十年了。我不覺得。在 3192 提示的保護下，考官提供參考答案有任何實質上的問題。我覺得沒有，考試因為太保守了。

\tc{1:28:18}
% cue 148 5298.240 5340.040
\hypertarget{c08:cue:148}{}現在，閱覽或複印試卷，這個可以同意。請問你們，為什麼你不可以去看看你的卷子到底有沒有被改到？上面有沒有批改月的痕跡。可以看嗎？那你說我要複印，他說你不能複印。你們覺得考試院在怕什麼？用你們白話回答。怕什麼？為什麼考試院不答應？複印不能複印。閱覽也不能複印。閱覽有很多種。你們想到的閱覽是眼睛去看的。但現在閱覽可以做到。透過電子攝影眼。
% cue 149 5340.040 5383.280
\hypertarget{c08:cue:149}{}你看吧，這有什麼困難。這不能都會啦。我都會啦。所以閱覽，你說閱覽，會傷害到卷子。閱覽會造成視覺獨感。這也是可憐的。以目前的科技。對不對？那複印。複印更不可能造成試卷有什麼傷害。那複印到底有什麼理由呢？不讓你複印。講話。你們都是考生。你們那麼沉默。難怪考試院他欺負你了。他欺負得扁扁的。因為你們太好欺負了。想到到底聽聽看。為什麼不可以啊？都沒有顧慮啊。

\tc{1:29:43}
% cue 150 5383.280 5441.770
\hypertarget{c08:cue:150}{}好。那我就不講了。那就讓他繼續不行了。因為你們都沒有顧慮嘛。我幹嘛。我幹嘛替你不憂慮呢？我都當不壓嘴。像我這樣的老師都討厭。他們很害得我。大家怎麼能走。我害其他老師很辛苦。好。我們現在回來看這個解釋文。那一組哪一個同學名大聲地念。到了句號停下來。大聲。所以看起來他在真的時刻要求再行評。那不能再行評。他的前提是只要彌封開拆。
% cue 151 5441.770 5477.660
\hypertarget{c08:cue:151}{}你現在在拆後就不能再行評。現在問你這件事情的流程講給我聽聽看。你們進考場的時候，你們的卷子是不是有一塊一塊蓋起來，考生的名字那塊看不見。對不對？那在閱卷的過程中也看不見。但是最後都改完了。算完分數了。要不要拆開啊？不然怎麼計分？誰計分？那個誰要跟那個分數對起來啊？這時候會拆開。拆開以後，才公布放榜。所以現在是你們放榜以後才去申請，
% cue 152 5477.660 5517.110
\hypertarget{c08:cue:152}{}事件複雜，這樣子程序清楚了。所以他說，只要你彌封開拆了，那原則上你就不可以要求再行評閱。除非，下面這句話解釋給我聽。什麼叫依形式觀察，即顯然發現有顯然之錯誤。舉例子，我聽聽看，這個事情你們不應該不懂。你們從小到大考試這麼多年，我講一件跟你們切身相關的事情，對你們有這麼困難，就可見我們這個教育把你的頭腦變成什麼頭腦？

\tc{1:31:57}
% cue 153 5517.110 5552.970
\hypertarget{c08:cue:153}{}分數的計算。好，那你講個道理。說我聽聽看，幾個例子，怎麼計算，怎麼計算，有幾種情況，慢慢講，我聽聽看，你說，我讓你講來，對，我沒有阻止你，表示這個方向是對的。比如說評位，他在加總成績的時候，加錯了。對，加錯了。好，加總成績加錯，有幾種可能。譬如說，5 加 5 加 5 加 5，加起來是不是應該等於 15，但是我寫錯了，我寫成 45，可能。先不要講錯，是故意。
% cue 154 5552.970 5587.430
\hypertarget{c08:cue:154}{}但是有沒有這種可能。三個客觀加強應該是 15，現在變成 45。這是不是處形式觀察？我不必看內容，我不必看你五分，對不對？我只看形式，這是不是顯然錯誤。第二種情形最明顯，我現在說每題配分百分之十，這比百分，所以每題都是十分。結果這題是配分十分，我給了你 15 分，這算不算顯然錯誤。四個題目加起來，我漏加一題，這算不算顯然錯誤。我講話有聽不懂嗎？那這種情況，
% cue 155 5587.430 5625.700
\hypertarget{c08:cue:155}{}為什麼？他說可以。除這種情形外，不得。這種情形，那就要再行評。有聽不懂嗎？這種情形是不是顯然違法？我說十分配分，我怎麼可能給你配到 15?5+5+5=15，任何小學生都會算。那怎麼會變成 45？那這種還能夠賣嗎？那如果是這種情況，那你就應該重新評。確認它是 5+5，應該是 15，還是 40？這一題應該是幾分？有聽懂嗎？其他的情形不可以。

\tc{1:33:45}
% cue 156 5625.700 5674.900
\hypertarget{c08:cue:156}{}為什麼？他需要維持考試的客觀公平。注意這句話。為了考試的客觀公平。所以你一公開拆貨，除非有這種情況，一律不能要求再行評。好了，到這裡，有清楚了嗎？好，請問，本件解釋有沒有不同意見？那幾個人寫了，翁岳生、楊日然。這三個都是學界出身的。三個都是臺大。翁先生當大法官 30 年，他那個時候是可以連任的，他要當月照。他總共只有三個意見。

\tc{1:34:34}
% cue 157 5674.900 5748.510
\hypertarget{c08:cue:157}{}一個個別的意見。這個不同意的只有這個。好啦。那我們來看看他講什麼。一、大聲的念啊。這句專屬圈起來。這個評分這件事情沒有別人可以期待的。就是典試委員的權利。繼續，不管你是哪一個理論。繼續，句話。這樣這句話有聽懂嗎？為什麼老師選這個集式？因為這個集式是最能代表考試權。考試權就是出題、閱卷、分的權利。這個權利是專屬於典試委員。
% cue 158 5748.510 5785.080
\hypertarget{c08:cue:158}{}你這個例子做例子。這一題我配分 10 分。你們可以跟我來挑戰說：老師，這個配分太少，應該 20 分。那你來當典試委員。我當典試委員。我說 10 分就 10 分。第二，這個 10 分裡面呢，我認為你答得很好。8 分。另外同學 7 分，還有同學 2 分。拿 2 分的人可不可以跟我挑戰，說老師，為什麼我不是 8 分。那個 8 分的人跟我挑戰說：我應該 9 分還是 10 分。我讀了你的卷子，我閱過你的卷子以後，我綜合的判斷，我認為是這樣子。
% cue 159 5785.080 5841.500
\hypertarget{c08:cue:159}{}就是這樣。這個權利是無專屬的。法院也不能推翻。不然法院就變成典試委員。這是不是最能說明考試權的，什麼比這個更努力。這樣有聽懂了嗎？好，二。所以剛才講的各種情況，客觀上如何如何，這都是違法行為。好，繼續。這就是我們講的訴訟權。訴訟權在保障當你的基本權遭受侵害時，你一定可以找到一個地方，他是公正的可以去向他拯救。好，繼續。因而尊重。

\tc{1:37:21}
% cue 160 5841.500 5896.210
\hypertarget{c08:cue:160}{}但。如，但如，那你們聽懂嗎？你說要尊重他，我不能干涉他行使，但是他的行使是覺得違法了，難道我要認嗎？不是。在這裡你看到每個權力都不自合，你的權力行使不能違法作為前提。你違法，對不起，就有人要來制裁你，要來糾正你。好，繼續。法院，代替評分，我不能替他決定你這題該得幾分。違，但是可以。有沒有程序違法，你看他像什麼？例如，定什麼？錯。

\tc{1:38:16}
% cue 161 5896.210 5916.660
\hypertarget{c08:cue:161}{}法定什麼？對。看你的電腦出了什麼問題，你就這麼相信電腦。我現在要問你，我們對於考試委員有資格要求，但有一個人根本是小學畢業，他能夠當考試委員嗎？
% cue 162 5916.660 5937.870
\hypertarget{c08:cue:162}{}那你有聽懂嗎？回到剛才那一句。你們不是告訴我說，我們去略覺的時候不可以知道這個題是誰出的，那我怎麼知道這個出題人有沒有符合？老師剛剛念的那句，你現在聽不懂我為什麼要念了吧？所以這兩個是互相矛盾的，雞生蛋還是蛋生雞。
% cue 163 5937.870 5972.410
\hypertarget{c08:cue:163}{}你說我們人民有權利知道考試委員有沒有符合法定資格，如果不符合法定資格，這個考試就是違法的，對不對？但是你又說考試委員的名稱，名姓名，要保密，你們要求事件評略，不能要求國蓋這個，這個資訊不得揭露，那不得揭露我怎麼知道他有沒有有資格？從而怎麼認定這個考試處沒有違法？我講話有聽懂嗎？聽不懂的舉手，我是可以再講一遍，那你現在說，
% cue 164 5972.410 5997.270
\hypertarget{c08:cue:164}{}聽懂你就真的要聽懂好這是第一種違法第二種例如對像剛才講的漏了或者算錯了好繼續一期三十分而給分欲三十分你為什麼念字都念的都念漏掉？
% cue 165 5997.270 6012.350
\hypertarget{c08:cue:165}{}怎麼會是啊？聽懂了嗎？這一題配分十分我給了十五分這樣是不是超過十分呢？這些東東都是錯還有或或還有啊，
% cue 166 6012.350 6022.670
\hypertarget{c08:cue:166}{}舉例子這是女生寫的卷就從寬男生寫的我就從嚴。

\tc{1:40:22}
% cue 167 6022.670 6027.780
\hypertarget{c08:cue:167}{}但是有一種情況能不能叫做濫用權利，
% cue 168 6027.780 6031.750
\hypertarget{c08:cue:168}{}以前我去當出題委員 閱卷委員，
% cue 169 6031.750 6033.050
\hypertarget{c08:cue:169}{}我最討厭一種卷子，
% cue 170 6033.050 6035.180
\hypertarget{c08:cue:170}{}不知道該怎麼形容，
% cue 171 6035.180 6040.370
\hypertarget{c08:cue:171}{}那個字寫著密密麻麻都念在一堆，
% cue 172 6040.370 6043.410
\hypertarget{c08:cue:172}{}自然就給他掃一點，
% cue 173 6043.410 6046.920
\hypertarget{c08:cue:173}{}這個算不算的濫用權利，
% cue 174 6046.920 6048.190
\hypertarget{c08:cue:174}{}因為你讓我看不錯，
% cue 175 6048.190 6051.230
\hypertarget{c08:cue:175}{}不知道考試委員是怎麼，
% cue 176 6051.230 6054.260
\hypertarget{c08:cue:176}{}沒有一個簽字筆會發，
% cue 177 6054.260 6056.510
\hypertarget{c08:cue:177}{}現在我們電腦閱卷 電腦閱卷，
% cue 178 6056.510 6059.790
\hypertarget{c08:cue:178}{}你聽到那閱卷是一個大辦公室，
% cue 179 6059.790 6061.470
\hypertarget{c08:cue:179}{}比這個大頭幾倍，
% cue 180 6061.470 6062.970
\hypertarget{c08:cue:180}{}桌子都這麼木頭桌子，
% cue 181 6062.970 6107.990
\hypertarget{c08:cue:181}{}一句句都都都都都都都都都都都都都都都都都都都都都都都都都都都都都都都都都都都都都都都都。因為要求考試員要求我們，閱卷委員要跟著學生的答卷，動。你才會寫成密成一團，我這個點要點到哪？中國是一個取辦考試，簽名的國。我們非常有辦法對付考官。我怎麼去保考官一名？跟著俱動啊。跟著俱動啊。跟著俱動你今天要不要點，不然你亂點，通啊。所以回到剛剛的問題。如果現在學生可以要求要閱卷，即便是以，可以確定你的卷是有未評過的。
% cue 182 6107.990 6153.790
\hypertarget{c08:cue:182}{}但是念東山吹的，完好無處，吹的原地而已。我講話也聽不懂啊。所以我剛才講的，私院面對現在的心情，他不應該這樣顧病之後，他需要解決問題。這些事情在我們大法官都已經告訴你的，不能挑戰，等下來念完就知道，都不能挑戰。好，濫用權利的。好，若拆角的這些，是不是都從形式上就可以判斷？有沒有一個需要重新讀一遍，然後再確認說，這題應該八分，應該六分，有沒有這樣的？

\tc{1:42:33}
% cue 183 6153.790 6193.410
\hypertarget{c08:cue:183}{}這些部分就是考試院月尾卷評分的算數，他說了就算了，來，我們講到這裡講一件事，這個學習考試你們怎麼考我不知道，但是我深表憂鬱，我看你們班上讀書情況非常不良，你們等一下不要來跟我講說，老師，我要申請閱卷啊，要不要給我打這個分數啊，我要來看一下啊，我知道我會主張 319，你們其他老師也跟著有 319 的保護啊，到這裡可以了，然後呢，再往下，你們看一下呢，這件解釋呢，
% cue 184 6193.410 6267.060
\hypertarget{c08:cue:184}{}第四點，他為什麼要發表一部不同意見，這是關鍵之所，你聽到目前為止，聽起來都一樣，那為什麼另外寫個意見呢，第四，說我們見解相同，故，任意啊，是他圈起來，我沒有任意，我是有客觀上合理的理由，我才要求的，唯，關鍵，本件的關鍵，來，借彌封有沒有開拆不是要點，要點是在有無客觀上，明顯違法之處，所以才可以要求，認為，這是多數，
% cue 185 6267.080 6307.690
\hypertarget{c08:cue:185}{}一天，現在這三個人不同意，其實，這樣有聽懂嗎？我要重新評論的時候，是不是要再封起來？所以我一樣不知道考生是誰，開拆以後可以再封回去，再要求重新評，這不行嗎？這個動作有很難嗎？就你們多數人在講什麼？一旦開拆彌封就不能要求了，這不是很奇怪嗎？所以我們三個人不同意，這樣有聽懂了嗎？好了，所以 319 號解釋的細節我就不再說了，總之 319 號是確實、確保、確定，閱卷評分。

\section{釋字第 461 號與參謀總長至立法院備詢}

\tc{1:45:07}
% cue 186 6307.690 6348.780
\hypertarget{c08:cue:186}{}法事委員或者典試委員，具有獨立的判斷預定。他說這樣就這樣，不能挑戰。319 號，來了。好，現在講 461，來。461 在哪裡？461 內主，好。請問 461 的爭點是什麼？用白話講。我練習印法院背前。好，到這裡就好了，到這裡就好了。那我現在請問你，參謀總長是什麼地位，你確到嗎？國防部部長的幕僚長，他是不是總統的幕僚長？有沒有提？在這個解釋裡有沒有提到？有還是沒有？

\tc{1:45:48}
% cue 187 6348.780 6389.060
\hypertarget{c08:cue:187}{}好，這個有點複雜，你要耐心的、專心的聽我講，就快要下課。我們在過去威權時期，蔣中正是代理軍隊出身的。蔣中正一天到晚都要求要反攻大陸，所以他希望繼續維持對軍隊的領導，但是我們憲法沒有這個設計，所以怎麼辦呢？所以當時就發明了一套制度，叫軍令軍政二元化，聽過這個名詞的句子，那聽過這個名詞的人請你們回答我，這怎麼二元化？就是軍令跟軍政參加。
% cue 188 6389.060 6428.510
\hypertarget{c08:cue:188}{}拆開來，簡單講是這樣，還沒有講完，軍政屬於誰，軍令屬於誰，軍令屬於總統，軍政屬於行政院長、國防部長，向立法院負責。這樣有聽懂嗎，參謀總長為什麼重要，那兩條線拆開來，但是有個交匯點，交匯點就是參謀總長，他是軍令系統總統的幕僚長，也是軍政系統行政院長、國防部長的幕僚長，所以這個人重不重要，這個權力最大的時候是在誰，向立法院，這個權力最大的時候是在誰，向立法院，這個權力是聽受到，叫國政參立法院，
% cue 189 6428.510 6477.710
\hypertarget{c08:cue:189}{}立法院以聽民選，定期改選，這時候立法委員不會再當相比賭章，尤其是最有名的立委叫什麼名，後來當總統的，陳水扁，陳水扁最有名的動作是什麼，拿一本六法全書，殺到哪裡殺一下，好聽懂，你如果要當立委，你要許阿別，那你當立委是有一套辦法，他告訴所有的記者，他叫過末雅，老闆今天會有動作，所以老闆進場的時候，攝影機都打開了，然後上臺他質詢，質詢了，你看。

\tc{1:47:57}
% cue 190 6477.710 6514.710
\hypertarget{c08:cue:190}{}你可以想像，郝柏村的答覆，不滿意啊。這時候郝柏村平是行政院長，他不滿意啊，三字經丟過去，怎樣，這樣才能一舉成名啊。我不是要教你們做壞事，他自在後面有很多的報應，那是他的事情，我不管。我只是講說，但當時的立法院的氛圍是這樣的，現在這個事情是發生在郝柏村當參謀總長，還沒當行政院長之前，大家都說，你是藏鏡人，你手上會有最大的權力，總統要經過你，國防部長行政院長要經過你，
% cue 191 6514.710 6527.380
\hypertarget{c08:cue:191}{}你手上會有 叫你來，你都不來，那你覺得這樣可以吧。好，大法官怎麼解釋，來，第一段，第年了。

\tc{1:48:47}
% cue 192 6527.380 6548.910
\hypertarget{c08:cue:192}{}判定是第九次，第七、第四次憲法增修，剛才問了 86 年第 997 年 7 月 21 公布的，好，增修條文第 3 條，第 2 項，第 1 款，規定，繼續。

\tc{1:49:08}
\speaker{學生}
% cue 193 6548.910 6566.040
\hypertarget{c08:cue:193}{}議員：向立法院提出施政方針及施政報告。\par\speaker{老師}好，這句話，原來憲法第 17 條，第 1 款，一模一樣，一個字都沒改，跳過這一句。OK，立法委員在開會時，由？

\tc{1:49:26}
\speaker{學生}
% cue 194 6566.040 6589.330
\hypertarget{c08:cue:194}{}議員：行政院院長及部會首長。\par\speaker{老師}對，此為？\par\speaker{學生}議員：憲法是民意政治、責任政治主義。\par\speaker{老師}好，你看，這時候憲法都已經修改，現在是總統制了，他還說，基於民意政治、責任政治所為的制度性設計，他還沉醉在當年的行政院。

\tc{1:49:49}
% cue 195 6589.330 6628.250
\hypertarget{c08:cue:195}{}當年的行政向立法，立法代表民意，所以基於民意政治、責任政治，所以你要來立法院報告。現在已經不是這樣了，他還用同樣的話，那這句話要怎麼理解？現在民意在哪裡？立法院也有民意，總統也有民意，都是人民權。人民要求總統負責，可以罷免總統。人民要立法院負責，人民可以罷免立法院。所以這條線是平行線，跟以前是不一樣的。我現在講的這個圖形，你去看老師的 YouTube，在 113 憲判 9,YouTube 的第一講，老師有畫來的圖。
% cue 196 6600.230 6638.140
\hypertarget{c08:cue:196}{}那這句話要怎麼理解？現在民意在哪裡？立法院有民意，總統也有民意，都是人民權。人民要求總統負責，可以罷免總統。人民要立法院負責，人民可以罷免立法院。所以這條線是平行線，跟以前是不一樣的。我現在講的這個圖形，你去看老師的 youtube，在 113 憲判 9 youtube 的第一講，老師有畫來圖，比較內閣制跟總統制，裡面所謂民意政治跟責任政治，究竟是在說什麼。兩個民是交同同樣的罰。
% cue 197 6638.160 6666.250
\hypertarget{c08:cue:197}{}這個八個字，他講的是不一樣的意思。我真的忍耐，我的大法官太神奇了。這八個字像魔術一樣，誰來都可以用。你相不相信，有一件變成半總統制，他一樣用這八個字。但意思一樣不一樣不一樣。好，繼續，國防部。

\tc{1:51:06}
\speaker{學生}
% cue 198 6666.250 6672.010
\hypertarget{c08:cue:198}{}關於國防是有其他政策的國防無關。自得大聲點。
% cue 199 6672.010 6675.710
\hypertarget{c08:cue:199}{}向行政院院長、國防部部長的請求。

\tc{1:51:15}
\speaker{老師}
% cue 200 6675.710 6690.370
\hypertarget{c08:cue:200}{}好，到這裡應該很清楚嗎？至於。

\tc{1:51:30}
% cue 201 6690.370 6711.250
\hypertarget{c08:cue:201}{}好，這句話講到這裡，什麼意思？用你的白話講。參謀總長要不要來院會備詢？要還是不要？要！他哪裡說要？自無上開條？自使用他不要來啊！他還告訴你理由，你還沒資格來，你只是幕僚長，你不是部會首長。你怎麼還說要？
% cue 202 6711.250 6752.750
\hypertarget{c08:cue:202}{}譬如，你自己看看你的中文，這句話你看不懂，你應該不是臺灣人。你給八十字道路開來，你深入湯老師的課。不值得滲透，我為什麼要滲透？啊！不是這個意思哦！他的意思，你看懂哦！他的意思說，沒有上開條文，就是你沒有資格用增修條文第三、第二項、第一項規定，到立法院的院會報告。理由，我們只對部會首長，你不是部會首長，你是幕僚長，你還不夠格。再問一個問題。

\tc{1:52:32}
% cue 203 6752.750 6785.850
\hypertarget{c08:cue:203}{}你們認為大法官這樣解釋，對還是不對？為什麼對？為什麼不對？講理由又定。站著為什麼？為什麼覺得是對？為什麼？這跟我剛才講的軍政、軍令、二元化有關。我問的問題，每個都是環環相扣。我沒有浪費你們的時間講些不相關的事情。在哪種情況下，你們不聽你們的叫什麼？在什麼情況下？

\tc{1:53:05}
\speaker{學生}
% cue 204 6785.850 6798.970
\hypertarget{c08:cue:204}{}憲法規定的部會首長、國防部也是常任一個部長和幕僚長，有責任再大一點，但也不應該去接受這一次。

\tc{1:53:18}
\speaker{老師}
% cue 205 6798.970 6823.710
\hypertarget{c08:cue:205}{}對，好，我可以簡單的回應你剛才的講法。等一下我們會看下一段，我們先看這一段。他至少確保了國防部的部長才是幕僚首長。你的參謀總長作為一個藏鏡人，我已經給你定性，叫做幕僚長。幕僚長你再怎麼強勢，你不可能成為部會首長。
% cue 206 6823.710 6861.690
\hypertarget{c08:cue:206}{}你沒有辦法到立法院的會學院會。我講這個話的意思，你們該聽懂，在當時軍政軍令二元化，長時期，因為威權時期，兩個威權強人的領導。這個時候你發現，參謀總長地位非常突出，國防部長很弱。如果大法官不做這個解釋，然後說參謀總長也可以來報告，那我告訴你以後都沒有任何國防部長，你是 Gay，我沒有羅友，我問那個真正就好。這計畫至少懂路了，國防部長作為，
% cue 207 6861.690 6881.020
\hypertarget{c08:cue:207}{}部會首長的，你無法造次，這是第一步。但第二步問題沒解決，剛才我問的問題是什麼，參謀總長要不要來，那現在你給一半的答案，你說院會不要來，院會不要來，就永遠不要來嗎？No。請念第二段。

\tc{1:54:58}
% cue 208 6898.890 6924.750
\hypertarget{c08:cue:208}{}哦，到這裡你聽懂了，院會不必來，誒，我沒有說你委員會也不要來，委員會是政府人員，還有社會上有關的人，我都可以邀請你來備詢，你們有沒有去過立法院備詢呢？還沒有啊，我下次請往上邀請你，你只要告訴他，我手上有票，就找你，來，建築。

\tc{1:55:24}
% cue 209 6924.750 6939.300
\hypertarget{c08:cue:209}{}建國行政院對立法院負責，是不是？行政院對立法院負責。好，我請問一下，419 號是什麼時候做的？
% cue 210 6939.300 6971.940
\hypertarget{c08:cue:210}{}27 年 7 月 24 號，修憲了沒？97 修憲了沒？修了！修了一年多了！我們大法官還在睡覺，還在用內閣制的，找問他講話。還在講行政對立法負責。我也真不知道該說什麼。可見了，李登輝真是很厲害。他把你們耍得團團轉，大法官也一樣跳脫不了他的手掌心。好，繼續。

\tc{1:56:20}
% cue 211 6980.290 7019.550
\hypertarget{c08:cue:211}{}除非，有應邀，好，有應邀，好，那現在講參謀總長，繼續。好，所以到這裡意思已經很清楚。你是不是部會首長，所以院會備詢你沒份，但是你是政府人員，委員會要你來，你是跑不掉的。
% cue 212 7019.550 7071.680
\hypertarget{c08:cue:212}{}你會叫你來，你可以不來嗎？除非，比方說，我問你，我們現在 F16 到底有幾架？同事人，可以隨時起飛？有多少？百分之多少？講！換向 16，幻象 2000，現在還能飛？不能飛？這話你會回答嗎？現在講話你都聽不懂了，沒關係啦，這一天算贈的事很怪就來了，證明裡面叫做 Day of Reckoning，來，不得，他不能，不可以來，唯！到這裡，問題已經全部結束了，但是，
% cue 213 7071.680 7116.120
\hypertarget{c08:cue:213}{}本案最重要不是這一句，是下一句，好下一句，至於，我們有問你嗎？本案有問這個嗎？沒有啊！大法官就有這個好習慣，他喜歡免費奉獸，他說，至於這三院院長呢？繼續，本於五院間互相互尊重，職務令長，政慣例，互受獨，消停備。這句話重不重要，我們這樣問參謀總長，他也回答我是另外三院院長，然後你的理由也非常有趣，權力分立，除了相當中，
% cue 214 7116.120 7151.160
\hypertarget{c08:cue:214}{}相互制衡，還有相互尊重，這不是我說的，這大法官說的，你不喜歡也不能否定。第二，依循憲憲政慣例，我們過去到現在為止，這三院院長有到立法院備詢過嗎？如果他們來立法院備詢，那請問司法考試監察如何獨立性實質？憲法只有說行政院向立法院負責，有沒有說司法建長考試三院向立法院負責？沒有吧！所以我已經跟你們講了，五權分立裡面，表面上該是五個，實際上最重要的就是這兩個，他們兩個關係，決定的是總統之內科。

\tc{1:59:11}
% cue 215 7151.160 7194.040
\hypertarget{c08:cue:215}{}其他三個都是依法、獨立行政院，那現在就告訴你這三個呢，不必到立法院負責。他講的是得。那我請問你，如果有一天，有一個司法院長，他說我要推動司法改革，我現在要向立法院爭取入憑，有史以來最龐大的司法預算，我願意自動、主動到立法院來進行專案報告，並且接受質詢，問可可。真的嗎？注意看這個字啊，得不受邀，你可以拒絕不來。我沒有說，你不來。
% cue 216 7194.040 7243.200
\hypertarget{c08:cue:216}{}我沒有說，你陰不受邀。那我更沒有說，你不可以主動來啊。我現在講的，是把你們當成聰明的人在訓練。因為我覺得你們當中會有立委。但是如果你不想當立委，我講話你當放屁，就當中不知道，這套法式的課是很靈活。我們的立委，懂不懂憲法。如果我是彭博昌，我可以做的事情可多了。那現在最有興趣就開始了。好啦，往下，繼續啊，還有一句啊。比方說，三院都有秘書長，都有處長，都有什麼長，對不對？於，
% cue 217 7243.200 7279.520
\hypertarget{c08:cue:217}{}這是一個案例，有法律上的案設計，應有上開憲法規定之使用。應有上開規定之使用，表示他也要到委員會來備去，因為他們不是依法獨立行使，所以你現在是不是看到司法院的秘書長、司法院的廳署長，常常被司法委員會、司法及法制委員會質詢的尷尬，有沒有？所以你們要看得懂，不是要把它當作笑話看，你們一定要理解，我們是一個民主國家，而且是一個憲政國家，我們這些，
% cue 218 7279.520 7297.260
\hypertarget{c08:cue:218}{}這些規定是有規則的，它這規則的設計是有道理的。好了，461 號講完了，現在我已經改，一路追改，只剩下到 632，好，因為它不認真了，所以我們下禮拜見。好，單早有司庭要宣布嗎？

\tc{2:01:37}
\speaker{學生}
% cue 219 7297.260 7316.350
\hypertarget{c08:cue:219}{}點名，數字是 4490。好，剛才老師講的話要記得的話，我們是哪一天要下禮拜，對不對？下禮拜三嘛，對不對？新書發表會是不是下禮拜三？下禮拜，下禮拜，下禮拜，加油。
% cue 220 7316.350 7318.600
\hypertarget{c08:cue:220}{}下禮拜，下禮拜，哦，下禮拜，加油。
% cue 221 7318.600 7321.530
\hypertarget{c08:cue:221}{}4490
% cue 222 7321.530 7346.120
\hypertarget{c08:cue:222}{}唉，我跟妳講，妳的柯克丁學校講一下，他暫時不要把人跟他表演，當然覺得必須要走。
% cue 223 7346.120 7382.670
\hypertarget{c08:cue:223}{}每天，我們跟他表演，因為土耀正的變成形。
% cue 224 7382.670 7492.220
\hypertarget{c08:cue:224}{}因為，在柯克丁學校，他在柯克丁學校，他在柯克丁學校，他在柯克丁學校講一下，他暫時不要把人跟他表演，當然覺得必須要走。

% 自 Notizen/W11_1123/ 各表由 Skripte/gen_tex.py 產生；改字改表後重產，不手改本檔。
\chapter{釋字第 627 號與總統刑事豁免權（2026-11-23）}

2026 年 11 月 23 日下午課堂錄音（約 105 分鐘）之逐字稿。左邊界為段首時間碼；未標說話人者為湯德宗老師。


\section{助理招募與課程影片字幕校對}

\tc{0:00:00}
% cue 1 0.770 21.230
\hypertarget{c11:cue:1}{}我們不想跟人家分享。好，算了。那剛才舉手的這些人下了課，如果你們有興趣，想要跟老師做研究，幫老師當助理，我現在要養成助理班。我現在缺少一些工作的人，但是如果你有興趣，再來個別跟我講。好不好？結束了嗎？
% cue 2 21.230 55.790
\hypertarget{c11:cue:2}{}當老師做研究，有一個條件，四個字。打死不退。沒有這個決心的，就不要來了。OK，講完了。打死不退是什麼意思？老師，我能吃苦。我能學習。我能符合你的要求。說不上，絕對要找我。中文化叫做無誠誤實。沒有誠意，不要在死痛。我不勉強。好啦，聽懂啦。等一下下課見啦。有人，我有重人。另外一個事情是，一開始老師就講過，嶽老師用 YouTube。現在 YouTube 他們小組的人告訴我說，
% cue 3 55.790 89.150
\hypertarget{c11:cue:3}{}我們還蠻多人的，但是每天都要培養新血。如果你們對於怎麼樣錄影，怎麼樣剪輯影片，怎麼樣上字幕，怎麼樣上字卡，所有的動作你們大概還不行，叫你們做校正。你們中文字看得寫出來，那個字幕打錯了，你會叫得出來嗎？那看你們中文，看你們的法律跟作文。我覺得每一集都有一些錯字。我沒有這麼大的能力去打到這麼細純。但是我們的讀者會三、五、四、寫清楚。
% cue 4 89.150 123.150
\hypertarget{c11:cue:4}{}我看老師寫錯了。老師建議你換一個小編。他不知道我有五個小編。哪有辦法管這麼多事情啊。所以這是你們的國學人問題啊。老師講真的講出來的知識錯誤。看了知識，他會發現要他改過來。如果有人想要做這個事情，歡迎來報名。好了，講太多了。好了，到了現在 633 結束了。來，我們今天進步 627。然後下一個，下一組，上一個禮拜我們有講過說，我要你們做 114 憲。

\section{114 年憲判字第 1 號報告與期末考試、上課心得}

\tc{0:02:03}
% cue 5 123.150 155.830
\hypertarget{c11:cue:5}{}114 憲那一組在哪裡？114 憲那你還沒講完啊？我現在是在唱一個大高歌啊？是不是穿黃衣服妹妹？是不是你？是吧？你看，老師就憑座位來找人了。我沒有看眼鏡，你對我有一段距離，我不是很清楚。但是你要準備你的 114 憲那裡，老師說了，沒有結束。講過了嘛，對不對？你們不要來測驗一個老頭子的記憶。啊，我對於我的公式，我是記得很清楚。然後班長要麻煩你提醒我一件事啊，下 12 月開始，
% cue 6 155.830 188.270
\hypertarget{c11:cue:6}{}用一月考試嘛，12 月開始，你跟老師講，你發信給我講，你說老師要準備題，我們要開始。有禮拜公告一題，禮拜公告。那有一題，我現在就可以公告。跟去年一樣，我要你們寫這個學期的上課心得。你有什麼心得？這個心得是準備，如果老師覺得這個心得很好，寫得很出色，我要把它傳給你們的學弟妹，作為下學期上課的參考。你們有看到上學期學長的心得嗎？沒有。

\tc{0:03:08}
% cue 7 188.270 222.410
\hypertarget{c11:cue:7}{}再問了一個人有嗎？只有一個人。請問你在哪裡看到？那是不是發的太慢了？你是不是建議老師要放在一個有效的地方？也就是在你們選課之前，需要能夠看到上個學期的同學選課的心得，以便我決定這個學期要不要選啊。我現在覺得這個學校的系統，我不是很清楚電容的人，但是我學的商話，我的助理不夠強。你們一定有辦法找到一個地方，讓大家都可以看得到。為什麼是要老師來傷腦筋呢？我上課，我的腦筋只要放在我的課上。
% cue 8 222.410 253.650
\hypertarget{c11:cue:8}{}課程裡面。課程外面的事情是你們去講的。口耳相傳，還是口耳相罵，你去傳跟罵，我不去介入。因為這是你們的言論自由，也是你們的評審。但是不管怎樣，這是應該可以給下一屆的學弟做參考。我知道動物很多混的人，這個腦袋裡面都要在小時候。每天活工都要上學，活工都要回家。然後這樣就過了五年，然後告訴我說，我的人生就這樣結束了。是的，那如果你覺得不要到時候後悔，那你就要聽我的。

\tc{0:04:13}
% cue 9 253.650 271.580
\hypertarget{c11:cue:9}{}他聽說，這個老師上課很兇耶，好難及格啊，以兒傳合一大堆。每次我都會跟你們講，少學期成績如何，分不了，哪有這種情形。哪有重要的就是我們緣分不足，我們不必別講。

\section{釋字第 627 號之案件事實：國務機要費案}

\tc{0:04:31}
% cue 10 271.580 287.430
\hypertarget{c11:cue:10}{}好了，第一，事先安一，待著。好，來，627 開始。請問本案事實是學格不可以說明一下，是一定從事實開始才有什麼好玩意的啦。請問你們這個案件，司法院的官網，有沒有事實？
% cue 11 287.430 324.890
\hypertarget{c11:cue:11}{}是神山藥？有，那你剛講有的多學。大聲音給大家聽。你們就不得慢慢聽，你就聽懂吧啦，這就不是本來。來，大聲一點。\par\speaker{學生}臺灣高等高檢署。查緝黑金中心。檢察官。\par\speaker{老師}這個人是老師的學弟，叫陳瑞仁。這個人非常有勇氣。一個人敢。現在的檢察官要找這種人，很少啦。現在的檢察官只要靠邊站，就升官發財。他既不想升官，又不想發財。他對付的是總統，以及總統夫人。
% cue 12 324.890 363.770
\hypertarget{c11:cue:12}{}總統夫人有幾個。好，繼續。就，申請人是本案的申請人。聽懂？我們現在是 627。627 是第一個有史以來由總統本人聲請憲法解釋來。NEVER。OK，繼續。好，所以本本案牽涉到的是國務機要費。國務機要費已經對你們來講，時間很早了。你們大概還沒出生吧？還沒出生，已經出生了，還沒有，對不對？所以我現在給你們講個故事。看。

\tc{0:06:03}
% cue 13 363.770 399.740
\hypertarget{c11:cue:13}{}我們中華民國偉大的政府要不要教你，因為你的歷史課本已經改善他糊塗了。中國史不教、外國史不教、臺灣史也不教。這個是臺灣史。而且最近發誓呢？他當然沒有要教。因為講起來很難。陳推銘的國務機要費有非常大的爭議。這裡面有很多筆。其中有一個我記得印象最清楚的。被我們當時的審計長請問。在座的人有聽過審計長三個字的請舉手。好，那能不能找一個條文，
% cue 14 399.740 435.140
\hypertarget{c11:cue:14}{}在憲法裡告訴我審計長是屬於什麼憲法機關的。他應該如何產生。老師這樣上課是不是對你們有一點幫助。我們不要講一些玄的東西，我們講有點實在的。審計長你們都聽過了。好，請說。\par\speaker{學生}憲法 104 條。憲法 104 條。憲法審計長由總統提名。經立法院同意任命。\par\speaker{老師}好，很好。這樣子都聽懂了嗎？那審計長是幹什麼的？他就是在審查每一個機關包含立法院，
% cue 15 435.140 466.440
\hypertarget{c11:cue:15}{}他的預算是怎麼開支的。你要說我說這個預算我賣 100 億，要買 F16、34 架。結果我把它買去拉法葉戰艦或者我把它拿去蓋太陽能光電板。然後大家喜哈哈哈哈的把錢都分光。那你覺得這個賬不能爆嗎？那每一個人都會列預算來誆老百姓啊。哎我現在需要多少錢不夠？還要追加預算。如果你是民主國家，你是人民，你說我是主人。我派一個代表叫立法委員、立法院。
% cue 16 466.440 499.700
\hypertarget{c11:cue:16}{}你建讀好。立法委員好像也是。好，我一定幫你監督主人我聽話。你們會覺得這些僕人可靠啊。他們對後衛沆瀉一切。行政院說我多編了三十億，是我綠色的，我三十億也分你十億吧。不可以。你們忘記了，行政院跟立法院有黨派的因素可以串成一家人。行政院不一定會對付立法院。要搬小英進去裡面是怎麼產生的？只有小英不想通過的法律，沒有小英通過的法律。這句話怎麼會產生的？

\tc{0:08:19}
% cue 17 499.700 534.030
\hypertarget{c11:cue:17}{}你把新家人都是綠色的、綠色的、執政的都是綠色的，現在不是。所以現在都被拼了。那以前為什麼沒有？四家人吵什麼吵？我司機執政權就不能分了。你看一些武界在臺南，現在大和多新聞。他今天又出來看記者會。有個很大，有很大三條啊。他說做不到要如何如何如何啊？我跟你們講話，你們是聽懂還是聽不懂？沒關係，我知道你們不想當主人，你想順起來。聊得又就好。

\tc{0:08:54}
% cue 18 534.030 566.550
\hypertarget{c11:cue:18}{}反正這個國家我已經說過了。對老師這輩人來講，我們經歷了國家的黑暗。從戒嚴到解嚴。從我們含窮、含窮，我們小時候的窮是你們聽不懂的。你得問爺爺那輩，看他會不會跟你講。我猜想你跟你們爺爺，大概也不太講話。更不會問他這麼難過。那如果你爺爺還有經歷，他更得跟你說一點。我沒有吃虧。我以前沒有手機，我現在有手機，我有雲板。我可以說，所有的世界上的奇怪的事情，我都經歷了。
% cue 19 566.550 600.860
\hypertarget{c11:cue:19}{}我都經歷過。剩下的是你們，你們現在的國家負債一大堆。你們現在看到這個市場每天變化這麼大。全球就不只有臺灣。臺灣今天也正前，也跌得很慘。你們自己想一想吧，這個國家交到你們的手裡是一個什麼局面。你繼續用這種態度當國家主人，那你最後就是，怎麼死呢？不要說我老師沒有跟你講，都不知道。這是很現實的問題。每一個當國家主人的人，一定要有權利者或者是主權者的意識。你要 sense 到，或是，我都不知道。
% cue 20 600.860 636.770
\hypertarget{c11:cue:20}{}哦，我是國家主人，你要的 sense，你沒有得到我同意，而你到底要不要安。你都沒有這種感覺。那個在商店的人，不管那叫什麼名，他一定欺負，他一定互動，一定玩弄你，把你耍的團團轉，所有的媒體，單一的聲音，看你怎麼樣。那你聽到的消息全部是裡面，你能知道什麼呢。好吧，講太多了。來，繼續，所以越來我說呢，最重要的一個意思，我記得的是，陳先生遠拿國務機要費去買他孫子的尿布。

\tc{0:10:36}
% cue 21 636.770 671.410
\hypertarget{c11:cue:21}{}他審計長報出來的，他不準，他有大概一千還是兩千萬左右的金額，不予核銷。這個按照我們現在的審計法的規定，如果這個機關的預算，國務機要費是編在總統府的預算，如果審計部不予核銷，會產生什麼後果？第一，機關首長跟機關兩個人要連帶負賠償責任。這個兩千萬，講兩千萬比較好記憶。誰出武？總統跟總統不要出武，相關人員還要議處。

\tc{0:11:11}
% cue 22 671.410 709.260
\hypertarget{c11:cue:22}{}你所審計長做不重要。美國的審計長任期多久，你們猜猜看，這是憲政常識。美國除了聯邦的法官是終身職，他幾乎憲法跟他病，就為他死掉了。其實第二個懲命的人，那就是審計長。請問你審計長按照美國的法律規定任期多久？15 年，夠不夠久？這樣聽懂嗎？有聽懂嗎？很好。那你去對比一下臺灣跟美國，你就知道你在講什麼。審計長這個職務，為什麼要這樣保障他？
% cue 23 709.260 743.090
\hypertarget{c11:cue:23}{}因為他非常重要。他是我們的 watchdog。中文的話叫做看門狗。我們要他看緊這些政府計劃，也不要亂花錢了。有沒有拿國務機要費，號稱叫國務機要，結果去給孫子買藥物的？那你覺得這樣可以嗎？還是你 care？如果你不 care，那我就不要討論了。反正所謂，都不 care，有司主人都不 care，那有什麼關係呢？好，就不要來幫我抱怨說我們年輕人都買不起房子。為了錢給人家用掉，當然買不起啦。好啦，繼續。

\tc{0:12:23}
\speaker{學生}
% cue 24 743.090 794.480
\hypertarget{c11:cue:24}{}有關於國務機要費的使用問題。\par\speaker{老師}好，有一個檢察官在黑金、查緝、黑金中心去分案偵查了。好，繼續。申請人是總統哦，他對總統下手了。好，後來，是，後來，所以有沒有起訴，偵查以後起訴了。起訴的人是吳淑珍和陳水扁，兩個人列為共同貪汙的共同正犯。但是陳水扁因為是憲法 52 條，
% cue 25 794.480 831.010
\hypertarget{c11:cue:25}{}所保障的總統豁免權。所以不敢列拉，就把他太太列了。其實講起來很有趣，你要兩個人才能整理了犯罪，所以叫共同正犯。他有一個不能列，就列另外一個。照理講，那個不能列的，會不會影響這個可以列？這個有解釋的問題。我問你，吳淑珍何德何能去看？如果他老公不是總統。所以真正的關鍵是在總統啊。人家是看了總統面子送給吳淑珍錢，有那麼笨的人喜歡吳淑珍去送給他錢的嗎？那真正的這個最痛的這個，
% cue 26 831.120 872.870
\hypertarget{c11:cue:26}{}他不能列，那你居然把另外一個跟他列成共同罪犯，這是可以討論的。但我不去車子的這個事情，你可以說這個檢察官已經用啦，要不怕，要已經做了，來本案繼續。所以案件呢，進入臺北地院審理啦。我們這個，我們的法官、我們的書記官寫這個案例事實，就會寫到那條文謅謅。我請問你字號有什麼重要？講個故事，你會講嗎？所以難怪人一看到法律出來，你給他關掉。

\tc{0:14:32}
% cue 27 872.870 906.350
\hypertarget{c11:cue:27}{}他電視器關掉，又在講官話。這是我們老師訓練我們法律系的學生要負最大的責任。老師上課從來不努力，你們這樣講話。你如果這樣講話，你就是得不到人家人民的欣賞。你公開完全的意思嗎？什麼字號？這個字號多難背。誰有辦法進出？我怎麼在意字號？你怎樣一句話說？本案進入臺北地方法院審理。不就好了嗎？他請出了嘛，要不審理？審理是不是到地略？不然直接跳高院嗎？還是跳最高院嗎？好，繼續。

\tc{0:15:06}
% cue 28 906.350 965.290
\hypertarget{c11:cue:28}{}申請人，這段話簡單地講，就是聲請意旨。陳水扁權為什麼來申請？他認為這是違憲的。現場就告訴你，羅勿社會條享有或民權，你還把我跟我太太認為共同對方，你討論壞了。我討論壞了，你討論壞了。所以我要來告。所以陳水扁做的事情是很有趣的，他是法律人，他認為他是義症的名詞，他要來告。那也不在乎以前沒有限令，沒有限令怎麼樣，總統就要告。好，那就來告了。繼續，用。

\tc{0:16:11}
% cue 29 971.820 1005.380
\hypertarget{c11:cue:29}{}好，所以這句話很重要。這句話為什麼很重要？這句話是說本案呢，好像是爭議，而不是一個疑義。假如陳水扁那裡他認為我不可以列為共同正犯，叫做疑義的話，現在這裡已經是爭議了。因為呢，臺北地院傳他來問事情。這個等一下我們就會看到，臺北地院是以證人的身份傳他來作證。
% cue 30 1005.380 1036.140
\hypertarget{c11:cue:30}{}做證，你太太做的事情怎樣怎樣？我問你，你為答，你是證人，這是刑案。刑案如果我傳你作證，證人是有違作證之義務，民事案件是沒有，你可以不去。刑案非得也來不可。因為刑案我們認為是發現犯罪的真實，它的公益性比較高，比民事訴訟判決是不是敗訴，是誰贏，誰多財產的問題不一樣。好，所以他就去了。去了以後呢，問他，可是這個總統說，對，
% cue 31 1036.140 1070.610
\hypertarget{c11:cue:31}{}我有拒絕回答他的權利，他說這是國家機密特權，這幾個字你不要相信他，他是不是真的這樣子，你要去查他的身體數，換出他牆，陳水扁，除非委任了非常高明的律師，不然他不一定知道什麼叫國家機密特權，這是一個專有的意思，我們憲法裡面沒有這個字，我相信在文靠解釋之前，我們的訴訟也沒有國家機密特權這個玩意兒，國家機密保，我們有國家機密保，但是我沒有特權這兩字，什麼叫國家機密特權？

\tc{0:17:50}
% cue 32 1070.610 1102.460
\hypertarget{c11:cue:32}{}我就光講這熱個字給你唸聽，如果我除非你今天這幾個有看，其他人應該單看國家機密特權也看不懂，不知道他什麼意思，好，他是這麼複雜，所以你要問我，我說我有特權，所以呢，有趣哦，這個是大法官界定的，大法官說這件事情，臺北地方法院和陳水扁總統兩個人發生爭議，很好，那我請問你，當你發生了這種爭議，現在這個算是我們大法官審理案件法，
% cue 33 1102.460 1140.710
\hypertarget{c11:cue:33}{}現在的憲法訴訟法也有的，這個算是大法官審判權當中的哪一個類型，這個像 520，憲法機關行使職權、適用憲法，與其他憲法機關行使職權、適用憲法發生爭議。我到現在為止講話，你如果聽不懂，你跟不上，讓我只能跟你講，這個學期還能夠用過。你們碰到老師，在你這個疑問當中大概就這一次，那你這一次錯過了，你以後再要後悔，那是你 50 歲的時候，不愛一起。

\tc{0:19:00}
% cue 34 1140.710 1174.590
\hypertarget{c11:cue:34}{}我們在 520，我一直跟你強調，這是一個爭議，因為這裡會具體的事件。然後呢，行政院說，我說停建就停建，立法院說 Junton Model，不是你說停建，是要我同意才能停建。是不是針對同一件事情核四要不要停建，或者當時具體的話說，停止滋用、停止使用、核四預算。這件事情兩個機關，都認為有憲法上的確率，做好要做的確率，但是兩個人間相處，所以這是爭議。這跟一個你法律條文看不懂。
% cue 35 1174.590 1210.910
\hypertarget{c11:cue:35}{}或者看到一個法律，你覺得他有沒有違憲，你自己有懷疑，沒有人跟你有爭議啊，誰跟你爭嗎？那不一樣，本案呢，如果照他的說法，那是很心靜的。你什麼時候看過法院行使職權和其他政府機關，包含總統產生爭議呢？什麼時候看過？你決定怎麼平平大？這個不知道誰寫，也許是申請人寫的，也許是我們大法官書記書寫的，這個事實，這樣子描述這件事情和偉大人。我們通常，你到法院去，
% cue 36 1210.930 1246.960
\hypertarget{c11:cue:36}{}你跟法院意見不一樣，你可以說法院跟我發生爭議，那我們要進行憲法訴訟，還是法院說我說了算，該裁定的，該裁定，該裁定的，不服去上訴，只不過我跟你爭，老子跟你爭嘛，我站在上面，底下的，但是本來不一樣，是總統被傳去要作證，他叫他講，他說我不講，誒有趣啊，這樣能不能構成那兩個憲法機關，行使，適用憲法，產生爭議，你們聽懂嗎？

\tc{0:20:46}
% cue 37 1246.960 1280.550
\hypertarget{c11:cue:37}{}就法院來講可以，因為我在傳訊證人，但是就總統而言，他有沒有這樣行使之，還是你們認為行使之，包含拒絕證言的國家機密特權，這也是我行使之，我講話有聽懂嗎？有聽懂嗎？我希望你真的是有聽懂，我不知道我講話是不是太誇張，但是我講話有層次，我一層一層的講，前面那層聽不懂，會不會這樣聽不懂，但前面聽懂，這一層不見得聽得懂，好啦，那不管怎樣，他認為有爭議，有爭議就有爭議，讓他收完。

\tc{0:21:20}
% cue 38 1280.550 1331.230
\hypertarget{c11:cue:38}{}繼續，原始就是因此啦，繼續，這一句話寫到這裡，你們做這個案子的同學，就要去後面理由書，一定在前面，理由書開始在前面，去找找看，大法官一定受理本案，不然他不會做成解釋，他受理本案的依據跟現在講的依據，是真的機關行使是會發生爭議嗎？如果是的話，那我們就剛才講，法院跟憲法機關行使是會發生爭議，這種案件極為少見。
% cue 39 1331.230 1363.910
\hypertarget{c11:cue:39}{}這是我們案的特色之一。這是認定問題。這是大法官要認定。我如果你們憲法訴訟法也好，當時的大法官、審理案件法也好，你規定大法官只有在下列情況下，下面的人通過這樣的事情，才能來申請解釋。這就是我的審判權。我的範圍就這麼大，法條一條條一條寫好，這裡有五個門還是六個門開在這裡。你要主打根本進來。我上次說過了，一定要先程序後實體。你這個由在為道。
% cue 40 1363.910 1396.650
\hypertarget{c11:cue:40}{}你走錯了門，你沒有找到證券門。對不起，你只能在門外。因為司法卻不告不理，司法不能主動親戚。我不能路見不平拔刀相助，那是檢察官的事。檢察官要起訴誰去調查誰去調查誰去調查誰。行政立法愛歸民服務，隨時可以主動親戚。他要定什麼法，要做什麼事情。這個不需要人民來申請。但是司法不行，司法的權力本質上就是被動的。所以記不記得美國憲法條文，
% cue 41 1396.650 1427.550
\hypertarget{c11:cue:41}{}第三條第一項第一句，說聯邦司法權僅及於個案或爭議。沒有案件沒有爭議就沒有司法權。案件或爭議是人家來招我才有。我不可能主動去抓招。好了，到這裡可以了。案件的事實已經說明完畢了。可以了吧。好，來，現在我們進行到解釋文。來一段一段的念，老師會一段一段的跟你們討論。那我希望你們都已經看懂吧。有看最好，沒有看的話，現在有點來不及了。

\section{解釋文第一段：總統刑事豁免權之範圍}

\tc{0:23:47}
% cue 42 1427.550 1473.160
\hypertarget{c11:cue:42}{}但也沒辦法跟直播這樣了。來，627。所以他先處理第一個議題。總統的刑事豁免權是怎麼樣的事情。好，一段段念，請繼續。好，這是原文的句子。那你再看一下，這個句子，是不是總統享有完全的刑事豁免？沒有哦。你如果犯內亂或外患，不能豁免。上次我們在 388 號解釋有講過嗎？388 可能今天沒有講。你如果內亂外患，
% cue 43 1473.160 1507.430
\hypertarget{c11:cue:43}{}外患，會有事多少條啊？還給你豁免啊？內亂是什麼意思？你要推翻我們的憲政秩序，跟毛澤東亞選舉選書嘛，還借先桌子嘛。賴清德如果這次選書給你先桌子，你覺得可以嗎？這不是典型的內亂罪嗎？或者你說我去勾結中共、勾結美國、勾結日本小帝國，推翻臺灣政府，這可以嗎？這叫外患。那我們選一個總統，我們允許他內亂、外患，他可以主張豁免。這種孩子不能扮他嗎？那我不討他太子了。是不是？
% cue 44 1507.430 1512.810
\hypertarget{c11:cue:44}{}你是國家主人，你會對這種總統家裡容忍，你一定頭把壞掉。所以這個我不答應。
% cue 45 1512.810 1519.580
\hypertarget{c11:cue:45}{}那內亂外患，以外的罪，會有什麼事情，要這麼說。這樣子聽懂了嗎？繼續。

\section{講義檔案之搜尋與復原}

\tc{0:25:19}
% cue 46 1519.580 1548.350
\hypertarget{c11:cue:46}{}稍微等我一下，我的檔案不知道什麼地方出問題。今天做的所有的註記，現在全部不見。
% cue 47 1548.350 1562.270
\hypertarget{c11:cue:47}{}就是有人跟老師開玩笑啦。有人可以幫忙找到這個新檔案嗎？電腦比較好的來找找看。你知道怎麼開我剛才開過的 Word 的檔案嗎？
% cue 48 1562.270 1579.010
\hypertarget{c11:cue:48}{}對，這個是因為我剛開的是一個 Word 檔。Word 檔有一個存儲的系列你知道在哪裡嗎？你知道嗎？最原始的檔案，最原始的檔案。你找給我看吧，在哪裡？
% cue 49 1579.010 1611.240
\hypertarget{c11:cue:49}{}對，你就告訴我啊，在哪裡啊？為什麼會？我為什麼會？不會？因為我不會，我不會去存到雲端的。我沒有這個習慣。雲端呢，很多人都會看到就是這個檔案。是不習慣到人家很小的檔案。是嗎？哪裡？沒有啊？這是我的雲端啊？沒有啊！不是這個，沒有啊！

\tc{0:26:51}
% cue 50 1611.240 1647.930
\hypertarget{c11:cue:50}{}沒有啊！不是在這裡。你們誰的電腦好？誰的 Word 的牆？不是這樣子？我沒有，我沒有放到。我現在只是覺得很奇怪。剛才有一個鈕瓦還是什麼？不小心，它就不見了。那我也不懂為什麼。來，你要，你要說什麼？來，你幫忙找一下。這裡，換個位置。來，請說。我要去找一下。我跟妳講了，通常我們在這裡底下呢，它會有一個地方顯示我們剛才開過的檔案。Word 呢，現在我在 Word 裡啊。
% cue 51 1647.930 1685.420
\hypertarget{c11:cue:51}{}然後呢，等下，讓我看一下這裡。這裡應該有啊。你看這裡有啊。144 個嗎？三分鐘前。應該不是這個。應該是更早的前。那更早的前是我沒有儲存的問題嗎？我在這裡做了很多的註記。我每一段都做了註記。我非常奇怪。你們有沒有人可以告訴老師，這個是為什麼？是有人在老師的電腦裡面做了手腳嗎？那為什麼我做的註記我寫了非常多？
% cue 52 1685.830 1724.040
\hypertarget{c11:cue:52}{}我又沒有把它刪掉。而且我每次都會處處。那你現在如果你說找剛才的記錄，剛才的記錄是應該可以看得到的。那你幫你找下來，我是問你，你找哪個記錄給我。這是在哪裡？在這裡就不會問你是要檔案的。你是要檔案嗎？檔案。那是這個嗎？這是剛剛，並不是這個，你再往下的。沒有啊，沒有，再往下，我在看。這裡沒有 627 啦，只有一次 627。這就是剛才啦，這兩分鐘前啦，兩分鐘前。沒有啊，我住氣的都在我底下，我有。

\tc{0:28:44}
% cue 53 1724.120 1731.530
\hypertarget{c11:cue:53}{}我不是用註解的方法，我是直接寫在上面的，
% cue 54 1731.530 1734.720
\hypertarget{c11:cue:54}{}這跟我們註解不關，我沒有在這裡做註解，
% cue 55 1734.720 1738.100
\hypertarget{c11:cue:55}{}這個是我在背景當中，不曉得出了什麼問題，
% cue 56 1738.100 1748.480
\hypertarget{c11:cue:56}{}那你告訴我在哪裡有這樣的草稿，哪裡有這樣的草稿，
% cue 57 1748.480 1752.370
\hypertarget{c11:cue:57}{}可以告訴我在哪裡，你可以找一找看在什麼地方，
% cue 58 1752.370 1759.910
\hypertarget{c11:cue:58}{}那你找給我看，那是在 Wallet，你在這裡，
% cue 59 1759.910 1768.630
\hypertarget{c11:cue:59}{}那你是要在哪裡去找還原，
% cue 60 1768.630 1771.620
\hypertarget{c11:cue:60}{}是啊，你可以找找看還原，
% cue 61 1771.620 1775.860
\hypertarget{c11:cue:61}{}我不知道他是什麼頁面，
% cue 62 1775.860 1778.750
\hypertarget{c11:cue:62}{}不想要提望是你頁面，
% cue 63 1778.750 1781.020
\hypertarget{c11:cue:63}{}不想你這個頁面在哪裡，
% cue 64 1781.020 1784.230
\hypertarget{c11:cue:64}{}我的打扮給你打扮是你要打，
% cue 65 1784.230 1786.380
\hypertarget{c11:cue:65}{}這個 627 就是剛剛，
% cue 66 1786.380 1788.530
\hypertarget{c11:cue:66}{}這不是不是我講的，
% cue 67 1788.530 1791.160
\hypertarget{c11:cue:67}{}我要有一個更早的 627，
% cue 68 1791.160 1795.460
\hypertarget{c11:cue:68}{}今天老師在辦公室裡面花了一個多小時，
% cue 69 1795.460 1796.460
\hypertarget{c11:cue:69}{}你要你瀏覽看看可能，
% cue 70 1796.460 1801.170
\hypertarget{c11:cue:70}{}627。

\tc{0:30:00}
% cue 71 1800.710 1862.740
\hypertarget{c11:cue:71}{}就是沒有在這裡啊。這裡哪裡有 627。3 點多啊。3 點多。現在已經在上課了。我沒有註解。我就在底下寫。真的註解。我就直接寫在底下。沒有。我有 OneDrive 剛才檢查過了。沒有。我沒有在 OneDrive 裡面。是嗎？我剛才就是為檔案的版本紀錄啊。等一下，等一下。你看嘛。現在檔案嘛。檔案的資訊嘛。這是資訊嘛。對不對。好。然後資訊呢哪裡。
% cue 72 1862.740 1899.520
\hypertarget{c11:cue:72}{}哪一個叫檔案？好。檔案定只能記錄這個嗎？那你說上傳的檔案啊。我沒有上傳，它為什麼會有一個問題？沒有啊。沒有，我沒有。我沒有在那裡啊。我沒有，我沒有，我沒有上傳。我沒有，剛才點啦。我沒有用 OneDrive，所以不會在這裡面認為。那我們可以看一下有 627 嗎？沒有啊。

\tc{0:31:39}
% cue 73 1899.520 1944.030
\hypertarget{c11:cue:73}{}然後，完，資訊，哪邊有關係，這邊，來，沒關係，反正內界我根本沒有在，這個是即時復原，會不會是嗎？我覺得是比較有關係，對啊，下肩下，首先看來是這個還是這個，有兩個，現在要找時間點，這個是下午 3 點 43,3 點 14,2 點 43,
% cue 74 1944.030 1947.420
\hypertarget{c11:cue:74}{}那這個是 3.43，
% cue 75 1947.420 1950.890
\hypertarget{c11:cue:75}{}所以剛才對，
% cue 76 1950.890 1952.310
\hypertarget{c11:cue:76}{}應該是這個，
% cue 77 1952.310 1954.330
\hypertarget{c11:cue:77}{}這個是 3.44，
% cue 78 1954.330 1958.720
\hypertarget{c11:cue:78}{}你交款了，
% cue 79 1958.720 1960.520
\hypertarget{c11:cue:79}{}所以我打開了，
% cue 80 1960.520 1962.760
\hypertarget{c11:cue:80}{}我打開關係上不對啊，
% cue 81 1962.760 1964.320
\hypertarget{c11:cue:81}{}所以不可能是這個，
% cue 82 1964.320 1965.360
\hypertarget{c11:cue:82}{}而是他下一個，
% cue 83 1965.360 1967.940
\hypertarget{c11:cue:83}{}你說這個嗎，
% cue 84 1967.940 1973.600
\hypertarget{c11:cue:84}{}所以這個才對，
% cue 85 1973.600 1976.700
\hypertarget{c11:cue:85}{}好 ok 好可以，
% cue 86 1976.700 1982.700
\hypertarget{c11:cue:86}{}我們再次謝謝，
% cue 87 1982.700 1985.500
\hypertarget{c11:cue:87}{}你等一下再不行，
% cue 88 1985.500 1986.880
\hypertarget{c11:cue:88}{}好 現在，
% cue 89 1986.880 1988.610
\hypertarget{c11:cue:89}{}你大家走到了，
% cue 90 1988.610 1991.810
\hypertarget{c11:cue:90}{}所以老師就變得很清楚，
% cue 91 1991.810 1992.890
\hypertarget{c11:cue:91}{}我變得很熱悶。

\tc{0:33:12}
% cue 92 1992.890 1994.980
\hypertarget{c11:cue:92}{}好 其他人來，
% cue 93 1994.980 1996.620
\hypertarget{c11:cue:93}{}我們剛才在第一段，
% cue 94 1996.620 1998.020
\hypertarget{c11:cue:94}{}第一段已經念完了，
% cue 95 1998.020 1998.880
\hypertarget{c11:cue:95}{}第一句，
% cue 96 1998.880 2000.350
\hypertarget{c11:cue:96}{}好 今天第二句，
% cue 97 2000.350 2006.050
\hypertarget{c11:cue:97}{}第一句剛才不是句號了嗎，
% cue 98 2006.050 2007.200
\hypertarget{c11:cue:98}{}念過了嗎，
% cue 99 2007.200 2009.920
\hypertarget{c11:cue:99}{}你剛才念你你忘記了，
% cue 100 2009.920 2011.710
\hypertarget{c11:cue:100}{}你比老師記性還短暫，
% cue 101 2011.710 2013.890
\hypertarget{c11:cue:101}{}一句念完了嗎，
% cue 102 2013.890 2017.200
\hypertarget{c11:cue:102}{}第二句念完了嗎，
% cue 103 2017.200 2018.830
\hypertarget{c11:cue:103}{}第二句我念到 388 嗎，
% cue 104 2018.830 2019.600
\hypertarget{c11:cue:104}{}第二句念完，
% cue 105 2019.600 2020.700
\hypertarget{c11:cue:105}{}好好好 ok，
% cue 106 2020.700 2022.530
\hypertarget{c11:cue:106}{}現在我的問題，
% cue 107 2022.530 2023.610
\hypertarget{c11:cue:107}{}那你們可以回答，
% cue 108 2023.610 2027.390
\hypertarget{c11:cue:108}{}現在我們規範的這個刑事豁免特權，
% cue 109 2027.390 2029.210
\hypertarget{c11:cue:109}{}那法官第二句說，
% cue 110 2029.210 2033.140
\hypertarget{c11:cue:110}{}這是憲法基於總統作為國家元首，
% cue 111 2033.140 2035.490
\hypertarget{c11:cue:111}{}我們給他的尊崇跟保障，
% cue 112 2035.490 2036.470
\hypertarget{c11:cue:112}{}好。

\tc{0:33:56}
% cue 113 2036.470 2037.690
\hypertarget{c11:cue:113}{}那我的問題是，
% cue 114 2037.690 2040.410
\hypertarget{c11:cue:114}{}如果這句話的意思到現在連凱，
% cue 115 2040.410 2042.330
\hypertarget{c11:cue:115}{}那就是只給國家元首，
% cue 116 2042.330 2044.730
\hypertarget{c11:cue:116}{}那是不是表示給政府元首，
% cue 117 2044.730 2048.990
\hypertarget{c11:cue:117}{}上次講過啦，
% cue 118 2048.990 2050.430
\hypertarget{c11:cue:118}{}一個國家有兩個頭，
% cue 119 2050.430 2054.080
\hypertarget{c11:cue:119}{}總統制的國家元首跟政府元首是同一人，
% cue 120 2054.080 2056.990
\hypertarget{c11:cue:120}{}但是內閣制的國家元首是虛位的總統，
% cue 121 2056.990 2058.940
\hypertarget{c11:cue:121}{}內閣制的政府元首是行政院長，
% cue 122 2058.940 2062.370
\hypertarget{c11:cue:122}{}或者叫總理，
% cue 123 2062.370 2064.770
\hypertarget{c11:cue:123}{}我們現在把 52 條寫的時候是內閣制，
% cue 124 2064.770 2066.320
\hypertarget{c11:cue:124}{}現在已經改成總統制了，
% cue 125 2066.320 2067.860
\hypertarget{c11:cue:125}{}那你要注意這句話，
% cue 126 2067.860 2069.860
\hypertarget{c11:cue:126}{}它說我們給總統保護，
% cue 127 2069.860 2073.170
\hypertarget{c11:cue:127}{}這個總統是把它當成國家元首給它尊榮，
% cue 128 2073.170 2077.090
\hypertarget{c11:cue:128}{}那意思就是不及於行政，
% cue 129 2077.090 2078.590
\hypertarget{c11:cue:129}{}不及於政府元首，
% cue 130 2078.590 2080.850
\hypertarget{c11:cue:130}{}所以 52 條只有總統，
% cue 131 2080.850 2084.500
\hypertarget{c11:cue:131}{}52 條不及於行政院長，
% cue 132 2084.500 2087.410
\hypertarget{c11:cue:132}{}我這樣講應該合理吧，
% cue 133 2087.410 2089.180
\hypertarget{c11:cue:133}{}這是我們看到這個句子，
% cue 134 2089.180 2093.380
\hypertarget{c11:cue:134}{}我們本身對它應該有的請問，
% cue 135 2093.380 2094.740
\hypertarget{c11:cue:135}{}我講話有聽懂嗎？

\tc{0:34:54}
% cue 136 2094.740 2096.710
\hypertarget{c11:cue:136}{}大法官講這個話說，
% cue 137 2096.710 2099.770
\hypertarget{c11:cue:137}{}大法官應該跟我們同樣的請問，
% cue 138 2099.770 2104.790
\hypertarget{c11:cue:138}{}那你的意思是不是真的這個意思，
% cue 139 2100.480 2134.540
\hypertarget{c11:cue:139}{}那你的意思是不是真的這個意思，如果不是你就不能上寫。你說是基於總統作為國家元首，答答答答答。那意思就是我沒有在保障這位的首長。那這個人他有行政真正的當權，我也不保障他有豁免特權。這個講起來其實，你們真的覺得我們的豁免是給一個虛位元首用語，沒辦法回答。他要給他回應反對。好，他是這麼講的。以上，這是第一段。好，來，你們現在可以留給他的意思。來，第二段。

\tc{0:35:34}
\speaker{學生}
% cue 140 2134.560 2176.680
\hypertarget{c11:cue:140}{}那一段或外患罪以外之罪者暫時不能以為刑事上訴究，並非完全不適用刑法或相關法律之刑罰，而定固為一種暫時性之程序障礙，而非總統就其犯罪行為享有實體之免責權。\par\speaker{老師}好，這句話重不重要？這句話是在告訴你 388 很早很早以前的解釋。說這個 52 條要講的這種所謂的刑事豁免權是暫時性的。

\tc{0:36:16}
% cue 141 2176.680 2214.090
\hypertarget{c11:cue:141}{}暫時性的什麼意思？那總統如果卸任呢？照樣要追究。第二件事情我問你們，如果是總統就任以前犯的罪呢？有沒有豁免權？我的問題有點奇怪了，再看一遍這個句子。他說這個豁免權是暫時性的程序障礙，程序障礙，而不是享有實體免責權。這個意思就是說你做了不對的事情，這件事情不會抹殺掉，不會把被子白死，不會。那現在是暫時性的罪，我現在暫時不能
% cue 142 2214.090 2240.790
\hypertarget{c11:cue:142}{}究辦你，那你卸任總統我可以究辦你，那罪不會消滅。那我問的問題是，如果他在就任以前，就任總統以前做的事情，現在他是總統，你現在不能究辦他。那就任以前，做總統以前，他犯了罪，那你覺得可不可以究辦？

\tc{0:37:20}
% cue 143 2240.790 2251.770
\hypertarget{c11:cue:143}{}選上到就任，選上到就任，你的意思是說，他還沒就任，他還沒就任，他就不是總統，他就不會有總統特權的。
% cue 144 2251.770 2274.240
\hypertarget{c11:cue:144}{}我現在講的是比方說川普，川普在第一任的犯罪，第一任，中間隔了四年是拜登，現在川普二任，那二任的時候我要追究一任的事情，你已經卸任了，你現在已經不是總統了，那一任做的事情任不任，要不要追究。

\tc{0:37:54}
% cue 145 2274.240 2297.470
\hypertarget{c11:cue:145}{}像具體的例子，2021 年 1 月，他們叫做 conversional capital resource risk，你可以把想像成內外，他們是在這個 Essalvania 上面演講的。
% cue 146 2297.470 2307.390
\hypertarget{c11:cue:146}{}那現在那個不一任政府啊，把我們的學區偷走了，把我們的票偷走了啊，大家現在想請進，我稍後加入你們，加了一組，就衝進了國會山莊，就暴動了。
% cue 147 2307.390 2347.040
\hypertarget{c11:cue:147}{}這件事情，Trump 在電視上也講這是有錄影的，那你認為他有沒有煽動暴動罪？舉例說啊，那是他在總統一任的總統任內，他現在要二任了，我們可以會追究他，有聽懂嗎？好，我跟你們講的案子已經記起來，Trump vs. US，那這個案子您是記住了嗎？Trump vs. US，這個案子我簡單地講一下，美國聯邦最高法院，他說按照我們美國憲法這個制度，是實施權力分立的。

\tc{0:39:07}
% cue 148 2347.040 2368.480
\hypertarget{c11:cue:148}{}所以總統這個權利的本質，就賦予他享有前任總統，絕對免於刑事訴追，就其有關於他專屬的職權。
% cue 149 2368.480 2378.650
\hypertarget{c11:cue:149}{}然後在這件事情按照憲法，只有他能做，他是前任總統，前任總統也是在豁免的範圍，這是美國的意思。
% cue 150 2378.650 2384.660
\hypertarget{c11:cue:150}{}第二，這件事情如果是你按照憲法，只有你能做，別人都不能做。
% cue 151 2384.660 2415.620
\hypertarget{c11:cue:151}{}那這件事情你有絕對的免責權，絕對。那如果你是說我現在以總統的官方身份、官方行為，那這個時候你也應該享有推定的免責權，就推定你有免責權。只有在一種情況下，你沒有免責權，就是你所做的非官方行為。
% cue 152 2415.620 2423.690
\hypertarget{c11:cue:152}{}這邊有兩個意思，第一個就是針對前任總統，前任總統以前當總統的時候的行為，他也不能追究。
% cue 153 2423.690 2442.460
\hypertarget{c11:cue:153}{}然後他的行為分成三處，而是第二次意思。依照憲法只有總統能做的，那是絕對免責。那如果不是，是總統的官方行為，那這官方行為推定免責。如果完全不是官方行為，那是完全沒有免責權。三層。

\tc{0:40:42}
% cue 154 2442.460 2459.580
\hypertarget{c11:cue:154}{}這個案子裡面最有名的就是 Sotomayor 這個大法官，他發表了一個非常有趣的意見，他說他堅決地反對，給予總統絕對豁免跟推定豁免的權利。他說本人試舉以下的例子，請你們
% cue 155 2459.600 2496.800
\hypertarget{c11:cue:155}{}多數大統合回答我。總統今天以三軍統帥的身份，下令美國的海豹部隊軍人啊，去暗殺他的政治上的政敵，因為他跟我選舉，他跟我選一波，他殺了，照你們的說法，他應該想往後面，因為他是官方型的，這是不是太方便了。第二個，以川普來講，坦白大事，特色，他很想 Pardon,Pardon,Pardon。好，你各位一筆錢五百萬，就 Pardon，公然行賄。那也是他的官方行為，因為總統才有赦免權。

\tc{0:41:36}
% cue 156 2496.800 2531.070
\hypertarget{c11:cue:156}{}別人他才有赦免權，別人沒有赦免權，他就用錢來交易赦免，這樣你是不是也要不許追究？同時你說前任總統他過去的行為，只要他是總統，你通通都不追究，那這怎麼得了？聽得懂啊？好，我正在問的第二個問題，388 號解釋跟本號解釋，本號解釋沒有推改，沒有推翻 388 號解釋的見解，對不對？他說 388 號解釋的意思，表示他尊重，他沒有去改他，好，好，繼續，那我現在，
% cue 157 2531.070 2572.430
\hypertarget{c11:cue:157}{}那個問題是 388 號，還有本號解釋，如果是看到了 Trump vs. US，你認為我們大陸會不會維持現在這種限縮跟保守的見解，現在的見解很清楚，是暫時性的，同時不是實質的免責權，所以如果照這樣來講的話，能不能給予前任總統的行為，可能引號，因此說總統看起來好像是現任總統，現任總統以前當總統的行為，
% cue 158 2572.430 2606.740
\hypertarget{c11:cue:158}{}有沒有免贖？我剛剛問的問題是比這個狂，我說他在就任總統以前，不是總統的行為，聽懂得懂嗎？有一個殺人犯，他殺了一部雪上總統，那殺人罪要不要救我？要不要追究我？我現在當總統，如果照三八八號和本號解釋的意思，我在做總統任內，你不能追究，我卸任當然可以追究，而且這不是我總統任內犯的罪，是我總統任內以外犯的罪，那你當然可以追究，可是到美國那就不一定好了，只要是你的官方行為，大概就不能追究。

\tc{0:43:26}
% cue 159 2606.740 2646.650
\hypertarget{c11:cue:159}{}私人行為是個一個，這個解釋，Trump vs. US,2024 年，一般認為是美國最高法院，向當權者低頭的開始，有人說開始，低頭的明顯行為，要屈服啊，這個解釋被批評得很厲害，有人說這個解釋是美國憲政墮落的開始，你們聽聽就好，美國的大法官不是像一個想像中這麼神聖，他在歷史上多次向當權者低頭，沒什麼，所以臺灣的大法官向當權者低頭，你不要覺得有什麼，好，往下。

\tc{0:44:06}
\speaker{學生}
% cue 160 2646.650 2690.080
\hypertarget{c11:cue:160}{}第五十二條，規定：不受刑事上之訴究、係指、刑事偵查及審判機關於總統任職期間，就總統涉犯內亂或外患罪以外之罪者，暫時不得以總統犯罪嫌疑人或被告而進行偵查判程序。\par\speaker{老師}好，這個也沒有問題，但是現在有一句不懂，最後一句，對於總統行使職權無直接關涉的措施。

\tc{0:44:50}
% cue 161 2690.080 2726.490
\hypertarget{c11:cue:161}{}我不知道，他意有所指，他是在講什麼，不在此限，不在此限的意思，就是說你可以偵查，也可以究辦，我沒有在我的特權保障範圍，你沒有被豁免，你沒有被豁免，那你就要被犯了，該怎麼辦，該怎麼辦，該怎麼辦，該怎麼辦，該怎麼辦，該怎麼辦，該怎麼辦，該怎麼辦，就得認了，有一個是清楚的，對於犯罪現場的即時勘察，總統殺了人，舉例說，有個犯罪現場，你現在過去採訪，
% cue 162 2726.490 2761.910
\hypertarget{c11:cue:162}{}以後找不到證據，不是說你現在當總統4年，我不能究辦，那你當了 4 年以後，我也就辦不成啊，所以證據要先保留下來，所以這個不算實際，我只是不能起訴你，我不能將對調查，但是不是表示犯罪現場的即時勘察，立即的勘察，我不可以做，可是還有一個，無直接關涉、官設置措施，劫空相愛，酌食、廢戒，是吧，你們若看得懂，請你們指教，什麼意思，這段什麼意思？

\tc{0:46:01}
% cue 163 2761.910 2778.570
\hypertarget{c11:cue:163}{}大法官，這是我最不喜歡的，你講話不要那麼文謅謅，你說清楚你在講什麼，如果你意有所指，請你講清楚，不然這句話下次又要解釋了，還沒有講清楚，把它又要解釋了，解釋了解釋了解釋了，好，下一位再換一段，\par\speaker{學生}總統。
% cue 164 2778.570 2799.820
\hypertarget{c11:cue:164}{}不及於因他人刑事案件而對總統的證據調查與證據保全。\par\speaker{老師}好，再停一下，請問你這句跟上句有什麼差別，上一段的末句，我剛講的那句，跟這段的第一句有什麼差別，上一段看起來是總統的案子，總統殺人，我剛，
% cue 165 2799.820 2834.960
\hypertarget{c11:cue:165}{}總統殺人，那你要即時調查這個，現場保全證據，這是 OK 的。現在我講的是另外一件事，他說給予他人，他人這是吳淑珍，吳淑珍的案子，總統可不可以去，對總統可不可以做證據調查和證據保全。這個案子其實很微妙，其實他們倆是共同正犯，是一起有犯意，聯絡一起犯的事情，他不寫這句，因為總統豁免掉，他只能講對他人，這個他人不是普通人。
% cue 166 2834.960 2873.620
\hypertarget{c11:cue:166}{}是跟他有犯意聯絡的共同正犯，但是不管他現在講是他人不是總統，這樣有聽懂嗎？這是本案的重點，吳淑珍被起訴了，吳淑珍被審判了，他現在傳他總統來作證，可不可以？他現在是不是在調查證據跟保全證據？不然他傳他來看什麼？萬一總統死了咧？就死無對證了，那這個共同正犯就成立不了了，算了，我贏會吧，好，唯，繼續。

\tc{0:47:53}
% cue 167 2873.620 2914.250
\hypertarget{c11:cue:167}{}好，到這裡暫停，這句話寫的很爛，有沒有看得懂吧？
% cue 168 2914.250 2951.440
\hypertarget{c11:cue:168}{}我又亂拼拼嗎？他文中中的句子，一句這麼長，牛頭不對馬嘴。你的意思其實是要告訴我很簡單的事情。對於他人的刑事案件，你要要求總統做證據調查和證據保存，不屬於總統52條免責、豁免的範圍。這麼一句話就好了，要講這麼多。好，不屬於嘛。所以怎麼樣？可以做吧？然後，下一句就不通。不屬於。刑事豁免。

\tc{0:49:11}
% cue 169 2951.440 2991.800
\hypertarget{c11:cue:169}{}你所做的證據調查跟保全處分。往下居。句話是不是很莫名其妙。你剛才在講證據調查跟證據保存。你現在馬上跳到，人身自由。你的意思是說證人不可以把他抓起來嗎？人是有可能的，你們訴訟法比我好。怕證人逃跑，所以把證人給他們狙擊起來。在你沒有完成作證之前，不讓你逃走。或者證人有種可能，有些證人會被暗殺掉。但是如果一旦作證了，應該，
% cue 170 2991.800 3028.930
\hypertarget{c11:cue:170}{}你不怕你被殺，因為證據都已經不在乎，證詞都已經取得了，證據證物也拿到了，我還怕你被殺了，你要是架子，你死了好了。但是他說均不包含，均不得限制總統人身自由，例如，大大大大，亦不得妨礙總統職權的政策。到最後，我請問你們，這句話是什麼？你覺得這句話是什麼？這句話是說，總統你可以去作證，但是總統作證有一個範圍，你的人身自由不能被剝奪，你不能因為總統，
% cue 171 3028.930 3067.770
\hypertarget{c11:cue:171}{}總統作證，然後妨礙到你接待外賓，舉行閱兵典禮，所以法院要配合你，這個跟總統有豁免，給他豁免的理由，是要尊崇以及不妨礙他行使之間相關，如果你妨礙總統行使職權，就不可以，這樣才有聽懂嗎？這是目的解釋，我為什麼給你豁免？因為我要給你能夠自由行使權的空間，不要動不動就被法院，把你接整個上我就不配了，那如果居住你人身自由的位置，
% cue 172 3067.770 3112.500
\hypertarget{c11:cue:172}{}抓去關起來，你當總統，說不行，證人都不行，所以我要排除。雖然你有作證的義務，但是證人的動作還是不能給你做，因為他牴觸了我們授予你這個豁免特權的目的。他有聽懂嗎，我最後一句是大法官沒講的，他不講這句話就造成你的理解的困難，大聲一點，可以拒絕啊，等下你往下看就知道了，我們的職業員是很有趣的，他為了證明他的清白，他主動到場，他不是被傳，不是被拘提到場，他收到了傳票，
% cue 173 3112.500 3133.520
\hypertarget{c11:cue:173}{}叫他去作證，他就去了，他主動出席的，那現在這句話有一個問題是，剛才講，你去當證人，不是你豁免的範圍，你不可以主張，我有豁免權，那不是這個意思嗎？好，繼續。

\tc{0:52:38}
% cue 174 3158.400 3189.730
\hypertarget{c11:cue:174}{}因為我們的立法，相關的法律，特別是刑事訴訟法，沒有想到會有這一招啦，我們約總統豁免就豁免了，通通沒有了，通通不必了，但他現在有一種情況，有沒有搜索是他人的案子，不是，是吳淑珍的案子，不是陳水扁的案子，但是因為吳淑珍跟他是共同正犯，所以搜索吳淑珍就會發現他犯罪的事實，對不對？那如果要來搜索，大法官就想到，那搜索可不可以逮捕特定人，這個現在新竹現場在抗議什麼？

\tc{0:53:09}
% cue 175 3189.730 3225.800
\hypertarget{c11:cue:175}{}他在絕食，今天進入第四天，不是搜索就把人帶走了嗎？他有放人，那是以後的事，那他把人帶走，他有扣留特定物件，把他的電腦還有他的手機都扣留。現在警方最厲害就是這個，那個人是不是最厲害就是放在你的手機裡還有電腦裡。我抓走了這兩樣，我們的行銷網絡，一般人是行銷網絡，或者個人交往網絡，選舉的時候是競選總部的來往部線，不知道，現在警察，
% cue 176 3225.800 3265.760
\hypertarget{c11:cue:176}{}警察官會把這個資訊透露給他的行政對手。你有人知道說，偵查不公開，我們偵查每天在公開，禁報每天都有消息，這就是我們臺灣的問題，這就是王國昌的問題，可以你還是學術手法。這麼大的漏洞，你不去抹起來，你不會去修法，那你活。你現在已經看到了他的濫用，你現在就抗議絕食有個屁用，笨蛋國民黨，你決一死了，你絕食就算死了，於是也無補，他的條文長得這樣，他都沒有這個限制，你們覺得是可不可怕，動不動就來搜索你。

\tc{0:54:25}
% cue 177 3265.820 3301.810
\hypertarget{c11:cue:177}{}搜索呢，原本抓一段時間，然後時間到了放出來，他可以不放，他可以申請羈押，看看法院要不要配合，柯文哲不是被羈押了很多天嗎，柯文哲現在有什麼證據啊，柯文哲一通話說他是怨言不是嗎，那這個是另外一件事，你把我的東西扣走，這才可怕，扣走我的手機，扣走我的電話，我怎麼選，那要看他扣到的是不是關鍵能力，小孩很底層的，那就沒關係，總幹事你就夠夠幹，徐欣營都不知道他到底跟什麼人有聯絡，他就不信，
% cue 178 3301.830 3340.930
\hypertarget{c11:cue:178}{}這個選舉公平，然後你現在更發現有趣的事情是，民進黨只要一提出申請，法院馬上就懂，該查的查，該抓的查，國務機場提出申請了，慢慢等，那這個可以這樣嗎，這就是現在檢方被詬病的，所以我一直主張，檢察官的這個問題，要去做一個根本處理，他現在是敗壞國家法制基礎，重要來源，動工就罵司法不公，法院有沒有那麼不公，可以討論，但是我覺得綁在一起，是檢察官非法官，檢察官。

\tc{0:55:40}
% cue 179 3340.930 3378.240
\hypertarget{c11:cue:179}{}法官至少會有駁回，柯文哲請要申請，法官有駁回啊，是到了上訴以後又來一次那次輸掉啦，那次輸掉的原因，我們要看題喔，現在是細節我騙講的，聽說律師都很清楚，要什麼時候去申請戰役，因為戰役駁回他可以算時間，算時間還要知道，今天值班的法官是誰，法官會動，這也算是秘密，這件事又很有趣，這件事有沒有讓你感覺內神通外鬼？

\tc{0:56:18}
% cue 180 3378.240 3413.810
\hypertarget{c11:cue:180}{}照理講，檢審應該分開，我們院方的事情，檢審應該不知道，你不覺得是不是笑話，你是地方法院的院長，你今天法院法官輪值的名單，檢審怎麼會知道，從而他可以挑選一個他喜歡的人當法官的時候，給他來申請羈押，他比較會通過，這是笑話嗎，不知道這個笑話是什麼，我也看把今天司法院有什麼作為，總之呢，你們現在發現，你們要當國家的主人，更難處處，現在你們學了法律，恭喜你們，但是你們一定，
% cue 181 3413.810 3447.220
\hypertarget{c11:cue:181}{}一定要跟老師一樣，碰到了這個問題，要這樣去問，為什麼為什麼為什麼，這個沒咖咖，沒咖咖在哪裡，那你如果找不到這個，你就沒有辦法回答，你下次有可能等你被抓了，你可能被誣陷，我不是說你一定是壞人，你可能被誣陷了，我的朋友當中是有很成功的律師，被其他的律師誣告，他現在也慘了，他今天也有麻煩，這是一個複雜的故事，好啦，想太多了，來，我們現在往下看。

\tc{0:57:27}
% cue 182 3447.220 3467.880
\hypertarget{c11:cue:182}{}你有搜索是必要的，立法要做這些規定，好，到這裡合理吧，立法沒有規定嘛，那本案現在到底怎麼辦，來，請繼續，還沒有嘛，還沒有承認，之前呢，好，這是總統同意的哦，好，無論，

\tc{0:58:05}
\speaker{學生}
% cue 183 3485.350 3518.190
\hypertarget{c11:cue:183}{}搜索扣押之適當性與必要性，非經該特別合議庭裁定准許不得為之，但搜索之處所，應避免總統執行職務及居住之處所，及抗告程序，適用刑事訴訟法相關規定。\par\speaker{老師}好，來，我第一個問題是，這個地方說，可以經過高等法院五人法官組成的特別合議庭的准許，就他對總統進行證據，
% cue 184 3518.190 3557.210
\hypertarget{c11:cue:184}{}去調查及保全。這件事情就是大法官的發明。用法律這樣寫。我們法律沒有規定，他說你趕快去規定。在沒有規定之前，就是這五個人組成特別規定。這是不是大法官發明的。這件事情也可以來印證 114憲判1。這個算不算是大法官的程序自主權的一種解釋。他就沒有規定嘛。他沒有規定，本命法者規定要等到哪天天荒地老，他也不會規定嘛。那這時候怎麼辦，本案要不要進行。所以最高級的法官，最高級的法官，最高級的法官，
% cue 185 3557.210 3593.650
\hypertarget{c11:cue:185}{}最高級的法官，他就說，哦，我們規定組成一個特別合議，你們來處理。這樣子有聽懂了嗎？這可以回應我上次講的，一一四憲判一，那五個大法官會說，現在的法律規定很難，我們不要造法律規定，我們要程序自主，我們要鬆綁。看這個，這個解釋是最接近程序自主，但這個程序自主，是不是只有補充現行法制不足，而沒有辦法推翻現行法制拘束。民事訴訟法、刑事訴訟法，說你沒有辦法推翻現行法制拘束。

\tc{0:59:53}
% cue 186 3593.650 3628.380
\hypertarget{c11:cue:186}{}或許補充現行法制拘束。你高興你就說我程序自主我要還，那不得了，那會天下大亂。講話又聽懂了嗎？我又再講一次意見、拜託。我現在盡可能的，從委要關屁。我這裡講到，我想到，哎，把別案子跟這個有關，我也順便跟你講一下。好，以上是第一個問題。第二個問題，跟這個有關。那麼大法官做的這個發明，算不算，訴外裁判，請表示意見。你們都知道，法院的法官、審理案件，原則上，不能說原則上，基本上也不能說建設上，
% cue 187 3628.380 3656.800
\hypertarget{c11:cue:187}{}是禁止訴外裁判的強調，不可以的。對吧？那明文大法官這番話是他的發明，那意思是說這是訴外裁判咯，那大法官可不可以做訴外裁判咯。舉手回答我，有聽懂，我希望你們有聽懂，難嗎？

\tc{1:00:56}
% cue 188 3656.800 3673.850
\hypertarget{c11:cue:188}{}對比今天這個學期沒想到，老師的書有，現在我正在能力地改，希望在學期期末考試寫書可以出來讓你們有東西可以看可以操。這個訴外裁判，對於抽象解釋，是不適用的。
% cue 189 3673.850 3710.470
\hypertarget{c11:cue:189}{}因為大法官在很多個解釋裡面說，我們大法官要解釋什麼東西，就是訴之標的，我要處理哪些問題，不受申請人訴之聲明的拘束。申請人申請一直問我三點，但這三點我聽不懂，但這三點讓我看到一點，我覺得那一點很有價值，我就解釋那一點，這三點我 ignore completely，小朋友聽懂啊。所以訴外裁判對大法官不適用。理由，因為我們是抽象解釋，不是具體個案審查，具體個案審查。

\tc{1:01:50}
% cue 190 3710.470 3751.530
\hypertarget{c11:cue:190}{}可能，不能講可能，沒有辦法再說什麼。下面有清楚了，來，然後往下，等一下等一下等一下，還有第二個問題。為什麼，剛才最後一句話說，如果你要做搜索扣押，要避免在總統執行職務跟居住的處所，總統府不能搜索扣押，居住的地方也不能搜索扣押，為什麼？你們不是承認可以為他人的案件，去要求總統做證據調查跟證據保全嗎？如果這個東西正好就放在，
% cue 191 3751.530 3789.220
\hypertarget{c11:cue:191}{}在總統府，是不是？因為證據最安全了，誰敢來搜索？結果你大法官居然跟我說不能搜索，那我請問你這個證據，湮滅了怎麼辦？怎麼證據保存？我講的話聽不懂的舉手，你不要看我，你看文字，文字上說什麼？你要看得懂，我再問的是哪個句話？這裡是不是矛盾？你前面說總統的豁免權不及於在他案中要求總統做證據調查跟證據保存，那是他案，你有義務去配合，那只不過我不能抓你的人。

\tc{1:03:09}
% cue 192 3789.220 3825.310
\hypertarget{c11:cue:192}{}現在我要搜索扣押的東西，你還有處所的限制，你說不能在總統執行職務的處所，就是總統府吧，總統的官邸就是他居住的地方，這兩個都不能搜索，那我問你，如果這兩件證物正好在總統和官邸，或者總統的辦公室，那你怎麼辦？不能搜索，那不能搜索怎麼保全，讓他滅失嗎？所以這個地方沒有代表，對矛盾，你這兩個目的哪個比較重要？再也看，是證據保全比較重要的，不然你就不要另外允許對總統做證據，
% cue 193 3825.310 3860.300
\hypertarget{c11:cue:193}{}既然允許，你現在就是好人做一半，說不能到那裡搜索，不能到那裡扣押。這又是為什麼？那難道你要讓這兩樣東西，就在那個辦公室，或者在他家裡面，關地裡面，消失嗎？這不是一個笑話嗎？總統一目的落敗，就把長山關地，總統搜不到。總統他人，他太安，或者你跟總統有關係，就把這個證物放在總統家，那你也安全，搜不到。這很奇怪的原因，很奇怪的其實。就是這個大法官，向前總一步又算退後退半分。我希望是立平答案。

\tc{1:04:20}
% cue 194 3860.300 3902.440
\hypertarget{c11:cue:194}{}這句話是不需要寫的，多寫多找啊。你就看到這個大法官，多麼怕得罪訪。我們的大法官現在勇敢，因為美國勇敢一點。美國來不敢，美國是絕對保護跟推定保護，推定豁免。好，再下一段，為證人之義務，也不及於。剛才是證據保全。證據保全跟為證人，還是有一點距離，對不對？搜索扣押是證物啦，但是證人是叫證人講話才有證典啦，證言也可以做的，這個他說也不及於。
% cue 195 3902.440 3930.140
\hypertarget{c11:cue:195}{}好，唯！所以你看我在哪裡，我在辦公室，就到辦公室來，你不可以傳我到法院去，妨礙我的公務進行。好，這樣可以嗎？以示對總統的尊崇。好，再下一段。

\tc{1:05:30}
% cue 196 3930.140 3956.700
\hypertarget{c11:cue:196}{}所以總統的夫人、總統的孩子、總統的親戚，沒有這個問題講得好，叫原則不得拋棄，表示有可拋棄之例外嗎？
% cue 197 3956.700 3967.660
\hypertarget{c11:cue:197}{}今天看我這樣的問題，很奇怪嗎？你說原則不得拋棄，這表示你承認有可拋棄之例外嗎？
% cue 198 3967.660 3972.510
\hypertarget{c11:cue:198}{}那你告訴我，這句話是什麼意思？聽不懂，哪一種是原則，哪一種是例外？
% cue 199 3972.510 3979.880
\hypertarget{c11:cue:199}{}原則上不能拋棄，但是有例外可以拋棄，那你告訴我例外可以拋棄的情形。
% cue 200 3979.880 3990.780
\hypertarget{c11:cue:200}{}下法官，每一個地方都留下一些浮履，這些浮履就會造成很多的肆肉混蛋，因為你沒有辦法說清楚。好，以上是一大段，第二大段，請繼續。

\section{總統之國家機密特權}

\tc{1:06:30}
% cue 201 3990.780 4006.990
\hypertarget{c11:cue:201}{}現在你標題漏掉了，標題。剛才講的是刑事豁免權，現在講的是另外一件事叫做機密特權，又是一個特權。好，剛才 52 條
% cue 202 4007.080 4038.720
\hypertarget{c11:cue:202}{}有刑事豁免權。現在憲法哪一條講總統有國家機密特權？沒有啊。所以這個二，全部是大法官的發明。造法造給你看，你不要問我為什麼我不知道。但是我就跟你講，他造了法，這是事實。憲法沒有這條。往下，行政權三個字圈起來。繼續。
% cue 203 4038.720 4059.550
\hypertarget{c11:cue:203}{}到這裡，你要不要跟我們交代一下，國家機密特權是根據憲法哪一條？哪有個大法官解釋可以不要有依據的？他就不跟你依據啊。我都無得無了。

\tc{1:07:39}
% cue 204 4059.550 4071.480
\hypertarget{c11:cue:204}{}一條多。那你不覺得很奇怪嗎？第二點，這個話中有話。他說憲法或者增修條文賦予總統行政權的範圍。
% cue 205 4071.480 4080.780
\hypertarget{c11:cue:205}{}搞不好。那我請問你，如果憲法沒有賦予總統行政權，那是不是表示總統就不能主張國家機密特權？
% cue 206 4080.780 4093.880
\hypertarget{c11:cue:206}{}我就要講有錯嗎？你的前提是這個前提嗎？好，再下來更難問。那我請問你，我們憲法或憲法增修條文有賦予總統行政權嗎？是根據哪一條？
% cue 207 4093.880 4132.940
\hypertarget{c11:cue:207}{}公然會換來所有的說明嗎？你的前提是，如果有賦予他行政權，在那個範圍內，他可以主張國家機密。再下來一個問題，國家機密特權一定是限於國家安全、國防、外交有關的資訊嗎？一定嗎？不是國防、不是外交、不是國家安全，就不會影響國家安全的利益。他說可能影響國家安全與國家利益者，這個範圍很大誒，我懂你了，看來是特別的。如果今天我公布臺灣怎麼樣製造去，
% cue 208 4132.940 4173.710
\hypertarget{c11:cue:208}{}我們的良率可以這麼高，人民幣有沒有影響國家安全，看看你的臺積電線要跌多少，臺積電線就是抓這些東西，這個 know-how，我不會設計，我去設計，但是我非常會知道，我的良率可以到 90\%幾\%，你叫別人營料，他可能做 100 片有 60 片是壞掉的，所以你永遠競爭不過我，我賣給你好，這就是我的 know-how，但我可不可以把它公布呢？你們覺得我現在講話是開玩笑嗎？如果有一個總統跟川普一樣愛塞，那就是交易的問題嗎？我舉個例子。

\tc{1:09:33}
% cue 209 4173.710 4199.490
\hypertarget{c11:cue:209}{}你們美國，Intel，現在要擠棋子。你要花三億，才有可能銀號 MD，可能追到臺積電。三億要花多少億？三百億？舉例說，我是臺灣總統，我跟你說，小慧麻煩呀。你如果給我兩百億，那你覺得是算不算好。所以你剛才才這樣的定義，是不是太小？

\tc{1:09:59}
% cue 210 4199.490 4213.980
\hypertarget{c11:cue:210}{}太小了，太小了。好，那現在依照什麼總統的國家機密特權？他沒說，再問第二個問題。你們有沒有印象？我幫我，
% cue 211 4213.980 4252.130
\hypertarget{c11:cue:211}{}585 有講到行政特權。585 講行政特權，他說是行政固有的特權。所以不需要法律授權，不需要憲法授權。這一句話，為什麼現在沒講，表示說這個是哪裡來的？我聽不出。我講話有聽懂嗎？兩個很像。一個是總統585號有行政特權。行政特權是對付了一百月以外的調查中不公開的。現在有國家機密特權，還是要來對付法院。你要問我，我不公開。一樣的，都是不公開。
% cue 212 4252.130 4289.840
\hypertarget{c11:cue:212}{}五八五號說，是行政固有的特權。所以這不帶規定。固有就是我固有的。哪裡需要規定？你至少像這一句。那你這個案子沒有講他是固有的。那請問國家機密特權是哪裡來的？再來一個問題。五八五講的行政特權，跟本號解釋講的國家機密特權，兩者是什麼關係？是合一是同一？還是包含跟包含瑜的關係？如果是後者是哪個包含那個？不要跟我一樣，你多這麼多，我是不想這麼責任。

\tc{1:11:29}
% cue 213 4289.840 4332.390
\hypertarget{c11:cue:213}{}但是老師不是不可以考一個這樣的考試。我這個考試也很簡單。我這針對 627，我出 4 個小題，每題 5 分卻合。而且這項考試卻非常明瞞。你讀懂都沒有讀懂，一目瞭然，無所遁忌於天地間。我的題目是包吾四。但是我要這麼考，會很希望。不過好在我的題目都是私卻過不如。今天大概有 1/3 的沒有來，所以我們也可以考慮一下。好了，再唸下一段。我剛問的題目我沒有給答案了，你們自己要去想。

\tc{1:12:12}
% cue 214 4332.390 4386.480
\hypertarget{c11:cue:214}{}我問的問題都是很合理的問題，不是我算起來，是大法官自己做的。他前面不去看後面，後面忘記了前面。這回頭有什麼辦法？往下，發聲音。這才是他要講的。如果牽涉到國家機密特權，總統你即使做證人，你也可以拒絕證言。拒絕證言當然跟證物要綁在一起，如果有關的話。好，到這裡。現在顯然是沒有規定。好，往下，要件及程序。我們法律是不是就是要件跟程序。

\tc{1:13:06}
% cue 215 4386.480 4449.490
\hypertarget{c11:cue:215}{}要件是實體要件，程序就是程序要件。好，兩個，應該就這個，繼續。沒有啊，聽到沒有。釋明兩個字圈起來。現在老師請問一下，釋明跟證明有什麼差別？確信其為真實。確信其為真實的機率要大於確信其為不真實的機率，這樣才要完成決定責任。如果不然，那會輸掉。那釋明呢，你現在講的都很對，你有去查，這個釋明是訴訟，
% cue 216 4449.490 4487.890
\hypertarget{c11:cue:216}{}辦理司法院發布的法官辦理民事訴訟事件。應行注意事項裡面的第幾點。八十九點吧，應該是點，好像沒有講。記得哦，八十九，去自己看一下。掃詩上次有講過，這個不管你將來考不考法官，這個東西是法官辦案你拿在手上，看看你就會不會當法官了。民事訴訟法、民事訴訟，沒有講的細則，就用這些來的。司法院依照 530 號解釋，
% cue 217 4487.890 4527.520
\hypertarget{c11:cue:217}{}作為國家最高司法機關所享有的司法規則，自己權可以下級法院用的。但是，不能干涉下級法院、法官等個案的審判獨立。那是抽象一般的規定，不是針對個案的指示。我現在條例不懂嗎？沒有啊。好，請繼續。請，誰判斷有沒有合理釋明？這就把主持審掉了。當然是法院啊，或者是檢察官啊，我傳你來，我要你當證人，那就看法院檢察官，尤其是法院，他認為沒有合理釋明。

\tc{1:15:27}
% cue 218 4527.520 4584.130
\hypertarget{c11:cue:218}{}沒有合理釋明。你講的話，只要讓以為，大概是這樣子，連大概都做不到，好，繼續，跳過，這就是條文的規定，為，對啊，那你要做一個處分，或者做一個裁定啊，因為在訴訟進行中，我們通常是用裁定啊，那麼，如果總統對這個裁定有不服呢，繼續，總統，總統不服，不敗，相關法律沒有針對總統特別規定，但是大法官認為，長情嘛，你是一個裁定，那我當然可以抗告啊，
% cue 219 4584.130 4633.200
\hypertarget{c11:cue:219}{}剛才一樣，所以剛才創造的五人合議庭，可以管，到底要不要對他做這個，必可以管他如果抗告聲明，也要怎麼說。好，繼續，請問為什麼這麼規定，何以應停止，因為這個不是開玩笑的，總統告訴你說，公開會影響國家安全、國家利益，那你公開就見光死了，所以呢，原來裁定前，那你原處分說叫我公開，這個原處分要暫停，早看見哦，這樣你聽懂嗎，這是目的的考慮。

\tc{1:17:13}
% cue 220 4633.200 4681.300
\hypertarget{c11:cue:220}{}那你執行啦，那就公開了，公開了就沒有秘密可言了嘛，公開了還什麼秘密，就保不住了嘛，所以要先，在沒有確定以前，要先把它保住，不讓它公開，繼續，所以也是用其他規定，如總統如，好，剛才是沒有釋明，現在如果合理釋明的話，繼續，這句話什麼意思，什麼叫予以尊重，予以尊重就是，可以不提出，不證言，不提出證物，是這個意思嗎，那我再問一下，那如果不提出證物，不提出證言，
% cue 221 4681.300 4726.180
\hypertarget{c11:cue:221}{}他本案因此辦不下去，不管說檢察官沒辦法結案，起訴或不起訴，他沒有證物，他怎麼起訴，關鍵證物，那法院也判不下去了，那怎麼辦，這個地方其實是利益衝突，是國家機密，泛泛的國家機密比較優先，還是個案要發現真實做成裁判，正確裁判的公共利益比較優先，這兩個利益是有專注的，這裡大法官應該要做一個衡量，並且給個答案，大法官沒有講，好，繼續。

\tc{1:18:46}
% cue 222 4726.180 4760.600
\hypertarget{c11:cue:222}{}要做一個衡量程序，我也沒有當過法官，我也不知道，但是我的理解，法官很不喜歡處理這個案件，等下下面你在理由書有看到他有個名字，因為你處理到這種保密案件，那你要付保密義務，因為你摸過，你看過，你讀過，你破過這種卷宗，上次老師講過嘛，引進風之所以死，是因為引進風知道太多名，這個拉法葉戰艦到底是怎麼回事，現在全世界最奇怪的就是這個拉法葉的秘密，在法國，知道拉法葉的秘密的人，
% cue 223 4760.600 4802.370
\hypertarget{c11:cue:223}{}不是跳樓、刺殺，就是被推下去突然死亡，還有什麼刺殺的，你進風也好，所以你司法官，你喜歡審這種案子，總統告訴你說這個人關下的機密，你還堅持要審，好，依照一切依照保密程序來進行，依照保密程序，你碰過這些機密，所以接下來，你可能幾年不能出國，司法官，你喜歡控訴案子，現在就是要去看保密規定是怎麼規定，我現在講的是個大概，它下面有理由書裡有講的，全文，這個保密規定的全文。

\tc{1:20:02}
% cue 224 4802.370 4839.900
\hypertarget{c11:cue:224}{}但我的理解是，法官都很不喜歡沾到這個，沾到這個東西很討厭和麻煩。講難聽一點，法官也不能申請，他有嫌疑的時候有人威脅他生命都是可以的，但是如果人要去比生命，隨時都可以去啦。就算警察跟著你，有什麼披露？你覺得警察保護你的命就還在嗎？你相信嗎？你真的這麼相信嗎？還是你最好，吉古、無知、八古、無代，你都沒事。一旦知道你的麻煩就開始了，什麼時候發生你都不知道，你又看過電影情節你就曉得了。

\tc{1:20:39}
% cue 225 4839.900 4901.380
\hypertarget{c11:cue:225}{}好啦，講太多了。總統，唯一字開始的最後這一句，什麼意思？這句話你是法官，你是檢察官，你讀到這裡，請問你這個案子，你犯不犯的話，你講話啦，然後你也依照保密程序做啦，然後才加了最後一句。你如果知道這些東西，本來做成裁判或者是不起訴處分的基礎，它就會公開啊，裁判要公開啊，那就會有危害國家安全喔，那你應認為有妨礙國家安全利益之虞，最後這句話是聽不懂的。所以呢，所以本案是，
% cue 226 4901.380 4943.130
\hypertarget{c11:cue:226}{}本案是不偵察終結，還是不起訴，還是不做判決，那本案怎麼結，當事而我該到這裡，會發現我們當法官生氣吧，跟 520 一樣生氣，他不寫折疑，越寫的問題越多，多道理，不曉得他在講事。這裡其實又是一個利益衝突，即便是這樣子，我認為法官應該還是有裁量權，他要來判斷，這件事情雖然是有國家利益，但是本案的爭解，或者本案的審判所實現的個案公平正義，
% cue 227 4943.130 4979.630
\hypertarget{c11:cue:227}{}高於他，所以我還是要走，該起訴的起訴，該判決的判決。我現在是問你一個問題，尼克森總統當年去竊聽對手、民主黨、競選總部，他們設在水門，Watergate 是一個旅館的名字，在布套曼，波多馬克克皇帝，他去竊聽他們總部，然後最後被人家告他了，報紙告他，或者是 New York Times，還是 Washington Post 的英文，告他了。大家可以跟洋洋，越查越證，越查越證。他開始的時候，哪一個人要查我？

\tc{1:22:59}
% cue 228 4979.630 5017.810
\hypertarget{c11:cue:228}{}哪一個人要查我就 fire，哪一個人要查我就 fire，他也 fire 了好幾個。最後連聯邦的檢察總長去相當於我們的檢察總長，法官叫法務部長，是同一人。他也要傻，他又要 fire，但是你越 fire，他就越 fire 越嚴重了。好，最後呢，法院要做這個事情呢，特別檢察官，任命特別檢察官，再傻。特別檢察官命令他交出來，他不交，法師訴訟。訴訟的結果，美國聯邦最高法院，最高法院說，這個案子，也許行政特權。
% cue 229 5017.810 5059.930
\hypertarget{c11:cue:229}{}行政特權。也許有所謂的人。但是，這案子看起來，你公開他，可能造證定的所害，但是不至於妨礙總統，做為一個總統。你欠別人家的東西，把你公開，你的領域沒有面子，你還是總統，你哪一塊權力少掉了。所以，他一個是具體的，一個是抽象的。衡量的結果，他認為本案定刻，還是公開。公開如果不到兩個月，尼克森就下臺。尼克森，如果不下臺，一般預計，尼克森會是第一個被彈劾下臺的。

\tc{1:24:19}
% cue 230 5059.930 5100.470
\hypertarget{c11:cue:230}{}在這個之前，還發生過一件有趣的事情，我也是最近才知道。他出國去訪問的，風風風風風風。你們覺得，移轉外交注意力，就可以拯救你在國內的麻煩嗎？臺灣我不知道。現在看到美國，我今天講了很多美國的例子，我希望對你們思考問題有幫助。我不是說美國一定對，但你去看看人家是怎麼處理這樣的事情。等到臺灣，我在我目前判斷，很快，很快就會碰到。
% cue 231 5100.470 5111.330
\hypertarget{c11:cue:231}{}那這種情形，我們應該怎麼處理。你不要說它一定不能處理。更不要說它是政治問題。不予受理。好，來，再下一段。

\section{國家機密之核定與法院審理程序}

\tc{1:25:11}
\speaker{學生}
% cue 232 5111.330 5122.810
\hypertarget{c11:cue:232}{}總統刑事用國家機密保護法，法院辦理，涉及國家機密案件，保密作業辦。

\tc{1:25:22}
\speaker{老師}
% cue 233 5122.810 5130.460
\hypertarget{c11:cue:233}{}來，我剛剛講這個辦法的名稱全名在這裡。你們有請求人可以去查一查。你願意當法官，你會碰到怎麼處理。繼續。

\tc{1:25:30}
\speaker{學生}
% cue 234 5130.460 5137.540
\hypertarget{c11:cue:234}{}總統刑事總統是否已依國家機密保護法。

\tc{1:25:37}
\speaker{老師}
% cue 235 5137.540 5140.980
\hypertarget{c11:cue:235}{}好，跳跳跳，這麼多條文。依照這些法院規定，
% cue 236 5140.980 5149.130
\hypertarget{c11:cue:236}{}怎樣？核定。依照這些條文的規定，做什麼？核定。

\tc{1:25:49}
% cue 237 5149.130 5161.920
\hypertarget{c11:cue:237}{}這個意思就是說，總統你要主張說，我現在這一件不能公開，因為它牽涉到國家機密。
% cue 238 5161.920 5165.790
\hypertarget{c11:cue:238}{}你一定要問你自己，國家機密你核定了嗎？
% cue 239 5165.790 5170.270
\hypertarget{c11:cue:239}{}核定的等級是什麼？是極機密，還是機密，還是密？
% cue 240 5170.270 5178.730
\hypertarget{c11:cue:240}{}那保護期、保密期限不一樣。你如果都沒有核定為國家機密，你可以隨便說這是國家機密嗎？
% cue 241 5178.730 5191.860
\hypertarget{c11:cue:241}{}只因為我不爽，我都要拿出來。這句話到目前為止寫的是對的。一般國家處理國家機密也是這樣的。但是下面就有問題了。如，

\tc{1:26:31}
\speaker{學生}
% cue 242 5191.860 5208.540
\hypertarget{c11:cue:242}{}如尚未依法核定為國家機密者，無從適用上開規定及新冠資訊。\par\speaker{老師}是啊，你沒有核定，當然不構成國家機密。對不對？這到這裡應該可以懂啦。但是唯以下就不懂了。來，請繼續。

\tc{1:26:48}
\speaker{學生}
% cue 243 5208.540 5226.400
\hypertarget{c11:cue:243}{}惟訴訟程序進行中，總統如將系爭資訊依法改核定為國家機密，或另行提出其他已核定之國家機密者。法院應改依上開規定之相關程序，續行其審理程序。

\tc{1:27:06}
\speaker{老師}
% cue 244 5226.400 5237.890
\hypertarget{c11:cue:244}{}好，這句話馬上就出現一個問題啦。前面那句，你剛才念的前面那句，說如果你沒有核定為國家機密，你不能主張他是國家機密。
% cue 245 5237.890 5253.820
\hypertarget{c11:cue:245}{}你不能要求法院要適用保密程序，不公開審理。不可以。但是，注授中華一半，陳追悲說，誒，我現在主張這是國家機密。那我依照國家機密保護法，把它核定為極機密。
% cue 246 5253.820 5264.410
\hypertarget{c11:cue:246}{}我記得極機密是保護二十個保護，保密期限是二十五年。要去看國家機密保護，我現在相信不得，我記得是二十五年。好，那二十五年內都不能公開。
% cue 247 5264.410 5270.720
\hypertarget{c11:cue:247}{}然後他講價，接著呢，法院就慘了，法院立刻要改，依照相關程序去辦理。

\tc{1:27:50}
% cue 248 5270.720 5284.190
\hypertarget{c11:cue:248}{}那你覺得這樣規定合理嗎？有這樣規定合理嗎？為什麼訴訟程序當中，可以追加核定的國家機密？
% cue 249 5284.190 5296.060
\hypertarget{c11:cue:249}{}剛才不是說，沒有核定的不算嗎？說沒有關係，我臨時現在發現了，我馬上追加，那你也可以。追加以後就被人家一照著賣法去處理。這裡面是不是又有一個價值判斷？
% cue 250 5296.060 5310.370
\hypertarget{c11:cue:250}{}這裡告訴你是說，原則上不是有提及我，但是國家機密這件事情太特別了。我為了要保護國家機密，你只要一講他是國家機密，好，就能照國家機密的規定來處理。避免妨礙國家機密。
% cue 251 5310.370 5322.980
\hypertarget{c11:cue:251}{}所以這裡的價值判斷，國家機密好像永遠高一級，好，到這裡還沒有完，到這裡如果你看到，我覺得已經很不錯了。最後一句，比較難，來，起立。

\tc{1:28:42}
\speaker{學生}
% cue 252 5322.980 5338.530
\hypertarget{c11:cue:252}{}整項並不因而違反國家機密保護法及法院辦理涉及國家機密案件，法院辦法的程序規定。\par\speaker{老師}好，他不違反相關程序規定。這句話什麼意思？
% cue 253 5338.530 5372.360
\hypertarget{c11:cue:253}{}你當年很長啊，我現在幫你摘要一下已經進行的程序。在你核定為國家機密之前，我已經進行的程序，幫他以程序免資源來做證，我取得了證言。好，諸如此類，或者你提交了證據，或已經入存了證據。好，那這種情況，他說不因為，因為你後來追加，然後我不因為，當時沒有，當時沒有，也不產生索機箱，不違反程序規定。這樣子講人，聽懂嗎？所以這裡我有一個問題。

\tc{1:29:32}
% cue 254 5372.360 5407.660
\hypertarget{c11:cue:254}{}所謂已進行之間，是指，在本號解釋之前已進行者。像本案，本案在我做朋友解釋之前，已經進行了很多了。這個算不算，這個應該既往不咎吧，不然法官很難了。但法官一解釋，法官都犯罪了。要不然，由任愛同驚刑，總統核定追加回國家機密之前已經進行了，誰叫你後來補罪家。在你補罪家之前，你不能怪我嘛。這兩種解釋都有可能。那是不違反相關的程序，意思就是說，他取得了證言，取得了證物，任務，任務，任務。
% cue 255 5407.680 5438.020
\hypertarget{c11:cue:255}{}有效的證物證言，可以作為起訴或裁判的基礎嗎？是不是這個意思。後來這句話寫來幹什麼？有聽懂了嗎？暫時處分的部分就不講了，反正沒什麼。好，627 號解釋到這裡為止。算不算說完了？來，我們往下看。解釋理由書第一段。第一段呢，你只需要念最後一句。

\section{解釋理由書：總統刑事豁免權與中央政府體制}

\tc{1:30:38}
\speaker{學生}
% cue 256 5438.020 5448.080
\hypertarget{c11:cue:256}{}內容與範圍。內容與範圍。
% cue 257 5448.080 5459.680
\hypertarget{c11:cue:257}{}尚非憲法法理上之必然，而屬各國憲法政策之決定。\par\speaker{老師}好，請問你，大法官一開封民意，講這一段寫的寫出來。
% cue 258 5459.680 5474.500
\hypertarget{c11:cue:258}{}能讀到這段文字，是不是跟我有同樣的疑問。這段文字講跟不講，不影響，完全不影響，前面我們唸過的解釋文。那他為什麼要講這一段呢？
% cue 259 5474.500 5485.060
\hypertarget{c11:cue:259}{}解釋大法官很有學問嘛。必要的開場呢？沒有辦法回答，很有趣嘛。所以這是政策。

\tc{1:31:25}
% cue 260 5485.060 5522.560
\hypertarget{c11:cue:260}{}換句話講，憲法要不要規定國家的元首有刑事豁免權。depends。愛規定就規定，不規定就不規定。那我們有規定，就直接從 52 條開始。而且前面已經有 388，你就從 388 開始。這句話是，這段話是不是廢話？這段話寫就不行。完全不影響本案的討論。對吧？好，388 號解釋。幾乎全文，主要的部分都出現在第三段。有吧？那你們就自己看一下。現在比較有趣的是，
% cue 261 5522.560 5561.260
\hypertarget{c11:cue:261}{}從我們這一節課的觀點，我們是在講我們的政府體制。講權力分立。你看大法官，這是最接近的一次。他以前都是說，介於民意、政治、責任、政治、所以如何。永遠都這樣講。這是八個字，是三人。是三人。是金剛照。沒有問題講這八個字。然後講完了，他答案就出來了。那這一次他沒有講這八個字。來，念一下第四段。字。10 月 27 多次生日。好，暫停一下。他這個訊息有沒有錯誤？

\tc{1:32:41}
% cue 262 5561.260 5592.700
\hypertarget{c11:cue:262}{}我們修憲是 84 年才修了嗎？喂，上課課講過了，我們白講。大老闆也不知道，哪一個是清楚。我們第一次憲法增修。你就算不算，要至少第二次，他卻是 84 年算。好啦，我不跟他計較了啦。反正大口愛怎麼講都可以。你看他切斷跟我們切斷，時段不一樣。好，經歷自憲法增修，來。

\tc{1:33:12}
\speaker{學生}
% cue 263 5592.700 5598.310
\hypertarget{c11:cue:263}{}國尊政府體制雖有所更動。雖有所更動。\par\speaker{老師}好，聽起來講得很輕鬆。
% cue 264 5598.310 5629.080
\hypertarget{c11:cue:264}{}這是剛才的憲法律修正，他說雖然有修正，雖然有改變，然，但是重點是然。

\tc{1:33:49}
\speaker{學生}
% cue 265 5629.080 5659.200
\hypertarget{c11:cue:265}{}「憲法觀之」，總統人僅僅享有憲法庭、憲法、憲法、增修條文、所列之權限。而行政權人依憲法第 53 條規定概括授予行政院。\par\speaker{老師}好，我說是最接近的意思。所以，如果照這樣的理解，那我現在是什麼體制呢？他還是不回答，他還是堵開了。那沒有關係，回到這句話最關鍵的部分，總統只享有憲法跟憲法、增修條文、
% cue 266 5659.200 5691.750
\hypertarget{c11:cue:266}{}所列舉之權限。剛才不是講，總統如果要行使國家機密特權，他必須是依照憲法和憲法增修條文有賦予行政權的範圍內。那你是不是在這裡講，列舉的規定、列舉的權限裡面，有沒有給予總統行政權？什麼行政權？所以在這個行政權範圍內，他才能有國家機密特權，這才是重點嘛。要不然你講這句話幹什麼？廢話嘛。然後他說，其他的行政權還是概括了社會行政院。從他這樣子看起來，他好像，
% cue 267 5691.770 5724.830
\hypertarget{c11:cue:267}{}暗示認為，好像是暗示我們是雙首長制或者半總統制。總統先把賦予一部分的行政權，所以我剛才承認他有行政特權，有國家機密不予公開的特權。另一方面，行政院還是會有全大部分概括的行政權。是這個意思嗎？那如果是這樣子，那就非常像是半總統制啦。但是他不用這個工用，沒有關係。那你從這裡可以看到，他認為這個案子是在半總統制的背景下來所提示的。好，再往下。

\tc{1:35:24}
% cue 268 5724.830 5770.690
\hypertarget{c11:cue:268}{}好，叫我幾個字就找不到了。好，這也是他的認定。他認為總統豁免權就是給國家元首用，跟政府元首不關。那你現在什麼政府體制也沒有必然的觀點。這句話是對的不對。好，然後這句話是下面這句話。因此。二條規定，因憲法，已經多次修正而需另做他解。本院釋字第 388 號解釋並無變更。
% cue 269 5770.690 5803.810
\hypertarget{c11:cue:269}{}對，因為當時也有主張。說 388 號解釋要重來。我要申請補充解釋。我要推翻 388。但訪問是多一事不如少一事。沒有必要。就這個意思。好，大家可以了嗎？下面還有很多很多很多很多。那我要想講的分戰。第八段。原則上不得拋棄。我剛剛說，你們看得懂嗎？他有例外嗎？那你們這小組人要回答我啊。你們這次都老師幫你們做了，你們都沒有在做嗎？

\tc{1:36:43}
\speaker{學生}
% cue 270 5803.810 5806.640
\hypertarget{c11:cue:270}{}我們這次都老師幫你們聽說，我們幫你們聽一下。
% cue 271 5806.640 5814.680
\hypertarget{c11:cue:271}{}文言 14 第 388 號，解釋應職位的規範而定。
% cue 272 5814.680 5824.650
\hypertarget{c11:cue:272}{}因此擔任於總統職位之個人、就總統之豁免權保障、範圍內之各項特權、原則上不得拋棄。
% cue 273 5824.650 5840.490
\hypertarget{c11:cue:273}{}其子總統原則上不得事前概括拋棄其豁免權而言，以免刑事偵查、審判、程序對總統之尊崇與職權、
% cue 274 5840.490 5843.490
\hypertarget{c11:cue:274}{}有效行使，造成無可預見之干擾。

\tc{1:37:23}
\speaker{老師}
% cue 275 5843.490 5854.590
\hypertarget{c11:cue:275}{}好，所以是他怕有干擾。那你既然是豁免就豁免，不可以拋棄。那如果隨便拋棄，那就會干擾。那我也不懂了，那到底有沒有干擾。

\tc{1:37:34}
% cue 276 5854.590 5866.250
\hypertarget{c11:cue:276}{}好，但你看，這更有趣的是，\par\speaker{學生}唯，但總統之刑事豁免權、本質上為、總統之憲法上特權、行使總統職權者。
% cue 277 5866.250 5877.240
\hypertarget{c11:cue:277}{}只有個別證據調查的機會。事實上是否造成總統尊崇與職權行使之損傷或妨礙，應有其判斷餘地。

\tc{1:37:57}
\speaker{老師}
% cue 278 5877.240 5883.370
\hypertarget{c11:cue:278}{}好，那這句話是不是矛盾？你說總統可以判斷，然後你又說原則上不得厚免。
% cue 279 5883.370 5890.780
\hypertarget{c11:cue:279}{}那這兩句話是不是矛盾？你幫我可以判斷啊？你幫我當笨蛋啊？怎麼說我不能厚免？
% cue 280 5890.780 5899.290
\hypertarget{c11:cue:280}{}通啊，你知道不通啊。然後他現在講這些話，都是要圓一件事、一個事情。
% cue 281 5899.290 5910.290
\hypertarget{c11:cue:281}{}你幫他只是一個自動放棄，那你該怎麼辦？簡單講就是這樣。我們大法官真可憐，就是鐵零插拼，往下。

\tc{1:38:30}
\speaker{學生}
% cue 282 5910.290 5926.640
\hypertarget{c11:cue:282}{}故除以總統為被告之刑事起訴與審判程序，或其他客觀上足認必然造成尊崇之損傷與職權行使之妨礙者外，其餘個別證據調查行為，縱為總統刑事豁免權所及。

\tc{1:38:46}
\speaker{老師}
% cue 283 5926.640 5933.050
\hypertarget{c11:cue:283}{}縱為他所及，表示他就是可以豁免啦。剛才你說不能，不能夠概括，
% cue 284 5933.050 5939.970
\hypertarget{c11:cue:284}{}概括了，原則上不能拋棄，事前概括了拋棄。好，但是你現在是已經說了，然後呢，違。
% cue 285 5939.970 5965.850
\hypertarget{c11:cue:285}{}好，這句話是不是，實際上是憋屈的論文在憋屈呢？

\tc{1:39:25}
% cue 286 5965.850 6004.380
\hypertarget{c11:cue:286}{}我覺得這裡有幾種可能，有一種可能，你就不要去處理這一段，你不要講原則上不能拋棄，他已經拋棄了，你就不要去處理他拋棄，就讓他拋棄了，然後你現在要說原則上不能拋棄，不過你個人如果拋棄了，拋棄你也沒有關係，那是廢話，那以後每個總統都要拋棄，你還有什麼原則上不能拋棄，聽不懂，這個就是沒有必要處理，我再講一個比較白的話，請問一個有權力的人，能不能拋棄權力，權力可不可以放棄，可不可以不行使？
% cue 287 6004.380 6046.000
\hypertarget{c11:cue:287}{}同理，我有特權，特權可不可以不行使？真說，原則上不能拋棄，但是你已經做了，但我就認為你也可以。講這個話，你要是要做什麼呢？維持法律的贏念嗎？還是維持什麼東西的尊嚴，我看不懂啊。要是我，我這一段我根本不行，這一段都寫了。此地無銀三百兩，要來解釋，你要常常問你自己，就像我剛開始，我第一段要不要寫一樣的。好，那現在我們唸一下第九段，唸完文就結束了。

\section{解釋理由書：總統之國家機密特權}

\tc{1:40:46}
\speaker{學生}
% cue 288 6046.000 6051.890
\hypertarget{c11:cue:288}{}國家機密特權，憲法並未明文規定。\par\speaker{老師}沒有規定。繼續。

\tc{1:40:51}
\speaker{學生}
% cue 289 6051.890 6062.910
\hypertarget{c11:cue:289}{}國家機密特權，惟依權力分立與制衡原則，行政首長依其固有之權能。\par\speaker{老師}又來了哦，依其固有之權能哦，有哦，跟 585 一樣哦，繼續。

\tc{1:41:02}
\speaker{學生}
% cue 290 6062.910 6079.130
\hypertarget{c11:cue:290}{}就有關國家安全、國防及外交之國家機密事項，有決定不予公開的權力，屬行政首長行政特權、本院釋字第 585 號，解釋足資參照。

\tc{1:41:19}
\speaker{老師}
% cue 291 6079.130 6094.320
\hypertarget{c11:cue:291}{}好！這句話有很多矛盾。第一，剛才說總統的行政特權是給予國家元首，不是給予行政首長的。那你現在自己說行政首長就有國家機密特權。
% cue 292 6094.320 6098.450
\hypertarget{c11:cue:292}{}第二個，你看這段話，回不回，能不能回答我剛才的問題。
% cue 293 6098.450 6110.310
\hypertarget{c11:cue:293}{}585 號講的行政特權，跟現在 627 號講的國家機密特權，一而二、二而一，兩個根本就是一件事。

\tc{1:41:50}
% cue 294 6110.310 6123.890
\hypertarget{c11:cue:294}{}還是兩個是誰包含與誰？選的是哪一個？這就是我們大法官的問題。你不是不可以寫，但是你寫了，你要考慮到你以前講過什麼話。
% cue 295 6123.890 6139.370
\hypertarget{c11:cue:295}{}那前面講的 585 號講的行政特權，跟你本號解釋 627 講的國家機密特權，什麼關係？要不要講清楚？你看這段話就講這樣子，這是什麼意思？
% cue 296 6139.370 6151.850
\hypertarget{c11:cue:296}{}接下來呢，它要解釋憲法什麼地方賦予它權利。這個解釋呢，非常的恐怖。我認為這個解釋呢，都非常不及格，就是下面這一段。

\tc{1:42:31}
% cue 297 6151.850 6169.000
\hypertarget{c11:cue:297}{}但是這個，其實吧，如果我要考，我會考一個很難的節目的話，只是唯一格。請你，就你所收集到的或者你所知道的大法官解釋，來論述一下，我們大法官，究竟認為我們現在是哪一種政府體制？
% cue 298 6169.000 6219.550
\hypertarget{c11:cue:298}{}這裡有很多其實，387、419、520，這個就分得寫了，蠻有得寫的。我其實考試只要考這題，但是我認為這樣考下來，但這個題目好。如果真的這麼考，你又考得好，至少恭喜，那你真的就這個學習對的主題，這個學習對的主題。今天時間到了，在他關燈以前，主要結束了。請問一下，那個 114 線單一的，你繼續準備哦，我總會改到時間抓到你的哦，來，我的班長，有什麼事情要宣布？好，其他同學還有問題嗎？

\tc{1:43:39}
% cue 299 6219.550 6224.720
\hypertarget{c11:cue:299}{}啊，如果有趁老師下課之前，你還可以問，有嗎？
% cue 300 6224.720 6264.510
\hypertarget{c11:cue:300}{}剛才來念的同學啊，哪一個同學剛才教我那個，我現在要把它關掉，所以你再來教我，把 627 找出來，我要儲存到正確的位置。我現在不知道 627 跑到哪裡去了，我已經把它關掉了，希望我來找到 TOMYTOL。先點完你的名字。
% cue 301 6264.510 6293.440 刪：whisper 幻覺字幕（「中文字幕 李宗盛」）

\backmatter
\printbibliography[heading=bibintoc,title=參考文獻]
\end{document}
